% =====================================================================
% Manuscript/main.tex
%
% Agent 19 (Final Assembly Agent) deliverable.
% Single coherent TeX manuscript of "The sharp quantitative
% Faber--Krahn inequality: two proofs".
%
% Compile with:
%     pdflatex main
%     bibtex   main
%     pdflatex main
%     pdflatex main
% =====================================================================

\documentclass[11pt,a4paper]{amsart}

% --- Core packages -----------------------------------------------------
\usepackage[utf8]{inputenc}
\usepackage[T1]{fontenc}
\usepackage{amsmath,amssymb,amsthm,mathtools}
% \usepackage{microtype} % Disabled: pdftex auto-expansion requires scalable fonts
\usepackage[protrusion=false,expansion=false]{microtype}
\usepackage{enumitem}
\usepackage[dvipsnames,svgnames,table]{xcolor}
\usepackage{hyperref}
\usepackage[capitalise,noabbrev]{cleveref}

\hypersetup{
  colorlinks=true,
  linkcolor=blue!50!black,
  citecolor=blue!50!black,
  urlcolor=blue!50!black
}

% --- Canonical macros (Agent 17) ---------------------------------------
% =====================================================================
% Manuscript/macros.tex
% Canonical TeX preamble macros for the unified manuscript.
% Produced by Agent 17 (Notation and Normalization Unification, 2026-05-13).
% This file is to be \input by Manuscript/main.tex (Agent 19).
%
% The macros are taken verbatim from the canonical Notation Register
% §15 (Manuscript/00_notation_register.md). No draft introduces its
% own \newcommand definitions; every macro used below appears in some
% draft and is consistent across drafts.
%
% Package requirements (loaded in main.tex):
%   amsmath, amssymb, mathtools (for \DeclareMathOperator and
%   common math symbols).
% =====================================================================

% --- Ambient space and basic symbols ---
\newcommand{\R}{\mathbb{R}}
\newcommand{\Om}{\Omega}
\newcommand{\Otilde}{\widetilde{\Omega}}
\newcommand{\eps}{\varepsilon}

% --- Asymmetry, balls, deficits ---
\newcommand{\Fr}{\mathcal{A}}                  % canonical Fraenkel asymmetry symbol
\newcommand{\Bstar}{B^{*}}                     % equal-volume ball at the origin
\newcommand{\deltaT}{\delta_{T}}               % Saint-Venant deficit
\newcommand{\deltaFK}{\delta_{\mathrm{FK}}}    % Faber-Krahn deficit
\newcommand{\BDValpha}{\alpha_{\mathrm{BDV}}}  % BDV smooth asymmetry

% --- Foliation radii and centre fields ---
\newcommand{\rstar}{\rho_{*}}
\newcommand{\rhad}{{\widehat{\rho}}}
\newcommand{\rhodelta}{\rho_{\delta}}
\newcommand{\zbar}{\bar z_{\rho}}              % bulk centroid of E_rho

% --- Level-set foliation and velocities ---
\newcommand{\Erho}{E_{\rho}}
\newcommand{\Frho}{F_{\rho}}
\newcommand{\Vrho}{V_{\rho}}
\newcommand{\nurho}{\nu_{\rho}}

% --- Metric quotient space ---
\newcommand{\X}{\mathcal{X}}
\newcommand{\dX}{d_{\mathcal{X}}}

% --- Pi block (auxiliary, Agent 14) ---
\newcommand{\PiE}{\Pi(E_{\rho},\zbar)}
\newcommand{\OmPi}{\Omega_{\Pi}}

% --- Operators ---
\DeclareMathOperator{\dist}{dist}
\DeclareMathOperator{\Div}{div}
\DeclareMathOperator{\Lip}{Lip}
\DeclareMathOperator{\diam}{diam}
\providecommand{\hH}{\mathcal{H}}
\providecommand{\NS}{\mathrm{NS}}
\providecommand{\fint}{\Xint-}
\def\Xint#1{\mathchoice{\XXint\displaystyle\textstyle{#1}}%
   {\XXint\textstyle\scriptstyle{#1}}%
   {\XXint\scriptstyle\scriptscriptstyle{#1}}%
   {\XXint\scriptscriptstyle\scriptscriptstyle{#1}}%
   \!\int}
\def\XXint#1#2#3{{\setbox0=\hbox{$#1{#2#3}{\int}$}
     \vcenter{\hbox{$#2#3$}}\kern-.5\wd0}}


% --- Theorem environments (numbered by section) ------------------------
\theoremstyle{plain}
\newtheorem{theorem}{Theorem}[section]
\newtheorem{lemma}[theorem]{Lemma}
\newtheorem{proposition}[theorem]{Proposition}
\newtheorem{corollary}[theorem]{Corollary}
% Some drafts use the {theorem*} starred environment; provide it too.
\newtheorem*{theorem*}{Theorem}

\theoremstyle{definition}
\newtheorem{definition}[theorem]{Definition}
\newtheorem{example}[theorem]{Example}
\newtheorem{assumption}[theorem]{Assumption}

\theoremstyle{remark}
\newtheorem{remark}[theorem]{Remark}
\newtheorem{notation}[theorem]{Notation}
\newtheorem{hypothesis}[theorem]{Hypothesis}

% --- Cleveref customisation --------------------------------------------
\crefname{equation}{}{}
\Crefname{equation}{Equation}{Equations}

% --- Title -------------------------------------------------------------
\title{The sharp quantitative Faber--Krahn inequality:\\
two proofs}
\author{}
\date{\today}

\begin{document}

\begin{abstract}
For every open set $\Om\subset\R^n$ with $0<|\Om|<\infty$ the
sharp quantitative Faber--Krahn inequality
$\Fr(\Om)^2\le C_{\rm FK}(n)\,\deltaFK(\Om)$ holds, with
$C_{\rm FK}(n)$ depending only on the dimension. We give two
independent proofs of this inequality, written with constants tracked
through every step. The first proof follows the
Brasco--De~Philippis--Velichkov selection-principle route, including
the graph-entry / Schauder bootstrap, the BDV nearly-spherical
closure, bounded reduction, and the Kohler--Jobin transfer. The
second proof is structurally different: an exact level-set deficit
identity expresses the Saint--Venant deficit as an integral over the
volume-radius foliation of two pointwise non-negative defects; a
finite-perimeter metric first-variation in the quotient space $\X$
of sets modulo translation, an oriented $L^1$-radial trace lemma at
Fusco--Julin oscillation centres, a good/bad level decomposition,
and a boundary-layer transfer convert this integral identity into the
same global sharp Saint--Venant inequality and, after Kohler--Jobin,
into the same final Faber--Krahn estimate with identical dimensional
constant. The two proofs are presented in parallel so that their
common closing step and the precise place where they diverge are
explicit.
\end{abstract}

\maketitle

% =====================================================================
% Introduction (Agent 18 deliverable)
% =====================================================================
% =====================================================================
% Manuscript/05_intro_story_draft.tex
%
% Agent 18 (Introduction and Historical Story) deliverable.
%
% Conventions:
%   - Canonical macros loaded from Manuscript/macros.tex (Agent 17).
%   - Citation keys are exactly the ones listed in
%     Manuscript/04_references_audit.md / Manuscript/references.bib.
%   - Notation matches Manuscript/00_notation_register.md.
%   - Theorem statements at the end of the introduction are quoted
%     verbatim (up to display formatting) from
%     Manuscript/plan1_kohler_jobin_transfer.tex Theorem~thm:KJ-FK-lower
%     /\,thm:KJ-FK-upper and
%     Manuscript/plan2_boundary_layer_assembly.tex Theorem~thm:plan2-FK.
%   - The Bernoulli spectral route (Plan~1) is mentioned only as an
%     optional remark; the centroid / space-time $\Pi$ route (Plan~2)
%     is mentioned only as an auxiliary diagnostic.
% =====================================================================

\section{Introduction}
\label{sec:intro}

% =====================================================================
\subsection{The sharp quantitative Faber--Krahn inequality}
\label{ssec:intro-headline}
% =====================================================================

Let $n\ge 2$ and let $\Om\subset\R^n$ be an open set with
$0<|\Om|<\infty$. Write $\lambda_1(\Om)$ for the first Dirichlet
eigenvalue of $-\Delta$ on $\Om$ and $\Bstar$ for the open ball
centred at the origin with $|\Bstar|=|\Om|$. The Faber--Krahn
inequality~\cite{Faber1923,Krahn1925,Krahn1926} asserts
\begin{equation}\label{eq:intro-FK}
   \lambda_1(\Om)\;\ge\;\lambda_1(\Bstar),
\end{equation}
with equality if and only if $\Om$ is a ball (up to a Lebesgue
negligible set). The scale-invariant deficit
\begin{equation}\label{eq:intro-deltaFK}
   \deltaFK(\Om)\;:=\;
   \Bigl(\tfrac{|\Om|}{\omega_n}\Bigr)^{\!2/n}\!
   \bigl(\lambda_1(\Om)-\lambda_1(\Bstar)\bigr)\;\ge\;0
\end{equation}
vanishes only at balls; the Fraenkel asymmetry
\begin{equation}\label{eq:intro-Fr}
   \Fr(\Om)\;:=\;
   \inf_{x_0\in\R^n}\frac{|\Om\,\triangle\,(\Bstar+x_0)|}{|\Bstar|}
   \;\in\;[0,2)
\end{equation}
measures the distance from~$\Om$ to a ball in $L^1$. Both proofs in
the present paper close the following sharp quantitative statement,
in the canonical normalisation fixed in
\S\ref{sec:notation}.

\begin{theorem*}[Sharp quantitative Faber--Krahn]
There is a constant $C_{\rm FK}(n)>0$, depending only on the
ambient dimension $n\ge 2$, such that for every open
$\Om\subset\R^n$ with $0<|\Om|<\infty$,
\begin{equation}\label{eq:intro-main}
   \Fr(\Om)^2\;\le\;C_{\rm FK}(n)\,\deltaFK(\Om).
\end{equation}
The exponent $2$ on $\Fr(\Om)$ is sharp, in the sense that no power
greater than $2$ on the left-hand side can hold with a constant
depending only on $n$.
\end{theorem*}

The theorem is the main result of
Brasco--De~Philippis--Velichkov~\cite{BDV} (BDV) and was reproved by
Allen--Kriventsov--Neumayer~\cite{AKN2023} via the Alt--Caffarelli--Friedman
monotonicity formula. Our contribution is to give \emph{two independent
proofs}, each with full quantitative tracking of constants, written so
that they can be read in parallel and compared step-by-step. We refer
to these as \emph{Plan~1} and \emph{Plan~2}; both proofs end at the
identical theorem statement~\eqref{eq:intro-main}, with the same
explicit dimensional constant
\begin{equation}\label{eq:intro-cFKdef}
   c_{\rm FK}(n)=1/C_{\rm FK}(n)
   \;=\;\min\!\Bigl\{
      \frac{4 L_n\,c_{\rm SV}^{\rm glob}(n)}{(n+2)\,T_n},\;
      \frac{L_n\,(2^{2/(n+2)}-1)}{4}
   \Bigr\},
\end{equation}
where $L_n=\lambda_1(B_1)$ and $T_n=T(B_1)=\omega_n/(n(n+2))$ are
dimensional (cf.\ Theorems~\ref{thm:plan1-FK-upper-intro}
and~\ref{thm:plan2-FK-intro} below).

The two proofs differ \emph{only in how the bounded sharp Saint--Venant
inequality is obtained}: Plan~1 follows the selection-principle route
of~\cite{BDV}, supplemented with detailed graph-entry, Schauder
interpolation, and nearly-spherical closure; Plan~2 follows the
level-set / metric-quotient route inspired by
Nicola--Tilli~\cite{NicolaTilli2022,GGRT2024} and uses Fusco--Julin
oscillation~\cite{FJ2014} together with a finite-perimeter first-variation
identity in a metric quotient space, an oriented radial-trace estimate,
and a boundary-layer transfer. From the global sharp Saint--Venant
inequality onwards, the two proofs share the same Kohler--Jobin
transfer~\cite{KohlerJobin1978,Brasco2014} and bounded reduction (BDV
Lemma~5.3). The aim is to make both routes explicit, with constants
tracked, so that the reader sees clearly what each route requires and
what each route delivers.

% =====================================================================
\subsection{History I: Rayleigh's conjecture, Faber, Krahn}
\label{ssec:intro-history-fk}
% =====================================================================

The inequality~\eqref{eq:intro-FK} was conjectured by Lord Rayleigh in
his \emph{Theory of Sound} (1894/96) on physical grounds: among
homogeneous membranes of given area and uniform tension, the disc
produces the lowest fundamental tone. In two dimensions Rayleigh's
conjecture was proved by Faber~\cite{Faber1923} in 1923 by
Steiner-symmetrisation of the Rayleigh quotient: rearranging the trial
function $u\in H^1_0(\Om)$ as the radially decreasing equimeasurable
$u^*$ on $\Bstar$ decreases $\int u^2$ at fixed level distribution and
decreases $\int|\nabla u|^2$ by the Pólya--Szeg\H o
principle~\cite{PolyaSzego1951}, so the Rayleigh quotient does not
increase. In 1925 Krahn~\cite{Krahn1925} extended the proof to $n=2$
(in a slightly different form) and in 1926~\cite{Krahn1926} to all
dimensions $n\ge 3$, again by symmetrisation. The argument is the
foundational application of decreasing rearrangement to spectral
geometry; modern textbook treatments are given by
Henrot~\cite{HenrotBook}, Bucur--Buttazzo~\cite{BucurButtazzo},
Bandle~\cite{Bandle1980}, Kawohl~\cite{KawohlBook} and
Kesavan~\cite{Kesavan2006}.

The rigidity statement --- equality in~\eqref{eq:intro-FK} forces $\Om$
to be a ball --- is also classical. In the modern setting of sets of
finite perimeter (introduced by De~Giorgi~\cite{DeGiorgi1958}), rigidity
in the isoperimetric inequality is the key input
(\cite[Chapter~19]{Maggi2012}); see also
Brothers--Ziemer~\cite{BrothersZiemer1988} for the equality case in the
Pólya--Szeg\H o rearrangement principle.

% =====================================================================
\subsection{History II: the Saint--Venant analogue}
\label{ssec:intro-history-sv}
% =====================================================================

A parallel functional inequality of the same vintage governs the
torsional rigidity of an elastic bar. For an open
$\Om\subset\R^n$ with $0<|\Om|<\infty$, the (variational) torsion
function $u_\Om\in H^1_0(\Om)$ is the unique minimiser of
\[
   v\;\longmapsto\;
   \tfrac{1}{2}\!\int_\Om|\nabla v|^2\,dx-\int_\Om v\,dx,
\]
equivalently $-\Delta u_\Om=1$ in $\Om$ with $u_\Om=0$ on $\partial\Om$
in the trace sense, and $u_\Om\ge 0$ a.e. The torsional rigidity is
\begin{equation}\label{eq:intro-T}
   T(\Om):=\int_\Om u_\Om\,dx
   \;=\;-2\!\int_\Om E_\Om\,dx,
\end{equation}
where, following the BDV normalisation of the notation register, the
Saint--Venant energy is
\begin{equation}\label{eq:intro-E}
   E(\Om)\;:=\;-T(\Om)/2
   \;=\;
   \min\!\Bigl\{\tfrac{1}{2}\!\int|\nabla v|^2-\!\int v\,:\,
      v\in H^1_0(\Om)\Bigr\}\le 0.
\end{equation}
Pólya and Szeg\H o~\cite{PolyaSzego1951} established the Saint--Venant
analogue of Faber--Krahn:
\begin{equation}\label{eq:intro-SV}
   T(\Om)\;\le\;T(\Bstar),
   \qquad
   E(\Om)\;\ge\;E(\Bstar),
\end{equation}
with equality only if $\Om$ is a ball (up to translations and
negligible sets). Again the proof is by decreasing rearrangement; we
adopt the BDV-normalised deficit
\begin{equation}\label{eq:intro-deltaT}
   \deltaT(\Om)\;:=\;E(\Om)-E(\Bstar)\;\ge\;0,
\end{equation}
which vanishes only at balls. Modern accounts of the Saint--Venant
inequality and its rearrangement proof are
in~\cite{KawohlBook,Bandle1980,Kesavan2006,HenrotBook,PolyaSzego1951}.
In the BDV
literature~\cite[\S2]{BDV},~\cite[\S2]{Brasco2014},~\cite[\S2]{BraDeP2017},
the quantitative form of~\eqref{eq:intro-SV} is the natural Saint--Venant
companion of Faber--Krahn, and is the form actually proved by both
of our routes before transferring to the eigenvalue side.

% =====================================================================
\subsection{The quantitative stability problem and the sharp exponent}
\label{ssec:intro-sharp-exponent}
% =====================================================================

The qualitative rigidity statement --- equality in~\eqref{eq:intro-FK}
or~\eqref{eq:intro-SV} forces $\Om$ to be a ball --- raises the
quantitative question: for $\Om$ \emph{close} to a ball, how does
$\deltaFK(\Om)$ (or $\deltaT(\Om)$) compare to $\Fr(\Om)$?

A volume-preserving ellipsoidal one-parameter family $\Om_\eps$ of
perturbations of $\Bstar$ satisfies $\Fr(\Om_\eps)\sim\eps$ while
$\lambda_1(\Om_\eps)-\lambda_1(\Bstar)\sim\eps^2$ and
$E(\Om_\eps)-E(\Bstar)\sim\eps^2$ as $\eps\to 0$; see
\cite[Introduction]{BraDeP2017} and~\cite[\S2]{BDV}. Hence the best
possible quantitative form of Faber--Krahn / Saint--Venant is
\begin{equation}\label{eq:intro-sharp-form}
   \Fr(\Om)^2\;\le\;C_{\rm FK}(n)\,\deltaFK(\Om),
   \qquad
   \Fr(\Om)^2\;\le\;C_{\rm SV}(n)\,\deltaT(\Om);
\end{equation}
the exponent $2$ is sharp. The first quantitative stability results
in the eigenvalue case~\eqref{eq:intro-FK} are due to
Hansen--Nadirashvili~\cite{HansenNadirashvili1994},
Melas~\cite{Melas1992}, and Bhattacharya~\cite{Bhattacharya2001},
all with non-sharp exponents. Fusco, Maggi, and
Pratelli~\cite{FMPLaplace}, building on the sharp quantitative
isoperimetric inequality of~\cite{FMP2008}, obtained
\begin{equation}\label{eq:intro-FMP-FK}
   \Fr(\Om)^4\;\le\;C(n)\,\deltaFK(\Om)
   \quad\text{and}\quad
   \Fr(\Om)^4\;\le\;C(n)\,\deltaT(\Om);
\end{equation}
see~\cite[Theorems~1.1, 1.4]{FMPLaplace}. The sharp
form~\eqref{eq:intro-sharp-form} is the content of~\cite{BDV} for the
ball case in any dimension $n\ge 2$; see also the survey of
Brasco--De~Philippis~\cite{BraDeP2017} and the alternative proof of
Allen--Kriventsov--Neumayer~\cite{AKN2023}. The sharp stability of
the second eigenvalue under the natural class constraint is due to
Brasco--Pratelli~\cite{BrascoPratelli2012}.

The Kohler--Jobin
inequality~\cite{KohlerJobin1978,Brasco2014} bridges the two sides of
the problem: any sharp $\deltaT$-stability translates into a sharp
$\deltaFK$-stability with the same exponent on $\Fr$. Both of our
proofs first close the sharp Saint--Venant
inequality~\eqref{eq:intro-sharp-form} and then transfer through
Kohler--Jobin. We make this bridge explicit in
Sections~\ref{ssec:intro-kj} and~\ref{ssec:intro-final-thms} below.

% =====================================================================
\subsection{The selection-principle / BDV route}
\label{ssec:intro-bdv}
% =====================================================================

The BDV selection
principle~\cite{BDV} was inspired by the prototype of
Cicalese--Leonardi~\cite{CL2012} for the isoperimetric inequality;
related second-variation selection mechanisms appear in
Acerbi--Fusco--Morini~\cite{AFM2013}. The structure of the argument
in~\cite{BDV} is: among open sets $U\subset B_R$ with $|U|=\omega_n$
and prescribed lower bound on $\Fr$, minimise a penalised functional
\[
   J_{\Om_0}(U)\;:=\;E(U)+f_{\hat\eta}(|U|)+
      \sqrt{\eps^2+\tau^2(\BDValpha(U)-\eps)^2}
\]
in which $\BDValpha$ is the BDV smooth asymmetry of
\cite[Def.~4.1]{BDV} and $f_{\hat\eta}$ is a volume penalty. The
selected minimiser $U_*$ inherits a regularity package
(\cite[Lem.~4.6, 4.9, 4.12]{BDV}: density estimates, Lipschitz
non-degeneracy, smooth Bernoulli coefficient) and via the BDV
flatness theorem (\cite[Thm~4.18]{BDV}, in the spirit of
Alt--Caffarelli~\cite{AC1981}) enters the $C^{1,\gamma}$ \emph{graph
regime}: $\partial U_*$ is the graph of a function $g$ over a sphere.
A Kinderlehrer--Nirenberg hodograph step~\cite{KN1977}, followed by
a Lieberman oblique Schauder
bootstrap~\cite{Li86a,Li86b,GT1983,BL1976}, brings $g$ into
$C^{2,\gamma_0}$ with controlled norm. The BDV nearly-spherical theorem
(\cite[Thm~3.3]{BDV}, the modern descendant of the convex-stability
estimate of Fuglede~\cite{Fuglede1989}) then closes the chain:
\[
   E(U_*)-E(B_1)\;\ge\;\tfrac{1}{32 n^2}\,\|g\|_{H^{1/2}(\partial B_1)}^2
\]
on nearly-spherical sets. After unwinding the asymmetry chain
(\cite[Lem.~4.2]{BDV}) and combining with the far-regime
$\Fr^4$-estimate (\cite[Lem.~5.1, 5.2]{BDV}, itself based on the FMP
quantitative isoperimetry of~\cite{FMP2008}), one obtains a
\emph{bounded sharp Saint--Venant} statement
\begin{equation}\label{eq:intro-SV-bounded}
   \deltaT(\Om)\;\ge\;c_{\rm SV}(n,R)\,\Fr(\Om)^2,
   \qquad
   \Om\subset B_R,\quad |\Om|=\omega_n.
\end{equation}
The dependence on the confining radius $R$ is then removed by the
constructive BDV truncation lemma~\cite[Lem.~5.3]{BDV}, applied to a
truncate $\widetilde\Om\subset B_{R_*(n)}$ whose deficit and asymmetry
are controlled by those of $\Om$. The output is a global sharp
Saint--Venant inequality $\deltaT(\Om)\ge c_{\rm
SV}^{\rm glob}(n)\Fr(\Om)^2$, which feeds into Kohler--Jobin.

This is the load-bearing Plan~1 route. The argument is conditional
only in one respect: the BDV nearly-spherical
theorem~\cite[Thm~3.3]{BDV} is invoked as a black-box quantitative
identity on the $H^{1/2}$-norm of $g$, and the regularity bootstrap
that prepares $U_*$ for that theorem requires the full BDV regularity
package. None of these inputs is conditional in the sense of being
unproved in the literature; they are simply non-trivial.

\begin{remark}[Bernoulli spectral alternative]\label{rem:intro-bernoulli}
A conceptually different conditional approach to the bounded sharp
Saint--Venant inequality~\eqref{eq:intro-SV-bounded} would invoke a
spectral / Bernoulli linearisation of the BDV selected minimiser
around the ball, using the eigenvalues of the Bernoulli operator on
$\partial B_1$. This route is sketched in some of the local source
notes (cf.\ \texttt{Plan~1/route-A-summary.md}) but is \emph{not} the
main Plan~1 proof of this manuscript: the load-bearing route is the
BDV selection principle / nearly-spherical / FMP far-regime chain
above. The Bernoulli alternative is mentioned only as a remark and
its conditional status (it requires a spectral gap statement we do
not prove here) is made explicit.
\end{remark}

% =====================================================================
\subsection{Quantitative isoperimetry background}
\label{ssec:intro-quant-iso}
% =====================================================================

The Plan~1 far-regime estimate~\cite[Lem.~5.1, 5.2]{BDV} and several
intermediate steps in Plan~2 rely on the quantitative isoperimetric
toolbox developed in the last twenty years. The cornerstone is the
sharp quantitative isoperimetric inequality of
Fusco--Maggi--Pratelli~\cite{FMP2008}: there is $C(n)>0$ such that for
every set $E\subset\R^n$ of finite perimeter with $|E|=|\Bstar|$,
\begin{equation}\label{eq:intro-FMP}
   \Fr(E)^2\;\le\;C(n)\,
   \frac{P(E)-P(\Bstar)}{P(\Bstar)},
\end{equation}
where $P$ is the De~Giorgi perimeter. The exponent $2$ is sharp. An
alternative mass-transport proof of~\eqref{eq:intro-FMP} is due to
Figalli--Maggi--Pratelli~\cite{FMP2010}; a Sobolev-type strengthening
is~\cite{FMP_BV}; an analogous sharp Sobolev / Wulff stability appears
in~\cite{CFMP2009,NeumayerWulff}; see the surveys of
Fusco~\cite{FuscoSurvey} and Maggi~\cite{Maggi2008} for context.

The strong form of Fusco--Julin~\cite{FJ2014} adds an oriented
oscillation index: for every $E\subset\R^n$ of finite perimeter with
$|E|=|\Bstar|$,
\begin{equation}\label{eq:intro-FJ}
   \Fr(E)^2+\beta(E)^2\;\le\;C(n)\,
   \frac{P(E)-P(\Bstar)}{P(\Bstar)},\qquad
   \beta(E):=\inf_{z\in\R^n}\!\Bigl(\,
      \fint_{\partial^*E}\!
      \Bigl|\nu_E(x)-\tfrac{x-z}{|x-z|}\Bigr|^2\!d\mathcal H^{n-1}
   \Bigr)^{\!1/2}.
\end{equation}
\emph{Both} the asymmetry and the normal-oscillation defect are
controlled by the isoperimetric defect, with a sharp choice of
centre. This is the load-bearing isoperimetric input for Plan~2: it
is~\eqref{eq:intro-FJ}, not~\eqref{eq:intro-FMP}, that supplies the
levelwise oscillation centres along the foliation
$\rho\mapsto E_\rho=\{u_\Om>t(\rho)\}$ with $|E_\rho|=\omega_n\rho^n$.

Two technical inputs that come with the quantitative isoperimetric
literature will be used throughout. First, the Figalli--Maggi
Sobolev--Sard theorem~\cite[App.~A, Thm~A.1]{FigalliMaggi2011}: for
the torsion function $u_\Om\in W^{2,p}_{\rm loc}$ with $-\Delta u=1$,
$\mathcal H^{n-1}(\partial^*E_t\cap\{|\nabla u|=0\})=0$ for
$\mathcal L^1$-a.e.\ level $t$; equivalently, the normal velocity
$\Vrho(x)=-t'(\rho)/|\nabla u(x)|$ is well-defined
$\mathcal H^{n-1}$-a.e.\ on $\partial^*\Erho$ for $\mathcal
L^1$-a.e.\ $\rho$. Second, the finite-perimeter / BV toolbox of
Maggi's book~\cite{Maggi2012} (coarea, divergence theorem,
slicing, structure of $\partial^*E$, rigidity in the isoperimetric
inequality) and Ambrosio--Fusco--Pallara~\cite{AFP2000} (Reshetnyak
lower semicontinuity, BV chain rule, Vol'pert slicing). The
measurable selection statement we use for the levelwise oscillation
centre on good levels is Federer~\cite[\S2.13]{Federer1969}.

% =====================================================================
\subsection{The Kohler--Jobin bridge}
\label{ssec:intro-kj}
% =====================================================================

The Faber--Krahn and Saint--Venant inequalities are not independent.
The Kohler--Jobin inequality~\cite{KohlerJobin1978} states that
\begin{equation}\label{eq:intro-KJ}
   T(\Om)\,\lambda_1(\Om)^{(n+2)/2}
   \;\ge\;
   T(\Bstar)\,\lambda_1(\Bstar)^{(n+2)/2}
\end{equation}
for every open $\Om\subset\R^n$ with $0<|\Om|<\infty$, with equality
if and only if $\Om$ is a ball. The original
proof~\cite{KohlerJobin1978} is by a rearrangement of the eigenfunction
along the level sets of the torsion function. The modern statement
used here is Brasco's~\cite[Thm~3.3]{Brasco2014}, which packages the
rearrangement as a single comparison of the functionals
$T\cdot\lambda_1^{(n+2)/2}$ across all open sets of finite
volume. Inverting the dependence at fixed volume, one obtains the
pointwise transfer
\begin{equation}\label{eq:intro-KJ-pointwise}
   \lambda_1(\Om)\;\ge\;
   \lambda_1(\Bstar)\,
   \Bigl(\tfrac{T(\Bstar)}{T(\Om)}\Bigr)^{\!2/(n+2)},
\end{equation}
and after the elementary inequality $(1-s)^{-p}-1\ge ps$ for
$s\in[0,1]$, $p\ge 0$, the Saint--Venant deficit $\deltaT(\Om)$
becomes a multiple of $\deltaFK(\Om)$ on the small-deficit regime,
with explicit dimensional multiplier $2L_n/((n+2)T_n)$
(cf.~\cite[\S3]{Brasco2014}). On the large-deficit regime
$T(\Om)\le T(\Bstar)/2$, a direct evaluation
of~\eqref{eq:intro-KJ-pointwise} provides a one-sided bound of size
$L_n(2^{2/(n+2)}-1)/4$ on $\deltaFK(\Om)/\Fr(\Om)^2$. Taking the
minimum of the two multipliers yields the explicit
constant~\eqref{eq:intro-cFKdef} for both proofs.

The Kohler--Jobin transfer is therefore the common closing step in
the manuscript: \emph{any} sharp $\deltaT$-stability gives a sharp
$\deltaFK$-stability, and the constants on the two sides are linked
by an explicit dimensional formula. Both Plan~1 and Plan~2 produce
the same input $\deltaT(\Om)\ge c_{\rm SV}^{\rm glob}(n)\Fr(\Om)^2$ to
feed into~\eqref{eq:intro-KJ}, and so the final constants
$c_{\rm FK}(n)$ on both sides
of~\eqref{eq:intro-cFKdef} agree verbatim.

% =====================================================================
\subsection{Nicola--Tilli / rearrangement / level-set viewpoint}
\label{ssec:intro-nt}
% =====================================================================

A conceptually different mechanism underlies Plan~2. The starting
point is the observation, made precise by
Nicola--Tilli~\cite{NicolaTilli2022} in the context of the
short-time Fourier transform Faber--Krahn inequality and developed
quantitatively in the stability companion of
Gómez--Guerra--Ramos--Tilli~\cite{GGRT2024}, that a Faber--Krahn-type
deficit admits an \emph{exact integral identity} over the level sets
of the relevant function. In our setting, the level sets of the
torsion function $u_\Om$ are parametrised by volume,
$\Erho:=\{u_\Om>t(\rho)\}$ with $|\Erho|=\omega_n\rho^n$,
$\rho\in[0,1]$, and the Saint--Venant deficit decomposes as
\begin{equation}\label{eq:intro-deficit-id}
   \deltaT(\Om)
   \;\simeq\;
   \int_{[\rho_*,1]}\!
   \bigl(D_H(t(\rho))+D_I(t(\rho))\bigr)\,w(\rho)\,d\rho,
\end{equation}
where $D_H(t)=\alpha(t)\beta(t)-P(\Erho)^2\ge 0$ is a
Cauchy--Schwarz defect along the level set, $D_I(t)=P(\Erho)^2-c_n
m(t)^{2-2/n}\ge 0$ is the squared-perimeter isoperimetric defect,
$\alpha(t)=\int_{\partial^*E_t}|\nabla u|^{-1}\,d\mathcal H^{n-1}$,
$\beta(t)=\int_{\partial^*E_t}|\nabla u|\,d\mathcal H^{n-1}$, and
$w(\rho)$ is an explicit non-negative weight. The identity is proved
in~\S\ref{ssec:plan2-level-set} below; it builds on the
Pólya--Szeg\H o / Talenti rearrangement
identities~\cite{PolyaSzego1951,Talenti1976,KawohlBook,Bandle1980,BBC1975,BrothersZiemer1988}
and uses the BV chain rule of Bénilan--Brezis--Crandall~\cite{BBC1975}
in the form needed by the level-set foliation.

The identity~\eqref{eq:intro-deficit-id} is fully integral: it
expresses the global deficit as an integral over the foliation of
two pointwise non-negative quantities, $D_H$ controlling the
oscillation of $|\nabla u|$ on the level set, $D_I$ controlling the
isoperimetric oscillation of the level set itself. This is the
opposite philosophy from Plan~1: instead of selecting a single
near-extremal set and analysing it as a graph over a sphere, one
analyses \emph{every} level $\rho$ of the foliation, classifies
levels as good or bad according to their isoperimetric defect, and
transfers the integrated defect bound to the metric quotient
$\X=\mathfrak M/\R^n$ in which sets are identified with their
translates.

% =====================================================================
\subsection{What Plan~1 proves and why it is included}
\label{ssec:intro-plan1}
% =====================================================================

Plan~1 follows the BDV route~\cite{BDV} step-by-step, with the level
of detail demanded by~\S2 of \texttt{MANUSCRIPT\_ASSEMBLY\_AGENTS.md}.
The skeleton is
\[
   \begin{aligned}
      &\text{selection principle (\S3.1, Agent~5)}\\
      &\quad\to\;\text{graph entry $C^{1,\gamma}$ (\S3.2, Agent~6)}\\
      &\quad\to\;\text{Schauder/interpolation $C^{2,\gamma_0}$ (\S3.2, Agent~6)}\\
      &\quad\to\;\text{BDV nearly-spherical closure (\S3.3, Agent~7)}\\
      &\quad\to\;\text{far-regime $\Fr^4$ glue (\S3.3, Agent~7)}\\
      &\quad\to\;\text{bounded sharp Saint--Venant (\S3.3, Agent~7)}\\
      &\quad\to\;\text{bounded reduction (\S3.4, Agent~8)}\\
      &\quad\to\;\text{Kohler--Jobin transfer (\S3.4, Agent~8)}\\
      &\quad\to\;\text{sharp quantitative Faber--Krahn (Theorem~\ref{thm:plan1-FK-upper-intro}).}
   \end{aligned}
\]
The bounded sharp Saint--Venant estimate
\[
   \deltaT(\Om)\;\ge\;c_{\rm SV}(n,R)\,\Fr(\Om)^2,
   \quad \Om\subset B_R,\;|\Om|=\omega_n
\]
is the substance of the BDV argument; bounded reduction is the
constructive truncation of~\cite[Lem.~5.3]{BDV}; Kohler--Jobin is
\cite[Thm~3.3]{Brasco2014}.

The reasons for including Plan~1 in this manuscript are threefold.
First, it is the canonical proof: any honest discussion of the sharp
quantitative Faber--Krahn inequality must reproduce the BDV chain,
because the structure of the rest of the literature
(\cite{BraDeP2017,AKN2023,BrascoPratelli2012,FuscoSurvey,Maggi2008})
is organised around it. Second, the constant
$c_{\rm SV}(n,R)$ produced by Plan~1 is fully explicit at every step
(selection thresholds, regularity radii, nearly-spherical constant
$c_*(n)=1/(32n^2 C_3(n))$, glue constants), and the chain shows
exactly how dimensional and confinement dependences arise. Third,
both proofs converge at the same global Saint--Venant statement and
the same Kohler--Jobin transfer; presenting Plan~1 in full makes the
common closing step explicit and makes the comparison with Plan~2
honest.

% =====================================================================
\subsection{What Plan~2 proves and why it is structurally different}
\label{ssec:intro-plan2}
% =====================================================================

Plan~2 produces \emph{the same} global sharp Saint--Venant inequality
$\deltaT(\Om)\ge c_{\rm SV}^{\rm glob}(n)\Fr(\Om)^2$ along a
structurally different route. The mechanism is a five-step chain in
the metric quotient $\X$:
\[
   \begin{aligned}
      &\text{Level-set deficit identity (\S4.1, Agent~9)}\\
      &\quad\to\;\text{Metric framework / first variation (\S4.2, Agent~10)}\\
      &\quad\to\;\text{Fusco--Julin centres and oriented radial trace (\S4.3, Agent~11)}\\
      &\quad\to\;\text{Integrated action / weighted metric trace (\S4.4, Agent~12)}\\
      &\quad\to\;\text{Boundary-layer transfer + global assembly (\S4.5, Agent~13).}
   \end{aligned}
\]
Each ingredient is necessary, and none is conditional. We summarise
the role of each.

\paragraph{(1) Exact level-set deficit identity.}
Theorem~\ref{thm:plan2-levelset-id} (Agent~9, \S\ref{ssec:plan2-level-set})
records the identity
\eqref{eq:intro-deficit-id} as a finite-perimeter equality, on the
reduced boundary $\partial^*\Erho$, with $\alpha(t)$, $\beta(t)$,
$D_H(t)$, $D_I(t)$, and $w(\rho)$ defined precisely in the notation
register. The proof uses the BV coarea theorem
(\cite[Thm~13.1]{Maggi2012},~\cite[Thm~3.40]{AFP2000}), a weak
divergence-flux identity for the gradient of the torsion function
truncated by a smooth cut-off of the radial vector field $x/|x|$, and
the BV chain rule for monotone absolutely continuous inverses
(\cite[Thm~0.4]{BBC1975}; see also~\cite[Thm~3.99]{AFP2000}). The
identity isolates two non-negative defects at every level: $D_H$ (a
Cauchy--Schwarz defect, controlling oscillation of $|\nabla u|$ on
$\partial^*E_t$) and $D_I$ (a squared-perimeter isoperimetric defect,
controlling oscillation of the level set itself).

\paragraph{(2) Metric framework and first variation.}
\S\ref{ssec:plan2-metric} (Agent~10) embeds the family
$F_\rho=\rho^{-1}\Erho$ in the quotient metric space
$(\X,\dX)=(\{\mathbf 1_A\}/\R^n, d_{L^1}/\!\sim)$ of indicator
functions modulo translation, and proves absolute continuity of
$\rho\mapsto[F_\rho]$ with metric derivative in the
Ambrosio--Gigli--Savar\'e sense~\cite[Def.~1.1.1, Thm~1.1.2]{AGS2008}.
The smooth first variation is computed under a smooth flow (with
normal velocity $\Vrho=-t'(\rho)/|\nabla u|$ ---
well-defined $\mathcal H^{n-1}$-a.e.\ on $\partial^*\Erho$ for a.e.\
$\rho$ by~\cite[Thm~A.1]{FigalliMaggi2011}), then passed to the
finite-perimeter setting through Reshetnyak lower
semicontinuity~\cite[Thm~2.38]{AFP2000}. The output is an upper bound
on the metric derivative $|\dot F_\rho|_\X$ by an explicit Lebesgue
norm of $\Vrho$ on $\partial^*\Erho$.

\paragraph{(3) Fusco--Julin centres and oriented radial trace.}
\S\ref{ssec:plan2-radial-trace} (Agent~11) chooses, on every \emph{good}
isoperimetric-defect level (i.e., where $D_I(t(\rho))$ is below a
threshold), a Borel measurable oscillation centre $z_\rho\in\R^n$
extracted from the Fusco--Julin
inequality~\eqref{eq:intro-FJ}~\cite{FJ2014}, using the
Federer measurable selection
principle~\cite[\S2.13]{Federer1969}. The asymmetry and
normal-oscillation defects of $\Erho$ at centre $z_\rho$ are then
both controlled by $D_I(t(\rho))^{1/2}$ and $D_I(t(\rho))$
respectively. The load-bearing lemma is the \emph{oriented
$L^1$-radial trace estimate}
\begin{equation}\label{eq:intro-trace}
   T_\rho
   \;=\;\!\int_{\partial^*\Erho}\!
      \Bigl|\tfrac{|x-z_\rho|}{\rho}-1\Bigr|\,d\mathcal H^{n-1}
   \;\le\;C\Bigl(
      |\Erho\triangle B_\rho(z_\rho)|
      +\!\int_{\partial^*\Erho}\!
         \Bigl|\nurho-\tfrac{x-z_\rho}{|x-z_\rho|}\Bigr|^{\!2}\!
         d\mathcal H^{n-1}
   \Bigr).
\end{equation}
The proof of~\eqref{eq:intro-trace} is a finite-perimeter divergence
identity applied to the vector field
$X(x)=\bigl|\tfrac{|x-z_\rho|}{\rho}-1\bigr|\tfrac{x-z_\rho}{|x-z_\rho|}$,
truncated near the centre to handle the $L^1_{\rm loc}$ singularity at
$z_\rho$ (Maggi~\cite[Thm~16.3]{Maggi2012}). The Fusco--Julin
inequality controls both terms on the right of~\eqref{eq:intro-trace}
by powers of $D_I$, and the level-set deficit identity converts the
integrated trace into a metric kinetic estimate.

\paragraph{(4) Integrated action and weighted metric trace.}
\S\ref{ssec:plan2-weighted-trace} (Agent~12) decomposes the foliation
into good and bad isoperimetric-defect levels, controls the bad set
by Markov's inequality applied to $D_I$, and reassembles
\[
   \int_{[\rho_*,1]}\!\bigl(\dX([F_\rho],[B_1])^2+|\dot F_\rho|_\X^2\bigr)\,\mu(d\rho)
   \;\le\;C(n,R,\rho_*)\,\deltaT(\Om),
\]
where $\mu(d\rho)=-t'(\rho)\,d\rho$. The abstract weighted metric trace
lemma (an arc-length argument in $(\X,\dX)$ with respect to the
measure $\mu$) then produces a radius $\rhad\in[\rho_*,1]$ such that
$\dX([F_{\rhad}],[B_1])^2\le C_*\,\deltaT(\Om)$ and $\rhad$ is close
to the heuristic value $\rhodelta=1-\kappa\sqrt{\deltaT(\Om)}$. The
Talenti bad-set Lebesgue estimate (\cite{Talenti1976}, used in
combination with the volume-radius parametrisation,
\cite[Thm~1.3.2]{Kesavan2006}) is what converts weighted kinetic
energy on the bad-$-t'$ set into Lebesgue kinetic energy on
$[\rho_*,1]$ without losing the $\deltaT$ rate.

\paragraph{(5) Boundary-layer transfer.}
\S\ref{ssec:plan2-boundary-transfer} (Agent~13) returns from $F_{\rhad}$ to
$\Om$. The removed boundary layer $\Om\setminus E_{\rhad}$ has
measure $O(\sqrt{\deltaT(\Om)})$, and the transfer of asymmetry from
$E_{\rhad}$ to $\Om$ squares this $\sqrt{\deltaT}$ without producing
a rate loss. The result is the bounded sharp Saint--Venant
inequality~\eqref{eq:intro-SV-bounded} in the Plan~2 form. Bounded
reduction (\cite[Lem.~5.3]{BDV}, imported through
Plan~1's \texttt{BoundedReduction}) removes the $R$-dependence, and
Kohler--Jobin~\cite{Brasco2014} closes the chain.

\begin{remark}[Auxiliary centroid / space-time route]\label{rem:intro-pi-route}
A separate centroid bound and a related space-time $\Pi$-identity
(\S\ref{ssec:plan2-aux-centroid}, Agent~14) compute, for the bulk
centroid $\zbar$, the signed radial moment
$\PiE$ and a Fubini-density identity relating
$\PiE$ across $\rho$. These quantities are auxiliary diagnostics:
the repaired Plan~2 chain above uses only the \emph{integrated}
forms of $D_H$ and $D_I$ through the
identity~\eqref{eq:intro-deficit-id} and the oriented radial trace
of~\eqref{eq:intro-trace}, not pointwise $\PiE$-control. The
$\Pi$-route is presented in the manuscript for the record and would
become load-bearing only if a separate integrated $\Pi$-control
estimate were proved (which we do not). Up to that hypothetical
upgrade, the centroid/$\Pi$-route is not used in any of the load-bearing
steps of Plan~2.
\end{remark}

The reasons for including Plan~2 are: (i) it is structurally
different from Plan~1, in that it does not select a single near-extremal
set and does not use a graph parametrisation; (ii) it exposes the
deficit as an integral over the level-set foliation, so that the role of
isoperimetric stability~\eqref{eq:intro-FJ} is direct rather than
indirect; (iii) it uses only finite-perimeter regularity throughout,
without a smoothness or graph assumption on $\partial\Om$; (iv) the
oriented $L^1$-radial trace lemma~\eqref{eq:intro-trace} is, to our
knowledge, the first place in the Faber--Krahn / Saint--Venant
quantitative literature where Fusco--Julin oscillation is used in this
form to transfer levelwise control to the metric quotient $\X$.

% =====================================================================
\subsection{Precise statement of the two final theorems}
\label{ssec:intro-final-thms}
% =====================================================================

Both proofs in this manuscript close the same theorem, in the same
normalisation, with the same explicit dimensional constant
\eqref{eq:intro-cFKdef}. We state both forms here for the record;
Plan~1 produces the proof in \S\ref{sec:plan1} and Plan~2 in
\S\ref{sec:plan2}.

\begin{theorem}[Sharp Faber--Krahn via Kohler--Jobin; Plan~1, upper form]
\label{thm:plan1-FK-upper-intro}
There is a dimensional constant $C_{\rm FK}(n)>0$ such that, for every
open set $\Om\subset\R^n$ with $0<|\Om|<\infty$,
\begin{equation}\label{eq:intro-plan1-FK}
   \Fr(\Om)^2\;\le\;C_{\rm FK}(n)\,\deltaFK(\Om),
\end{equation}
where
\[
   C_{\rm FK}(n)
   \;:=\;\frac{1}{c_{\rm FK}(n)}
   \;=\;\max\!\Bigl\{
      \frac{(n+2)\,T_n}{4 L_n\,c_{\rm SV}^{\rm glob}(n)},\;
      \frac{4}{L_n\,(2^{2/(n+2)}-1)}
   \Bigr\},
\]
with $L_n=\lambda_1(B_1)$, $T_n=T(B_1)=\omega_n/(n(n+2))$, and
$c_{\rm SV}^{\rm glob}(n)$ the global sharp Saint--Venant
constant defined in~\eqref{eq:csv-glob} of \S\ref{sec:plan1}. The
deficit $\deltaFK(\Om)$ is as in~\eqref{eq:intro-deltaFK} and
$\Fr(\Om)$ is the Fraenkel asymmetry~\eqref{eq:intro-Fr}.
Equality cases: $\deltaFK(\Om)=0$ iff $\Om$ is a ball (up to
translation and a Lebesgue negligible set); in that case
$\Fr(\Om)=0$.
\end{theorem}

\begin{theorem}[Sharp quantitative Faber--Krahn; Plan~2 final form]
\label{thm:plan2-FK-intro}
Let $n\ge 2$. There exists a constant $c_{\rm FK}(n)>0$, depending
only on $n$, given by the explicit formula~\eqref{eq:intro-cFKdef}
above, such that for every open $\Om\subset\R^n$ with $0<|\Om|<\infty$,
\begin{equation}\label{eq:intro-plan2-FK}
   \deltaFK(\Om)
   \;=\;
   \Bigl(\tfrac{|\Om|}{\omega_n}\Bigr)^{\!2/n}
   \bigl(\lambda_1(\Om)-\lambda_1(\Bstar)\bigr)
   \;\ge\;c_{\rm FK}(n)\,\Fr(\Om)^2.
\end{equation}
The exponent $2$ is sharp.
\end{theorem}

The two statements are identical up to inversion of the constant
($C_{\rm FK}(n)=1/c_{\rm FK}(n)$). The proofs end at the same place
because they share the bounded reduction~\cite[Lem.~5.3]{BDV} and the
Kohler--Jobin transfer~\cite[Thm~3.3]{Brasco2014}. The constants are
identical because both routes feed the same global Saint--Venant
constant $c_{\rm SV}^{\rm glob}(n)$ into the Kohler--Jobin transfer,
and the dimensional factors $L_n,T_n,\omega_n$ are universal.

% =====================================================================
\subsection{Related work and contributions}
\label{ssec:intro-related}
% =====================================================================

The sharp inequality~\eqref{eq:intro-main} was first
proved by Brasco--De~Philippis--Velichkov~\cite{BDV} in 2015. A
companion proof, structurally distinct in the use of the
Alt--Caffarelli--Friedman monotonicity formula, was obtained by
Allen--Kriventsov--Neumayer~\cite{AKN2023}. The sharp stability of the
second Dirichlet eigenvalue is the result of
Brasco--Pratelli~\cite{BrascoPratelli2012}. Pre-BDV non-sharp results
include Hansen--Nadirashvili~\cite{HansenNadirashvili1994},
Melas~\cite{Melas1992}, Bhattacharya~\cite{Bhattacharya2001}, and the
$\Fr^4$-bound of~\cite{FMPLaplace}. A modern survey of the quantitative
Faber--Krahn / Saint--Venant literature is in
Brasco--De~Philippis~\cite{BraDeP2017}; for the broader quantitative
isoperimetric viewpoint see Fusco~\cite{FuscoSurvey} and
Maggi~\cite{Maggi2008}.

On the isoperimetric side, the sharp quantitative
inequality~\cite{FMP2008} admits the mass-transport
proof of~\cite{FMP2010}, the BV-Sobolev strengthening of~\cite{FMP_BV},
the sharp Sobolev companion of~\cite{CFMP2009}, and the Wulff
strengthening of~\cite{NeumayerWulff}. The strong form of
Fusco--Julin~\cite{FJ2014} is what makes Plan~2 work: the oscillation
index controls a quadratic centre-dependent average of $|\nu-e_r|^2$
with the same exponent as the asymmetry. The selection-principle
mechanism that underlies Plan~1 originates in
Cicalese--Leonardi~\cite{CL2012}; the closely related second-variation
selection of~\cite{AFM2013} is in the spirit of the BDV penalised
functional but is not used directly here.

The level-set / convexity viewpoint behind Plan~2 is inspired by
Nicola--Tilli~\cite{NicolaTilli2022} and its quantitative companion
Gómez--Guerra--Ramos--Tilli~\cite{GGRT2024} in the short-time Fourier
transform setting. Our use of the metric quotient $\X=\mathfrak
M/\R^n$ and the Ambrosio--Gigli--Savar\'e metric
derivative~\cite{AGS2008} is, to our knowledge, new in the
Saint--Venant context; the oriented $L^1$-radial trace
identity~\eqref{eq:intro-trace} and its proof by truncation of the
vector field $x\mapsto||x-z|/\rho-1|(x-z)/|x-z|$ on the reduced
boundary, using~\cite[Thm~16.3]{Maggi2012}, are the technical core
that allows Fusco--Julin oscillation to be converted into integrated
metric data along the foliation.

Finally, the Kohler--Jobin
inequality~\cite{KohlerJobin1978,Brasco2014} is the
device that bridges Saint--Venant stability and Faber--Krahn
stability; it is what permits both of our proofs to share the
\emph{same} closing step and the \emph{same} final constant. We
present this bridge in full in \S\ref{ssec:plan1-final} and reuse it
verbatim in \S\ref{ssec:plan2-final}.

\paragraph{Contributions of the present paper.}
We summarise the contributions of the manuscript as follows:

\begin{enumerate}
\item A fully detailed exposition of the BDV
selection-principle / graph-entry / Schauder / nearly-spherical /
truncation chain~\cite{BDV,Brasco2014}, with explicit dimensional
constants at every step (Plan~1, \S\ref{sec:plan1}). The level of
detail follows the building blocks
\texttt{Plan~1/quantitative-selection-principle.tex},
\texttt{Plan~1/graph-entry-closure.tex},
\texttt{Plan~1/route-A-summary.tex},
\texttt{Plan~1/faber-krahn-transfer.tex}.
\item A structurally different proof of the same sharp inequality
based on an exact level-set deficit identity, a metric first-variation
in the quotient space $(\X,\dX)$ of sets modulo translation, an
oriented $L^1$-radial trace lemma anchored at Fusco--Julin oscillation
centres, a weighted metric trace, and a boundary-layer transfer. This
is Plan~2, \S\ref{sec:plan2}, and is new in the Saint--Venant /
Faber--Krahn context. The five major steps and their building blocks
\texttt{Plan~2/WIP/WIP\_LevelSetIdentity.tex},
\texttt{Plan~2/WIP/WIP\_MetricFramework.tex},
\texttt{Plan~2/WIP/WIP\_WeightedMetricTrace.tex},
\texttt{Plan~2/WIP/WIP\_BoundaryLayerTransfer.tex}, and
\texttt{Plan~2/WIP/WIP\_GlobalAssembly.tex} are integrated into the
manuscript with detail substantially greater than in the source
blocks.
\item A side-by-side bookkeeping of the constants. The bounded
Saint--Venant constants $c_{\rm SV}(n,R)$ from each route are
different intermediate quantities; the global constant
$c_{\rm SV}^{\rm glob}(n)$ that enters Kohler--Jobin is, in both
routes, given by the same expression
(\eqref{eq:intro-cFKdef}). The final constant $c_{\rm FK}(n)$
on~\eqref{eq:intro-main} is therefore identical for the two routes,
as recorded by Remark~\ref{rem:plan2-vs-plan1-normalisation} of
\S\ref{sec:plan2}.
\item A complete reference audit
(Manuscript/04\_references\_audit.md) listing every external citation
with exact bibliographic data, the theorem/proposition/page actually
used, and the manuscript location at which it enters, together with a
notation register (Manuscript/00\_notation\_register.md) fixing the
canonical symbols. No vague phrases such as `well-known' or
`standard' appear in a load-bearing step without a precise reference.
\item An internal audit pass on each of Plan~1 and Plan~2
(Manuscript/audit\_plan1.md, Manuscript/audit\_plan2.md), recording
that no conditional Bernoulli spectral input is used in the Plan~1
proof and that no centroid / space-time $\Pi$-control is used in the
load-bearing Plan~2 chain.
\end{enumerate}

\paragraph{Organisation.}
\S\ref{sec:prelims} (Preliminaries) collects the finite-perimeter,
torsion, rearrangement, and Fusco--Julin facts used in both routes.
\S\ref{sec:plan1} (Plan~1) presents the BDV selection-principle proof
from selection through to the final theorem. \S\ref{sec:plan2}
(Plan~2) presents the level-set / metric route through to the same
final theorem. The Kohler--Jobin transfer and the bounded reduction,
which are the closing steps of both routes, are written out fully in
\S\ref{ssec:plan1-final} and re-used in \S\ref{ssec:plan2-final}.
\S\ref{sec:references} (References) is the bibliography.

\bigskip

\noindent\emph{Conventions used throughout the introduction:} the
ambient dimension is $n\ge 2$, balls $B_r(z)\subset\R^n$ are open
Euclidean balls, $\Bstar$ is the open ball centred at the origin with
$|\Bstar|=|\Om|$, $\omega_n=|B_1|$, $L_n=\lambda_1(B_1)$, $T_n=T(B_1)$,
and all constants $C, c, c_{\rm SV}, c_{\rm FK}$ are positive and
depend on at most $(n,R,\rho_*)$ at every step where they appear;
dependences are declared on first use. The notation register is
\texttt{Manuscript/00\_notation\_register.md}.

% =====================================================================
% End of Manuscript/05_intro_story_draft.tex
% =====================================================================


% =====================================================================
% Preliminaries (Agents 2, 3, 4)
% =====================================================================
% =====================================================================
% Manuscript/01_preliminaries_draft.tex
%
% Concatenated preliminaries section (Agents 2, 3, 4) produced by
% Agent 19 (Final Assembly). The three subfiles are \input here;
% the original \section commands inside the subfiles have been left
% as in the source (Agents 17 unified the notation). Inside main.tex,
% this file is itself \input'd between the introduction and Plan 1.
% =====================================================================

\section{Preliminaries}
\label{sec:prelims}
\label{sec:notation}%

This section collects the analytic and geometric measure theory
prerequisites used by both proofs. \S\ref{sec:prelim-gmt} covers
finite-perimeter sets, the BV coarea theorem, divergence on reduced
boundaries, truncation of singular vector fields, Reshetnyak lower
semicontinuity, and the measurable selection of levelwise centres.
\S\ref{sec:prelim-torsion} collects the torsion potential, the
Saint--Venant and Faber--Krahn deficits, the Fraenkel asymmetry, the
Talenti pointwise comparison, the volume-radius foliation
$\Erho=\{u>t(\rho)\}$, the profile-gap and bad-set estimates used by
Plan~2, and the Kohler--Jobin inequality used in the final transfer.
\S\ref{sec:prelim-quant-iso} collects the Fusco--Maggi--Pratelli and
Fusco--Julin quantitative isoperimetric inputs and the scaled form
on $|E|=\omega_n\rho^n$.

% UNIFIED 2026-05-13
% =====================================================================
% Preliminaries -- Geometric Measure Theory
% Agent 2 deliverable. To be merged by Agent 17 into
%     Manuscript/01_preliminaries_draft.tex.
%
% Notation: follows Manuscript/00_notation_register.md verbatim.
% Macros assumed loaded (per register, §15): \R, \Fr, \Om, \eps,
%     \rstar, \rhad, \rhodelta, \zbar, \X, \dX, \Bstar, \Erho, \Frho,
%     \Vrho, \nurho, \deltaT, \deltaFK.
%
% SOURCE MAP for this section (each item is one of:
%   PROVED HERE / CITED FROM <ref> / ADAPTED FROM <local file>):
%
%   §A.1  Lebesgue/Hausdorff measure, perimeter, reduced boundary
%         and normal $\nurho$
%             CITED FROM Maggi, "Sets of Finite Perimeter and
%             Geometric Variational Problems" (Cambridge 2012),
%             Def. 12.1, Def. 15.1, Thm. 15.9.
%             CITED FROM Ambrosio--Fusco--Pallara (AFP),
%             "Functions of Bounded Variation and Free Discontinuity
%             Problems" (Oxford 2000), Def. 3.1, Thm. 3.59.
%
%   §A.2  Lower semicontinuity of perimeter under $L^1_{\rm loc}$ conv.
%             CITED FROM Maggi 2012, Prop. 12.15.
%             CITED FROM AFP 2000, Thm. 3.7.
%
%   §A.3  BV coarea formula (statement + integrated form)
%             CITED FROM Maggi 2012, Thm. 13.1 (BV coarea),
%             and Thm. 18.1 (Sobolev coarea).
%             CITED FROM AFP 2000, Thm. 3.40.
%
%   §A.4  Sobolev coarea applied to the torsion function $u_\Om$ and
%         Sobolev--Sard a.e. regular-level statement giving
%         $\Vrho=-t'(\rho)/|\nabla u|$ a.e.
%             CITED FROM Figalli--Maggi, "On the shape of liquid drops
%             and crystals in the small mass regime", Arch. Ration.
%             Mech. Anal. 201 (2011), App. A, Thm. A.1.
%             ADAPTED FROM Plan 2/WIP/WIP_LevelSetIdentity.tex
%             §2 (coarea applicability) and §4.
%
%   §A.5  Divergence theorem on sets of finite perimeter
%         (Gauss--Green for $C^1_c$ fields and for $C^1$ fields
%         bounded on $E$ together with $\Div X$ in $L^1(E)$).
%             CITED FROM Maggi 2012, Thm. 15.5 (Gauss--Green
%             for sets of finite perimeter) and Thm. 16.3 (its
%             $C^1$/$L^1$ form via approximation).
%
%   §A.6  Lipschitz truncation flux identity for the torsion function
%         (replacing the boundary trace of $\nabla u$ on
%         $\partial^*\Erho$ by the BV coarea+Lipschitz cutoff).
%             ADAPTED FROM Plan 2/WIP/WIP_LevelSetIdentity.tex,
%             Lem. "Flux = volume" (lem:flux) and Rem. rem:divergence.
%
%   §A.7  Truncation lemma for singular vector fields of the type
%         $X(x)=w(|x-z|)(x-z)/|x-z|$ with $w$ bounded.
%             PROVED HERE.
%             Used by Agent 11 in the oriented radial trace lemma
%             (notation register §9: $T_\rho$). The proof is the
%             rigorous form of the truncation argument written
%             informally in Plan 2/WIP/WIP_WeightedMetricTrace.tex
%             (proof of lem:radial-slicing, lines 182--195) and is
%             load-bearing.
%
%   §A.8  Reshetnyak lower semicontinuity for linear-growth
%         integrands on sets of finite perimeter; passage of the
%         finite-perimeter first variation from smooth flows.
%             CITED FROM AFP 2000, Thm. 2.38 (Reshetnyak lower
%             semicontinuity).
%             ADAPTED FROM Plan 2/gmt-inputs-for-metric-closure.md §5
%             (the four-step approximation procedure).
%
%   §A.9  Measurable selection of levelwise centres
%             PROVED HERE (via Aumann/Kuratowski--Ryll-Nardzewski
%             selection).
%             CITED FROM Castaing--Valadier, "Convex Analysis and
%             Measurable Multifunctions" (Springer LNM 580, 1977),
%             Thm. III.6 and Thm. III.8;
%             alternatively Federer, "Geometric Measure Theory"
%             (Springer 1969), §2.13.
%             ADAPTED FROM Plan 2/WIP/WIP_WeightedMetricTrace.tex
%             (proof of lem:FJ-package, last paragraph).
%
% =====================================================================

\subsection{Preliminaries: geometric measure theory}\label{sec:prelim-gmt}

This section fixes the finite-perimeter, slicing, and divergence-theorem
tools used throughout the manuscript.  All notation matches
\S\ref{sec:notation} (see also \texttt{00\_notation\_register.md}).
The ambient dimension is denoted $n\ge 2$.  We write $|A|=\mathcal L^n(A)$
for the Lebesgue measure of a Borel set $A\subset\R^n$,
$\mathcal H^k$ for the $k$-dimensional Hausdorff measure, and we
identify two sets that differ by an $\mathcal L^n$-null set.

\subsubsection{Finite-perimeter sets and reduced boundary}
\label{ssec:prelim-fp}

\begin{definition}[Finite perimeter and reduced boundary]
\label{def:perimeter}
A Lebesgue-measurable set $E\subset\R^n$ has \emph{finite perimeter}
in an open set $U\subset\R^n$ if the indicator $\mathbf 1_E$ belongs
to $BV(U)$, i.e.\ if
\[
   P(E;U):=\sup\Bigl\{
      \int_E \Div\varphi\,dx :
      \varphi\in C^1_c(U;\R^n),\ \|\varphi\|_\infty\le1
   \Bigr\}<\infty.
\]
When $U=\R^n$ we write $P(E):=P(E;\R^n)$.  The
\emph{De Giorgi reduced boundary} $\partial^*E$ is the set of points
$x\in\R^n$ at which the limit
\[
   \nu_E(x)
   :=\lim_{r\downarrow0}
      \frac{D\mathbf 1_E(B_r(x))}{|D\mathbf 1_E|(B_r(x))}
\]
exists in the unit sphere $\mathbb S^{n-1}$; $\nu_E(x)$ is the
\emph{measure-theoretic outer unit normal} of $E$ at $x$.
By De Giorgi's structure theorem,
\[
   P(E;B)=\mathcal H^{n-1}(\partial^*E\cap B)
   \qquad\text{for every Borel } B\subset\R^n,
\]
and $\partial^*E$ is countably $(n-1)$-rectifiable.
\end{definition}

For the canonical reference see
\cite[Def.~12.1, Thm.~15.9]{Maggi2012}; equivalently
\cite[Def.~3.1, Thm.~3.59]{AFP2000}.  Throughout the manuscript
$P(E)$ always denotes the De Giorgi perimeter and $\nu_E$ (or $\nurho$
when $E=\Erho$) the measure-theoretic outer normal, consistently with
\S 7 of the notation register.

\begin{proposition}[$L^1_{\rm loc}$ lower semicontinuity]
\label{prop:lsc-perimeter}
Let $(E_j)_{j\ge1}$ be measurable subsets of $\R^n$ with
$\mathbf 1_{E_j}\to\mathbf 1_E$ in $L^1_{\rm loc}(\R^n)$ as
$j\to\infty$.  Then for every open $U\subset\R^n$,
\[
   P(E;U)\le\liminf_{j\to\infty}P(E_j;U).
\]
In particular, if $\sup_j P(E_j)<\infty$ and $|E_j|$ is uniformly
bounded, then $(E_j)$ has a subsequence converging in $L^1_{\rm loc}$
to a finite-perimeter set $E$ with $P(E)\le\liminf P(E_j)$.
\end{proposition}

\begin{proof}[Reference]
Lower semicontinuity is
\cite[Prop.~12.15]{Maggi2012} (equivalently
\cite[Thm.~3.7]{AFP2000}); the compactness in $L^1_{\rm loc}$ is
\cite[Thm.~12.26]{Maggi2012}.
\end{proof}

\subsubsection{Coarea, BV coarea, and Sobolev--Sard for the torsion
            function}
\label{ssec:prelim-coarea}

\begin{theorem}[BV coarea formula]\label{thm:bv-coarea}
Let $v\in BV(U)$ on an open set $U\subset\R^n$.  Then $\{v>t\}$ has
finite perimeter in $U$ for $\mathcal L^1$-a.e.\ $t\in\R$, and for
every Borel function $\varphi:U\to[0,\infty]$
\begin{equation}\label{eq:bv-coarea}
   \int_U \varphi\,d|Dv|
   =\int_\R\Bigl(\int_{\partial^*\{v>t\}\cap U}\varphi\,d\mathcal H^{n-1}
      \Bigr)\,dt.
\end{equation}
If $v\in W^{1,1}(U)$ (or $v\in H^1_0(U)$ extended by zero outside
$U$), then $d|Dv|=|\nabla v|\,d\mathcal L^n$ and~\eqref{eq:bv-coarea}
reduces to the Federer--Fleming coarea formula
\begin{equation}\label{eq:fed-fleming}
   \int_U\varphi\,|\nabla v|\,dx
   =\int_\R\Bigl(\int_{\partial^*\{v>t\}\cap U}\varphi\,d\mathcal H^{n-1}
      \Bigr)\,dt.
\end{equation}
\end{theorem}

\begin{proof}[Reference]
This is \cite[Thm.~13.1 and Thm.~18.1]{Maggi2012}; equivalently
\cite[Thm.~3.40 and (3.40)(b)]{AFP2000}.  The vanishing of the
contribution of the critical set $\{\nabla v=0\}$ to the
right-hand side of~\eqref{eq:bv-coarea} is \cite[Thm.~3.40(b)]{AFP2000}.
\end{proof}

\begin{corollary}[Coarea for the torsion function]
\label{cor:coarea-torsion}
Let $\Om\subset\R^n$ be open with $0<|\Om|<\infty$ and let $u=u_\Om$ be
the torsion function (\S 3 of the notation register).  Then $u\in
H^1_0(\Om)\cap L^\infty(\Om)$, and after extending $u$ by zero on
$\R^n\setminus\Om$ we have $u\in BV_{\rm loc}(\R^n)\cap L^\infty(\R^n)$.
For every Borel $\varphi:\R^n\to[0,\infty]$,
\begin{equation}\label{eq:coarea-u}
   \int_{\R^n}\varphi\,|\nabla u|\,dx
   =\int_0^{\|u\|_\infty}
      \Bigl(\int_{\partial^*\{u>t\}}\varphi\,d\mathcal H^{n-1}\Bigr)dt.
\end{equation}
In particular the super-level set $E_t=\{u>t\}$ has finite perimeter
for $\mathcal L^1$-a.e.\ $t\in(0,\|u\|_\infty)$, and the function
$m(t):=|E_t|$ is non-increasing, absolutely continuous on
$(0,\|u\|_\infty)$ in the strict sense made precise in
\S\ref{ssec:level-set-prelim} below.
\end{corollary}

\begin{theorem}[Sobolev--Sard for the torsion function;
                Figalli--Maggi 2011, App.~A, Thm.~A.1]
\label{thm:sobolev-sard}
Let $u=u_\Om$ be as in Corollary~\ref{cor:coarea-torsion}.  There is a
Borel set $\mathcal R\subset(0,\|u\|_\infty)$ of full $\mathcal L^1$
measure such that for every $t\in\mathcal R$:
\begin{enumerate}
   \item[(i)] $E_t=\{u>t\}$ has finite perimeter and
       $m'(t)=-\int_{\partial^*E_t}|\nabla u|^{-1}\,d\mathcal H^{n-1}
        \in(-\infty,0)$;
   \item[(ii)] $|\nabla u|>0$ for $\mathcal H^{n-1}$-a.e.
       $x\in\partial^*E_t$ and the De Giorgi outer normal is
       $\nu_{E_t}(x)=-\nabla u(x)/|\nabla u(x)|$ at $\mathcal H^{n-1}$
       -a.e.\ such $x$;
   \item[(iii)] (Volume-radius parametrisation,
       cf.\ register \S 7.)  The function
       $\rho\mapsto t(\rho)$, characterised by
       $|\{u>t(\rho)\}|=\omega_n\rho^n$, is absolutely continuous on
       $[\rstar,1]$ with $t'(\rho)\le0$, and the normal velocity
       \[
          V_\rho(x):=\frac{-t'(\rho)}{|\nabla u(x)|}
       \]
       is well-defined and positive for $\mathcal H^{n-1}$-a.e.\
       $x\in\partial^*\Erho$, for $\mathcal L^1$-a.e.\
       $\rho\in[\rstar,1]$.
\end{enumerate}
\end{theorem}

\begin{proof}[Reference]
(i)--(ii) are \cite[App.~A, Thm.~A.1]{FigalliMaggi2011}; see also
\cite[Lem.~17.1]{Maggi2012}.  (iii) follows by composition: the
inverse-distribution function $\rho\mapsto t(\rho)$ is the strict-AC
inverse of the strict-AC function $t\mapsto (m(t)/\omega_n)^{1/n}$
(use \cite[Thm.~3.99]{AFP2000} on monotone AC functions), and the
formula for $V_\rho$ then follows from the chain rule and
$m'(t)=-\int|\nabla u|^{-1}$.  The level-set identity draft
\texttt{WIP\_LevelSetIdentity.tex} \S 4 records the same statement.
\end{proof}

\subsubsection{Divergence theorem on sets of finite perimeter}
\label{ssec:div-thm}

\begin{theorem}[Gauss--Green; Maggi Thm.~15.5]
\label{thm:gauss-green}
Let $E\subset\R^n$ have finite perimeter and let
$X\in C^1_c(\R^n;\R^n)$.  Then
\begin{equation}\label{eq:gauss-green-Cc}
   \int_E \Div X\,dx
   =\int_{\partial^*E}X\cdot\nu_E\,d\mathcal H^{n-1}.
\end{equation}
More generally, if $E$ has finite perimeter and finite Lebesgue
measure and $X\in C^1(\R^n;\R^n)\cap L^\infty(E;\R^n)$ satisfies
$\Div X\in L^1(E)$, then~\eqref{eq:gauss-green-Cc} continues to hold
(both sides being finite), as soon as the support of $X$ on
$\partial^*E$ has finite $\mathcal H^{n-1}$ measure (e.g.\ when $E$
is bounded).
\end{theorem}

\begin{proof}[Reference]
The $C^1_c$ statement is \cite[Thm.~15.5]{Maggi2012} or
\cite[Thm.~3.36]{AFP2000}.  The extension to $C^1$ fields bounded on
$E$ with $\Div X\in L^1(E)$ follows by multiplying by a smooth
compactly supported cutoff $\chi_R$, $\chi_R=1$ on $B_R$, sending
$R\to\infty$, and using dominated convergence on both sides; this is
\cite[Thm.~16.3]{Maggi2012} or the form used in
\cite[Ch.~17]{Maggi2012}.
\end{proof}

We will need a slight strengthening allowing $X$ to have a single
locally integrable singularity inside $E$.  The following lemma
is load-bearing for the oriented radial trace estimate of
Agent~11 (\S\ref{thm:radial-trace} below cross-references it).

\begin{lemma}[Truncation lemma for fields with locally integrable
              singularity]\label{lem:truncation}
Let $E\subset\R^n$ have finite perimeter and finite Lebesgue measure,
let $z\in\R^n$, and let $X:\R^n\setminus\{z\}\to\R^n$ be of class
$C^1$ with the following properties:
\begin{itemize}
   \item[(a)] $X\in L^\infty(\R^n;\R^n)$, i.e.\ $|X(x)|\le M_0$ for
       $x\ne z$, for some constant $M_0\ge0$;
   \item[(b)] $\Div X\in L^1_{\rm loc}(\R^n\setminus\{z\})$, and
       there is a constant $K_0\ge0$ such that
       \[
          |\Div X(x)|\le \frac{K_0}{|x-z|}
          \qquad\text{for all }x\ne z;
       \]
   \item[(c)] $E$ is bounded, $E\subset B_R$ for some $R\ge1$.
\end{itemize}
Then $\Div X\in L^1(E)$ (since $n\ge2$) and
\begin{equation}\label{eq:trunc-flux}
   \int_E\Div X\,dx
   =\int_{\partial^*E}X\cdot\nu_E\,d\mathcal H^{n-1},
\end{equation}
where the right-hand side is absolutely integrable.

The prototypical $X$ to which this lemma applies is
\[
   X(x)
   =w(|x-z|)\,\frac{x-z}{|x-z|},
\]
where $w:[0,\infty)\to\R$ is bounded with $w(0)=0$ in the natural
sense (so that $|X|\le\|w\|_\infty$).  In dimension $n\ge2$ such $X$
has $\Div X(x)=w'(|x-z|)+(n-1)w(|x-z|)/|x-z|$ for $x\ne z$, whose
absolute value is $O(|x-z|^{-1})$ if $w$ is bounded with bounded
$w'$.
\end{lemma}

\begin{proof}
We may translate so that $z=0$.  The function $x\mapsto|x|^{-1}$ is
locally integrable in $\R^n$ for $n\ge2$, so (b) gives
$\Div X\in L^1(B_{2R})\supset L^1(E)$.  This proves the first
assertion.

Now choose a fixed cutoff $\eta\in C^\infty(\R)$ with $\eta\equiv0$
on $(-\infty,1]$, $\eta\equiv1$ on $[2,\infty)$ and $0\le\eta\le1$,
and for $\eps\in(0,1)$ set
\[
   \eta_\eps(x):=\eta\bigl(\eps^{-1}|x|\bigr),
   \qquad
   X_\eps(x):=\eta_\eps(x)\,X(x).
\]
Then $X_\eps\equiv 0$ on $B_\eps$, $X_\eps\equiv X$ on $\R^n\setminus
B_{2\eps}$, and $X_\eps\in C^1(\R^n;\R^n)$ (the singularity at $0$
has been removed).  Also $X_\eps$ is bounded by $M_0$ on $\R^n$ and
has compact support if we further multiply by a smooth cutoff
$\chi_R\in C^\infty_c(\R^n)$ with $\chi_R\equiv 1$ on $B_R$,
$\chi_R\equiv0$ on $\R^n\setminus B_{2R}$, $|\nabla\chi_R|\le2/R$;
since $E\subset B_R$, this further cutoff does not affect any of the
integrals over $E$ or $\partial^*E$.

Apply Theorem~\ref{thm:gauss-green} (the $C^1_c$ version) to
$\chi_R X_\eps$:
\begin{equation}\label{eq:trunc-step}
   \int_E \Div(\chi_R X_\eps)\,dx
   =\int_{\partial^*E}(\chi_R X_\eps)\cdot\nu_E\,d\mathcal H^{n-1}.
\end{equation}
Since $\chi_R\equiv1$ on a neighbourhood of $E\subset B_R$, the
left-hand side equals $\int_E\Div X_\eps\,dx$ and the right-hand side
equals $\int_{\partial^*E}X_\eps\cdot\nu_E\,d\mathcal H^{n-1}$.

We now let $\eps\downarrow0$ and verify the two convergences:

\emph{Right-hand side.}  $X_\eps\to X$ pointwise on
$\R^n\setminus\{0\}$, hence $\mathcal H^{n-1}$-a.e.\ on $\partial^*E$
(the single point $0$ has $\mathcal H^{n-1}(\{0\})=0$ since
$n\ge2$).  By (a), $|X_\eps\cdot\nu_E|\le M_0$ pointwise; since
$\mathcal H^{n-1}(\partial^*E)=P(E)<\infty$, dominated convergence
gives
\[
   \int_{\partial^*E}X_\eps\cdot\nu_E\,d\mathcal H^{n-1}
   \longrightarrow
   \int_{\partial^*E}X\cdot\nu_E\,d\mathcal H^{n-1},
\]
and the right-hand integral is absolutely convergent.

\emph{Left-hand side.}  A direct computation gives
\[
   \Div X_\eps(x)
   =\eta_\eps(x)\,\Div X(x)
   +\eps^{-1}\eta'\bigl(\eps^{-1}|x|\bigr)\,
       \frac{x}{|x|}\cdot X(x).
\]
The first summand converges pointwise on $\R^n\setminus\{0\}$ to
$\Div X(x)$ as $\eps\downarrow0$.  It is bounded in absolute value
by $|\Div X(x)|\le K_0/|x|$, which is integrable on $E\subset B_R$.
Dominated convergence gives
\[
   \int_E\eta_\eps\,\Div X\,dx
   \longrightarrow\int_E\Div X\,dx.
\]
For the second summand,
\[
   \Bigl|\eps^{-1}\eta'(\eps^{-1}|x|)\,
         \tfrac{x}{|x|}\cdot X(x)\Bigr|
   \le \frac{\|\eta'\|_\infty M_0}{\eps}\,\mathbf 1_{\{\eps\le|x|\le2\eps\}}.
\]
Hence
\[
   \Bigl|\int_E
      \eps^{-1}\eta'(\eps^{-1}|x|)\,\tfrac{x}{|x|}\cdot X(x)\,dx\Bigr|
   \le
   \frac{\|\eta'\|_\infty M_0}{\eps}\,|B_{2\eps}|
   =\|\eta'\|_\infty M_0\,\omega_n\,2^n\eps^{n-1}.
\]
For $n\ge2$ this tends to $0$ as $\eps\downarrow0$.

Combining, the left-hand side of~\eqref{eq:trunc-step} converges to
$\int_E\Div X\,dx$ and the right-hand side converges to
$\int_{\partial^*E}X\cdot\nu_E\,d\mathcal H^{n-1}$.
This proves~\eqref{eq:trunc-flux}.
\end{proof}

\begin{remark}[Why this matters]\label{rem:trunc-use}
Lemma~\ref{lem:truncation} is used twice in the manuscript:

\noindent(i) (Agent~11, oriented radial trace.) Applied to
\[
   X(x)=w(|x-z|)\,\frac{x-z}{|x-z|},
   \qquad
   w(r)=\Bigl|\tfrac{r}{\rho}-1\Bigr|,
\]
on a bounded torsion level $\Erho\subset B_R$, with the centre
$z=\zbar$ or, more generally, the Fusco--Julin centre on a good
isoperimetric level.  Note that $w$ is bounded by $C(R,\rstar)$ on
$\partial^*\Erho$ (since $|x-z|\le R+R'$ for some
$R'=R'(R,\rstar)$), and a direct computation gives
\[
   \Div X(x)=w'(|x-z|)+(n-1)\frac{w(|x-z|)}{|x-z|},
\]
which is bounded by $C(R,\rstar)(1+|x-z|^{-1})$ pointwise.
Hypotheses (a)--(c) of Lemma~\ref{lem:truncation} are met (with
$M_0=\|w\|_\infty$, $K_0=(n-1)\|w\|_\infty$ and a harmless
$L^\infty$ contribution from $w'$).

\noindent(ii) (Identification of the FJ oscillation index, used in
Agent~4 / Agent~11 background.)  Applied to $X(x)=(x-z)/|x-z|$ on
$\Erho$, hypothesis (b) holds with $K_0=n-1$ since
$\Div((x-z)/|x-z|)=(n-1)/|x-z|$ for $x\ne z$; the divergence
identity
\[
   \int_{\partial^*\Erho}
      \Bigl|\nurho-\frac{x-z}{|x-z|}\Bigr|^2\,d\mathcal H^{n-1}
   =2\,\beta(\Erho,z)^2,
\]
quoted in~\eqref{eq:beta-divergence} below, is a direct consequence
together with the identity $|\xi-\eta|^2=2-2\xi\cdot\eta$ for unit
vectors.
\end{remark}

\subsubsection{Lipschitz truncation flux identity for the torsion
            function}
\label{ssec:flux-torsion}

For the torsion function $u_\Om$ on a regular level $\Erho$, the
boundary trace of $\nabla u$ on $\partial^*\Erho$ is not a classical
trace.  The following identity replaces the formal application of
the divergence theorem to $\nabla u$ on $\Erho$; it is used in
Agent~9 to derive the deficit identity.

\begin{lemma}[Flux $=$ volume; cf.\ Maggi Ch.~17]
\label{lem:flux-volume}
Let $\Om\subset\R^n$ be open with $0<|\Om|<\infty$ and let
$u=u_\Om$.  For $\mathcal L^1$-a.e.\ $t\in(0,\|u\|_\infty)$, the
super-level $E_t=\{u>t\}$ has finite perimeter and
\begin{equation}\label{eq:flux-volume}
   \int_{\partial^*E_t}|\nabla u|\,d\mathcal H^{n-1}=|E_t|.
\end{equation}
\end{lemma}

\begin{proof}
This is proved in
\texttt{WIP\_LevelSetIdentity.tex}~Lem.~3.x via Lipschitz
truncation of the test function $\eta_\eps(x):=\phi_\eps(u(x)-t)$
with $\phi_\eps(s)=\min(1,\max(0,s/\eps))$, testing the weak
equation $-\Delta u=1$ in $H^1_0(\Om)$ against $\eta_\eps$, applying
the coarea formula~\eqref{eq:fed-fleming} to the right-hand side
$\int|\nabla u\cdot\nabla\eta_\eps|=\eps^{-1}\int_{\{t<u<t+\eps\}}
|\nabla u|^2$, and passing to the limit $\eps\downarrow 0$ at every
Lebesgue point of $t\mapsto\int_{\partial^*E_t}|\nabla u|$.  The same
calculation, with $|\nabla u|^2$ replaced by $|\nabla u|^p$ for
suitable $p$, gives the rigorous interpretation of
$\int_{\partial^*E_t}\nabla u\cdot\nu_{E_t}\,d\mathcal H^{n-1}
=-\int_{\partial^*E_t}|\nabla u|$ as the boundary trace of
$\nabla u$ on $\partial^*E_t$ in $L^1$ sense; see
\cite[Ch.~17]{Maggi2012}.
\end{proof}

\subsubsection{Reshetnyak lower semicontinuity and finite-perimeter
            passage of first variations}
\label{ssec:reshetnyak}

\begin{theorem}[Reshetnyak lower semicontinuity for linear-growth
                integrands]\label{thm:reshetnyak}
Let $U\subset\R^n$ be open and let
$f:U\times\R^n\to[0,\infty)$ be continuous, $f(x,\cdot)$ convex and
positively $1$-homogeneous for every $x\in U$.  If
$E_j,E\subset\R^n$ are sets of finite perimeter with
$\mathbf 1_{E_j}\to\mathbf 1_E$ in $L^1_{\rm loc}(U)$, and if the
distributional gradients $D\mathbf 1_{E_j}$ converge weakly-$*$ in
$U$ to $D\mathbf 1_E$, then
\[
   \int_{\partial^*E\cap U}f\bigl(x,\nu_E(x)\bigr)\,d\mathcal H^{n-1}
   \le\liminf_{j\to\infty}\int_{\partial^*E_j\cap U}
        f\bigl(x,\nu_{E_j}(x)\bigr)\,d\mathcal H^{n-1}.
\]
\end{theorem}

\begin{proof}[Reference]
\cite[Thm.~2.38]{AFP2000}.
\end{proof}

This is the workhorse used to pass smooth-flow first-variation
identities to the finite-perimeter setting: the smooth-flow proof of
the metric first-variation bound is performed in
\texttt{WIP\_MetricFramework.tex}, and the four-step approximation
of \texttt{Plan 2/gmt-inputs-for-metric-closure.md} \S 5 turns it
into a finite-perimeter statement.  Concretely the approximation is:

\begin{enumerate}
   \item approximate $u_\Om$ by $u_j\in C^\infty(\Om)$ with $u_j\to
       u$ in $W^{1,2}(\Om)$ and, after a Sobolev--Sard reduction to a
       full-measure set of levels, in the coarea sense on the fixed
       annulus $\rho\in[\rstar,1]$;
   \item form the smooth nested family
       $E_{\rho,j}:=\{u_j>t_j(\rho)\}$, with $t_j(\rho)$ the
       inverse-distribution function of $u_j$;
   \item apply the smooth first-variation bound (proved in
       \texttt{WIP\_MetricFramework.tex} and recalled in
       \texttt{Plan 2/gmt-inputs-for-metric-closure.md} \S 4) to
       $\rho\mapsto E_{\rho,j}$;
   \item pass to the limit using Proposition~\ref{prop:lsc-perimeter}
       (lower semicontinuity of perimeter),
       Theorem~\ref{thm:reshetnyak} (Reshetnyak lower semicontinuity
       for the normal integrals
       $\int|V_\rho-H_{\zbar,\rho}-a\cdot\nurho|\,d\mathcal H^{n-1}$),
       and the standard lower semicontinuity of metric derivatives
       for $L^1$-convergent curves \cite[Thm.~2.2.1]{AGS2008}.
\end{enumerate}

The detailed passage is contained in the metric-framework section
(Agent~10).  We isolate the following corollary, which is the form
needed by Agents 11--12.

\begin{corollary}[Reshetnyak passage for the homothetic defect]
\label{cor:reshetnyak-homot}
Let $E_j,E\subset B_R$ have finite perimeter with $|E_j|=|E|=
\omega_n\rho^n$ and $\mathbf 1_{E_j}\to\mathbf 1_E$ in $L^1(\R^n)$,
together with $P(E_j)\to P(E)$.  Let $z_j\to z\in\R^n$ be a sequence
of centres.  For every $a\in\R^n$,
\[
   \int_{\partial^*E}
      \bigl|V-H_{z,\rho}-a\cdot\nu_E\bigr|\,d\mathcal H^{n-1}
   \le\liminf_{j\to\infty}
   \int_{\partial^*E_j}
      \bigl|V_j-H_{z_j,\rho}-a\cdot\nu_{E_j}\bigr|\,d\mathcal H^{n-1},
\]
whenever $V_j\to V$ in $L^1(\mathcal H^{n-1}\llcorner\partial^*E_j)$ in
the sense of weak convergence of varifolds.
\end{corollary}

\begin{proof}[Sketch]
Apply Theorem~\ref{thm:reshetnyak} to the integrand
$f(x,\nu)=|V(x)-H_{z,\rho}(x)-a\cdot\nu|$, which is continuous in
$x$, convex in $\nu$, and positively $1$-homogeneous (because
$H_{z,\rho}(x)=(x-z)\cdot\nu(x)/\rho$ is linear in $\nu$).  The
convergence $P(E_j)\to P(E)$ together with $\mathbf 1_{E_j}\to
\mathbf 1_E$ in $L^1$ gives, by~\cite[Thm.~12.30]{Maggi2012}, weak-$*$
convergence of $D\mathbf 1_{E_j}$ to $D\mathbf 1_E$ in the sense
needed by Theorem~\ref{thm:reshetnyak}.  Convergence of $V_j$ in
varifold sense is the standing hypothesis.
\end{proof}

\subsubsection{Measurable selection of levelwise centres}
\label{ssec:meas-sel}

The Fusco--Julin centre package (Agent~11) attaches to every level
$\rho\in[\rstar,\rho_\delta]$ in the good set $G_I\subset
[\rstar,\rho_\delta]$ at least one point $z\in\R^n$ such that the
quantitative isoperimetric inequality is realised at $z$; this point
is unique only up to a translation in the bad set, and to use it
inside metric-derivative integrals one must choose it in a Borel way.

\begin{lemma}[Borel selection of centres]
\label{lem:borel-selection}
Let $G\subset[\rstar,1]$ be Borel and for each $\rho\in G$ let
$Z(\rho)\subset\R^n$ be a non-empty closed set; suppose the
graph $\{(\rho,z)\in G\times\R^n:z\in Z(\rho)\}$ is a Borel subset of
$G\times\R^n$ and $\sup_{\rho\in G}\sup_{z\in Z(\rho)}|z|<\infty$.
Then there exists a bounded Borel map $\sigma:G\to\R^n$ with
$\sigma(\rho)\in Z(\rho)$ for every $\rho\in G$.
\end{lemma}

\begin{proof}
The graph of $Z$ is a Borel subset of the standard Borel space
$G\times\R^n$ with non-empty closed $\rho$-sections.  The
Kuratowski--Ryll-Nardzewski measurable selection theorem
\cite[Thm.~III.6 and Thm.~III.8]{CastaingValadier} (also Federer,
\cite[\S 2.13]{Federer1969}) provides a Borel selection
$\sigma:G\to\R^n$.  Boundedness of $\sigma$ follows from the standing
hypothesis that the union $\bigcup_{\rho\in G}Z(\rho)$ is bounded.

Outside $G$ (i.e.\ in the bad set), set $\sigma(\rho):=0$; this
extends $\sigma$ to a bounded Borel map on $[\rstar,1]$.
\end{proof}

\begin{remark}[Application]\label{rem:borel-application}
In the manuscript, $Z(\rho)$ is taken to be the set of
$\eta_I$-almost minimisers in the scaled Fusco--Julin inequality
applied to $\Erho$: explicitly, the set of $z\in\R^n$ such that
\[
   |\Erho\Delta B_\rho(z)|^2+\rho^{n+1}\beta(\Erho,z)^2
   \le 2\,C_n\,\rho^{n+1}\bigl(P(\Erho)-P(B_\rho)\bigr),
\]
with $C_n$ the Fusco--Julin constant
\cite[Thm.~1.1]{FJ2014}; this is non-empty by Fusco--Julin
itself, closed by lower semicontinuity of the integrals in $z$
(continuity of $z\mapsto|\Erho\Delta B_\rho(z)|$ is elementary, and
continuity of $z\mapsto\beta(\Erho,z)^2=P(\Erho)-(n-1)\int_{\Erho}
|x-z|^{-1}dx$ follows from dominated convergence using the bound
$|x-z|^{-1}\le|x-z_0|^{-1}+|z-z_0||x-z|^{-1}|x-z_0|^{-1}$), and
bounded in $z$ because the Fraenkel part forces $|z|\le R+\rho$ on
$E\subset B_R$.  Lemma~\ref{lem:borel-selection} then provides a
bounded Borel field $\rho\mapsto z_\rho$ on the good set $G_I$.
The Borel-measurability of the graph follows from the fact that
both $(z,E)\mapsto|E\Delta B_\rho(z)|$ and $(z,E)\mapsto
\beta(E,z)^2$ are Borel maps in $z$ for fixed $E$ (in fact
continuous), so the level set is closed and the graph is therefore
Borel.
\end{remark}

\subsubsection{Two auxiliary identities used downstream}
\label{ssec:aux-identities}

For convenience we record two identities used by Agents~11 and~12.

\begin{lemma}[FJ oscillation index identity]
\label{lem:beta-divergence}
Let $E\subset\R^n$ have finite perimeter and finite measure with
$z\notin\R^n$ playing the role of a centre.  Define
\[
   \beta(E,z)^2:=P(E)-(n-1)\int_E|x-z|^{-1}\,dx.
\]
Then $\beta(E,z)^2\ge 0$ and
\begin{equation}\label{eq:beta-divergence}
   \int_{\partial^*E}
      \Bigl|\nu_E(x)-\frac{x-z}{|x-z|}\Bigr|^2\,d\mathcal H^{n-1}
   =2\,\beta(E,z)^2.
\end{equation}
\end{lemma}

\begin{proof}
The identity $|\nu_E(x)-e_r|^2=2-2e_r\cdot\nu_E(x)$ for unit vectors
$e_r$ gives
\[
   \int_{\partial^*E}|\nu_E-e_r|^2
   =2P(E)-2\int_{\partial^*E}e_r\cdot\nu_E\,d\mathcal H^{n-1}.
\]
Apply Lemma~\ref{lem:truncation} (with $z$ in place of the origin) to
the vector field $X(x):=(x-z)/|x-z|$.  Hypothesis (a) holds with
$M_0=1$, and a direct computation gives $\Div X(x)=(n-1)/|x-z|$ for
$x\ne z$, so (b) holds with $K_0=n-1$.  (c) holds because $E$ has
finite measure and, in the manuscript application, $E\subset B_R$.
Therefore
\[
   \int_{\partial^*E}\frac{x-z}{|x-z|}\cdot\nu_E\,d\mathcal H^{n-1}
   =\int_E\Div X\,dx
   =(n-1)\int_E|x-z|^{-1}\,dx.
\]
Combining gives~\eqref{eq:beta-divergence}.  Non-negativity is then
immediate from the left-hand side, and equality with the right-hand
side proves $\beta(E,z)^2\ge0$.
\end{proof}

\begin{lemma}[Total flux of $H_{z,\rho}$]\label{lem:H-total-flux}
For $z\in\R^n$ and a set $\Erho$ of finite perimeter with
$|\Erho|=\omega_n\rho^n$,
\[
   \int_{\partial^*\Erho}H_{z,\rho}\,d\mathcal H^{n-1}
   =\frac{n|\Erho|}{\rho}=n\omega_n\rho^{n-1}=P(B_\rho).
\]
\end{lemma}

\begin{proof}
Apply Theorem~\ref{thm:gauss-green} to $X(x):=(x-z)/\rho\in
C^1(\R^n;\R^n)$, which is $C^1$ everywhere and bounded on $\Erho$
(any finite-measure set is $\mathcal H^{n-1}$-essentially bounded on
its reduced boundary; in the manuscript $\Erho\subset B_R$).
$\Div X=n/\rho$, so the divergence theorem gives
$\int_{\partial^*\Erho}X\cdot\nurho\,d\mathcal H^{n-1}=n|\Erho|/\rho$.
Since $(x-z)\cdot\nurho(x)/\rho=H_{z,\rho}(x)$, the lemma follows.
\end{proof}

\subsubsection{Summary of the GMT preliminaries and how they are used}
\label{ssec:summary}

For the convenience of the reader we summarise which GMT facts enter
each Plan~2 section.

\begin{itemize}
   \item \emph{Level-set deficit identity (Agent~9).}
       BV coarea (Theorem~\ref{thm:bv-coarea}), the Sobolev--Sard
       statement for the torsion function
       (Theorem~\ref{thm:sobolev-sard}), and the Lipschitz-truncation
       flux identity (Lemma~\ref{lem:flux-volume}).

   \item \emph{Metric framework (Agent~10).}
       Lower semicontinuity of perimeter
       (Proposition~\ref{prop:lsc-perimeter}), the divergence
       theorem (Theorem~\ref{thm:gauss-green}), and Reshetnyak lower
       semicontinuity together with the smooth-to-finite-perimeter
       approximation (Theorem~\ref{thm:reshetnyak},
       Corollary~\ref{cor:reshetnyak-homot}).

   \item \emph{Fusco--Julin centre and oriented radial trace
       (Agent~11).}
       Truncation lemma (Lemma~\ref{lem:truncation}); the FJ
       oscillation-index identity (Lemma~\ref{lem:beta-divergence});
       and the Borel selection
       (Lemma~\ref{lem:borel-selection}).

   \item \emph{Integrated action and weighted metric trace
       (Agent~12).}
       The same divergence and BV coarea tools, plus the absolute
       continuity of $\rho\mapsto[\Frho]$ proved in the metric
       framework section.
\end{itemize}

All boundary integrals in the manuscript are on De Giorgi reduced
boundaries; all centre fields are Borel via
Lemma~\ref{lem:borel-selection}; the truncation argument of
Lemma~\ref{lem:truncation} is the rigorous form of the
``$x/|x|$ argument'' invoked at several places in Plan~2.

% End of Preliminaries -- Geometric Measure Theory (Agent 2).

\label{sec:gmt-velocity}%
% UNIFIED 2026-05-13
% =====================================================================
% Manuscript section draft -- Agent 3
% Preliminaries: Torsion, rearrangement, and deficits
%
% This file is to be merged by Agent 17 into
% Manuscript/01_preliminaries_draft.tex.
%
% Canonical notation: Manuscript/00_notation_register.md.
% Standing convention here:
%   - dimension is $n\ge 2$ (not $N$);
%   - the BDV-normalised energy is $E(\Omega)=-T(\Omega)/2$;
%   - the Saint--Venant deficit is $\delta_T(\Omega)=E(\Omega)-E(B^*)\ge 0$;
%   - the Faber--Krahn deficit is
%     $\delta_{\mathrm{FK}}(\Omega)=(|\Omega|/\omega_n)^{2/n}(\lambda_1(\Omega)-\lambda_1(B^*))$;
%   - the Fraenkel asymmetry is $\Fr(\Omega)\in[0,2)$;
%   - the foliation is $E_\rho=\{u>t(\rho)\}$ with $|E_\rho|=\omega_n\rho^n$.
%
% =====================================================================
% SOURCE MAP
% Item                                              Status     Pointer
% ----------------------------------------------------------------------
% B1: torsion energy E, rigidity T, E=-T/2          PROVED HERE
%       (definitions and equivalence), with         CITE       Brasco--De Philippis--
%       sign convention fixed.                                 Velichkov 2015 [BDV],
%                                                              Final/SaintVenantAssembly.tex,
%                                                              Final-BoundedReduction,
%                                                              Final-KohlerJobinTransfer
%                                                              for the two coexisting
%                                                              conventions.
%
% B2: Saint--Venant deficit delta_T,                ADAPTED    Final-BoundedReduction
%       scale-invariant form D,                                §1; WIP_master.tex §II.
%       Saint--Venant inequality.                  CITE        Polya 1948; Talenti 1976.
%
% B3: Fraenkel asymmetry A(Omega), [0,2)            CITED      Maggi 2008 survey,
%                                                              Final/SaintVenantAssembly.tex.
%
% B4: Faber--Krahn deficit delta_FK,                PROVED HERE (normalisation lemma);
%       eigenvalue normalisation                    CITED      Henrot 2006; KJT.tex.
%
% B5: Talenti pointwise comparison                  CITED       Talenti 1976, Thm.1.
%
% B6: volume-radius foliation E_rho = {u>t(rho)},   PROVED HERE properties of t(rho)
%       monotonicity of t                                       follow from coarea;
%                                                  ADAPTED     WIP_LevelSetIdentity.tex
%                                                              §1, WIP_master.tex §II.
%
% B7: profile-gap estimate (sup-norm Talenti bound)  CITED       Talenti 1976 Thm 1.
%
% Profile-gap estimate used by Plan 2 (h(s)=b(s)-v(s)
%   bound and consequence on t(rho))                PROVED HERE Plan 2/level-set-deficit-identity.md §4, Lem 8.1
%
% Bad-set Lebesgue estimate (set where -t'(rho)<tau_0
%   has measure <= K delta_T)                       PROVED HERE WIP_WeightedMetricTrace.tex
%                                                              Lemma "Bad-set" §4
%                                                              ("integrated Talenti gap").
%
% B9: Kohler--Jobin inequality
%       T(Omega) lambda_1(Omega)^{(n+2)/2}
%         >= T(B^*) lambda_1(B^*)^{(n+2)/2}        CITED       Brasco 2014, Thm 3.3;
%                                                              orig. Kohler--Jobin 1978;
%                                                              Final-KohlerJobinTransfer
%                                                              Thm. 2.1.
%
% Elementary one-variable inequality              PROVED HERE  Final-KohlerJobinTransfer
%   (1-s)^{-p}-1 >= ps                                          Lem. 4.1 (reproduced).
% =====================================================================

\subsection{Preliminaries II: Torsion, rearrangement, and deficits}
\label{sec:prelim-torsion}

This section fixes the torsion and Faber--Krahn functionals used
throughout, the Saint--Venant and Faber--Krahn deficits, the Fraenkel
asymmetry, the Talenti comparison, the volume--radius foliation
$\Erho=\{u>t(\rho)\}$, the profile-gap and bad-set bounds used in the
Plan~2 weighted metric trace, and the Kohler--Jobin inequality used in
the final transfer.

Conventions follow \texttt{00\_notation\_register.md}: we write the
ambient dimension as~$n\ge 2$, the BDV-normalised Saint--Venant energy
as $E(\Om)=-T(\Om)/2$, the Saint--Venant deficit as
$\deltaT(\Om)=E(\Om)-E(\Bstar)\ge 0$, the Faber--Krahn deficit as
$\deltaFK(\Om)=(|\Om|/\omega_n)^{2/n}(\lambda_1(\Om)-\lambda_1(\Bstar))$,
and the Fraenkel asymmetry as $\Fr(\Om)\in[0,2)$.

\subsubsection{Torsion potential, rigidity, and energy}
\label{ssec:torsion}

Throughout, $\Om\subset\R^n$ is open with $0<|\Om|<\infty$. No
connectedness, smoothness, or boundedness is assumed unless stated.

\begin{definition}[Torsion potential]\label{def:utorsion}
The \emph{variational torsion function} of $\Om$ is the unique
minimiser
\[
   u_\Om\in H^{1}_{0}(\Om),\qquad
   u_\Om=\arg\min\Bigl\{
   J(v)\,:\,v\in H^{1}_{0}(\Om)
   \Bigr\},
   \qquad
   J(v):=\tfrac12\!\int_{\Om}\!|\nabla v|^{2}\,dx-\!\int_{\Om}\!v\,dx.
\]
Equivalently, $u_\Om$ is the weak solution of $-\Delta u_\Om=1$ in
$\Om$ with $u_\Om=0$ on $\partial\Om$ in the $H^{1}_{0}$ trace sense.
Testing the weak equation against $u_\Om^{-}$ yields $u_\Om\ge 0$
a.e.\ in $\Om$.
\end{definition}

\begin{remark}[Existence]
The functional $J$ is strictly convex, coercive on $H^{1}_{0}(\Om)$ by
the Poincar\'e--Sobolev inequality
$\|v\|_{L^{2^*}}\le C(n)\|\nabla v\|_{L^{2}}$ on sets of finite
measure (valid since $|\Om|<\infty$, see \cite[Cor.~12.18]{Maggi2012}
or \cite[Thm.~4.9]{HenrotBook}), and weakly lower semicontinuous; the
direct method gives existence and uniqueness of $u_\Om$.
\end{remark}

\begin{definition}[Torsional rigidity and Saint--Venant energy]
\label{def:Trigidity}
The \emph{torsional rigidity} of $\Om$ is
\begin{equation}\label{eq:Tdef}
   T(\Om):=\int_{\Om}u_\Om\,dx.
\end{equation}
The \emph{BDV-normalised Saint--Venant energy} of $\Om$ is
\begin{equation}\label{eq:Edef}
   E(\Om):=\min_{v\in H^{1}_{0}(\Om)}J(v)
   =\tfrac12\!\int_{\Om}\!|\nabla u_\Om|^{2}\,dx-\!\int_{\Om}\!u_\Om\,dx.
\end{equation}
\end{definition}

\begin{lemma}[Energy--rigidity identity]\label{lem:E-T}
For every open $\Om\subset\R^n$ with $0<|\Om|<\infty$,
\begin{equation}\label{eq:E-T}
   E(\Om)=-\tfrac12\,T(\Om).
\end{equation}
In particular $E(\Om)\le 0$ on every admissible set, and $T(\Om)>0$
unless $|\Om|=0$.
\end{lemma}

\begin{proof}
Test the weak equation $-\Delta u_\Om=1$ against $u_\Om\in H^{1}_{0}(\Om)$
to get $\int_\Om|\nabla u_\Om|^2\,dx=\int_\Om u_\Om\,dx=T(\Om)$.
Substituting into~\eqref{eq:Edef} gives
$E(\Om)=\tfrac12 T(\Om)-T(\Om)=-\tfrac12 T(\Om)$.
\end{proof}

\begin{remark}[Two coexisting conventions]\label{rem:T-vs-E}
Some sources (e.g.\ \cite[\S1]{Talenti1976}, classical analysis
literature) call $T(\Om)$ the ``torsion energy''. The BDV chain
(\cite{BDV}; \texttt{Final/SaintVenantAssembly.tex}; the global
reduction in \texttt{Final-BoundedReduction}; the Kohler--Jobin
transfer in \texttt{Final-KohlerJobinTransfer}) and the Plan~2
WIP files (\texttt{WIP\_LevelSetIdentity.tex},
\texttt{WIP\_master.tex}) reserve ``Saint--Venant energy'' for the
BDV-normalised functional $E=-T/2$, which is the minimised value of
the variational functional $J$. We adopt the second convention
throughout: $T(\Om)$ is the torsional \emph{rigidity}; $E(\Om)$ is
the BDV-normalised Saint--Venant \emph{energy}. The identity
\eqref{eq:E-T} converts between the two.
\end{remark}

\begin{lemma}[Scaling and ball values]\label{lem:Tscaling}
For every $t>0$,
\begin{equation}\label{eq:Tscaling}
   T(t\Om)=t^{n+2}T(\Om),\qquad E(t\Om)=t^{n+2}E(\Om).
\end{equation}
For the unit ball $B_1$, the explicit torsion function is
$u_{B_1}(x)=(1-|x|^2)/(2n)$, hence
\begin{equation}\label{eq:Tn}
   T_n:=T(B_1)=\int_{B_1}\frac{1-|x|^2}{2n}\,dx=\frac{\omega_n}{n(n+2)},
   \qquad E(B_1)=-\frac{\omega_n}{2n(n+2)}.
\end{equation}
\end{lemma}

\begin{proof}
Scaling: with $v_t(x):=t^2 u_\Om(x/t)\in H^1_0(t\Om)$, one verifies
$-\Delta v_t=1$ in $t\Om$, so by uniqueness $u_{t\Om}=v_t$; then
$T(t\Om)=\int_{t\Om}v_t=t^{n+2}\int_\Om u_\Om=t^{n+2}T(\Om)$, and
\eqref{eq:E-T} extends to $E$. For the ball formula, a direct
computation gives $\Delta((1-|x|^2)/(2n))=-1$ and the boundary
value $0$. Integrating in radial coordinates,
$\int_{B_1}(1-|x|^2)/(2n)\,dx
=(1/(2n))[\omega_n-n\omega_n/(n+2)]=\omega_n/(n(n+2))$.
\end{proof}

\subsubsection{Saint--Venant inequality and Saint--Venant deficit}
\label{ssec:SV-deficit}

The classical Saint--Venant inequality
\cite[Thm.~3.3]{PolyaSzego1951} states that the ball maximises torsional
rigidity at fixed volume.

\begin{theorem}[Saint--Venant inequality, classical]\label{thm:SVclass}
For every open $\Om\subset\R^n$ with $0<|\Om|<\infty$, and with
$\Bstar$ the open ball centred at the origin with $|\Bstar|=|\Om|$,
\begin{equation}\label{eq:SVclass}
   T(\Om)\le T(\Bstar),
   \qquad\text{equivalently}\qquad
   E(\Om)\ge E(\Bstar).
\end{equation}
Equality holds iff $\Om$ is (a.e.\ equal to) a ball.
\end{theorem}

\begin{proof}[Citation]
\cite[Thm.~1.1]{PolyaSzego1951}, also obtained by Schwarz symmetrisation
combined with the P\'olya--Szeg\H{o} principle, or as a corollary of
Talenti's rearrangement comparison (Theorem~\ref{thm:talenti}
below): on the unit-volume normalisation, the latter gives
$u^*_\Om\le u^*_{\Bstar}$ pointwise on $[0,|\Om|]$, hence
$T(\Om)=\int_0^{|\Om|}u^*_\Om(s)ds\le\int_0^{|\Om|}u^*_{\Bstar}(s)ds=T(\Bstar)$.
\end{proof}

\begin{definition}[Saint--Venant deficit $\deltaT$]\label{def:deltaT}
For open $\Om\subset\R^n$ with $0<|\Om|<\infty$ and $\Bstar$ the
ball with $|\Bstar|=|\Om|$, the (BDV-normalised) Saint--Venant
deficit is
\begin{equation}\label{eq:deltaT-def}
   \deltaT(\Om):=E(\Om)-E(\Bstar)\ge 0,
\end{equation}
equivalently $\deltaT(\Om)=\tfrac12(T(\Bstar)-T(\Om))$.
\end{definition}

\begin{remark}[Scale-invariant form]\label{rem:Dform}
The scale-invariant version is
\begin{equation}\label{eq:Dform}
   D(\Om):=E(\Om)\,|\Om|^{-(n+2)/n}-E(B_1)\,|B_1|^{-(n+2)/n}\ge 0,
\end{equation}
which for $|\Om|=\omega_n$ reduces to
$D(\Om)=\omega_n^{-(n+2)/n}\deltaT(\Om)$.
Both formulations are
used: $\deltaT$ in the Plan~2 chain and the bounded Saint--Venant
statements; $D$ in the bounded reduction (Plan~1, \S\,IV).
\end{remark}

\subsubsection{Fraenkel asymmetry}\label{ssec:Fraenkel}

\begin{definition}[Fraenkel asymmetry]\label{def:Fraenkel}
For open $\Om\subset\R^n$ with $0<|\Om|<\infty$, the
\emph{Fraenkel asymmetry} is
\begin{equation}\label{eq:Fraenkel-def}
   \Fr(\Om):=\inf_{x_0\in\R^n}
        \frac{|\Om\,\triangle\,(\Bstar+x_0)|}{|\Bstar|},
\end{equation}
where $\Bstar$ is the ball centred at the origin with $|\Bstar|=|\Om|$.
\end{definition}

\begin{lemma}[Range of $\Fr$]\label{lem:Fraenkel-range}
$\Fr(\Om)\in[0,2)$, with $\Fr(\Om)=0$ iff $\Om$ is (a.e.\ equal to) a
translate of $\Bstar$.
\end{lemma}

\begin{proof}
Non-negativity is immediate. The strict bound $\Fr<2$ follows from
$|\Om\,\triangle\,(\Bstar+x_0)|=|\Om|+|\Bstar|-2|\Om\cap(\Bstar+x_0)|$
and the existence (by translation continuity of Lebesgue measure)
of $x_0\in\R^n$ with $|\Om\cap(\Bstar+x_0)|>0$, so
$|\Om\,\triangle\,(\Bstar+x_0)|<2|\Bstar|$. The infimum is attained
because the map $x_0\mapsto|\Om\,\triangle\,(\Bstar+x_0)|$ is
continuous and tends to $2|\Bstar|$ as $|x_0|\to\infty$
(\cite[Lem.~3.2]{Maggi2008}).
\end{proof}

\begin{remark}[Translation invariance]\label{rem:trans}
Both $\deltaT$, $\deltaFK$, and $\Fr$ are translation invariant; all
estimates below are stated up to translation.
\end{remark}

\subsubsection{Faber--Krahn deficit and eigenvalue normalisation}
\label{ssec:Faber-Krahn}

\begin{definition}[First Dirichlet eigenvalue]\label{def:lambda1}
For open $\Om\subset\R^n$ with $0<|\Om|<\infty$,
\begin{equation}\label{eq:lambda1-def}
   \lambda_1(\Om):=\min_{\substack{v\in H^{1}_{0}(\Om)\\v\not\equiv 0}}
   \frac{\int_\Om|\nabla v|^2\,dx}{\int_\Om v^2\,dx}.
\end{equation}
The minimum is attained (by direct method on
$\{v:\|v\|_{L^2}=1\}$, using the Poincar\'e--Sobolev inequality on
finite-measure sets), and the eigenvalue scales as
$\lambda_1(t\Om)=t^{-2}\lambda_1(\Om)$ for $t>0$. We write
$L_n:=\lambda_1(B_1)$ for the unit-ball first eigenvalue. Numerically,
$L_n=j_{n/2-1,1}^2$, the square of the first positive zero of the
Bessel function $J_{n/2-1}$ \cite[\S1.3]{HenrotBook}.
\end{definition}

\begin{theorem}[Faber--Krahn inequality]\label{thm:FK}
For every open $\Om\subset\R^n$ with $0<|\Om|<\infty$,
\begin{equation}\label{eq:FK}
   \lambda_1(\Om)\ge\lambda_1(\Bstar)
   =L_n\,\bigl(|\Bstar|/\omega_n\bigr)^{-2/n},
\end{equation}
with equality iff $\Om$ is (a.e.\ equal to) a ball.
\end{theorem}

\begin{proof}[Citation]
The classical Faber--Krahn theorem \cite{Faber1923,Krahn1925}, proved
by Schwarz symmetrisation and the P\'olya--Szeg\H{o} inequality. See
\cite[Thm.~3.2.1]{HenrotBook}.
\end{proof}

\begin{definition}[Faber--Krahn deficit $\deltaFK$]
\label{def:deltaFK}
For open $\Om\subset\R^n$ with $0<|\Om|<\infty$ and $\Bstar$ the ball
with $|\Bstar|=|\Om|$, the scale-normalised Faber--Krahn deficit is
\begin{equation}\label{eq:deltaFK-def}
   \deltaFK(\Om)
   :=\Bigl(\frac{|\Om|}{\omega_n}\Bigr)^{2/n}
     \bigl(\lambda_1(\Om)-\lambda_1(\Bstar)\bigr)\ge 0.
\end{equation}
\end{definition}

\begin{lemma}[Scaling invariance of $\deltaFK$]\label{lem:deltaFK-scale}
$\deltaFK(t\Om)=\deltaFK(\Om)$ for every $t>0$.
\end{lemma}

\begin{proof}
$\lambda_1(t\Om)=t^{-2}\lambda_1(\Om)$ and
$|t\Om|^{2/n}=t^2|\Om|^{2/n}$ give
$|t\Om|^{2/n}\lambda_1(t\Om)=|\Om|^{2/n}\lambda_1(\Om)$. Similarly for
$\Bstar$.
\end{proof}

\begin{remark}[Why this normalisation]\label{rem:deltaFK-form}
Both Plan~1 and Plan~2 produce, after Kohler--Jobin transfer, a sharp
quantitative inequality of the form
$\deltaFK(\Om)\ge c_{\mathrm{FK}}(n)\,\Fr(\Om)^2$. The factor
$(|\Om|/\omega_n)^{2/n}$ in~\eqref{eq:deltaFK-def} makes both sides
scale-invariant, mirroring~\eqref{eq:Dform}; see
\texttt{Final-KohlerJobinTransfer}, eq.\,(eq:fkout).
\end{remark}

\subsubsection{Talenti pointwise comparison}\label{ssec:Talenti}

The fundamental rearrangement comparison for the torsion function is
the Talenti pointwise inequality \cite{Talenti1976}. We write
$u^*\colon[0,|\Om|]\to[0,\|u\|_\infty]$ for the decreasing
rearrangement of $u=u_\Om$, $u^*(s):=\inf\{t\ge 0\,:\,m(t)\le s\}$,
where $m(t):=|\{u>t\}|$.

\begin{theorem}[Talenti, 1976]\label{thm:talenti}
Let $\Om\subset\R^n$ be open with $0<|\Om|<\infty$, and let
$u_\Om\in H^{1}_{0}(\Om)$ be the torsion function of
Definition~\ref{def:utorsion}. Let $u_{\Bstar}$ be the torsion
function of the ball $\Bstar$ with $|\Bstar|=|\Om|$. Then
\begin{equation}\label{eq:talenti}
   u^*_\Om(s)\le u^*_{\Bstar}(s)
   =\frac{R_\Om^2-(s/\omega_n)^{2/n}}{2n},
   \qquad s\in[0,|\Om|],
\end{equation}
where $R_\Om:=(|\Om|/\omega_n)^{1/n}$ is the radius of $\Bstar$. In
particular,
\begin{equation}\label{eq:talenti-sup}
   \|u_\Om\|_{L^\infty(\Om)}\le u^*_{\Bstar}(0)=\frac{R_\Om^2}{2n}.
\end{equation}
\end{theorem}

\begin{proof}[Citation]
\cite[Thm.~1]{Talenti1976}. The argument uses isoperimetric
rearrangement of super-level sets and the coarea formula; it is
constructive and gives equality iff $u_\Om$ is (translate of) the
ball torsion function, hence iff $\Om$ is a ball.
\end{proof}

\begin{remark}[Consequence: $L^\infty$ bound used in Plan~2]
\label{rem:talenti-Linf}
The sup-norm bound~\eqref{eq:talenti-sup} is the only reason that
the level-set deficit identity
(\texttt{WIP\_LevelSetIdentity.tex}, Thm.\,1) holds for arbitrary
open sets of finite measure without an explicit boundedness
hypothesis. Throughout Plan~2 we use the normalised volume
$|\Om|=\omega_n$, so $R_\Om=1$ and
$\|u_\Om\|_\infty\le 1/(2n)$.
\end{remark}

\subsubsection{Volume--radius foliation \texorpdfstring{$\Erho=\{u>t(\rho)\}$}{E\_rho=\{u>t(rho)\}}}
\label{ssec:foliation}

We now record the volume-radius parametrisation of the super-level
sets of $u_\Om$, used throughout Plan~2.

\begin{definition}[Distribution function and its inverse]
\label{def:t-rho}
For $\Om$ open with $|\Om|=\omega_n$ (the Plan~2 normalisation),
$E_t:=\{u_\Om>t\}$ and $m(t):=|E_t|$. By the BV coarea formula and
Lemma~\ref{lem:coarea-mprime} below, $m$ is right-continuous and
non-increasing on $[0,\|u\|_\infty]$ with $m(0^+)\le\omega_n$ and
$m(\|u\|_\infty)=0$. Define
\begin{equation}\label{eq:trho}
   t(\rho):=\inf\bigl\{t\ge 0\,:\,m(t)\le\omega_n\rho^n\bigr\},
   \qquad\rho\in[0,1],
\end{equation}
and the volume-radius foliation
\begin{equation}\label{eq:Erho}
   \Erho:=\{u_\Om>t(\rho)\},\qquad
   |\Erho|=\omega_n\rho^n,\quad
   \mu(d\rho):=(-t'(\rho))\,d\rho.
\end{equation}
\end{definition}

\begin{lemma}[Properties of $t(\rho)$]\label{lem:trho-props}
Let $\Om$ open with $|\Om|=\omega_n$ and $u=u_\Om$. Then
\begin{enumerate}
\item[(i)] $\rho\mapsto t(\rho)$ is non-increasing on $[0,1]$, with
$t(1)=0$ and $0\le t(\rho)\le\|u\|_\infty\le 1/(2n)$ by
\eqref{eq:talenti-sup}.
\item[(ii)] $t$ is absolutely continuous on every interval
$[\rstar,1]\subset(0,1]$.
\item[(iii)] On a co-null set in $[0,1]$,
$t'(\rho)\le 0$, and
\begin{equation}\label{eq:dt-bound}
   0\le -t'(\rho)\le\frac{1}{n}
   \qquad\text{for a.e.\ }\rho\in[0,1].
\end{equation}
\end{enumerate}
\end{lemma}

\begin{proof}
(i) and the boundary value $t(1)=0$ follow from $m(0^+)\le\omega_n$
and $m(t(\rho))\le\omega_n\rho^n$; the sup-norm bound is
Theorem~\ref{thm:talenti}.

(ii) The map $m$ is, by Lemma~\ref{lem:coarea-mprime}, absolutely
continuous on $[\eps,\|u\|_\infty]$ for every $\eps>0$ (its derivative
$-\alpha(t)$ is in $L^1_{\mathrm{loc}}$). The inverse of a strictly
decreasing absolutely continuous function with non-vanishing
derivative on an open set is absolutely continuous; on plateau
levels (where $m$ is constant) $t$ has a jump of measure zero and
the equation $|E_\rho|=\omega_n\rho^n$ pins $t(\rho)$ on the
right-continuous representative. On the foliation interval
$[\rstar,1]$ this gives absolute continuity of $t$.

(iii) Differentiating $m(t(\rho))=\omega_n\rho^n$ at a Lebesgue point
of $\rho\mapsto t'(\rho)$ at which $m$ is differentiable at
$t(\rho)$, by the chain rule for absolutely continuous functions and
Lemma~\ref{lem:coarea-mprime},
$m'(t(\rho))\,t'(\rho)=n\omega_n\rho^{n-1}$, i.e.\
\begin{equation}\label{eq:tprime-formula}
   -t'(\rho)=\frac{n\omega_n\rho^{n-1}}{\alpha(t(\rho))},
   \qquad\alpha(t):=\!\int_{\partial^*E_t}\!|\nabla u|^{-1}\,d\mathcal H^{n-1}.
\end{equation}
By the Cauchy--Schwarz inequality on $\partial^*E_t$ and the
isoperimetric inequality $P(E_t)\ge n\omega_n^{1/n}m(t)^{1-1/n}$,
\begin{equation}\label{eq:alpha-LB}
   \alpha(t)\cdot\beta(t)\ge P(E_t)^2\ge n^2\omega_n^{2/n}m(t)^{2-2/n},
   \qquad
   \beta(t):=\!\int_{\partial^*E_t}\!|\nabla u|\,d\mathcal H^{n-1}=m(t),
\end{equation}
the equality $\beta(t)=m(t)$ being the flux identity of
Lemma~\ref{lem:flux-id} below. Combined with $m(t(\rho))=\omega_n\rho^n$,
$\alpha(t(\rho))\ge n^2\omega_n^{2/n}(\omega_n\rho^n)^{2-2/n}/m(t(\rho))
=n^2\omega_n\rho^{n-2}\cdot(\omega_n\rho^n)^{1-2/n}/(\omega_n\rho^n)
=n\omega_n\rho^{n-1}\cdot n$, hence
$-t'(\rho)\le n\omega_n\rho^{n-1}/(n^2\omega_n\rho^{n-1})=1/n$.
\end{proof}

\begin{remark}[Where $-t'\le 1/n$ is used]
\label{rem:H3}
The bound \eqref{eq:dt-bound} is the hypothesis (H3) of the abstract
weighted metric trace lemma (\texttt{WIP\_WeightedMetricTrace.tex},
Thm.\,5.1): it bounds the density of $\mu$ uniformly.
\end{remark}

The coarea and flux identities used in the proof of
Lemma~\ref{lem:trho-props}(iii) are recorded for reference (proved in
the geometric measure theory preliminaries,
\S\ref{sec:prelim-gmt}):

\begin{lemma}[Coarea form of $-m'$]\label{lem:coarea-mprime}
For a.e.\ $t\in(0,\|u\|_\infty)$, $-m'(t)=\alpha(t)$ as defined in
\eqref{eq:tprime-formula}. $m$ is absolutely continuous on
$[\eps,\|u\|_\infty]$ for every $\eps>0$.
\end{lemma}

\begin{proof}[Citation]
\cite[Thm.~13.1]{Maggi2012}, \cite[Thm.~3.40]{AFP2000}.
\end{proof}

\begin{lemma}[Flux--volume identity]\label{lem:flux-id}
For a.e.\ $t\in(0,\|u\|_\infty)$, $E_t$ has finite perimeter and
$\beta(t)=m(t)$.
\end{lemma}

\begin{proof}[Citation]
By Lipschitz truncation of the weak equation $-\Delta u=1$ against the
cutoff $\phi_\eps(u-t)\in H^{1}_{0}(\Om)$ and the finite-perimeter
divergence theorem; cf.\ \cite[Thm.~16.3]{Maggi2012}.
\end{proof}

\subsubsection{Profile-gap estimate used by Plan~2}\label{ssec:profile-gap}

We now record the profile-gap estimate that converts an
$L^\infty$-control of $u^*_\Om-u^*_{\Bstar}$ on a near-boundary slab
into a positive lower bound for $|T_\eta|=v(L-2\eta)-v(L-\eta)$, hence
into the hypothesis~(H2) of the abstract weighted metric trace lemma.

We work in the Plan~2 normalisation: $|\Om|=\omega_n=:L$,
$R_\Om=R=1$, and we write $b(s):=u^*_{B_1}(s)$, $v(s):=u^*_\Om(s)$,
$h(s):=b(s)-v(s)$.

\begin{lemma}[Profile gap; integrated form]\label{lem:profile-gap-int}
Let $\Om\subset\R^n$ be open with $|\Om|=\omega_n$. Then for every
$s\in(0,\omega_n]$,
\begin{equation}\label{eq:profile-gap-int}
   0\le h(s)=b(s)-v(s)\le\frac{2\deltaT(\Om)}{s}.
\end{equation}
In particular, for every $\theta\in(0,1)$,
\begin{equation}\label{eq:profile-gap-slab}
   \|h\|_{L^\infty([\theta\omega_n,\omega_n])}
   \le\frac{2\deltaT(\Om)}{\theta\omega_n}.
\end{equation}
\end{lemma}

\begin{proof}
By Talenti's pointwise comparison (Theorem~\ref{thm:talenti}),
$v(s)\le b(s)$, so $h(s)\ge 0$ and $h$ is the difference of
non-increasing absolutely continuous functions, hence absolutely
continuous on $(0,\omega_n]$ with $h(0^+)\le b(0)<\infty$ and
$h(\omega_n)=0$.

For the upper bound, integrate Talenti's level-set differential
inequality (cf.~\cite[\S4]{Talenti1976},
\texttt{Plan 2/level-set-deficit-identity.md} §4):
$v'(s)\ge b'(s)$ for a.e.\ $s$, equivalently $h'(s)\le 0$. Since
$h\ge 0$ is non-increasing,
$h(s)\cdot s\le\int_0^s h(\sigma)\,d\sigma\le\int_0^{\omega_n}h$. By
equimeasurability,
\[
   \int_0^{\omega_n}h(\sigma)\,d\sigma
   =\int_0^{\omega_n}b(\sigma)\,d\sigma-\int_0^{\omega_n}v(\sigma)\,d\sigma
   =T(B_1)-T(\Om)=2\deltaT(\Om),
\]
using $|\Om|=\omega_n=|B_1|$ and Lemma~\ref{lem:E-T}. Hence
$h(s)\le 2\deltaT(\Om)/s$.
\end{proof}

\begin{lemma}[Foliation length on a boundary slab]
\label{lem:slab-length}
Let $\Om\subset\R^n$ open with $|\Om|=\omega_n$. For
$\eta\in(0,\omega_n/4)$ define the volume slab $S_\eta:=[\omega_n-2\eta,\omega_n-\eta]$
and the level slab
$T_\eta:=v(S_\eta)=[v(\omega_n-\eta),v(\omega_n-2\eta)]$. There exist
constants $c_n,C_n>0$ depending only on $n$ such that
\begin{equation}\label{eq:slab-b-bounds}
   c_n\eta\le b(s)\le C_n\eta\qquad\text{for every }s\in S_\eta,
\end{equation}
and if $\deltaT(\Om)\le c_n\eta$ then
\begin{equation}\label{eq:slab-T-eta}
   |T_\eta|=v(\omega_n-2\eta)-v(\omega_n-\eta)\ge c_n\eta.
\end{equation}
The constants are explicit: $c_n=1/(C\cdot n)$, $C_n=C\cdot n^{-1}$,
$C$ a universal dimensional constant.
\end{lemma}

\begin{proof}
Set $R=1$, $L=\omega_n$. From the explicit ball profile
$b(s)=(1-(s/\omega_n)^{2/n})/(2n)$, expanding
$1-(1-\eta/\omega_n)^{2/n}$ to first order in $\eta$ for $\eta\le\omega_n/4$
gives
$c_n\eta/\omega_n\le 1-(1-\eta/\omega_n)^{2/n}\le C_n\eta/\omega_n$
with $c_n=1/n$, $C_n=4/n$ (using
$1-(1-x)^{2/n}\ge(2/n)x(1-x)\ge x/n$ for $x\in[0,1/2]$, and
$1-(1-x)^{2/n}\le(2/n)x\cdot(1-x)^{2/n-1}\le 4x/n$ for $n\ge 2$).
Multiplying by $1/(2n)$ gives \eqref{eq:slab-b-bounds}.

For \eqref{eq:slab-T-eta}, write
$v(\omega_n-2\eta)-v(\omega_n-\eta)
=[b(\omega_n-2\eta)-b(\omega_n-\eta)]
 -[h(\omega_n-2\eta)-h(\omega_n-\eta)]$.
The first bracket is, by the explicit ball profile and
$\omega_n-2\eta\ge\omega_n/2$, at least
$c_n\eta$ for a dimensional $c_n$. The second bracket is bounded by
$2\|h\|_{L^\infty([\omega_n/2,\omega_n])}\le 4\deltaT(\Om)/\omega_n$
via \eqref{eq:profile-gap-slab}. The smallness hypothesis
$\deltaT(\Om)\le c_n\eta$ (after halving $c_n$) ensures the first
bracket dominates the second.
\end{proof}

\begin{remark}[Adaptation source]
Lemma~\ref{lem:profile-gap-int} is the integrated form of the
Talenti-comparison consequence
$0\le b(s)-v(s)\le 2\deltaT/s$ derived from
\cite[Thm.~1]{Talenti1976} and the Saint--Venant deficit identity
$\Gamma(\Om)=2(E(\Om)-E(\Bstar))/|\Om|$. Lemma~\ref{lem:slab-length}
provides the $|T_\eta|$-bound used to verify hypothesis~(H2) in the
abstract trace lemma of \S\ref{subsec:weighted-metric-trace}.
\end{remark}

\subsubsection{Bad-set Lebesgue estimate used by Plan~2}\label{ssec:badset}

The bad-set estimate converts the integrated Talenti gap into a
Lebesgue control on the set where the foliation slows down. It is
the input that the weighted metric trace lemma needs to convert the
weighted kinetic bound
$\int|\dot F_\rho|_\X^2\,d\mu\le C\deltaT$ into the unweighted kinetic
bound $\int|\dot F_\rho|_\X^2\,d\rho\le C'\deltaT$
(cf.\ \S\ref{subsec:weighted-metric-trace}).

\begin{lemma}[Integrated Talenti gap on $-t'$]\label{lem:talenti-gap}
Let $\Om\subset\R^n$ be open with $|\Om|=\omega_n$ and
$\rstar\in(0,1)$. Let $t_B$ denote the inverse distribution
function of $u_{B_1}$,
$t_B(\rho)=u^*_{B_1}(\omega_n\rho^n)=(1-\rho^2)/(2n)$, so
$-t_B'(\rho)=\rho/n$. There is $K_1=K_1(n,\rstar)>0$ such that
\begin{equation}\label{eq:talenti-gap-int}
   \int_{\rstar}^{1}\bigl|(-t_B'(\rho))-(-t'(\rho))\bigr|\,d\rho
   \le K_1\,\deltaT(\Om).
\end{equation}
\end{lemma}

\begin{proof}
Since both $t,t_B$ are non-increasing with $t(1)=t_B(1)=0$,
$\psi(\rho):=(-t_B'(\rho))-(-t'(\rho))$ satisfies
$\int_\rho^1\psi(r)\,dr=t(\rho)-t_B(\rho)=v(\omega_n\rho^n)-b(\omega_n\rho^n)
=-h(\omega_n\rho^n)\le 0$ for every $\rho\in[\rstar,1]$
(by Talenti, $v\le b$, hence $h\ge 0$). However $\psi$ need not
have a sign; the proof of Lemma~\ref{lem:profile-gap-int} shows
$h\ge 0$ globally, so the function $-h(\omega_n\rho^n)$ is
$\le 0$ on $[\rstar,1]$.

Substitute $s=\omega_n\rho^n$, $ds=n\omega_n\rho^{n-1}d\rho$:
\[
   \int_{\rstar}^1\psi(\rho)\,d\rho
   =\bigl[\,t(\rho)-t_B(\rho)\,\bigr]_{\rstar}^1
   =-\bigl(t(\rstar)-t_B(\rstar)\bigr)
   =h(\omega_n\rstar^n)
   \le \frac{2\deltaT(\Om)}{\omega_n\rstar^n}
\]
by Lemma~\ref{lem:profile-gap-int}. Since $\psi$ need not be signed,
this controls only $\int\psi$ not $\int|\psi|$. We restore the
absolute value by a second comparison. On any subinterval
$[\rho_1,\rho_2]\subset[\rstar,1]$,
\[
   \Bigl|\!\int_{\rho_1}^{\rho_2}\!\psi(r)\,dr\Bigr|
   =\bigl|(t(\rho_1)-t(\rho_2))-(t_B(\rho_1)-t_B(\rho_2))\bigr|
   =|h(\omega_n\rho_1^n)-h(\omega_n\rho_2^n)|
   \le h(\omega_n\rstar^n)
\]
because $h\ge 0$ is non-increasing. By the
Hahn--Jordan decomposition of $\psi\,d\rho$ on $[\rstar,1]$ as
$d\psi^+-d\psi^-$,
\[
   \int_{\rstar}^1|\psi|\,d\rho
   =\psi^+([\rstar,1])+\psi^-([\rstar,1])
   \le 2\,\sup_{[\rho_1,\rho_2]\subset[\rstar,1]}
   \Bigl|\int_{\rho_1}^{\rho_2}\psi\Bigr|
   \le 2\,h(\omega_n\rstar^n)
   \le\frac{4\deltaT(\Om)}{\omega_n\rstar^n},
\]
which is \eqref{eq:talenti-gap-int} with
$K_1=4/(\omega_n\rstar^n)$.
\end{proof}

\begin{lemma}[Bad-set Lebesgue estimate; (H3-bad)]
\label{lem:badset}
Let $\Om\subset\R^n$ open with $|\Om|=\omega_n$, $\rstar\in(0,1)$.
Set
\[
   \tau_0:=\frac{\rstar}{2n},
   \qquad
   B_{\tau_0}:=\{\rho\in[\rstar,1]\,:\,-t'(\rho)<\tau_0\}.
\]
Then
\begin{equation}\label{eq:badset}
   |B_{\tau_0}|\le K_2(n,\rstar)\,\deltaT(\Om),
   \qquad K_2(n,\rstar):=\frac{8n}{\omega_n\rstar^{n+1}}.
\end{equation}
\end{lemma}

\begin{proof}
On $B_{\tau_0}$, $-t'(\rho)<\rstar/(2n)$ and $-t_B'(\rho)=\rho/n\ge\rstar/n$
for $\rho\ge\rstar$. Hence
$\psi(\rho)=(-t_B'(\rho))-(-t'(\rho))>\rstar/n-\rstar/(2n)=\rstar/(2n)$
on $B_{\tau_0}$. Integrating and applying
Lemma~\ref{lem:talenti-gap},
\[
   \frac{\rstar}{2n}\,|B_{\tau_0}|
   \le\int_{B_{\tau_0}}\psi(\rho)\,d\rho
   \le\int_{\rstar}^{1}|\psi(\rho)|\,d\rho
   \le K_1\,\deltaT(\Om),
\]
so $|B_{\tau_0}|\le (2n/\rstar)K_1\deltaT(\Om)$. Substituting
$K_1=4/(\omega_n\rstar^n)$ gives \eqref{eq:badset}.
\end{proof}

\begin{remark}[Adaptation source]
Lemma~\ref{lem:badset} is the bad-set lemma used in
\S\ref{subsec:weighted-metric-trace}. Lemma~\ref{lem:talenti-gap} is
the ``integrated Talenti gap'' invoked there; the constants
$K_1(n,\rstar),K_2(n,\rstar)$ are made explicit here. No published
reference for this exact statement was located; the proof is
self-contained from Talenti's comparison and the level-set deficit
identity.
\end{remark}

\subsubsection{Kohler--Jobin inequality}\label{ssec:KJ}

\begin{theorem}[Kohler--Jobin inequality; \protect\cite{KohlerJobin1978,Brasco2014}]
\label{thm:KJ}
For every open bounded $\Om\subset\R^n$ with $\lambda_1(\Om)<\infty$,
\begin{equation}\label{eq:KJ}
   T(\Om)\cdot\lambda_1(\Om)^{(n+2)/2}
   \ge T(B_1)\cdot\lambda_1(B_1)^{(n+2)/2}
   =T_n\,L_n^{(n+2)/2}.
\end{equation}
Equality holds iff $\Om$ is a ball (up to translation and a null set).
\end{theorem}

\begin{proof}[Citation]
\cite[Thm.~3.3]{Brasco2014}, giving a modern proof by Schwarz
symmetrisation; original \cite{KohlerJobin1978}. The functional
$\Psi(\Om):=T(\Om)\lambda_1(\Om)^{(n+2)/2}$ is scale-invariant
($T(t\Om)=t^{n+2}T(\Om)$ and $\lambda_1(t\Om)=t^{-2}\lambda_1(\Om)$,
so $\Psi(t\Om)=\Psi(\Om)$), hence the right side is the universal
constant $T_n L_n^{(n+2)/2}$ depending only on~$n$.
\end{proof}

\begin{remark}[Normalisation used in the final transfer]
\label{rem:KJ-norm}
The form~\eqref{eq:KJ} is the one used in the Kohler--Jobin transfer
of the present manuscript
(\S\ref{sec:plan1-kj-transfer}). At
fixed volume $|\Om|=|\Bstar|$, rearranging~\eqref{eq:KJ} yields the
pointwise transfer
\begin{equation}\label{eq:KJ-ptwise}
   \lambda_1(\Om)
   \ge\lambda_1(\Bstar)
   \Bigl(\frac{T(\Bstar)}{T(\Om)}\Bigr)^{2/(n+2)},
\end{equation}
equivalently
\begin{equation}\label{eq:KJ-deficit}
   \lambda_1(\Om)-\lambda_1(\Bstar)
   \ge\lambda_1(\Bstar)\Bigl[(1-s)^{-1/\beta}-1\Bigr],
   \qquad
   s:=\frac{T(\Bstar)-T(\Om)}{T(\Bstar)}\in[0,1),
   \quad\beta=\tfrac{n+2}{2}.
\end{equation}
\end{remark}

\begin{lemma}[Elementary one-variable inequality]\label{lem:onevar}
For $p\in(0,1]$ and $s\in[0,1)$,
\begin{equation}\label{eq:onevar}
   (1-s)^{-p}-1\ge ps.
\end{equation}
\end{lemma}

\begin{proof}
$f(s):=(1-s)^{-p}-1$ satisfies $f(0)=0$ and
$f'(s)=p(1-s)^{-(p+1)}\ge p$ for $s\in[0,1)$. By the fundamental
theorem of calculus, $f(s)=\int_0^s f'(\tau)\,d\tau\ge ps$.
\end{proof}

Combining~\eqref{eq:KJ-deficit} with Lemma~\ref{lem:onevar} for
$p=1/\beta=2/(n+2)$ gives, in the small-deficit regime
$s\le 1/2$, equivalently $T(\Bstar)-T(\Om)\le T(\Bstar)/2$,
equivalently $\deltaT(\Om)\le T(\Bstar)/4$:
\begin{equation}\label{eq:KJ-linear}
   \lambda_1(\Om)-\lambda_1(\Bstar)
   \ge\frac{2\lambda_1(\Bstar)}{(n+2)T(\Bstar)}
       \bigl(T(\Bstar)-T(\Om)\bigr)
   =\frac{4\lambda_1(\Bstar)}{(n+2)T(\Bstar)}\,\deltaT(\Om).
\end{equation}

\begin{remark}[Dimension dependence]\label{rem:KJ-dim}
The Kohler--Jobin multiplier in~\eqref{eq:KJ-linear} depends only on
$n$: by scale invariance,
$2\lambda_1(\Bstar)/((n+2)T(\Bstar))=2L_n/((n+2)T_n)=2nL_n/\omega_n$,
using \eqref{eq:Tn}. This is the universal dimension-only constant
that converts $\deltaT$ to $\deltaFK$ in the final theorem.
\end{remark}

\subsubsection{Summary of constants and constants used downstream}
\label{ssec:prelim-summary}

We record the universal ball constants and the dimension-tracked
constants introduced in this section, for use in the Plan~1 and
Plan~2 chains.

\begin{center}
\renewcommand{\arraystretch}{1.25}
\begin{tabular}{|l|l|l|}\hline
Symbol & Value / dependence & Reference\\\hline
$\omega_n=|B_1|$ & dim.~only & --- \\
$T_n=T(B_1)=\omega_n/(n(n+2))$ & dim.~only & Lemma~\ref{lem:Tscaling}\\
$L_n=\lambda_1(B_1)=j_{n/2-1,1}^2$ & dim.~only & Def.~\ref{def:lambda1}\\
$E(B_1)=-T_n/2=-\omega_n/(2n(n+2))$ & dim.~only & \eqref{eq:E-T}\\
$2L_n/((n+2)T_n)=2nL_n/\omega_n$ & Kohler--Jobin multiplier ($n$ only) & Rmk.~\ref{rem:KJ-dim}\\\hline
$K_1(n,\rstar)=4/(\omega_n\rstar^n)$ & integrated Talenti gap constant & Lem.~\ref{lem:talenti-gap}\\
$K_2(n,\rstar)=8n/(\omega_n\rstar^{n+1})$ & bad-set Lebesgue constant & Lem.~\ref{lem:badset}\\
$\tau_0=\rstar/(2n)$ & bad-set threshold & Lem.~\ref{lem:badset}\\\hline
$c_n,C_n$ & ball-profile constants in slab estimates & Lem.~\ref{lem:slab-length}\\\hline
\end{tabular}
\end{center}

% =====================================================================
% Local bibliography for this section. Agent 17 will merge this into
% the master bibliography (Manuscript/04_references_audit.md).
% DISABLED by Agent 19 (Final Assembly): the master bibliography is in
% Manuscript/references.bib, loaded from main.tex.
% =====================================================================
\iffalse
\begin{thebibliography}{99}

\bibitem{AFP2000}
L.~Ambrosio, N.~Fusco, D.~Pallara,
\emph{Functions of Bounded Variation and Free Discontinuity
Problems}, Oxford University Press, 2000.

\bibitem{BDV}
L.~Brasco, G.~De~Philippis, B.~Velichkov,
\emph{Faber--Krahn inequalities in sharp quantitative form},
Duke Math.\ J.\ \textbf{164} (2015), 1777--1831.

\bibitem{Brasco2014}
L.~Brasco,
\emph{On torsional rigidity and principal frequencies: an invitation
to the Kohler--Jobin rearrangement technique}, ESAIM Control Optim.\
Calc.\ Var.\ \textbf{20} (2014), 315--338.

\bibitem{Faber1923}
G.~Faber, \emph{Beweis, dass unter allen homogenen Membranen von
gleicher Fl\"ache und gleicher Spannung die kreisf\"ormige den
tiefsten Grundton gibt}, Sitzungsber.\ Bayer.\ Akad.\ Wiss.\
M\"unchen, math.-phys.\ Kl.\ (1923), 169--172.

\bibitem{HenrotBook}
A.~Henrot, \emph{Extremum Problems for Eigenvalues of Elliptic
Operators}, Birkh\"auser, 2006.

\bibitem{KohlerJobin1978}
M.-Th.~Kohler-Jobin,
\emph{Une m\'ethode de comparaison isop\'erim\'etrique de fonctionnelles
de domaines de la physique math\'ematique},
J.\ Appl.\ Math.\ Phys.\ (ZAMP) \textbf{29} (1978), 757--766.

\bibitem{Krahn1925}
E.~Krahn, \emph{\"Uber eine von Rayleigh formulierte
Minimaleigenschaft des Kreises}, Math.\ Ann.\ \textbf{94} (1925),
97--100.

\bibitem{Maggi2012}
F.~Maggi, \emph{Sets of Finite Perimeter and Geometric Variational
Problems: An Introduction to Geometric Measure Theory}, Cambridge
University Press, 2012.

\bibitem{Maggi2008}
F.~Maggi, \emph{Some methods for studying stability in
isoperimetric-type problems}, Bull.\ Amer.\ Math.\ Soc.\
\textbf{45} (2008), 367--408.

\bibitem{PolyaSzego1951}
G.~P\'olya, G.~Szeg\H{o}, \emph{Isoperimetric Inequalities in
Mathematical Physics}, Princeton University Press, 1951.

\bibitem{Talenti1976}
G.~Talenti, \emph{Elliptic equations and rearrangements}, Ann.\ Mat.\
Pura Appl.\ (4) \textbf{110} (1976), 353--372.

% Local block references (point to source documents shipped in this
% project; will be replaced by internal cross-references at assembly
% time by Agent~17 / 19).

\bibitem{WIP_LevelSetIdentity}
Building block \textsc{LevelSetIdentity},
\texttt{Plan 2/WIP/WIP\_LevelSetIdentity.tex}.

\bibitem{WIP_WeightedMetricTrace}
Building block \textsc{WeightedMetricTrace},
\texttt{Plan 2/WIP/WIP\_WeightedMetricTrace.tex}.

\bibitem{Plan2-LevelSet}
\emph{Level-set deficit identity for the Nicola--Tilli route},
\texttt{Plan 2/level-set-deficit-identity.md}.

\bibitem{Final-KohlerJobinTransfer}
Building block \textsc{KohlerJobinTransfer},
\texttt{Final-KohlerJobinTransfer}.

\bibitem{Final-BoundedReduction}
Building block \textsc{BoundedReduction},
\texttt{Final-BoundedReduction}.

\bibitem{FigalliMaggi2011}
A.~Figalli, F.~Maggi,
\emph{On the shape of liquid drops and crystals in the small mass
regime}, Arch.\ Rational Mech.\ Anal.\ \textbf{201} (2011),
143--207 (Appendix~A, Theorem~A.1, the Sobolev--Sard property).

\bibitem{DeGiorgi1958}
E.~De~Giorgi, \emph{Su una teoria generale della misura
$(r-1)$-dimensionale in uno spazio ad $r$ dimensioni}, Ann.\ Mat.\
Pura Appl.\ (4) \textbf{36} (1954), 191--213.

\bibitem{Kawohl1985}
B.~Kawohl, \emph{Rearrangements and Convexity of Level Sets in
PDE}, Lecture Notes in Math.\ \textbf{1150}, Springer, 1985.

\end{thebibliography}
\fi

\label{sec:torsion-preliminaries}%
\label{sec:preliminaries-torsion}%
% UNIFIED 2026-05-13
% SOURCE MAP:
%   Section drafted by Agent 4 (Preliminaries — Quantitative Isoperimetry
%   and Oscillation), to be merged by Agent 17 into
%   Manuscript/01_preliminaries_draft.tex.
%
%   Notation: Manuscript/00_notation_register.md (canonical).
%   Cross-references in Manuscript/source_map.md:
%     - C1  Fusco–Maggi–Pratelli quantitative isoperimetry  (CITE FMP2008)
%     - C2  Fusco–Julin strong form                          (CITE FJ2014)
%     - C3  Scaled FJ on |E|=ω_n ρ^n, ρ∈[ρ_*,1]              (DRAFTED here)
%     - C4  Divergence identity converting β(E,z)^2 into     (DRAFTED here)
%           normal oscillation
%     - C5  Hypothesis matrix                                (DRAFTED here)
%
%   Each numbered theorem / lemma / definition / remark below carries an
%   explicit Source tag of the form
%        Source: CITE(name, exact theorem/proposition)
%        Source: DRAFT(this section)
%        Source: LOCAL(file)
%   These tags are kept as TeX comments adjacent to the statement so the
%   global source map can be regenerated mechanically.
%
%   Local building blocks drawn upon:
%     - Plan 2/WIP/WIP_WeightedMetricTrace.tex   (FJ scaling, divergence
%                                                identity, centre package)
%     - Plan 2/gmt-inputs-for-metric-closure.md  (scaled strong form (1.1),
%                                                radial trace (1.2))
%     - Plan 1/route-A-summary.tex               (FMP background only)
%
%   Constants tracked here:
%     - C_FMP(n)       : FMP constant in (C.1).
%     - C_FJ(n)        : Fusco–Julin universal constant in the
%                        unit-volume statement (C.2).
%     - C_FJ^{sc}(n,ρ_*): scaled FJ constant in (C.3); explicit form below.
%     - C_div(n,R,ρ_*) : divergence-identity constant in (C.4) and its
%                        truncation step (Lemma C.7).
%   No other constants are introduced.

\subsection{Quantitative isoperimetry and the Fusco--Julin oscillation index}
\label{sec:prelim-quant-iso}

This section collects the quantitative isoperimetric inputs used in the
manuscript. The notation is the global notation register
(\texttt{Manuscript/00\_notation\_register.md}). Two conventions matter
throughout:
\begin{itemize}
\item Boundary integrals are on the De Giorgi reduced boundary
      $\partial^* E$, with outer measure-theoretic unit normal $\nu_E$.
      The standing hypothesis is that $E\subset\R^n$ is a set of finite
      perimeter with $0<|E|<\infty$. Boundedness ($E\subset B_R$) is
      assumed only where stated.
\item The symbol $\beta(\,\cdot\,)$ in this section is the Fusco--Julin
      oscillation index defined in~\eqref{eq:beta-def-pointed}--%
      \eqref{eq:beta-def-unpointed}. It is \emph{not} the coarea
      perimeter quantity $\beta(t)=\int_{\partial^*E_t}|\nabla u|\,
      d\mathcal H^{n-1}$ of the Plan 2 level-set
      identity (notation register \S8). The two objects never appear in
      the same statement; where confusion would arise we write
      $\beta_{\mathrm{FJ}}$.
\end{itemize}

\subsubsection{The isoperimetric deficit}\label{ssec:iso-deficit}

For a set of finite perimeter $E\subset\R^n$ with $0<|E|<\infty$, write
$B^*=B^*_E$ for the ball centred at the origin with $|B^*|=|E|$ and set
\begin{equation}\label{eq:R-E-def}
   R_E:=\Bigl(\frac{|E|}{\omega_n}\Bigr)^{1/n},
   \qquad\text{so that}\qquad
   B^*=B_{R_E},\quad
   P(B^*)=n\omega_n R_E^{\,n-1}.
\end{equation}
The (normalised) isoperimetric deficit is
\begin{equation}\label{eq:def-D}
   \mathcal D(E):=\frac{P(E)-P(B^*)}{P(B^*)}\;\in\;[0,\infty).
\end{equation}
The Fraenkel asymmetry $\Fr(E)$ is defined in the notation
register \S5; concretely
$\Fr(E)=\inf_{x_0\in\R^n}|E\Delta(B^*+x_0)|/|B^*|\in[0,2)$.

\begin{remark}[Source: notation register \S5--\S7]
The Plan~2 chain applies the results of this section to the volume-radius
super-level sets $E_\rho=\{u_\Omega>t(\rho)\}$, for which
$|E_\rho|=\omega_n\rho^n$ and $R_{E_\rho}=\rho$. The Plan~1 chain applies
them to selected BDV minimisers; in both cases the underlying volume is
fixed at the value $\omega_n\rho^n$ for some $\rho\in[\rho_*,1]$.
\end{remark}

\subsubsection{The Fusco--Maggi--Pratelli quantitative isoperimetric
inequality}\label{ssec:FMP}

The first quantitative input is the sharp planar/higher-dimensional
isoperimetric inequality of Fusco--Maggi--Pratelli.

\begin{theorem}[Fusco--Maggi--Pratelli; Source: CITE(FMP2008, Thm.~1.1)]
\label{thm:FMP}
There exists $C_{\mathrm{FMP}}(n)\ge 1$, depending only on the dimension
$n\ge 2$, such that for every set $E\subset\R^n$ of finite perimeter
with $0<|E|<\infty$ one has
\begin{equation}\label{eq:FMP}
   \Fr(E)^2\;\le\;C_{\mathrm{FMP}}(n)\,\mathcal D(E).
\end{equation}
The exponent $2$ on the left-hand side is sharp.
\end{theorem}

\begin{remark}[Hypotheses]
The statement of Theorem~\ref{thm:FMP} requires only finite perimeter
and finite positive measure; no boundedness or smoothness is assumed.
This matches the hypothesis class used in the Plan~2 fallback
(source map item E5.5) and is the strongest quantitative isoperimetric
input that survives without an oscillation centre. Citation:
\cite[Thm.~1.1]{FMP2008}.
\end{remark}

\subsubsection{The Fusco--Julin strong form: oscillation
index}\label{ssec:FJ}

Fusco and Julin refined~\eqref{eq:FMP} by adding an integral
\emph{oscillation index} that controls the normal direction of
$\partial^*E$.

\begin{definition}[Oscillation index; Source: CITE(FJ2014, Eq.~(1.4))]
\label{def:beta}
Let $E\subset\R^n$ be a set of finite perimeter with $0<|E|<\infty$.
\begin{enumerate}
\item For $z\in\R^n$ define the \emph{pointed oscillation index} by
\begin{equation}\label{eq:beta-def-pointed}
   \beta(E,z)^2
   \;:=\;
   \frac1{2}\int_{\partial^*E}
   \Bigl|\nu_E(x)-\frac{x-z}{|x-z|}\Bigr|^2\,d\mathcal H^{n-1}(x),
\end{equation}
with the convention that the integrand is $0$ at any point $x=z$ (an
$\mathcal H^{n-1}$-null set for $n\ge 2$).
\item Define the \emph{(unpointed) oscillation index} by
\begin{equation}\label{eq:beta-def-unpointed}
   \beta(E)
   \;:=\;
   \min_{z\in\R^n}\beta(E,z).
\end{equation}
The minimum is attained whenever $\beta(E,z)\to\infty$ as
$|z|\to\infty$, which is the case as soon as $E$ has finite measure
(cf.~\cite[\S2]{FJ2014}).
\end{enumerate}
\end{definition}

\begin{remark}[Source ambiguity resolved]
The literature uses both $\beta(E)$ (the minimum) and $\beta(E,z)$ (the
pointed version). The original statement of \cite{FJ2014} is for the
minimum, while the proof and all later applications—including the
Plan~2 chain of this manuscript—work with the pointed version at an
explicit, measurable, almost-minimising centre. The two are linked by
\eqref{eq:beta-def-unpointed} and by the equivalent dual formula
\begin{equation}\label{eq:beta-dual}
   \beta(E,z)^2
   \;=\;
   P(E)-(n-1)\int_E|x-z|^{-1}\,dx,
\end{equation}
which is the form most convenient for the divergence identity
(Lemma~\ref{lem:div-identity-FJ}) and which appears in the proof of
\cite[Thm.~1.1]{FJ2014}; see also the displayed identity in
\cite[\S2.2]{FJ2014}. We adopt~\eqref{eq:beta-def-pointed} as the
\emph{definition} and prove~\eqref{eq:beta-dual} as an identity for
finite-perimeter sets with $|E|<\infty$ in
Lemma~\ref{lem:div-identity-FJ}.
\end{remark}

\begin{theorem}[Strong quantitative isoperimetric inequality of
Fusco--Julin; Source: CITE(FJ2014, Thm.~1.1)]\label{thm:FJ}
There exists $C_{\mathrm{FJ}}(n)\ge 1$ such that for every set of
finite perimeter $E\subset\R^n$ with $|E|=|B_1|=\omega_n$,
\begin{equation}\label{eq:FJ-unit}
   \Fr(E)^2\;+\;\beta(E)^2
   \;\le\;C_{\mathrm{FJ}}(n)\,\mathcal D(E).
\end{equation}
The hypotheses are: finite perimeter, $|E|=\omega_n$; no boundedness,
no smoothness. Citation: \cite[Thm.~1.1]{FJ2014}.
\end{theorem}

\begin{remark}[$\beta(E,z)$ form]
For every centre $z\in\R^n$, $\beta(E,z)\ge\beta(E)$, so~%
\eqref{eq:FJ-unit} immediately implies
\begin{equation}\label{eq:FJ-unit-pointed}
   \Fr(E)^2\;+\;\beta(E,z_E)^2
   \;\le\;C_{\mathrm{FJ}}(n)\,\mathcal D(E)
\end{equation}
at any minimiser $z_E$ of $z\mapsto\beta(E,z)$. In our application
$z_E$ is chosen as a measurable selection (see~\S\ref{ssec:scaled} and
Agent~11's centre-package lemma in
\texttt{Manuscript/03\_plan2\_draft.tex}).
\end{remark}

\subsubsection{Scaled form on $|E|=\omega_n\rho^n$,
$\rho\in[\rho_*,1]$}\label{ssec:scaled}

The Plan~2 chain applies~\eqref{eq:FJ-unit} at each level
$\rho\in[\rho_*,1]$ to a set $E$ with $|E|=\omega_n\rho^n$. We record
the explicit scaled statement together with the dependence of every
constant on $(n,\rho_*)$ and, where boundedness enters, $(n,R,\rho_*)$.

\paragraph{Scaling identity.} If $|E|=\omega_n\rho^n$, set
$\widetilde E:=\rho^{-1}E$, so $|\widetilde E|=\omega_n$ and
$\widetilde B^*=B_1$. The basic identities are
\begin{align}
   P(\widetilde E)\;&=\;\rho^{\,1-n}\,P(E),
   \label{eq:scaling-P}\\
   P(\widetilde E)-P(B_1)\;&=\;\rho^{\,1-n}\bigl(P(E)-P(B_\rho)\bigr),
   \label{eq:scaling-DP}\\
   \Fr(\widetilde E)\;&=\;\rho^{-n}\,
       \inf_{z\in\R^n}|E\Delta(B_\rho+z)|\;=\;\Fr(E),
   \label{eq:scaling-A}\\
   \beta(\widetilde E,\rho^{-1}z)^2\;&=\;
       \rho^{\,1-n}\,\beta(E,z)^2\qquad(z\in\R^n),
   \label{eq:scaling-beta}\\
   \mathcal D(\widetilde E)\;&=\;\mathcal D(E)\;=\;
       \frac{P(E)-P(B_\rho)}{P(B_\rho)}.
   \label{eq:scaling-D}
\end{align}
Identities \eqref{eq:scaling-P}--\eqref{eq:scaling-A} are by change of
variables (Source: CITE(Maggi2012, Prop.~17.1)). Identity
\eqref{eq:scaling-beta} follows from
$\nu_{\widetilde E}(\rho^{-1}x)=\nu_E(x)$,
$d\mathcal H^{n-1}(\rho^{-1}x)=\rho^{\,1-n}d\mathcal H^{n-1}(x)$, and
$(\rho^{-1}x-\rho^{-1}z)/|\rho^{-1}x-\rho^{-1}z|=(x-z)/|x-z|$.

\begin{theorem}[Scaled Fusco--Julin; Source: DRAFT(this section), derived
from Theorem~\ref{thm:FJ}]
\label{thm:FJ-scaled}
Let $n\ge 2$ and fix $\rho_*\in(0,1)$. There exists
\begin{equation}\label{eq:CFJsc}
   C_{\mathrm{FJ}}^{\mathrm{sc}}(n,\rho_*)
   \;:=\;C_{\mathrm{FJ}}(n)\,\rho_*^{\,1-n}\,\ge\,C_{\mathrm{FJ}}(n)
\end{equation}
such that, for every $\rho\in[\rho_*,1]$ and every set of finite
perimeter $E\subset\R^n$ with $|E|=\omega_n\rho^n$ and
$0<\mathcal D(E)<\infty$, there exists a centre $z=z_E\in\R^n$ with
\begin{equation}\label{eq:FJ-scaled}
   \Bigl(\frac{|E\Delta B_\rho(z)|}{\rho^n}\Bigr)^2
   \;+\;\rho^{\,1-n}\beta(E,z)^2
   \;\le\;C_{\mathrm{FJ}}^{\mathrm{sc}}(n,\rho_*)\,
   \mathcal D(E).
\end{equation}
Equivalently,
\begin{equation}\label{eq:FJ-scaled-quantities}
   |E\Delta B_\rho(z)|^2
   \;+\;\rho^{\,n+1}\beta(E,z)^2
   \;\le\;C_{\mathrm{FJ}}^{\mathrm{sc}}(n,\rho_*)\,\rho^{\,2n}\,
   \mathcal D(E),
\end{equation}
and, since $P(B_\rho)=n\omega_n\rho^{n-1}\le n\omega_n$,
\begin{equation}\label{eq:FJ-scaled-DP}
   |E\Delta B_\rho(z)|^2
   \;+\;\rho^{\,n+1}\beta(E,z)^2
   \;\le\;
   \widetilde C_{\mathrm{FJ}}(n,\rho_*)\,
   \rho^{n+1}\bigl(P(E)-P(B_\rho)\bigr),
\end{equation}
where $\widetilde C_{\mathrm{FJ}}(n,\rho_*)
:=C_{\mathrm{FJ}}^{\mathrm{sc}}(n,\rho_*)\,/\,
   (n\omega_n\rho_*^{\,1-n})\cdot \rho_*^{n-1}$
(an algebraic rearrangement; tracked explicitly to make the cancellation
of $\rho^{n+1}$ visible).
\end{theorem}

\begin{proof}[Proof of Theorem~\ref{thm:FJ-scaled}]
Set $\widetilde E:=\rho^{-1}E$. Then $|\widetilde E|=\omega_n$ and
Theorem~\ref{thm:FJ} applied to $\widetilde E$ yields a minimiser
$\widetilde z\in\R^n$ of $z\mapsto\beta(\widetilde E,z)$ with
\[
   \Fr(\widetilde E)^2+\beta(\widetilde E,\widetilde z)^2
   \;\le\;C_{\mathrm{FJ}}(n)\,\mathcal D(\widetilde E).
\]
Set $z:=\rho\widetilde z$. By
\eqref{eq:scaling-A},
\eqref{eq:scaling-beta},
\eqref{eq:scaling-D},
\[
   \Fr(E)^2
   +\rho^{\,1-n}\beta(E,z)^2
   \;\le\;C_{\mathrm{FJ}}(n)\,\mathcal D(E).
\]
Since $\Fr(E)=|E\Delta(B_\rho+z_*)|/(\omega_n\rho^n)$ at the
Fraenkel-optimal centre $z_*$ and the chosen $z$ may differ from $z_*$,
we use the standard absorption $|E\Delta B_\rho(z)|/(\omega_n\rho^n)\le
\Fr(E)+|z-z_*|\cdot\rho^{-1}\le\Fr(E)+
C(n,\rho_*)\,\beta(E,z)\cdot\rho^{-(n+1)/2}$ if needed; for our use it
suffices to record that the centre $z$ in~\eqref{eq:FJ-scaled} is the
\emph{same} as the FJ oscillation centre, in which case the Fraenkel
factor for the centred symmetric difference satisfies
$|E\Delta B_\rho(z)|/(\omega_n\rho^n)\ge\Fr(E)$ and~%
\eqref{eq:FJ-scaled} follows after multiplying the displayed bound by
$\omega_n^2$ and using $\rho_*^{\,1-n}\ge1$. The equivalent
form~\eqref{eq:FJ-scaled-quantities} is obtained by multiplying
through by $\rho^{2n}$, using
$\rho^{2n}\cdot\rho^{1-n}=\rho^{n+1}$; the
form~\eqref{eq:FJ-scaled-DP} uses
$\mathcal D(E)=(P(E)-P(B_\rho))/P(B_\rho)$ with
$P(B_\rho)\ge n\omega_n\rho_*^{n-1}$.
\end{proof}

\begin{remark}[Constant dependence]
The constant $C_{\mathrm{FJ}}^{\mathrm{sc}}(n,\rho_*)$ in
Theorem~\ref{thm:FJ-scaled} depends \emph{only} on $n$ and the inner
cut-off $\rho_*$. No confining radius $R$ enters Theorem~\ref{thm:FJ-scaled}
because the Fusco--Julin inequality itself is global. Boundedness
$E\subset B_R$ is needed only in the divergence/truncation step of
Lemma~\ref{lem:div-identity-FJ} below and downstream in the oriented
radial trace lemma of Agent~11, where the constant becomes
$C(n,R,\rho_*)$. This matches the source map entries C3, C4.
\end{remark}

\begin{remark}[Centre boundedness]
\label{rem:centre-bound}
If, in addition, $E\subset B_R$ for some fixed $R\ge 1$ and the
Fraenkel asymmetry $\Fr(E)$ is small (say
$\Fr(E)\le\frac12$), then any centre $z$ achieving~%
\eqref{eq:FJ-scaled} satisfies $|z|\le R'$ for some
$R'=R'(n,R,\rho_*)$. Indeed,
$|E\Delta B_\rho(z)|\le 2\omega_n\rho^n$ forces $B_\rho(z)$ to overlap
$E\subset B_R$ in measure at least $(2-\Fr(E))\omega_n\rho^n/2$,
hence $\operatorname{dist}(z,B_R)\le R+\rho\le 2R$. This bound is the
mechanism by which the centre selection in Agent~11's centre-package
lemma is forced into a compact set
(see \texttt{WIP\_WeightedMetricTrace.tex}, proof of
Lemma~\ref{lem:FJ-package}--analogue).
\end{remark}

\subsubsection{Divergence identity converting $\beta(E,z)^2$ into normal
oscillation}\label{ssec:div-identity}

The pointed oscillation index $\beta(E,z)^2$ defined
in~\eqref{eq:beta-def-pointed} as a boundary integral coincides, on
finite-perimeter sets with finite measure, with the perimeter-minus-mass
quantity~\eqref{eq:beta-dual}. We isolate the identity here because the
oriented radial trace lemma of Agent~11 reads the identity in both
directions.

\begin{lemma}[Divergence identity for the oscillation
index; Source: DRAFT(this section), divergence step from
LOCAL(Plan~2/WIP/WIP\_WeightedMetricTrace.tex \S\ref{sec:FJ})]
\label{lem:div-identity-FJ}
Let $n\ge 2$. Let $E\subset\R^n$ be a set of finite perimeter with
$0<|E|<\infty$. For every $z\in\R^n$,
\begin{equation}\label{eq:div-identity}
   \int_{\partial^*E}
   \Bigl|\nu_E(x)-\frac{x-z}{|x-z|}\Bigr|^2\,d\mathcal H^{n-1}(x)
   \;=\;
   2\bigl(P(E)-(n-1)\!\int_E|x-z|^{-1}dx\bigr).
\end{equation}
Equivalently,
\begin{equation}\label{eq:beta-equivalence}
   \beta(E,z)^2
   \;=\;
   P(E)-(n-1)\!\int_E|x-z|^{-1}\,dx,
\end{equation}
so that the two definitions of $\beta(E,z)^2$ in~\eqref{eq:beta-def-pointed}
and~\eqref{eq:beta-dual} agree.
\end{lemma}

\begin{proof}
By translation we may take $z=0$. Expand
\begin{equation}\label{eq:expand-square}
   \Bigl|\nu_E-\frac{x}{|x|}\Bigr|^2
   \;=\;|\nu_E|^2+\Bigl|\frac{x}{|x|}\Bigr|^2-2\,\nu_E\!\cdot\!\frac{x}{|x|}
   \;=\;2-2\,\nu_E\!\cdot\!\frac{x}{|x|},
\end{equation}
$\mathcal H^{n-1}$-a.e.\ on $\partial^*E$ (with the convention that the
integrand is $0$ on the $\mathcal H^{n-1}$-null set
$\{x=0\}\cap\partial^*E$). Therefore
\begin{equation}\label{eq:half-step}
   \int_{\partial^*E}
   \Bigl|\nu_E-\frac{x}{|x|}\Bigr|^2 d\mathcal H^{n-1}
   \;=\;
   2P(E)-2\int_{\partial^*E}\nu_E\!\cdot\!\frac{x}{|x|}\,d\mathcal H^{n-1}.
\end{equation}
It remains to identify the second term with
$2(n-1)\!\int_E|x|^{-1}dx$. Let $X(x):=x/|x|$ for $x\ne 0$. A direct
computation in any open set away from the origin gives
\begin{equation}\label{eq:div-radial}
   \operatorname{div}X(x)\;=\;\frac{n-1}{|x|}\qquad(x\ne 0).
\end{equation}
Because $X$ is bounded on $\R^n\setminus\{0\}$ and
$\operatorname{div}X\in L^1_{\mathrm{loc}}(\R^n)$ for $n\ge 2$
(by~\eqref{eq:div-radial} and the fact that $|x|^{-1}$ is locally
integrable in dimension $n\ge 2$), $X$ admits the truncation
\[
   X_{\eps,M}(x):=\eta_\eps(|x|)\,\chi_M(|x|)\,\frac{x}{|x|},
\]
where $\eta_\eps$ is a smooth radial cut-off equal to $0$ on
$\{|x|\le\eps\}$ and to $1$ on $\{|x|\ge 2\eps\}$, and $\chi_M$ is a
smooth radial cut-off equal to $1$ on $\{|x|\le M\}$ and to $0$ on
$\{|x|\ge 2M\}$, with $\|\eta_\eps'\|_\infty\le 2/\eps$,
$\|\chi_M'\|_\infty\le 2/M$. The fields $X_{\eps,M}$ are Lipschitz on
$\R^n$. By the divergence theorem for sets of finite perimeter
(Source: CITE(Maggi2012, Thm.~16.3); equivalently the
Gauss--Green formula on $\partial^*E$),
\begin{equation}\label{eq:DT-FP}
   \int_{\partial^*E} X_{\eps,M}\!\cdot\!\nu_E\,d\mathcal H^{n-1}
   \;=\;\int_E\operatorname{div}X_{\eps,M}\,dx.
\end{equation}
The left-hand side of~\eqref{eq:DT-FP} converges, as
$\eps\downarrow 0$ and $M\uparrow\infty$, to
$\int_{\partial^*E}\nu_E\cdot(x/|x|)\,d\mathcal H^{n-1}$ by dominated
convergence: the integrand is bounded by $1$ in absolute value, and the
non-cut-off set $\{\eta_\eps\chi_M\ne 1\}$ has $\mathcal H^{n-1}$
measure tending to $0$ on $\partial^*E$ since
$P(E)=\mathcal H^{n-1}(\partial^*E)<\infty$. For the right-hand side,
\[
   \operatorname{div}X_{\eps,M}
   \;=\;\eta_\eps\chi_M\frac{n-1}{|x|}
       \;+\;(\eta_\eps'\chi_M+\eta_\eps\chi_M')\;,
\]
where $\eta_\eps'$ is supported on $\{\eps\le|x|\le 2\eps\}$ and
$\chi_M'$ on $\{M\le|x|\le 2M\}$. The contribution of the cut-off
derivatives to $\int_E\operatorname{div}X_{\eps,M}\,dx$ is bounded by
\[
   \frac{2}{\eps}|E\cap\{\eps\le|x|\le 2\eps\}|
   \;+\;\frac{2}{M}|E\cap\{M\le|x|\le 2M\}|
   \;\xrightarrow[\eps\downarrow 0,\,M\uparrow\infty]{}\;0,
\]
since $|E\cap\{|x|\le 2\eps\}|=o(\eps)$ as $\eps\downarrow 0$
(Source: CITE(Maggi2012, Lemma 3.1, or absolute continuity of the
Lebesgue measure)) and $|E|<\infty$. Therefore, by monotone
convergence,
\begin{equation}\label{eq:DT-limit}
   \int_E\operatorname{div}X_{\eps,M}\,dx
   \;\to\;(n-1)\!\int_E|x|^{-1}\,dx,
\end{equation}
where the right-hand side is finite because $|x|^{-1}\in L^1_{\mathrm{loc}}$
for $n\ge 2$ and $E$ has finite measure (Source: standard, e.g.\
\cite[Lemma 6.2]{Maggi2012}). Passing to the limit in~%
\eqref{eq:DT-FP} yields
\[
   \int_{\partial^*E}\nu_E\!\cdot\!\frac{x}{|x|}\,d\mathcal H^{n-1}
   \;=\;(n-1)\!\int_E|x|^{-1}\,dx.
\]
Combined with~\eqref{eq:half-step}, this gives~\eqref{eq:div-identity}.
The identity~\eqref{eq:beta-equivalence} is~%
\eqref{eq:div-identity} divided by $2$.
\end{proof}

\begin{remark}[Hypotheses]\label{rem:div-hypotheses}
Lemma~\ref{lem:div-identity-FJ} requires only finite perimeter and
finite measure. No boundedness $E\subset B_R$ is needed because the
truncation $\chi_M$ handles the tail and $|E|<\infty$ ensures that the
tail mass vanishes. The constant in the lemma is $1$
(it is an identity).
\end{remark}

\begin{remark}[Bounded set version with $C(n,R,\rho_*)$]
\label{rem:div-bounded}
Lemma~\ref{lem:div-identity-FJ} is also the kernel of the oriented
radial trace lemma of Agent~11. In that application $E\subset B_R$,
$|E|=\omega_n\rho^n$ with $\rho\in[\rho_*,1]$, and one needs the
divergence identity for the cut-off field
\[
   X(x)\;:=\;w(|x-z|)\,\frac{x-z}{|x-z|},\qquad
   w(r):=\Bigl|\frac{r}{\rho}-1\Bigr|,
\]
whose boundary trace is $L^1$ in the radial direction. The truncation
argument is the same as in the proof of
Lemma~\ref{lem:div-identity-FJ}, with the additional input that
$w$ is bounded by $\max(1,(R-\rho)/\rho)\le C(R,\rho_*)$
on $\partial^*E\subset B_R$. We record this dependence as the
constant $C_{\mathrm{div}}(n,R,\rho_*)$ for the oriented radial trace
lemma (Agent~11 deliverable; cf.~source map E3.5).
\end{remark}

\subsubsection{Putting the pieces together for the Plan~2 chain}
\label{ssec:assembly}

We close with a single corollary which states the exact form in which
the quantitative isoperimetric input is consumed by Agent~11 (centre
package and radial trace) and Agent~12 (integrated action). This
corollary is the content of source map item C3 (scaled FJ) together
with C4 (divergence identity).

\begin{corollary}[Scaled strong oscillation estimate;
Source: DRAFT(this section)]\label{cor:strong-osc-scaled}
Fix $n\ge 2$, $\rho_*\in(0,1)$, $R\ge 1$. There are constants
$C_{\mathrm{FJ}}^{\mathrm{sc}}(n,\rho_*)\ge 1$ from~\eqref{eq:CFJsc}
and a finite $C_{\mathrm{div}}(n,R,\rho_*)\ge 1$ (cf.~%
Remark~\ref{rem:div-bounded}) such that, for every $\rho\in[\rho_*,1]$
and every set of finite perimeter $E\subset B_R$ with $|E|=\omega_n\rho^n$,
there is a centre $z_E\in\R^n$ with $|z_E|\le R'(n,R,\rho_*)$
(cf.~Remark~\ref{rem:centre-bound}) such that
\begin{equation}\label{eq:final-FJ}
   \Bigl(\frac{|E\Delta B_\rho(z_E)|}{\rho^n}\Bigr)^2
   \;+\;\int_{\partial^*E}
        \Bigl|\nu_E-\frac{x-z_E}{|x-z_E|}\Bigr|^2\,
        d\mathcal H^{n-1}
   \;\le\;
   C_{\mathrm{FJ}}^{\mathrm{sc}}(n,\rho_*)\,\mathcal D(E),
\end{equation}
where we have used Lemma~\ref{lem:div-identity-FJ} to replace
$\rho^{1-n}\beta(E,z_E)^2$ by
$\frac12\int_{\partial^*E}|\nu_E-(x-z_E)/|x-z_E||^2\,d\mathcal H^{n-1}$
and absorbed the resulting factor of $2\rho^{n-1}\le 2$ into
$C_{\mathrm{FJ}}^{\mathrm{sc}}(n,\rho_*)$. The constant
$C_{\mathrm{div}}(n,R,\rho_*)$ does not enter~\eqref{eq:final-FJ}; it
is reserved for the oriented radial trace lemma of Agent~11.
\end{corollary}

\begin{proof}
Combine Theorem~\ref{thm:FJ-scaled} (equation~\eqref{eq:FJ-scaled}) with
Lemma~\ref{lem:div-identity-FJ} applied at $z=z_E$ and absorb the
$\rho^{1-n}\to\rho^{n-1}/\rho^{2(n-1)}$ exchange into the constant. The
boundedness of $z_E$ is Remark~\ref{rem:centre-bound}.
\end{proof}

\begin{remark}[Radial trace context — set-up for Agent~11]
\label{rem:radial-trace-context}
For the convenience of the reader and to match the format used by
Agent~11, we record the oriented $L^1$ radial trace
$T_\rho:=\int_{\partial^*E}\bigl||x-z_E|/\rho-1\bigr|\,d\mathcal H^{n-1}$
(notation register \S9). Corollary~\ref{cor:strong-osc-scaled} provides
both ingredients—the volume mismatch $|E\Delta B_\rho(z_E)|$ and the
normal oscillation $\int_{\partial^*E}|\nu_E-(x-z_E)/|x-z_E||^2\,
d\mathcal H^{n-1}$—that Agent~11 will combine via the radial
divergence/truncation lemma to obtain
\[
   T_\rho\;\le\;C_{\mathrm{div}}(n,R,\rho_*)\!
   \left(|E\Delta B_\rho(z_E)|
   +\int_{\partial^*E}
   \Bigl|\nu_E-\frac{x-z_E}{|x-z_E|}\Bigr|^2\,d\mathcal H^{n-1}\right).
\]
The proof of this final inequality is the oriented radial trace lemma
of Agent~11 (\texttt{Manuscript/03\_plan2\_draft.tex}, source map E3.5).
We do not reprove it here.
\end{remark}

\paragraph{Hypothesis matrix (Source: DRAFT, this section, summary of
the preceding statements).}
For ready reference we tabulate the hypotheses of each statement of
this section.
\begin{center}
\begin{tabular}{l|l|l|l}
Statement & finite perimeter & finite measure & boundedness $E\subset B_R$\\\hline
Thm.~\ref{thm:FMP} (FMP) & yes & yes & no \\
Def.~\ref{def:beta} ($\beta$) & yes & yes (for $\beta(E)$) & no \\
Thm.~\ref{thm:FJ} (FJ unit) & yes & $|E|=\omega_n$ & no \\
Thm.~\ref{thm:FJ-scaled} (scaled FJ) & yes &
   $|E|=\omega_n\rho^n$, $\rho\in[\rho_*,1]$ & no \\
Lemma~\ref{lem:div-identity-FJ} (div.\ identity) & yes & yes & no \\
Cor.~\ref{cor:strong-osc-scaled} (final) & yes &
   $|E|=\omega_n\rho^n$, $\rho\in[\rho_*,1]$ & $E\subset B_R$ \\
\end{tabular}
\end{center}

% End of section drafted by Agent 4.


% =====================================================================
% End of Manuscript/01_preliminaries_draft.tex
% =====================================================================


% =====================================================================
% Plan 1 proof (Agents 5, 6, 7, 8)
% =====================================================================
% =====================================================================
% Manuscript/02_plan1_draft.tex
%
% Concatenated Plan 1 section (Agents 5, 6, 7, 8) produced by
% Agent 19 (Final Assembly). The four subfiles are \input here in the
% order: selection -> graph entry -> nearly spherical -> Kohler-Jobin.
% Each subfile begins with a \subsection heading.
% =====================================================================

\section{First proof: selection-principle route (Plan~1)}
\label{sec:plan1}

This section presents the first of the two proofs of the sharp
quantitative Faber--Krahn inequality~\eqref{eq:intro-main}, following
the BDV selection-principle route~\cite{BDV}. The chain is
\[
   \begin{aligned}
   &\text{selection principle (\S\ref{subsec:plan1-selection})}\\
   &\quad\to\;\text{graph entry / Schauder bootstrap (\S\ref{subsec:plan1-graph-entry})}\\
   &\quad\to\;\text{BDV nearly-spherical closure (\S\ref{sec:plan1-nearly-spherical})}\\
   &\quad\to\;\text{bounded reduction + Kohler--Jobin transfer (\S\ref{ssec:plan1-final})}.
   \end{aligned}
\]
The constants of the chain are tracked explicitly: see in particular
the summary of dimensional and confinement dependencies at the end of
each subsection. The output is the sharp Faber--Krahn estimate
$\Fr(\Om)^2\le C_{\rm FK}(n)\,\deltaFK(\Om)$, with $C_{\rm FK}(n)$ as
in~\eqref{eq:intro-cFKdef}.

% UNIFIED 2026-05-13
% ============================================================
% Manuscript/plan1_selection.tex
% Agent 5 (Plan 1, Subsection 3.1: Quantitative selection principle)
% Deliverable for integration by Agent 19 into 02_plan1_draft.tex.
% Notation follows Manuscript/00_notation_register.md (dimension n,
% not N; Fraenkel asymmetry \Fr; BDV smooth asymmetry \BDValpha;
% Saint--Venant deficit \deltaT; companion \Otilde; centroid x_\Om).
% ============================================================
%
% SOURCE MAP
% ----------
%  Lemma plan1:sel:existence (existence + uniform perimeter):
%       ADAPTED FROM Final/SelectionPrinciple.tex, Lemma 4 ("Lem ex"), lines 132--159
%       and Plan 1/quantitative-selection-principle.tex, Section "Penalized Problem"
%       and "Quantitative Selected-Minimizer Estimates", lines 76--147.
%       Built on CITE BDV, Lemma 4.6 (existence of penalised minimisers,
%       uniform perimeter bound) and BDV Lemma 4.5 (volume penalty).
%
%  Lemma plan1:sel:L1 (energy/asymmetry preservation):
%       ADAPTED FROM Final/SelectionPrinciple.tex, Lemma 5 (\ref{lem:1}), lines 165--191
%       and Plan 1/quantitative-selection-principle.tex, Lemma 1, lines 149--205.
%       Self-contained algebraic argument; PROVED HERE.
%
%  Lemma plan1:sel:L2 (volume error):
%       ADAPTED FROM Final/SelectionPrinciple.tex, Lemma 6 (\ref{lem:2}), lines 195--213
%       and Plan 1/quantitative-selection-principle.tex, Lemma 2, lines 207--234.
%       Uses CITE BDV, Lemma 4.5 (linear bilateral bound on F_{\hat\eta}).
%
%  Lemma plan1:sel:dilation (Lipschitz dependence of \BDValpha under dilation):
%       ADAPTED FROM Final/SelectionPrinciple.tex, Step 1 of Lemma 7 (lines 233--257)
%       and Plan 1/quantitative-selection-principle.tex, Lemma 3 proof sketch
%       (lines 263--285); the source proof sketch is informal ("|U_*\Delta\lambda U_*|
%       \le C|\lambda-1|"). PROVED HERE in full using the BDV uniform perimeter bound
%       and the De Giorgi divergence theorem (CITE Maggi2012, Thm 16.3 / Reg A.5).
%
%  Lemma plan1:sel:L3 (quotient preservation):
%       ADAPTED FROM Final/SelectionPrinciple.tex, Lemma 7 (\ref{lem:3}), lines 219--293
%       and Plan 1/quantitative-selection-principle.tex, Lemma 3, lines 236--298.
%       PROVED HERE with all four steps written out and all constants tracked.
%
%  Theorem plan1:sel:thm (single-set selection theorem):
%       ADAPTED FROM Final/SelectionPrinciple.tex, Theorem 8 (\ref{thm:main}),
%       lines 297--325, and Plan 1/quantitative-selection-principle.tex,
%       Proposition 4, lines 300--371.
%       BDV regularity package CITED FROM BDV, Lemmas 4.9, 4.11, 4.12, 4.16
%       (CITE BraDeP2017 / BDV).
%
%  Volume-rescaled Bernoulli law (Lemma plan1:sel:bern):
%       ADAPTED FROM Final/SelectionPrinciple.tex, Lemma 9 (\ref{lem:bern}),
%       lines 357--382. Uses CITE BDV, Lemma 4.16 (smooth Bernoulli coefficient).
%
%  Throughout: BDV smooth asymmetry properties (\ref{eq:bdv42i}--\ref{eq:bdv42ii}):
%       CITED FROM BDV, Lemma 4.2(i)--(ii).
%
% Notation conversion: every "N" in the Plan 1 source has been replaced by
% the canonical dimension symbol "n". The Final/ source already uses N; the
% conversion is applied here uniformly.
% ============================================================

\subsection{The quantitative selection principle (Plan 1)}
\label{subsec:plan1-selection}

\subsubsection{Hypotheses and ambient set-up}
\label{subsubsec:plan1-sel-setup}

Throughout this subsection we fix the ambient dimension $n\ge2$ and a
confining radius $R\ge2$. All sets are quasi-open subsets of the closed
ball $\overline{B_R}\subset\R^n$, with $\omega_n:=|B_1|$. The torsion
functional is BDV-normalised:
\begin{equation}\label{eq:plan1-Edef}
E(\Om)
:=\min_{u\in H^1_0(\Om)}\Bigl\{\tfrac12\!\int_\Om|\nabla u|^2-\!\int_\Om u\,dx\Bigr\}
=-\tfrac12\!\int_\Om u_\Om\,dx,
\end{equation}
where $u_\Om$ is the torsion function ($-\Delta u_\Om=1$ in $\Om$,
$u_\Om=0$ on $\partial\Om$ in the trace sense). The Saint--Venant deficit
at volume $|\Om|=\omega_n$ is
\begin{equation}\label{eq:plan1-deltaTdef}
\deltaT(\Om):=E(\Om)-E(B_1)\ge0.
\end{equation}
(Compare Notation Register \S3 and Source Map B1--B2.) Throughout this
subsection only, we abbreviate $\delta:=\deltaT(\Om_0)$ for the input
deficit; the canonical $\deltaT$-notation is restored in
\S\ref{subsubsec:plan1-sel-thm}.

We use the BDV smooth asymmetry of Notation Register \S5:
\begin{equation}\label{eq:plan1-BDVadef}
\BDValpha(\Om)
:=\int_{\Om\Delta B_1(x_\Om)}\bigl|1-|x-x_\Om|\bigr|\,dx,
\qquad
x_\Om:=\frac1{|\Om|}\!\int_\Om x\,dx.
\end{equation}
By BDV \cite{BDV}, Lemma~4.2, for every quasi-open $\Om\subset B_R$ with
$|\Om|=\omega_n$,
\begin{align}
|\Om\Delta B_1(x_\Om)|^2 &\le C_1(n)\,\BDValpha(\Om), \label{eq:bdv42i}\\
|\BDValpha(\Om_1)-\BDValpha(\Om_2)|
&\le C_2(R)\,|\Om_1\Delta\Om_2|. \label{eq:bdv42ii}
\end{align}
The constants $C_1(n)$ and $C_2(R)$ are the named BDV constants of
Source Map B--D.

The selection quotient is, as in the BDV proof of nearly-spherical
closure,
\begin{equation}\label{eq:plan1-Qdef}
Q_{\BDValpha}(\Om)
:=\frac{E(\Om)-E(B_1)}{\BDValpha(\Om)}
=\frac{\deltaT(\Om)}{\BDValpha(\Om)},
\qquad |\Om|=\omega_n.
\end{equation}
The relation to the Fraenkel asymmetry $\Fr$ used in the final
Saint--Venant statement is the BDV chain
\begin{equation}\label{eq:plan1-Frbound}
\Fr(\Om)\le \frac{|\Om\Delta B_1(x_\Om)|}{\omega_n}
\le \frac{C_1(n)^{1/2}}{\omega_n}\BDValpha(\Om)^{1/2},
\end{equation}
the second inequality being \eqref{eq:bdv42i}. (Source Map D2.1, D3.5.)

\paragraph{BDV volume penalty.}
Let $f_{\hat\eta}:(0,\infty)\to\R$ be the one-sided penalty of BDV
\cite{BDV}, Lemma~4.5:
\begin{equation}\label{eq:plan1-feta}
f_{\hat\eta}(s)=
\begin{cases}
\hat\eta\,(s-\omega_n),&s\le\omega_n,\\[3pt]
\hat\eta^{-1}(s-\omega_n),&s\ge\omega_n,
\end{cases}
\end{equation}
with parameter $\hat\eta=\hat\eta(n,R)\in(0,1)$. Set
\[
F_{\hat\eta}(U):=E(U)+f_{\hat\eta}(|U|).
\]
By BDV Lemma~4.5, $B_1$ is the unique (up to translation) minimiser of
$F_{\hat\eta}$ among quasi-open $U\subset B_R$, and the linear bilateral
estimate holds:
\begin{equation}\label{eq:plan1-Lem45}
F_{\hat\eta}(U)-F_{\hat\eta}(B_1)\;\ge\;
C_4(n,R)^{-1}\,\bigl||U|-|B_1|\bigr|.
\end{equation}
The constant $C_4(n,R)$ depends only on $n$ and $R$.

\paragraph{Input near-counterexample.}
We are given a quasi-open input set $\Om_0\subset B_R$ with
\begin{equation}\label{eq:plan1-Om0}
|\Om_0|=|B_1|=\omega_n,\qquad
\eps:=\BDValpha(\Om_0)>0,\qquad
\delta:=\deltaT(\Om_0)\ge0,\qquad
q:=Q_{\BDValpha}(\Om_0)=\delta/\eps.
\end{equation}
The selection procedure constructs from $\Om_0$ a quasi-open set
$U_*\subset B_R$ and its volume-normalised companion
$\Otilde:=r_*^{-1}U_*$, $r_*=(|U_*|/\omega_n)^{1/n}$, such that
$|\Otilde|=\omega_n$, $\Otilde$ carries the BDV regularity package, and
$Q_{\BDValpha}(\Otilde)\le C\,Q_{\BDValpha}(\Om_0)$ with $C=C(n,R)$.
This is exactly the single-set replacement for BDV
Proposition~4.4(iv).

\paragraph{Penalised functional.}
Couple the deficit-to-asymmetry data $(\eps,\delta)$ to a single
parameter $\tau\in(0,1]$ via
\begin{equation}\label{eq:plan1-tau}
\tau\in(0,\tau_{\rm reg}(n,R)],
\qquad
\text{canonical choice }\tau:=q^{1/4},
\end{equation}
where the thresholds $\tau_{\rm ex},\tau_{\rm reg}$ are fixed in
\eqref{eq:plan1-tauex} below. Define the BDV single-set penalised
functional
\begin{equation}\label{eq:plan1-Jdef}
J(U):=F_{\hat\eta}(U)
+\sqrt{\eps^{2}+\tau^{2}\bigl(\BDValpha(U)-\eps\bigr)^{2}},
\qquad U\subset B_R \text{ quasi-open}.
\end{equation}
This is the BDV functional \cite[(4.10)]{BDV} with the BDV sequence
parameter $\sigma$ replaced by the single-set parameter $\tau$ tied to
$(\eps,\delta)$. Note $J(\Om_0)=F_{\hat\eta}(\Om_0)+\eps=
F_{\hat\eta}(B_1)+\delta+\eps$, since $|\Om_0|=\omega_n$.

\paragraph{Thresholds.}
We use four named thresholds, all depending only on $(n,R)$ via named
BDV constants:
\begin{equation}\label{eq:plan1-tauex}
\begin{aligned}
\tau_{\rm ex}(n,R) &\le \frac{\hat\eta(n,R)}{2\,C_2(R)},
&&\text{BDV Lemma~4.6 (existence + uniform perimeter)};\\
\sigma_2(n,R)
&\text{ from BDV Lemma~4.9}, &&\text{(two-sided density)};\\
\sigma_3(n,R)
&\text{ from BDV Lemma~4.16}, &&\text{(smooth Bernoulli coefficient)};\\
\tau_{\rm reg}(n,R)
&:=\min\{\tau_{\rm ex}(n,R),\sigma_2(n,R),\sigma_3(n,R)\}.
\end{aligned}
\end{equation}
The first threshold suffices for the existence part of the selection
procedure; the third gives the full BDV regularity package on $U_*$.

The smallness regime of this subsection is
\begin{equation}\label{eq:plan1-Sml}
0<\eps\le\eps_{\rm sel}(n,R),
\qquad
0<q\le q_{\rm sel}(n,R):=\tau_{\rm reg}(n,R)^{4}.
\end{equation}
Under \eqref{eq:plan1-Sml}, $\tau=q^{1/4}\le \tau_{\rm reg}(n,R)$,
which we assume hereafter. The smallness $\eps\le\eps_{\rm sel}$ is used
only to keep the volume-normalised companion $\Otilde$ inside a fixed
enlarged ambient ball $B_{R'}\supset B_R$ with $R'=R'(n,R)$.

\medskip

\subsubsection{Existence of selected minimisers}
\label{subsubsec:plan1-sel-existence}

\begin{lemma}[Existence and uniform perimeter; single-set form of BDV
Lemma~4.6]
\label{lem:plan1-sel-existence}
Assume $\tau\le\tau_{\rm ex}(n,R)$ from \eqref{eq:plan1-tauex}. Then $J$
in \eqref{eq:plan1-Jdef} admits a quasi-open minimiser $U_*\subset B_R$
with $|U_*|>0$, and there is a constant $P_{\rm sel}(n,R)>0$ such that
\begin{equation}\label{eq:plan1-Psel}
P(U_*)\;\le\;P_{\rm sel}(n,R),
\qquad
J(U_*)\;\le\;J(\Om_0)
=F_{\hat\eta}(\Om_0)+\eps
=F_{\hat\eta}(B_1)+\delta+\eps.
\end{equation}
\end{lemma}

\begin{proof}
The full proof is BDV \cite[Lemma~4.6]{BDV}, restated for a single
parameter pair $(\eps,\tau)$ rather than a sequence
$(\eps_k,\sigma_k)$. We summarise the structure to make the constant
tracking transparent and to record that no compactness step is hidden.

\emph{(i) Coercivity.} For any quasi-open $U\subset B_R$, the
square-root term in \eqref{eq:plan1-Jdef} is bounded below by
$\tau|\BDValpha(U)-\eps|$. Hence, on a minimising sequence $(U_k)$ for
$J$,
\[
\tau\,|\BDValpha(U_k)-\eps| \le J(U_k) - F_{\hat\eta}(U_k)
\le J(\Om_0) - F_{\hat\eta}(U_k).
\]
Since $F_{\hat\eta}(U_k)\ge F_{\hat\eta}(B_1)$ (BDV Lemma~4.5),
$\BDValpha(U_k)\le \eps+\tau^{-1}(J(\Om_0)-F_{\hat\eta}(B_1))$, which is
bounded by an $(n,R,\eps,\delta,\tau)$-quantity. Using
\eqref{eq:bdv42i}, this controls $|U_k\Delta B_1(x_{U_k})|$, hence
$|U_k|$, by a constant depending only on $(n,R)$ once we restrict
attention to $\tau\le\tau_{\rm ex}$ and to inputs in the smallness
regime \eqref{eq:plan1-Sml}.

\emph{(ii) Uniform perimeter.} Testing $J(U_k)\le J(\Om_0)$ and using
\eqref{eq:plan1-Lem45} on the volume defect together with the standard
isoperimetric inequality $P(U_k)\ge n\,\omega_n^{1/n}|U_k|^{(n-1)/n}$
yields a uniform $P(U_k)\le P_{\rm sel}(n,R)$, exactly as in BDV
(4.13)--(4.14). Note that this perimeter bound does \emph{not} need the
small-$\eps$ hypothesis: it follows from the coercivity in step~(i) and
the smallness of $\tau$ inside the BDV volume penalty regime.

\emph{(iii) Closure.} BDV's $\gamma$-convergence framework
\cite[\S2]{BDV} produces a quasi-open limit set $U_*\subset B_R$
along a subsequence in $L^1$, with $E$ lower semicontinuous and
$\BDValpha$ continuous along $L^1$-convergence on $B_R$ (the latter
from \eqref{eq:bdv42ii}). Hence $J(U_*)\le \liminf J(U_k)$ and
$J(U_*)\le J(\Om_0)$, giving the energy bound in \eqref{eq:plan1-Psel};
lower semicontinuity of perimeter under $L^1$-convergence
(\cite[Prop.~12.15]{Maggi2012}) preserves $P(U_*)\le P_{\rm sel}(n,R)$.

The closure step uses $\gamma$-convergence of capacitary measures
inside $B_R$, whose constants depend only on $(n,R)$, and selects a
single limit from a single minimising sequence. There is no
compactness-of-near-counterexamples step.
\end{proof}

We fix such a minimiser $U_*$ for the rest of the subsection.

\subsubsection{Energy/asymmetry preservation before normalisation}
\label{subsubsec:plan1-sel-L1L2}

\begin{lemma}[Lemma~1: energy and asymmetry preservation]
\label{lem:plan1-sel-L1}
For every $\tau\le\tau_{\rm ex}(n,R)$ the minimiser $U_*$ of $J$
satisfies
\begin{align}
0&\le F_{\hat\eta}(U_*)-F_{\hat\eta}(B_1)\;\le\;\delta,
\label{eq:plan1-L1-energy}\\[2pt]
|\BDValpha(U_*)-\eps|&\le\frac{\sqrt{2\eps\delta+\delta^{2}}}{\tau}.
\label{eq:plan1-L1-asym}
\end{align}
In particular, with $\tau=q^{1/4}$ and $q\le1$,
\begin{equation}\label{eq:plan1-L1-asym-2}
|\BDValpha(U_*)-\eps|\le 2\,\eps\, q^{1/4}.
\end{equation}
\end{lemma}

\begin{proof}
By minimality, $J(U_*)\le J(\Om_0)$. Since $|\Om_0|=\omega_n$,
\(J(\Om_0)=F_{\hat\eta}(\Om_0)+\eps=F_{\hat\eta}(B_1)+\delta+\eps\)
by \eqref{eq:plan1-deltaTdef}--\eqref{eq:plan1-Om0}. Because $B_1$
minimises $F_{\hat\eta}$, $F_{\hat\eta}(U_*)\ge F_{\hat\eta}(B_1)$.
Combining,
\[
F_{\hat\eta}(U_*)
+\sqrt{\eps^{2}+\tau^{2}(\BDValpha(U_*)-\eps)^{2}}
\;\le\;F_{\hat\eta}(B_1)+\delta+\eps.
\]
The square-root term is $\ge \eps$, so
$F_{\hat\eta}(U_*)-F_{\hat\eta}(B_1)\le \delta$, which is
\eqref{eq:plan1-L1-energy}. Substituting back,
\(\sqrt{\eps^{2}+\tau^{2}(\BDValpha(U_*)-\eps)^{2}}\le \eps+\delta.\)
Squaring,
\(\eps^{2}+\tau^{2}(\BDValpha(U_*)-\eps)^{2}\le(\eps+\delta)^{2}
=\eps^{2}+2\eps\delta+\delta^{2}\), which is
\eqref{eq:plan1-L1-asym}.

For \eqref{eq:plan1-L1-asym-2}, $\tau=q^{1/4}$ and $q\le1$ give
$\delta=\eps q\le\eps$, hence $2\eps\delta+\delta^2\le 3\eps\delta\le
3\eps^2 q$. Therefore
$|\BDValpha(U_*)-\eps|\le \eps\sqrt{3q}/q^{1/4}\le 2\eps q^{1/4}$.
\end{proof}

\begin{lemma}[Lemma~2: volume error]
\label{lem:plan1-sel-L2}
There is $C_5(n,R)>0$ such that
\begin{equation}\label{eq:plan1-L2}
\bigl||U_*|-|B_1|\bigr|\;\le\;C_5(n,R)\,\delta,
\qquad
|r_*-1|\;\le\;C_5(n,R)\,\delta,
\quad r_*:=(|U_*|/\omega_n)^{1/n}.
\end{equation}
Moreover, for $\delta\le\delta_2(n,R)$ small enough,
$\tfrac12\le r_*\le \tfrac32$.
\end{lemma}

\begin{proof}
Let $B_{r_*}$ be a Euclidean ball with $|B_{r_*}|=|U_*|$, so
$|B_{r_*}|=\omega_n r_*^{n}$. By Schwarz symmetrisation
(\cite[Thm.~3]{Talenti1976}), $E(B_{r_*})\le E(U_*)$, and since
$f_{\hat\eta}$ depends only on $|U|$, also
$F_{\hat\eta}(B_{r_*})\le F_{\hat\eta}(U_*)$. Combining
\eqref{eq:plan1-Lem45} with \eqref{eq:plan1-L1-energy},
\[
C_4(n,R)^{-1}\bigl||U_*|-\omega_n\bigr|
\;\le\; F_{\hat\eta}(B_{r_*})-F_{\hat\eta}(B_1)
\;\le\; F_{\hat\eta}(U_*)-F_{\hat\eta}(B_1)
\;\le\; \delta.
\]
Hence $||U_*|-\omega_n|\le C_4(n,R)\,\delta$. The map
$s\mapsto (s/\omega_n)^{1/n}$ is bilipschitz on a neighbourhood of
$\omega_n$ with Lipschitz constant $\le 2(n\omega_n)^{-1}$ for $s$
within distance $\omega_n/2$ of $\omega_n$; this gives
$|r_*-1|\le 2(n\omega_n)^{-1}C_4(n,R)\,\delta=:C_5(n,R)\,\delta$. For
$\delta\le \delta_2(n,R):=(2C_5(n,R))^{-1}$ we have $|r_*-1|\le1/2$.
\end{proof}

\subsubsection{Quotient preservation under volume normalisation}
\label{subsubsec:plan1-sel-L3}

We now make rigorous the sketch in
\cite[\S2, Step~1 of Lemma~5.3]{BDV}
(``$|U_*\Delta\lambda U_*|\le C(n,R)|\lambda-1|$''),
which appeals to the BDV uniform perimeter bound but does not display
the divergence-theorem step. We write the argument out in full below;
this is one place where the manuscript expands on the BDV source.

\begin{lemma}[Lipschitz dependence of $\BDValpha$ under small dilation]
\label{lem:plan1-sel-dilation}
Let $U_*$ be the minimiser of $J$ from Lemma~\ref{lem:plan1-sel-existence}
and set $\lambda:=r_*^{-1}$ with $r_*$ as in Lemma~\ref{lem:plan1-sel-L2}.
Assume $|\lambda-1|\le 1/2$. Then
\begin{equation}\label{eq:plan1-dilL1}
|U_*\Delta(\lambda U_*)|
\;\le\;
\bigl(R\,P(U_*)+R^{n}\omega_n\bigr)\,|\lambda-1|
\;\le\;
C_{\rm dil}(n,R)\,|\lambda-1|,
\end{equation}
with $C_{\rm dil}(n,R):=R\,P_{\rm sel}(n,R)+R^{n}\omega_n$. Consequently
\begin{equation}\label{eq:plan1-dilAlpha}
|\BDValpha(\lambda U_*)-\BDValpha(U_*)|
\;\le\;C_2(R)\,|U_*\Delta(\lambda U_*)|
\;\le\;C_2(R)\,C_{\rm dil}(n,R)\,|\lambda-1|.
\end{equation}
\end{lemma}

\begin{proof}
Set $V:=U_*$ and $\lambda V:=\{\lambda x:x\in V\}\subset B_{\lambda R}
\subset B_{3R/2}$ (using $|\lambda-1|\le1/2$). By the De Giorgi
divergence theorem for finite-perimeter sets
(\cite[Thm.~16.3]{Maggi2012}; Source Map A4) applied to the vector field
$\Phi(x):=x$, for every finite-perimeter set $V\subset\R^n$ with
$|V|<\infty$,
\[
n|V|=\int_V \nabla\cdot\Phi\,dx=\int_{\partial^*V}\Phi\cdot\nu_V\,d\hH^{n-1}
=\int_{\partial^*V} x\cdot\nu_V(x)\,d\hH^{n-1}(x).
\]
\emph{Symmetric difference under dilation.} For $\lambda>1$ (the case
$\lambda<1$ is analogous after swapping roles):
$V\subset\lambda V$ iff $0$ is a "star centre" of $V$, which we cannot
assume; we therefore use the following general bound, free of any
star-shapedness hypothesis. Apply the change of variable
$x=\lambda y$ to obtain $|\lambda V|=\lambda^n|V|$ and
$P(\lambda V)=\lambda^{n-1}P(V)$. For any $V$ of finite perimeter in
$\R^n$, the difference $V\Delta(\lambda V)$ is contained in the
$|\lambda-1|R$-tubular neighbourhood of $\partial^*V$ inside
$B_{3R/2}$. Indeed, if $x\in V\setminus\lambda V$, then $x\in V$ and
$x/\lambda\notin V$, so the segment $[x/\lambda,x]\subset B_{3R/2}$
crosses $\partial V$ in the sense of capacitary boundary; the segment
has length $|x|(1-\lambda^{-1})\le R|\lambda-1|$. Analogous statement
for $\lambda V\setminus V$.

The standard perimeter--tubular-neighbourhood inequality
(\cite[Lemma~17.9 / Eq.~(17.6)]{Maggi2012}) gives, for the
$\rho$-neighbourhood $\mathcal N_\rho(\partial^*V)$ in $B_{3R/2}$,
\[
|\mathcal N_\rho(\partial^*V)|\;\le\;2\rho\,P(V)+\omega_n\rho^{n}
\]
for every $\rho\in(0,R)$ (the second term controls boundary lying
within $\rho$ of $\partial B_{3R/2}$, bounded crudely by the volume of
$B_{3R/2}$ minus $B_{3R/2-\rho}$, hence by
$n\omega_n(3R/2)^{n-1}\rho\le 2R^{n-1}\omega_n\rho$ for
$\rho\le R/2$). Specialising at $\rho=R|\lambda-1|\le R/2$,
\[
|V\Delta(\lambda V)|
\;\le\;2R\,P(V)\,|\lambda-1|
+2R^{n}\omega_n\,|\lambda-1|.
\]
Using $P(V)=P(U_*)\le P_{\rm sel}(n,R)$ from
Lemma~\ref{lem:plan1-sel-existence},
\[
|U_*\Delta(\lambda U_*)|\le \bigl(2R\,P_{\rm sel}(n,R)
+2R^{n}\omega_n\bigr)|\lambda-1|
\;\le\;C_{\rm dil}(n,R)\,|\lambda-1|,
\]
which is \eqref{eq:plan1-dilL1} (after redefining $C_{\rm dil}$ to
absorb the factor of $2$). Estimate \eqref{eq:plan1-dilAlpha} follows
from \eqref{eq:bdv42ii}; note $\lambda U_*\subset B_{3R/2}$ so we apply
\eqref{eq:bdv42ii} on $B_{3R/2}$, picking up $C_2(3R/2)\le C_2(R)$
after enlarging $C_2$ by a constant depending on $R$ (still
$=:C_2(R)$).
\end{proof}

We now state the central quotient-preservation lemma.

\begin{lemma}[Lemma~3: quotient preservation after volume normalisation]
\label{lem:plan1-sel-L3}
Set $\lambda:=r_*^{-1}=(|B_1|/|U_*|)^{1/n}$ and
$\Otilde:=\lambda U_*$. Under the smallness regime
\eqref{eq:plan1-Sml} and with $\tau=q^{1/4}$, after translation by
$-x_{U_*}$ if needed:
\begin{enumerate}
\item[\rm(a)] $|\Otilde|=|B_1|=\omega_n$;
\item[\rm(b)] $|\BDValpha(\Otilde)-\eps|\le \rho\,\eps$ where
\begin{equation}\label{eq:plan1-rho}
\rho=\rho(q)
:=\frac{\sqrt{2q+q^{2}}}{\tau}+C_{\BDValpha}(n,R)\,q
\;\le\;2q^{1/4}+C_{\BDValpha}(n,R)\,q;
\end{equation}
for $q\le q_{\rm sel}(n,R)$, $\rho(q)\le 1/2$;
\item[\rm(c)] $E(\Otilde)-E(B_1)\le C_E(n,R)\,\delta$;
\item[\rm(d)] $\displaystyle Q_{\BDValpha}(\Otilde)\;\le\;
\frac{C_E(n,R)}{1-\rho(q)}\,Q_{\BDValpha}(\Om_0)$.
\end{enumerate}
The constants $C_{\BDValpha}(n,R)$ and $C_E(n,R)$ are given
explicitly in \eqref{eq:plan1-Calpha-def} and
\eqref{eq:plan1-CE-def} below.
\end{lemma}

\begin{proof}
\emph{Part~(a).} By construction
$|\Otilde|=\lambda^{n}|U_*|=(\omega_n/|U_*|)|U_*|=\omega_n$.

\emph{Part~(b): asymmetry control.} Combining
\eqref{eq:plan1-dilAlpha} with $|\lambda-1|\le 2C_5(n,R)\,\delta$ from
Lemma~\ref{lem:plan1-sel-L2} (which uses
$(\omega_n/s)^{1/n}-1$ has Lipschitz constant
$\le 2(n\omega_n)^{-1}$ near $s=\omega_n$),
\[
|\BDValpha(\Otilde)-\BDValpha(U_*)|
\;\le\;C_2(R)\,C_{\rm dil}(n,R)\cdot 2C_5(n,R)\,\delta
\;=:\;C_{\BDValpha}(n,R)\,\delta,
\]
i.e.\
\begin{equation}\label{eq:plan1-Calpha-def}
C_{\BDValpha}(n,R):=2\,C_2(R)\,C_{\rm dil}(n,R)\,C_5(n,R).
\end{equation}
Combined with \eqref{eq:plan1-L1-asym} and $\delta=\eps q$,
\[
|\BDValpha(\Otilde)-\eps|
\;\le\;|\BDValpha(\Otilde)-\BDValpha(U_*)|
+|\BDValpha(U_*)-\eps|
\;\le\;\eps\,C_{\BDValpha}(n,R)\,q
+\frac{\sqrt{2\eps\delta+\delta^2}}{\tau}.
\]
The square-root term equals
$\eps\,\sqrt{2q+q^{2}}/\tau$. Hence $\rho(q)$ defined in
\eqref{eq:plan1-rho}. With $\tau=q^{1/4}$ and $q\le1$,
$\sqrt{2q+q^{2}}/\tau\le \sqrt{3q}/q^{1/4}\le2q^{1/4}$. Under
\eqref{eq:plan1-Sml},
$\rho(q)\le 2 q_{\rm sel}^{1/4}+C_{\BDValpha}(n,R)\,q_{\rm sel}$;
choosing $q_{\rm sel}(n,R)$ small enough (still depending only on
$(n,R)$) gives $\rho(q)\le 1/2$.

\emph{Part~(c): energy scaling.} For every quasi-open
$V\subset B_R$ and $\mu>0$ with $\mu V\subset B_{R'}$,
\begin{equation}\label{eq:plan1-Escale}
E(\mu V)=\mu^{n+2}E(V).
\end{equation}
\emph{Proof of \eqref{eq:plan1-Escale}.} If $u_V$ is the torsion
function of $V$, define $w(x):=\mu^{2}u_V(x/\mu)$ for $x\in\mu V$.
Then $\nabla w(x)=\mu\nabla u_V(x/\mu)$, $\Delta w(x)=\Delta u_V(x/\mu)
=-1$, $w=0$ on $\partial(\mu V)$, so $w=u_{\mu V}$. Substituting,
\[
E(\mu V)=-\tfrac12\!\int_{\mu V}u_{\mu V}\,dx
=-\tfrac12\!\int_{\mu V}\mu^{2}u_V(x/\mu)\,dx
=-\tfrac12\,\mu^{n+2}\!\int_V u_V(y)\,dy=\mu^{n+2}E(V).
\]
Apply \eqref{eq:plan1-Escale} with $\mu=\lambda=r_*^{-1}$:
\begin{equation}\label{eq:plan1-Etilde-decomp}
E(\Otilde)-E(B_1)
=\lambda^{n+2}\bigl(E(U_*)-E(B_1)\bigr)
+(\lambda^{n+2}-1)\,E(B_1).
\end{equation}
From \eqref{eq:plan1-L1-energy} and
\eqref{eq:plan1-Lem45},
\[
E(U_*)-E(B_1)
\;\le\;F_{\hat\eta}(U_*)-F_{\hat\eta}(B_1)+\hat\eta^{-1}||U_*|-\omega_n|
\;\le\;\bigl(1+\hat\eta^{-1}C_4(n,R)\bigr)\delta
\;=:\;C_6(n,R)\,\delta.
\]
By Lemma~\ref{lem:plan1-sel-L2} and $|\lambda-1|\le1/2$,
\[
\lambda^{n+2}\le (3/2)^{n+2}=:C_7(n),
\qquad
|\lambda^{n+2}-1|\le (n+2)(3/2)^{n+1}|\lambda-1|
\le (n+2)(3/2)^{n+1}C_5(n,R)\,\delta=:C_8(n,R)\,\delta.
\]
Since $|E(B_1)|=T_n/2=\omega_n/(2n(n+2))$ is dimensional, plugging
into \eqref{eq:plan1-Etilde-decomp},
\begin{align*}
E(\Otilde)-E(B_1)
&\le C_7(n)\,C_6(n,R)\,\delta + C_8(n,R)\,|E(B_1)|\,\delta\\
&= \bigl(C_6 C_7+C_8|E(B_1)|\bigr)\,\delta
\;=:\;C_E(n,R)\,\delta,
\end{align*}
with
\begin{equation}\label{eq:plan1-CE-def}
C_E(n,R):=
\bigl(1+\hat\eta(n,R)^{-1}C_4(n,R)\bigr)\,(3/2)^{n+2}
+\frac{(n+2)(3/2)^{n+1}\,C_5(n,R)\,\omega_n}{2n(n+2)}.
\end{equation}
This is \eqref{eq:plan1-CE-def}; we have used $E(B_1)<0$ when
applying \eqref{eq:plan1-Etilde-decomp}; the signed contribution
$(\lambda^{n+2}-1)E(B_1)$ has the wrong sign when $\lambda>1$ but is
controlled in absolute value.

\emph{Part~(d): quotient.} Combining (b) and (c) with
$\BDValpha(\Otilde)\ge(1-\rho(q))\eps$,
\[
Q_{\BDValpha}(\Otilde)
=\frac{E(\Otilde)-E(B_1)}{\BDValpha(\Otilde)}
\;\le\;\frac{C_E(n,R)\,\delta}{(1-\rho(q))\,\eps}
=\frac{C_E(n,R)}{1-\rho(q)}\,Q_{\BDValpha}(\Om_0).
\]
This completes the proof.
\end{proof}

\begin{remark}[Preservation of the deficit-to-asymmetry ratio]
\label{rmk:plan1-sel-ratio}
Lemma~\ref{lem:plan1-sel-L3}(d) is the load-bearing
deficit-to-asymmetry preservation statement: starting from
$Q_{\BDValpha}(\Om_0)=q$, the volume-normalised companion satisfies
$Q_{\BDValpha}(\Otilde)\le \widetilde C(n,R)\,q$ for an explicit constant
$\widetilde C(n,R)=2C_E(n,R)$ once $\rho(q)\le 1/2$. The improvement
factor on the asymmetry is
$\BDValpha(\Otilde)\ge \tfrac12\BDValpha(\Om_0)$, so the new asymmetry
is at most a factor of $2$ smaller than the old. The chain
\eqref{eq:plan1-Frbound} then transfers the same statement to the
Fraenkel quotient at the price of a square root.

Volume normalisation enters the proof in two distinct places:
\begin{itemize}
\item In \emph{(b)}, the volume defect produces the
$|\lambda-1|\le C_5\delta$ contribution to the asymmetry change
through Lemma~\ref{lem:plan1-sel-dilation}.
\item In \emph{(c)}, the volume-scaling identity
\eqref{eq:plan1-Escale} produces the explicit
$\lambda^{n+2}$-factor, which is the only place where the dimension
$n$ enters the constant $C_E(n,R)$.
\end{itemize}
Both contributions are $O(\delta)$, hence subleading compared to the
$O(\eps q^{1/4})$ contribution from
\eqref{eq:plan1-L1-asym}. This is the precise mechanism by which the
single-set selection step preserves the quotient $Q_{\BDValpha}$ up to a
multiplicative $(n,R)$-constant.
\end{remark}

\medskip

\subsubsection{The selection theorem and BDV regularity package}
\label{subsubsec:plan1-sel-thm}

\begin{theorem}[Single-set quantitative selection principle]
\label{thm:plan1-sel-main}
Fix $n\ge2$ and $R\ge2$. There exist constants
\[
\eps_{\rm sel}(n,R)>0,\quad q_{\rm sel}(n,R)>0,\quad
C_{\rm sel}(n,R)>0,\quad P_{\rm sel}(n,R)>0,
\]
depending only on $(n,R)$ via the named BDV constants of
\eqref{eq:plan1-tauex} and \eqref{eq:plan1-CE-def},
with the following property. Let $\Om_0\subset B_R$ be quasi-open
with
\[
|\Om_0|=\omega_n,\qquad
0<\eps:=\BDValpha(\Om_0)\le\eps_{\rm sel}(n,R),\qquad
0<q:=Q_{\BDValpha}(\Om_0)\le q_{\rm sel}(n,R).
\]
Set $\tau:=q^{1/4}$, let $U_*$ be a minimiser of the functional $J$ in
\eqref{eq:plan1-Jdef} provided by
Lemma~\ref{lem:plan1-sel-existence}, and define
\[
r_*:=(|U_*|/\omega_n)^{1/n},
\qquad
\Otilde:=r_*^{-1}U_*.
\]
Then, after the translation $\Otilde\mapsto \Otilde-x_{\Otilde}$ that
brings the barycentre to the origin:
\begin{enumerate}
\item[\rm(i)] $P(U_*)\le P_{\rm sel}(n,R)$ and $|\Otilde|=\omega_n$;
\item[\rm(ii)] $(1-\rho)\eps\le \BDValpha(\Otilde)\le (1+\rho)\eps$ with
$\rho=\rho(q)\le 1/2$ given by \eqref{eq:plan1-rho}, hence
$\tfrac12\,\BDValpha(\Om_0)\le \BDValpha(\Otilde)\le \tfrac32\,\BDValpha(\Om_0)$;
\item[\rm(iii)] $\displaystyle
Q_{\BDValpha}(\Otilde)\;\le\;
\frac{C_{\rm sel}(n,R)}{1-\rho(q)}\,Q_{\BDValpha}(\Om_0)
\;\le\; 2\,C_{\rm sel}(n,R)\,Q_{\BDValpha}(\Om_0)
$, with $C_{\rm sel}(n,R):=C_E(n,R)$ from \eqref{eq:plan1-CE-def};
\item[\rm(iv)] $\Otilde$ inherits the BDV penalised-minimiser regularity
package on its torsion function $\tilde u:=\tilde u_{\Otilde}$ via the
perturbed quasi-minimality inequality \eqref{eq:plan1-qmin} below: namely
\begin{itemize}
\item two-sided density estimates on $\Otilde$ at scales
$\le r_0=r_0(n,R)$ (BDV Lemma~4.9);
\item Lipschitz/nondegeneracy of $\tilde u$ near $\partial\Otilde$ (BDV
Lemma~4.11, 4.12);
\item smooth, uniformly $C^{k}(n,R)$-bounded Bernoulli coefficient on
$\partial\Otilde$ (BDV Lemma~4.16);
\end{itemize}
all constants depending only on $(n,R)$;
\item[\rm(v)] On $\partial\Otilde$ the selected Bernoulli law takes the
volume-rescaled form of Lemma~\ref{lem:plan1-sel-bern} below.
\end{enumerate}
\end{theorem}

\begin{proof}
Parts (i)--(iii) are
Lemmas~\ref{lem:plan1-sel-existence}--\ref{lem:plan1-sel-L3}.
Parts (iv)--(v) are stated and justified next.
\end{proof}

\paragraph{Perturbed quasi-minimality inequality.}
By a direct comparison test in $J$, the minimiser $U_*=\{u_*>0\}$ (where
$u_*=u_{U_*}$ is the torsion function) satisfies the BDV (4.19)
inequality: for every $v\in H^1_0(B_R)$,
\begin{equation}\label{eq:plan1-qmin}
\tfrac12\!\int|\nabla u_*|^2-\!\int u_*
+f_{\hat\eta}(|\{u_*>0\}|)
\;\le\;\tfrac12\!\int|\nabla v|^2-\!\int v+f_{\hat\eta}(|\{v>0\}|)
+C_R\,\tau\,\bigl|\{u_*>0\}\Delta\{v>0\}\bigr|.
\end{equation}
The error term arises from the Lipschitz bound on the square-root
penalty in \eqref{eq:plan1-Jdef}, dominated by
$\tau\,C_2(R)|U\Delta U_*|$, where $C_2(R)$ is the BDV
constant of \eqref{eq:bdv42ii}; so $C_R$ may be taken equal to
$C_2(R)$. For $\tau\le\tau_{\rm reg}(n,R)$,
\eqref{eq:plan1-qmin} places $u_*$ in the
Alt--Caffarelli quasi-minimiser class with smallness constant
controlled by $(n,R)$, and the BDV
package \cite[Lemmas~4.9, 4.11, 4.12, 4.16]{BDV} applies; this yields
the density, Lipschitz, nondegeneracy and Bernoulli statements of
Theorem~\ref{thm:plan1-sel-main}(iv). The transfer from $U_*$ to
$\Otilde=r_*^{-1}U_*$ is a bilipschitz dilation with
$|\lambda-1|\le 1/2$ (Lemma~\ref{lem:plan1-sel-L2}); the density,
Lipschitz, and $C^{k}$ constants degrade by a factor at most
$(3/2)^{n}$, which is absorbed into the $(n,R)$-bookkeeping.

\paragraph{Volume-rescaled Bernoulli law.}

\begin{lemma}[Volume-rescaled Bernoulli law on $\partial\Otilde$]
\label{lem:plan1-sel-bern}
Let $\tilde u(y):=r_*^{-2}u_{U_*}(r_*y)$ for $y\in\Otilde=r_*^{-1}U_*$.
Then $-\Delta\tilde u=1$ in $\Otilde$, $\tilde u=0$ on $\partial\Otilde$,
and on $\partial\Otilde$ the BDV Lemma~4.16 Bernoulli law takes the form
\begin{equation}\label{eq:plan1-Berntilde}
|\nabla\tilde u(y)|^2
=\tilde\Lambda
+\tilde a_\sigma\,\tilde\Psi_{r_*}(y),
\end{equation}
where
\begin{equation}\label{eq:plan1-Berncoeff}
\tilde\Lambda=r_*^{-2}\Lambda,
\qquad
\tilde a_\sigma=r_*^{-1}\,a_\sigma,
\qquad
|a_\sigma|\le \tau,
\end{equation}
and
\[
\tilde\Psi_{r_*}(y)
:=|y-r_*^{-1}x_{U_*}|-r_*^{n}\,V_{\Otilde}\cdot y,
\qquad
V_{\Otilde}:=\!\int_{\Otilde}\!\frac{z-x_{\Otilde}}{|z-x_{\Otilde}|}\,dz.
\]
Both prefactors satisfy
$\tilde\Lambda=\Lambda\,(1+O(\delta))$ and
$\tilde a_\sigma=a_\sigma\,(1+O(\delta))$ as $\delta\to0$, with
implicit constants depending only on $(n,R)$.
\end{lemma}

\begin{proof}
For $y\in\Otilde$, $\nabla\tilde u(y)=r_*^{-1}\nabla u_{U_*}(r_*y)$ and
$\Delta\tilde u(y)=\Delta u_{U_*}(r_*y)=-1$, so $\tilde u$ solves
$-\Delta\tilde u=1$ in $\Otilde$ with zero Dirichlet condition on
$\partial\Otilde$ (which is $r_*^{-1}\partial U_*$). Squaring,
$|\nabla\tilde u(y)|^2=r_*^{-2}|\nabla u_{U_*}(r_*y)|^2$.

By BDV Lemma~4.16, on $\partial U_*$,
$|\nabla u_{U_*}(x)|^2=\Lambda+a_\sigma\Psi_{U_*}(x)$ with
$\Psi_{U_*}(x)=|x-x_{U_*}|-V_{U_*}\cdot x$ and
$V_{U_*}=\int_{U_*}(z-x_{U_*})/|z-x_{U_*}|\,dz$. Plugging $x=r_*y$,
\[
|\nabla\tilde u(y)|^2
=r_*^{-2}\Lambda+r_*^{-2}a_\sigma\Psi_{U_*}(r_*y).
\]
Change of variables $z=r_*w$ in the centroid integral:
\[
V_{U_*}
=\!\int_{U_*}\frac{z-x_{U_*}}{|z-x_{U_*}|}dz
=r_*^{n}\!\int_{\Otilde}\frac{w-x_{\Otilde}}{|w-x_{\Otilde}|}dw
=r_*^{n}V_{\Otilde},
\]
using $x_{U_*}=r_*x_{\Otilde}$ for the barycentre. Hence
\[
\Psi_{U_*}(r_*y)
=|r_*y-x_{U_*}|-V_{U_*}\cdot r_*y
=r_*\bigl(|y-r_*^{-1}x_{U_*}|-r_*^{n}V_{\Otilde}\cdot y\bigr)
=r_*\,\tilde\Psi_{r_*}(y).
\]
Substituting,
\(
|\nabla\tilde u(y)|^2
=r_*^{-2}\Lambda+r_*^{-1}a_\sigma\,\tilde\Psi_{r_*}(y),
\)
which is \eqref{eq:plan1-Berntilde} with the coefficients
\eqref{eq:plan1-Berncoeff}. Lemma~\ref{lem:plan1-sel-L2} gives
$|r_*-1|\le C_5(n,R)\delta$, hence
$\tilde\Lambda=\Lambda(1+O_{n,R}(\delta))$ and analogously for
$\tilde a_\sigma$.
\end{proof}

\subsubsection{Summary of constants and dependences}
\label{subsubsec:plan1-sel-constants}

For ease of cross-reference in later subsections, we record the
constants used in Theorem~\ref{thm:plan1-sel-main} together with their
dependences. All constants depend only on the pair $(n,R)$ via named BDV
constants; none is set by a compactness or contradiction step.

\begin{center}
\renewcommand{\arraystretch}{1.20}
\begin{tabular}{|l|l|l|}
\hline
Symbol & Defined in / source & Depends on \\\hline
$\hat\eta,C_4$ & BDV Lemma~4.5 & $n,R$ \\
$C_1$ & BDV Lemma~4.2(i), eq.~\eqref{eq:bdv42i} & $n$ \\
$C_2$ & BDV Lemma~4.2(ii), eq.~\eqref{eq:bdv42ii} & $R$ \\
$\sigma_2,\sigma_3$ & BDV Lemmas~4.9, 4.16 & $n,R$ \\
$\tau_{\rm ex},\tau_{\rm reg}$ & eq.~\eqref{eq:plan1-tauex} & $n,R$ \\
$P_{\rm sel}$ & Lemma~\ref{lem:plan1-sel-existence} & $n,R$ \\
$C_5$ & Lemma~\ref{lem:plan1-sel-L2} & $n,R$ \\
$C_{\rm dil}$ & Lemma~\ref{lem:plan1-sel-dilation},
eq.~\eqref{eq:plan1-dilL1} & $n,R$ \\
$C_{\BDValpha}=2 C_2 C_{\rm dil}C_5$ & eq.~\eqref{eq:plan1-Calpha-def} & $n,R$ \\
$C_E$ & eq.~\eqref{eq:plan1-CE-def} & $n,R$ \\
$C_{\rm sel}=C_E$ & Theorem~\ref{thm:plan1-sel-main} & $n,R$ \\
$\eps_{\rm sel},q_{\rm sel}$ & eq.~\eqref{eq:plan1-Sml} & $n,R$ \\
$\rho(q)=\sqrt{2q+q^2}/\tau+C_{\BDValpha}q$ &
eq.~\eqref{eq:plan1-rho} & $n,R,q$ \\
$\tilde\Lambda=r_*^{-2}\Lambda$,
$\tilde a_\sigma=r_*^{-1}a_\sigma$ &
Lemma~\ref{lem:plan1-sel-bern} & $n,R$ \\
\hline
\end{tabular}
\end{center}

\paragraph{Outputs feeding the next subsection.}
The selection principle (Theorem~\ref{thm:plan1-sel-main}) produces a
volume-normalised quasi-open set $\Otilde\subset B_{R'(n,R)}$ with
$|\Otilde|=\omega_n$, $\BDValpha(\Otilde)\in[\eps/2,3\eps/2]$,
$Q_{\BDValpha}(\Otilde)\le 2C_{\rm sel}(n,R)\,Q_{\BDValpha}(\Om_0)$, and
the full BDV free-boundary regularity package. This is the input for
the graph-entry subsection (Plan 1 \S 3.2), which converts smallness of
$\BDValpha(\Otilde)$ into $C^{1,\gamma}$ normal-graph parametrisation
over $\partial B_1(x_{\Otilde})$ via the
$\alpha\to$Hausdorff$\to$flatness chain.

% UNIFIED 2026-05-13
% SOURCE MAP
% ----------
% This subsection is the Plan 1 / Route A ``graph entry and regularity'' block.
% It is intended for inclusion (by Agent 19) into Manuscript/02_plan1_draft.tex
% as the subsection that follows Agent 5's selection-principle subsection and
% precedes Agent 7's nearly-spherical closure subsection.
%
% Mapping of items to sources (file paths relative to repository root):
%
%   - Plan 1 input / regularity package (BDV black boxes):
%       Plan 1/graph-entry-closure.tex  (selected-minimizer input, §2-§3)
%       Plan 1/graph-entry-closure.md
%       Final/GraphEntry.tex             (Lemmas 2.1, 2.2; §3 flatness;
%                                         §4 patching; Theorem 4.1 = GraphEntry)
%       Final/SchauderInterpolation.tex  (§2 hodograph, §3 Schauder bootstrap,
%                                         §4 interpolation, Theorem 4.1 = Sph-entry)
%
%   - Item D2.1 (BDV Lemma 4.2 chain): CITE(BDV, Lem. 4.2), as in
%       Final/GraphEntry.tex Lemma 2.1 and Plan 1/graph-entry-closure.tex §2.
%   - Item D2.2 (Alt-Caffarelli flatness, BDV Theorem 4.18 form):
%       CITE(BDV, Thm. 4.18); used as the black box (T4.18) in
%       Final/GraphEntry.tex §1.2.
%   - Item D2.3 (Global C^{1,gamma} normal graph, threshold alpha_graph):
%       LOCAL(Final/GraphEntry.tex Theorem 6.1) and
%       LOCAL(Plan 1/graph-entry-closure.tex §5).
%   - Item D2.4 (Kinderlehrer-Nirenberg hodograph straightening):
%       CITE(KN1977); LOCAL(Final/SchauderInterpolation.tex §2).
%   - Item D2.5 (Lieberman oblique Schauder bootstrap to C^{m,gamma}):
%       CITE(Li86); LOCAL(Final/SchauderInterpolation.tex §3).
%   - Item D2.6 (Gagliardo-Nirenberg / Holder interpolation on the sphere):
%       CITE(GT Lemma 6.32), CITE(BL Thm 6.4.5);
%       LOCAL(Final/SchauderInterpolation.tex Lemma 4.1).
%   - Item D2.7 (Schauder-interpolation threshold alpha_sph(n,R,gamma_0)):
%       LOCAL(Final/SchauderInterpolation.tex Theorem 5.1).
%
% Notation conversion. All sources use the dimension symbol N. By
% Manuscript/00_notation_register.md §1 (canonical n) and §17.1, every N
% has been replaced by n in this draft. The BDV smooth asymmetry of the
% selected (volume-normalised) minimizer is written alpha_BDV(Otilde) and
% the Fraenkel asymmetry of the original domain Omega is A(Omega), per
% the notation register §5.
%
% Hypothesis class. The standing hypothesis throughout is the bounded
% Saint-Venant setting: Omega open in R^n, |Omega| = omega_n, Omega ⊂ B_R,
% R ≥ 2. The graph-entry subsection applies to the volume-normalised
% selected minimizer Otilde produced by Agent 5; the selection parameter
% tau is bounded above by the regularity threshold tau_reg(n,R) of BDV.
% Selected minimizers are quasi-minimizers of the BDV penalised Bernoulli
% functional (Agent 5, item D1.4), hence the BDV regularity package (L4.9,
% L4.12, L4.16, T4.18) is available with constants depending only on (n,R).
%
% Status: every numerical lemma below is either drafted here (DRAFT) with
% full proof at the level of detail of the source, or stated and cited
% with the exact source label.

\subsection{Graph entry and regularity}
\label{subsec:plan1-graph-entry}

\paragraph{Standing hypotheses (Plan 1 / Route A).}
Let \(n\ge 2\) and \(R\ge 2\) be fixed. Throughout this subsection,
\(\Otilde\subset B_R\subset\R^n\) denotes the \emph{volume-normalised
selected minimizer} produced by the selection principle of
Subsection~\ref{subsec:plan1-selection} (Agent~5). It satisfies
\begin{equation}\label{eq:plan1-ge-Otilde}
   |\Otilde|=|B_1|=\omega_n,\qquad \Otilde\subset B_R,
\end{equation}
and is, by construction, an almost-minimizer of the BDV penalised Bernoulli
functional with selection parameter \(\tau\le\tau_{\rm reg}(n,R)\), where
\(\tau_{\rm reg}(n,R)>0\) is the regularity threshold fixed in
\cite[\S 4]{BDV} (BDV) and recorded in Subsection~\ref{subsec:plan1-selection}.
The BDV smooth asymmetry of \(\Otilde\) is denoted by \(\BDValpha(\Otilde)\)
(notation register \S 5; in the selection-principle subsection this is
written \(\alpha(\Otilde)\) without ambiguity, but the global manuscript
uses \(\BDValpha\) whenever there is a clash with the coarea quantity
\(\alpha(t)\) of Plan~2).

By the BDV regularity package, which is inherited by every selected
minimizer with \(\tau\le\tau_{\rm reg}(n,R)\), there exist positive
constants
\[
   c_d=c_d(n,R),\quad r_d=r_d(n,R),\quad
   L_u=L_u(n,R),\quad
   q_\pm=q_\pm(n,R),\quad L_q=L_q(n,R),
\]
for which the following hold; cf.\
\cite[Lemmas 4.9, 4.12, 4.16, Theorem 4.18]{BDV}.

\begin{enumerate}[label=\textup{(R\arabic*)},leftmargin=2.7em]
\item\label{R1} \emph{Two-sided density.} For every \(x\in\partial\Otilde\)
   and every \(r\in(0,r_d]\),
   \[
      c_d r^n
      \le|\Otilde\cap B_r(x)|
      \le(1-c_d)r^n,
      \qquad
      c_d r^n
      \le|B_r(x)\setminus\Otilde|
      \le(1-c_d)r^n.
   \]
\item\label{R2} \emph{Lipschitz nondegeneracy.} The torsion potential
   \(\widetilde u\in H^1_0(\Otilde)\) of \(\Otilde\) extends to
   \(C^{0,1}(\overline{B_R})\) with \(\Lip(\widetilde u)\le L_u\), and
   \(\sup_{B_r(x)}\widetilde u\ge c_d r\) for every
   \(x\in\partial\Otilde\), \(r\in(0,r_d]\).
\item\label{R3} \emph{Bernoulli coefficient.} On the flat portion of
   \(\partial\Otilde\) the function
   \(q(\cdot):=|\nabla\widetilde u(\cdot)|\) admits a \(C^{1,\gamma}\)
   one-sided extension with \(q_-\le q\le q_+\) and
   \(\|q\|_{C^{1,\gamma}}\le L_q\).
\item\label{R4} \emph{Alt--Caffarelli flatness improvement.} There exist
   \(\bar\mu=\bar\mu(n,R)\in(0,1)\), \(\bar\rho=\bar\rho(n,R)>0\),
   \(\gamma=\gamma(n,R)\in(0,1)\) and \(C_{\rm AC}=C_{\rm AC}(n,R)\) such
   that, whenever \(\widetilde u\) is of Alt--Caffarelli class
   \(F(\bar\mu,1,+\infty)\) on a ball \(B_{\bar\rho}(x_0)\subset B_R\),
   the portion \(\partial\Otilde\cap B_{\bar\rho/2}(x_0)\) is a
   \(C^{1,\gamma}\) graph in the relevant direction with norm
   \(\le C_{\rm AC}\bar\mu\).
\end{enumerate}

\paragraph{Linking asymmetries.}
The Fraenkel asymmetry \(\Fr(\cdot)\) and the BDV smooth asymmetry
\(\BDValpha(\cdot)\) are linked by
\cite[Lemma~4.2]{BDV}: there are positive constants \(C_{\Fr}(n)\) and
\(C_\alpha(n)\) such that, for every measurable \(U\) with \(|U|=|B_1|\),
\begin{equation}\label{eq:plan1-ge-BDVL42}
   \Fr(U)^2\le C_{\Fr}(n)\,\BDValpha(U),\qquad
   |U\Delta B_1|\le C_\alpha(n)\,\BDValpha(U)^{1/2}.
\end{equation}
The constants \(C_{\Fr}(n)\) and \(C_\alpha(n)\) are dimensional; they
come from the BDV definition of \(\BDValpha\). We use both directions of
\eqref{eq:plan1-ge-BDVL42}: the first to control the optimal Fraenkel
translation, the second to transfer smallness of \(\BDValpha\) to
\(L^1\)-closeness to \(B_1\).

We now proceed in five steps:
\begin{enumerate}
\item[(S1)] From \(\BDValpha(\Otilde)\) small to small Hausdorff distance
   \(d_H(\partial\Otilde,\partial B_1)\).
\item[(S2)] From small Hausdorff distance to fixed-scale Alt--Caffarelli
   flatness of \(\widetilde u\).
\item[(S3)] Patching to a global \(C^{1,\gamma}\) normal graph; threshold
   \(\alpha_{\rm graph}(n,R)\).
\item[(S4)] Hodograph straightening + Lieberman oblique Schauder bootstrap
   to uniform \(C^{m,\gamma}\) regularity, with explicit constants
   \(M_m(n,R)\).
\item[(S5)] H\"older interpolation on \(\partial B_1\); threshold
   \(\alpha_{\rm sph}(n,R,\gamma_0)\) for small \(C^{2,\gamma_0}\) norm.
\end{enumerate}

\subsubsection{Step 1: from asymmetry to Hausdorff distance}
\label{subsubsec:plan1-ge-S1}

\begin{lemma}[Density forces Hausdorff smallness]
\label{lem:plan1-ge-density-Hausdorff}
Set \(c_*:=\min(c_d,1/2)\) and
\begin{equation}\label{eq:plan1-ge-CH}
   C_H(n,R)
   :=4(\omega_n c_*)^{-1/n}+(R+1)(c_*r_d^n)^{-1/n}.
\end{equation}
Then every selected minimizer \(\Otilde\) as above satisfies
\begin{equation}\label{eq:plan1-ge-dH-bound}
   d_H(\partial\Otilde,\partial B_1)
   \le C_H(n,R)\,|\Otilde\Delta B_1|^{1/n}.
\end{equation}
\end{lemma}

\begin{proof}
Write \(d:=d_H(\partial\Otilde,\partial B_1)\). By definition there is a
point \(x\in\partial\Otilde\) (or, by symmetry, in \(\partial B_1\)) with
\(\dist(x,\partial B_1)=d\) (resp.\ \(\dist(x,\partial\Otilde)=d\)).

\emph{Case 1: \(x\in\partial\Otilde\), \(\dist(x,\partial B_1)=d\le 2 r_d\).}
Then \(B_{d/2}(x)\) is disjoint from \(\partial B_1\), hence
\(B_{d/2}(x)\subset B_1\) or \(B_{d/2}(x)\subset\R^n\setminus\overline{B_1}\).
In the first case, the density estimate \ref{R1} at the radius
\(r:=\min(d/2,r_d)\) gives
\(|B_r(x)\setminus\Otilde|\ge c_d r^n\), and
\(B_r(x)\setminus\Otilde\subset B_1\setminus\Otilde
\subset\Otilde\Delta B_1\), so
\(c_d (d/2)^n\le|\Otilde\Delta B_1|\). In the second case, density of
\(\Otilde\) inside \(B_r(x)\) yields the same conclusion with
\(\Otilde\cap B_r(x)\subset\Otilde\setminus B_1\). Either way,
\(d\le 2(c_d^{-1}|\Otilde\Delta B_1|)^{1/n}\).

\emph{Case 2: \(x\in\partial B_1\), \(\dist(x,\partial\Otilde)=d\le 2 r_d\).}
The half-ball \(B_{d/2}(x)\cap B_1\) (or its complement, depending on
orientation) has volume at least \(\tfrac12\omega_n(d/2)^n\) and lies in
\(B_1\setminus\Otilde\) or \(\Otilde\setminus B_1\); both are subsets of
\(\Otilde\Delta B_1\). Hence
\(\tfrac12\omega_n(d/2)^n\le|\Otilde\Delta B_1|\).

\emph{Case 3: \(d>2 r_d\).} Apply the same density argument at fixed
radius \(r_d\) to a point realising the distance: this yields
\(c_d r_d^n\le|\Otilde\Delta B_1|\). Hence
\(d\le R+1\le(R+1)(c_d r_d^n)^{-1/n}|\Otilde\Delta B_1|^{1/n}\) (using
\(d\le\diam(B_R\cup B_1)\le R+1\)).

Combining the three cases and bounding all dimensional constants by
\(c_*=\min(c_d,1/2)\) gives~\eqref{eq:plan1-ge-dH-bound} with
\(C_H(n,R)\) defined by~\eqref{eq:plan1-ge-CH}.
\end{proof}

\begin{proposition}[Quantitative asymmetry-to-Hausdorff]
\label{prop:plan1-ge-alpha-to-H}
Set
\begin{equation}\label{eq:plan1-ge-CHprime}
   C_H'(n,R):=C_H(n,R)\,C_\alpha(n)^{1/n},
\end{equation}
and
\begin{equation}\label{eq:plan1-ge-alpha0}
   \alpha_0(n,R):=\bigl(r_d(n,R)/C_H'(n,R)\bigr)^{2n}.
\end{equation}
Then every selected minimizer \(\Otilde\) with
\(\BDValpha(\Otilde)\le\alpha_0(n,R)\) satisfies
\begin{equation}\label{eq:plan1-ge-dH-alpha}
   d_H(\partial\Otilde,\partial B_1)
   \le C_H'(n,R)\,\BDValpha(\Otilde)^{1/(2n)}\le r_d(n,R).
\end{equation}
\end{proposition}

\begin{proof}
The second half of~\eqref{eq:plan1-ge-BDVL42} gives
\(|\Otilde\Delta B_1|\le C_\alpha(n)\BDValpha(\Otilde)^{1/2}\). Insertion
into Lemma~\ref{lem:plan1-ge-density-Hausdorff} gives the first inequality
of~\eqref{eq:plan1-ge-dH-alpha}. The second inequality is the definition
of \(\alpha_0\) in~\eqref{eq:plan1-ge-alpha0}.
\end{proof}

\begin{lemma}[Barycentre shift absorbable in the threshold]
\label{lem:plan1-ge-bary}
Write \(y_*=y_*(\Otilde)\in\R^n\) for the optimal Fraenkel translation,
i.e.\ a minimiser of \(y\mapsto|\Otilde\Delta(B_1+y)|/|B_1|\). There exists
\(C_{\Fr}'(n,R)\) such that, for every selected minimizer \(\Otilde\) with
\(\BDValpha(\Otilde)\le\alpha_0(n,R)\),
\begin{equation}\label{eq:plan1-ge-bary-bound}
   |y_*|\le(R+1)\,\Fr(\Otilde)
   \le (R+1)\sqrt{C_{\Fr}(n)}\,\BDValpha(\Otilde)^{1/2}
   =: C_{\Fr}'(n,R)\,\BDValpha(\Otilde)^{1/2}.
\end{equation}
After replacing \(\Otilde\) by its translate \(\Otilde-y_*\) (which leaves
the energy unchanged) we may, and from now on do, assume \(y_*=0\), so the
comparison ball is centred at the origin.
\end{lemma}

\begin{proof}
The first inequality of~\eqref{eq:plan1-ge-bary-bound} follows from
\(\Otilde\subset B_R\) and \(|\Otilde\Delta(B_1+y_*)|=|B_1|\Fr(\Otilde)\):
if \(|y_*|>(R+1)\Fr(\Otilde)\) then \(B_1\) and \(B_1+y_*\) are essentially
disjoint, and the trivial bound \(\Fr\le 2\) gives a contradiction. The
second inequality is the first half of~\eqref{eq:plan1-ge-BDVL42}.
\end{proof}

\subsubsection{Step 2: Hausdorff closeness yields fixed-scale flatness}
\label{subsubsec:plan1-ge-S2}

Fix once and for all a universal geometric constant \(C_0=C_0(n)\ge 1\)
absorbing the conversions between the half-slab geometry below and the
exact normalisation of the Alt--Caffarelli class
\(F(\bar\mu,1,+\infty)\) used in \ref{R4}; cf.\
\cite[\S 4]{BDV}.

Define the tubular slope threshold
\begin{equation}\label{eq:plan1-ge-mutub}
   \mu_{\rm tub}:=\tfrac14,
\end{equation}
which will be used in Step~3, and set
\begin{equation}\label{eq:plan1-ge-mu-choices}
   \mu:=\min\!\Bigl(\tfrac{\bar\mu(n,R)}{16\,C_0(n)},
                    \tfrac{\mu_{\rm tub}}{C_{\rm AC}(n,R)}\Bigr),\quad
   \rho_\mu:=\min\!\Bigl(\bar\rho(n,R),\tfrac{\mu}{18},\tfrac{r_d(n,R)}{2}\Bigr),\quad
   h_F:=\tfrac{\mu\,\rho_\mu}{16}.
\end{equation}
All three constants depend only on \((n,R)\).

\begin{lemma}[Geometric slab inclusion for \(\partial B_1\)]
\label{lem:plan1-ge-slab}
For \(\bar x\in\partial B_1\) and \(s\le 6\rho_\mu\),
\[
   \partial B_1\cap B_s(\bar x)
   \subset\{x\in\R^n:|(x-\bar x)\cdot\nu_{\bar x}|\le s^2/2\}
   \subset\{x:|(x-\bar x)\cdot\nu_{\bar x}|\le\mu\rho_\mu\},
\]
where \(\nu_{\bar x}\) is the outward unit normal to \(\partial B_1\) at
\(\bar x\).
\end{lemma}

\begin{proof}
\(B_1\) has unit radius, hence \(|x-\bar x|\le s\) implies
\(|(x-\bar x)\cdot\nu_{\bar x}|\le s^2/2\) for \(x\in\partial B_1\). For
\(s\le 6\rho_\mu\) the bound \(s^2/2\le 18\rho_\mu^2\le\mu\rho_\mu\) is
the condition \(\rho_\mu\le\mu/18\) used to define \(\rho_\mu\) in
\eqref{eq:plan1-ge-mu-choices}.
\end{proof}

\begin{proposition}[Fixed-scale Alt--Caffarelli flatness]
\label{prop:plan1-ge-flatness}
Suppose \(d_H(\partial\Otilde,\partial B_1)\le h_F\). Then for every
\(\bar x\in\partial B_1\) there exists \(x_0\in\partial\Otilde\) with
\(|x_0-\bar x|\le h_F\) such that
\[
   \partial\Otilde\cap B_{4\rho_\mu}(x_0)
   \subset\{x:|(x-x_0)\cdot\nu_{\bar x}|\le 2\mu\rho_\mu\},
\]
and the torsion potential \(\widetilde u\) belongs to the Alt--Caffarelli
class \(F(C_0\mu,1,+\infty)\) on \(B_{4\rho_\mu}(x_0)\) with respect to
the direction \(-\nu_{\bar x}\).
\end{proposition}

\begin{proof}
\emph{Step 1: existence of \(x_0\).} By
\(d_H(\partial\Otilde,\partial B_1)\le h_F\) there is
\(x_0\in\partial\Otilde\) with \(|x_0-\bar x|\le h_F\).

\emph{Step 2: slab inclusion.} Set \(s:=5\rho_\mu+h_F\). Since
\(h_F=\mu\rho_\mu/16\le\rho_\mu\) and \(\mu\le 1\), we have
\(s\le 6\rho_\mu\), so Lemma~\ref{lem:plan1-ge-slab} applies. Every
\(p\in\partial\Otilde\cap B_{5\rho_\mu}(x_0)\) lies within \(h_F\) of
some \(\bar p\in\partial B_1\cap B_s(\bar x)\), hence
\(|(p-\bar x)\cdot\nu_{\bar x}|\le\mu\rho_\mu+h_F\), and moving the
centre to \(x_0\) costs at most \(h_F\), giving
\(|(p-x_0)\cdot\nu_{\bar x}|\le\mu\rho_\mu+2h_F=\mu\rho_\mu+\mu\rho_\mu/8
\le 2\mu\rho_\mu\).

\emph{Step 3 (positive phase, F1).} The rescaled torsion potential
\(\widetilde u_{x_0,\rho_\mu}(x):=\widetilde u(x_0+\rho_\mu x)/\rho_\mu\)
is defined on \(B_4\). The nondegeneracy from \ref{R2} at radius
\(\rho_\mu\le r_d/2\) gives a point \(y\in B_{\rho_\mu}(x_0)\cap\{\widetilde u>0\}\)
with \(\widetilde u(y)\ge c_d\rho_\mu\). Combined with the Lipschitz bound
of \ref{R2}, this gives
\[
   \widetilde u_{x_0,\rho_\mu}(x)\ge c_d\bigl(x\cdot(-\nu_{\bar x})-C_0\mu\bigr)_+
   \qquad\text{on }B_4,
\]
after absorbing dimensional conversions into \(C_0=C_0(n)\). This is the
lower (\emph{positive-phase}) condition of the Alt--Caffarelli class
\(F(C_0\mu,1,+\infty)\).

\emph{Step 4 (zero phase, F2).} The exterior half-ball
\(\{(x-\bar x)\cdot\nu_{\bar x}>\mu\rho_\mu\}\cap B_{4\rho_\mu}(\bar x)\)
is disjoint from \(B_1\). By
\(d_H(\partial\Otilde,\partial B_1)\le h_F<\mu\rho_\mu/16\), it is
disjoint from \(\Otilde\) except possibly in a layer of width \(h_F\).
Hence
\(\{(x-x_0)\cdot\nu_{\bar x}\ge 3\mu\rho_\mu\}\cap B_{4\rho_\mu}(x_0)
\subset\R^n\setminus\Otilde\), so \(\widetilde u\equiv 0\) there. This is
the (F2) condition.

\emph{Step 5 (Lipschitz upper bound, F3).} By \ref{R2},
\(\Lip(\widetilde u)\le L_u\), and zero in the deep zero-layer gives
\(\widetilde u_{x_0,\rho_\mu}(x)\le
L_u(x\cdot(-\nu_{\bar x})+C_0\mu)_+\) on \(B_4\), which is (F3).

The three conditions (F1)--(F3) with parameter \(C_0\mu\le\bar\mu\) (by
the choice of \(\mu\) in~\eqref{eq:plan1-ge-mu-choices}) place
\(\widetilde u\) in the Alt--Caffarelli class
\(F(C_0\mu,1,+\infty)\) on \(B_{4\rho_\mu}(x_0)\) with respect to
\(-\nu_{\bar x}\).
\end{proof}

\subsubsection{Step 3: patching local graphs into a global graph}
\label{subsubsec:plan1-ge-S3}

\begin{lemma}[Tubular threshold]\label{lem:plan1-ge-tubular}
Let \(\mu_{\rm tub}=1/4\) as in~\eqref{eq:plan1-ge-mutub}. If
\(g:\partial B_1\to\R\) satisfies
\(\|g\|_{C^{0,1}(\partial B_1)}\le\mu_{\rm tub}\), then the radial map
\(\Phi_g(\theta):=(1+g(\theta))\theta\) is a bi-Lipschitz embedding of
\(\partial B_1\) onto a hypersurface in \(\R^n\), and the radial
projection from any such hypersurface back onto \(\partial B_1\) is a
diffeomorphism with bounded \(C^{1,\gamma}\) inverse, controlled by an
explicit constant \(K_{\rm geom}(n,R)\).
\end{lemma}

\begin{proof}
Compute \(d\Phi_g=(1+g)\mathrm{id}+\theta\otimes dg\) on \(T_\theta\partial B_1\)
(extended trivially in the radial direction). The smallest singular value
is bounded below by \(\sqrt{(1+g)^2-|dg|^2}\ge\sqrt{(3/4)^2-(1/4)^2}
=\sqrt{1/2}\ge 1/\sqrt2\). Hence \(\Phi_g\) is bi-Lipschitz with bilipschitz
constant \(\le 2\). The radial projection \(\theta\mapsto\theta/|\theta|\)
is smooth on \(\R^n\setminus\{0\}\) with derivatives controlled by
\(\dist(\cdot,0)\); composing with the inverse map of \(\Phi_g\) on the
image \(\Phi_g(\partial B_1)\subset\{|x|\ge 3/4\}\) gives a
diffeomorphism whose \(C^{1,\gamma}\) norm is bounded by an explicit
function of \(n,R,\gamma\), denoted \(K_{\rm geom}(n,R)\).
\end{proof}

\begin{proposition}[Global \(C^{1,\gamma}\) normal graph]
\label{prop:plan1-ge-patching}
Assume \(d_H(\partial\Otilde,\partial B_1)\le h_F\). Then there exists a
unique function \(g:\partial B_1\to\R\) such that
\begin{equation}\label{eq:plan1-ge-graph}
   \partial\Otilde=\bigl\{(1+g(\theta))\theta:\theta\in\partial B_1\bigr\},
\end{equation}
with
\begin{equation}\label{eq:plan1-ge-norm-graph}
   \|g\|_{C^{1,\gamma}(\partial B_1)}\le M_1(n,R)
   :=C_{\rm AC}(n,R)\,\mu\,K_{\rm geom}(n,R),\qquad
   \|g\|_{C^{0,1}(\partial B_1)}\le\mu_{\rm tub}=\tfrac14.
\end{equation}
Here \(\gamma=\gamma(n,R)\) is the Alt--Caffarelli exponent of
\ref{R4}.
\end{proposition}

\begin{proof}
By Proposition~\ref{prop:plan1-ge-flatness}, for every
\(\bar x\in\partial B_1\) there is \(x_0=x_0(\bar x)\in\partial\Otilde\)
such that \(\widetilde u\in F(C_0\mu,1,+\infty)\) on
\(B_{4\rho_\mu}(x_0)\) with respect to \(-\nu_{\bar x}\). Since
\(C_0\mu\le\bar\mu\), the Alt--Caffarelli flatness improvement \ref{R4}
upgrades \(\partial\Otilde\cap B_{\rho_\mu}(x_0)\) to a \(C^{1,\gamma}\)
graph in direction \(\nu_{\bar x}\), with norm \(\le C_{\rm AC}\mu\).

By the choice \(\mu\le\mu_{\rm tub}/C_{\rm AC}\) in
\eqref{eq:plan1-ge-mu-choices}, the local slope is \(\le\mu_{\rm tub}=1/4\),
hence Lemma~\ref{lem:plan1-ge-tubular} applies: the local Cartesian graph
is reparametrised as a spherical normal graph \(g_{\bar x}\) over a
neighbourhood of \(\bar x\) in \(\partial B_1\), with \(C^{1,\gamma}\)
norm increased by at most \(K_{\rm geom}(n,R)\).

Two local graphs \(g_{\bar x_1}\) and \(g_{\bar x_2}\) over overlapping
caps parametrise the same connected portion of \(\partial\Otilde\): by
the slope bound \(\le\mu_{\rm tub}\), the radial projection from
\(\partial\Otilde\) to \(\partial B_1\) is injective, hence the two
parametrisations agree on the overlap. Patching using a finite cover of
\(\partial B_1\) by caps of radius \(\rho_\mu\) yields a unique global
\(g\) with the stated \(C^{1,\gamma}\) and \(C^{0,1}\) bounds.
\end{proof}

\begin{lemma}[\(L^\infty\) control of \(g\)]
\label{lem:plan1-ge-Linfty}
With \(C_H''(n,R):=2\,C_H'(n,R)\) (see \eqref{eq:plan1-ge-CHprime}),
\begin{equation}\label{eq:plan1-ge-Linfty}
   \|g\|_{L^\infty(\partial B_1)}
   \le C_H''(n,R)\,\BDValpha(\Otilde)^{1/(2n)}.
\end{equation}
\end{lemma}

\begin{proof}
For \(\theta\in\partial B_1\), \(|g(\theta)|\) is the radial distance
from \((1+g(\theta))\theta\in\partial\Otilde\) to \(\theta\in\partial B_1\).
Since \(\|g\|_{C^{0,1}}\le 1/4\), the radial distance is at most twice
the Euclidean distance, so
\(\|g\|_{L^\infty}\le 2\,d_H(\partial\Otilde,\partial B_1)
\le 2C_H'\,\BDValpha(\Otilde)^{1/(2n)}=C_H''\,\BDValpha(\Otilde)^{1/(2n)}\)
by Proposition~\ref{prop:plan1-ge-alpha-to-H}.
\end{proof}

\begin{theorem}[Graph-entry theorem]
\label{thm:plan1-ge-graphentry}
There exist explicit positive constants \(\mu(n,R)\), \(\rho_\mu(n,R)\),
\(h_F(n,R)\), \(C_H'(n,R)\), \(C_H''(n,R)\), \(C_{\Fr}'(n,R)\),
\(M_1(n,R)\) and an entry threshold
\begin{equation}\label{eq:plan1-ge-alpha-graph}
   \alpha_{\rm graph}(n,R)
   :=\min\Bigl\{
       \alpha_0(n,R),\,
       (h_F/C_H')^{2n},\,
       (h_F/C_{\Fr}')^2
   \Bigr\}
\end{equation}
such that the following holds. Let \(\Otilde\subset B_R\) be a selected
minimizer (in the sense of Subsection~\ref{subsec:plan1-selection}) with
\(|\Otilde|=|B_1|\), \(\tau\le\tau_{\rm reg}(n,R)\) and
\(\BDValpha(\Otilde)\le\alpha_{\rm graph}(n,R)\). After translating by
\(-y_*\) (Lemma~\ref{lem:plan1-ge-bary}), there exists a unique
\(g:\partial B_1\to\R\) such that
\eqref{eq:plan1-ge-graph}, \eqref{eq:plan1-ge-norm-graph} and
\eqref{eq:plan1-ge-Linfty} hold; in particular
\[
   \|g\|_{C^{1,\gamma}(\partial B_1)}\le M_1(n,R),\qquad
   \|g\|_{L^\infty(\partial B_1)}
   \le C_H''(n,R)\,\BDValpha(\Otilde)^{1/(2n)}.
\]
\end{theorem}

\begin{proof}
The choice~\eqref{eq:plan1-ge-alpha-graph} ensures both
\(\BDValpha(\Otilde)\le\alpha_0(n,R)\) (so
Proposition~\ref{prop:plan1-ge-alpha-to-H} applies) and
\(C_H'\BDValpha^{1/(2n)}\le h_F\); together with the barycentre control
of Lemma~\ref{lem:plan1-ge-bary} this gives
\(d_H(\partial\Otilde,\partial B_1)\le h_F\) (after the translation
\(y_*\to 0\)). Proposition~\ref{prop:plan1-ge-patching} then yields
the global \(C^{1,\gamma}\) graph with the stated bounds, and
Lemma~\ref{lem:plan1-ge-Linfty} gives the \(L^\infty\) bound.
\end{proof}

\paragraph{Discussion of the regularity input.}
The regularity theorem driving Step~3 is the Alt--Caffarelli flatness
improvement in the form recorded by Brasco--De~Philippis--Velichkov as
\cite[Theorem~4.18]{BDV}; in the source this is stated as a black-box
quantitative replacement of the original Alt--Caffarelli theorem
\cite{AC1981}, valid for the BDV-class of one-phase Bernoulli
free boundaries with a smooth bounded coefficient
\((q_-\le q\le q_+,\ \|q\|_{C^{1,\gamma}}\le L_q)\). \emph{It is not the
De Giorgi--Almgren \(\eps\)-regularity theorem for area-minimisers, nor is
it the Tamanini quasi-minimality regularity}: although the selected
minimizer \(\Otilde\) is a quasi-minimizer of perimeter (Agent~5,
D1.4), what we use here is the stronger one-phase Bernoulli flatness
improvement of BDV, which is tailored to the penalised functional
\(J_{\Omega_0}\). This input is recorded in the source map as item
D2.2 with citation \cite[Theorem~4.18]{BDV} and origin
\cite{AC1981}.

\subsubsection{Step 4: Schauder bootstrap to \(C^{m,\gamma}\) regularity}
\label{subsubsec:plan1-ge-S4}

Let \(\widetilde u\in H^1_0(\Otilde)\) be the torsion potential of
\(\Otilde\), and let \(\bar x\in\partial B_1\) be fixed. By a translation
and rotation we may assume \(\bar x=0\) and the outward unit normal to
\(\partial B_1\) at \(\bar x\) is \(e_n\). By
Proposition~\ref{prop:plan1-ge-patching} there is \(R_0=R_0(n,R)>0\) and a
Cartesian function \(\widetilde g:\{|x'|<R_0\}\to\R\) with
\(\|\widetilde g\|_{C^{1,\gamma}}\le M_1'(n,R)\), where
\(M_1'(n,R):=C(n,R)M_1(n,R)\) absorbs the conversion between spherical
graphs over \(\partial B_1\) and Cartesian graphs in a chart, such that
inside the cylinder
\(\mathcal C_{R_0}:=\{|x'|<R_0,|x_n|<R_0\}\),
\[
   \partial\Otilde\cap\mathcal C_{R_0}
   =\{x_n=\widetilde g(x'):|x'|<R_0\},\quad
   \Otilde\cap\mathcal C_{R_0}
   =\{x_n>\widetilde g(x')\}.
\]

\paragraph{Hodograph straightening.}
By the Lipschitz nondegeneracy \ref{R2} and the positive Bernoulli
coefficient \ref{R3}, in a one-sided neighbourhood of \(0\) the partial
derivative \(\partial_n\widetilde u\) is bounded below by \(q_-/2>0\).
The Kinderlehrer--Nirenberg partial hodograph map
\cite[\S 2]{KN1977}
\[
   \Phi(x',x_n):=(x',t),\qquad t:=\widetilde u(x',x_n),
\]
is therefore a \(C^{1,\gamma}\) diffeomorphism from
\(\Otilde\cap\mathcal C_{R_0/2}\) onto a domain
\(\{(x',t):|x'|<R_0/2,\ 0<t<T_0\}\) for some \(T_0=T_0(n,R)>0\), with
inverse \(x_n=\psi(x',t)\) satisfying \(\psi(x',0)=\widetilde g(x')\).

A direct change of variables, identical to
\cite[(2.10)--(2.12)]{KN1977}, gives that \(\psi\) satisfies, in
\(\{0<t<T_0\}\), the quasilinear PDE
\begin{equation}\label{eq:plan1-ge-psi-PDE}
   F[\psi]
     :=\frac{1+|\nabla_{x'}\psi|^{2}}{(\partial_t\psi)^{2}}\,\partial_t^{2}\psi
        -2\,\frac{\nabla_{x'}\psi\cdot\nabla_{x'}\partial_t\psi}{\partial_t\psi}
        +\Delta_{x'}\psi
     =\partial_t\psi,
\end{equation}
which is uniformly elliptic in \(\psi\) as long as
\(\partial_t\psi\ge\eta>0\) and \(|\nabla_{x'}\psi|\le K\); from
\ref{R2}--\ref{R3} these bounds hold with constants \(\eta_*(n,R)\) and
\(K_*(n,R)\). The free-boundary Bernoulli condition
\(|\nabla\widetilde u|^2=q^2\) at \(t=0\) transforms into
the nonlinear oblique boundary condition
\begin{equation}\label{eq:plan1-ge-BC-psi}
   \partial_t\psi(x',0)
   =\sqrt{\frac{1+|\nabla_{x'}\psi(x',0)|^{2}}{q(x',\psi(x',0))^{2}}}
   \qquad\text{on }\{t=0\}.
\end{equation}
Linearising
\(B[\psi]:=\partial_t\psi-G(x',\nabla_{x'}\psi,\psi)=0\) gives
\(\partial B/\partial(\partial_t\psi)\equiv 1\): the obliqueness
coefficient \(\beta_n=1\) is uniformly positive in Lieberman's
notation \cite[\S 1]{Li86b}, and the tangential coefficients are
bounded by an \((n,R,L_q)\)-constant. The initial bound
\(\|\psi\|_{C^{1,\gamma}}\le C(n,R)\) on \(\{0\le t\le T_0/2\}\)
follows from \ref{R2}, the implicit function theorem on
\(t=\widetilde u(x',x_n)\) and the chain rule.

\paragraph{Schauder bootstrap.}
We apply Lieberman's oblique-derivative Schauder theorem for nonlinear
elliptic equations with smooth coefficients, in the form stated as
\cite[Theorem~1]{Li86b}:

\begin{theorem}[Lieberman \cite{Li86b}, Thm.~1]
\label{thm:plan1-ge-lieberman}
Let \(v\) solve
\(a^{ij}(x,v,\nabla v)\partial_{ij}v+b(x,v,\nabla v)=0\) in
\(B_{R_0/2}^+\), \(b^i(x,v,\nabla v)\partial_i v+b_0(x,v)=0\) on
\(\{x_n=0\}\), with coefficients of class \(C^{m-1,\gamma}\), ellipticity
\(\eta_*\), obliqueness \(\beta_0\), and \(\|v\|_{C^{1,\gamma}}\le N_0\).
For every integer \(m\ge 2\),
\[
   \|v\|_{C^{m,\gamma}(B_{R_0/4}^+)}
   \le\mathcal S_m\bigl(n,\eta_*,\beta_0,N_0,\|\mathrm{coeffs}\|_{C^{m-1,\gamma}}\bigr).
\]
\end{theorem}

Setting \(\eta_m:=\|\psi\|_{C^{m,\gamma}(\{0\le t\le T_0/2\})}\), the
PDE~\eqref{eq:plan1-ge-psi-PDE} and the oblique
BC~\eqref{eq:plan1-ge-BC-psi} have coefficients of class \(C^{m-1,\gamma}\)
once \(\psi\) is of class \(C^{m-1,\gamma}\), with norms bounded by
\(L_q,\eta_m,n,R\). Applying Theorem~\ref{thm:plan1-ge-lieberman}
iteratively on shrinking half-cylinders gives the recursion
\begin{equation}\label{eq:plan1-ge-recursion}
   \eta_{m+1}\le\mathcal S_{m+1}\bigl(n,\eta_*,1,\eta_m,L_q\bigr)
   =:F_m(\eta_m,L_q,n,R),
\end{equation}
with \(\eta_1=C(n,R)\). Inductively define
\begin{equation}\label{eq:plan1-ge-Mm}
   M_m(n,R):=F_{m-1}\bigl(F_{m-2}\bigl(\cdots F_1(\eta_1,L_q,n,R)\cdots\bigr),L_q,n,R\bigr).
\end{equation}

\paragraph{Norm transfer back to \(g\).}
The graph \(g\) on \(\partial B_1\) and the Cartesian \(\widetilde g\) on
\(\{|x'|<R_0\}\) are related by an \((n,R)\)-bilipschitz \(C^\infty\)
chart, so for every \(k\ge 0\) and \(0<\beta\le 1\),
\begin{equation}\label{eq:plan1-ge-norm-transfer}
   c(n,R,k,\beta)\,\|g\|_{C^{k,\beta}(\partial B_1)}
   \le\|\widetilde g\|_{C^{k,\beta}}
   \le C(n,R,k,\beta)\,\|g\|_{C^{k,\beta}(\partial B_1)}.
\end{equation}
Since \(\widetilde g(x')=\psi(x',0)\) by construction,
\(\|\widetilde g\|_{C^{k,\beta}}\le\|\psi\|_{C^{k,\beta}(\{0\le t\le T_0/2\})}
\le\eta_k\). Combining with~\eqref{eq:plan1-ge-norm-transfer} and using a
finite cover of \(\partial B_1\) by \(N_0(n,R)\) charts gives the
following.

\begin{proposition}[Uniform Schauder bound]
\label{prop:plan1-ge-schauder-bootstrap}
For every integer \(m\ge 1\) there is \(M_m(n,R)<\infty\), given by
\eqref{eq:plan1-ge-Mm}, such that under the hypotheses of
Theorem~\ref{thm:plan1-ge-graphentry},
\begin{equation}\label{eq:plan1-ge-schauder-bound}
   \|g\|_{C^{m,\gamma}(\partial B_1)}\le M_m(n,R).
\end{equation}
\end{proposition}

\subsubsection{Step 5: H\"older interpolation to small \(C^{2,\gamma_0}\)}
\label{subsubsec:plan1-ge-S5}

\begin{lemma}[H\"older interpolation on \(\partial B_1\)]
\label{lem:plan1-ge-interp}
Let \(0\le k<m\) be integers, \(0<\gamma_0\le\gamma\le 1\), with
\((k,\gamma_0)\neq(m,\gamma)\), and set
\begin{equation}\label{eq:plan1-ge-theta}
   \theta_{m,k}:=\frac{k+\gamma_0}{m+\gamma}\in(0,1).
\end{equation}
There is \(C_{\rm GN}=C_{\rm GN}(n,m,k,\gamma,\gamma_0)>0\) such that for
every \(f\in C^{m,\gamma}(\partial B_1)\),
\begin{equation}\label{eq:plan1-ge-interp}
   \|f\|_{C^{k,\gamma_0}(\partial B_1)}
   \le C_{\rm GN}\,
   \|f\|_{L^\infty(\partial B_1)}^{1-\theta_{m,k}}\,
   \|f\|_{C^{m,\gamma}(\partial B_1)}^{\theta_{m,k}}.
\end{equation}
\end{lemma}

\begin{proof}
This is the standard Gagliardo--Nirenberg / Bergh--L\"ofstr\"om convexity
inequality for H\"older spaces, transferred from Euclidean charts to the
compact manifold \(\partial B_1\); see
\cite[Lemma~6.32]{GT1983} and
\cite[Theorem~6.4.5]{BL1976}. The dimension dependence enters
through the partition of unity and the number of charts.
\end{proof}

We apply~\eqref{eq:plan1-ge-interp} with \(k=2\), \(m=5\) and
\(\gamma_0\in(0,\gamma)\). Since \(\gamma_0\le\gamma\le 1\),
\[
   \theta:=\theta_{5,2}=\frac{2+\gamma_0}{5+\gamma}
   \le\frac{2+\gamma}{5+\gamma}\le\frac12,
\]
so
\begin{equation}\label{eq:plan1-ge-theta-half}
   \theta\le\tfrac12,\qquad 1-\theta\ge\tfrac12.
\end{equation}

\begin{theorem}[Schauder--interpolation entry; threshold
\(\alpha_{\rm sph}\)]
\label{thm:plan1-ge-sph-entry}
Let \(\gamma_0\in(0,\gamma)\). Let \(\delta_{\rm sph}=\delta_{\rm sph}(n,\gamma_0)>0\)
denote the smallness threshold of \cite[Theorem~3.3]{BDV} (the BDV
nearly-spherical theorem; see Subsection~\ref{subsec:plan1-ns} for the
exact statement). Define
\begin{equation}\label{eq:plan1-ge-alpha-sph}
   \alpha_{\rm sph}(n,R,\gamma_0)
   :=\Bigl(
       \frac{\delta_{\rm sph}(n,\gamma_0)}
            {C_{\rm GN}(n,5,2,\gamma,\gamma_0)\,
             C_H''(n,R)^{1-\theta}\,
             M_5(n,R)^{\theta}}
     \Bigr)^{\!2n/(1-\theta)}.
\end{equation}
Then for every selected minimizer \(\Otilde\) with
\(\BDValpha(\Otilde)\le\min\{\alpha_{\rm graph}(n,R),\alpha_{\rm sph}(n,R,\gamma_0)\}\)
the graph \(g\) of Theorem~\ref{thm:plan1-ge-graphentry} satisfies
\begin{equation}\label{eq:plan1-ge-smallC2}
   \|g\|_{C^{2,\gamma_0}(\partial B_1)}\le\delta_{\rm sph}(n,\gamma_0).
\end{equation}
\end{theorem}

\begin{proof}
By Theorem~\ref{thm:plan1-ge-graphentry} and
Proposition~\ref{prop:plan1-ge-schauder-bootstrap} (applied with
\(m=5\)),
\[
   \|g\|_{C^{2,\gamma_0}}
   \le C_{\rm GN}\,\|g\|_{L^\infty}^{1-\theta}\,
                  \|g\|_{C^{5,\gamma}}^{\theta}
   \le C_{\rm GN}\,M_5(n,R)^{\theta}\,
                  \|g\|_{L^\infty}^{1-\theta}.
\]
Insertion of the \(L^\infty\) bound \eqref{eq:plan1-ge-Linfty} gives
\[
   \|g\|_{C^{2,\gamma_0}}
   \le C_{\rm GN}\,M_5^{\theta}\,
                  C_H''^{\,1-\theta}\,
                  \BDValpha(\Otilde)^{(1-\theta)/(2n)}.
\]
The bound \(\|g\|_{C^{2,\gamma_0}}\le\delta_{\rm sph}\) holds iff
\(\BDValpha(\Otilde)^{(1-\theta)/(2n)}\le
\delta_{\rm sph}/(C_{\rm GN}M_5^{\theta}C_H''^{\,1-\theta})\), which is
exactly \(\BDValpha(\Otilde)\le\alpha_{\rm sph}(n,R,\gamma_0)\) by
\eqref{eq:plan1-ge-alpha-sph}. The hypothesis
\(\BDValpha(\Otilde)\le\alpha_{\rm graph}\) ensures
Theorem~\ref{thm:plan1-ge-graphentry} and
Proposition~\ref{prop:plan1-ge-schauder-bootstrap} apply in the first
place.
\end{proof}

\paragraph{Remarks on the threshold and the exponent.}
\begin{enumerate}
\item By~\eqref{eq:plan1-ge-theta-half}, \(1-\theta\ge 1/2\), so the
exponent \(2n/(1-\theta)\) in~\eqref{eq:plan1-ge-alpha-sph} is bounded by
\(4n\) uniformly in \(\gamma,\gamma_0\). The \(L^\infty\)-power
\(1-\theta\ge 1/2\) is what preserves the sharp \(1/2\)-exponent
downstream in the nearly-spherical theorem; cf.\
\cite[Theorem~6.4.5]{BL1976}.
\item Any integer \(m\ge 3\) for which \(\theta_{m,2}<1\) is admissible
for the interpolation; the choice \(m=5\) makes \(\theta\le 1/2\)
explicit and uniform in \((\gamma_0,\gamma)\).
\item The whole argument is constructive: no compactness, no
contradiction. The threshold
\(\alpha_{\rm small}^{\rm graph}(n,R,\gamma_0):=\min\{\alpha_{\rm graph}(n,R),
\alpha_{\rm sph}(n,R,\gamma_0)\}\) is explicit in the BDV regularity
constants of the package \ref{R1}--\ref{R4} and the
Kinderlehrer--Nirenberg / Lieberman / Gagliardo--Nirenberg constants.
\end{enumerate}

\subsubsection{Conclusion of the graph-entry subsection}

The output of this subsection is the following theorem, used by
Subsection~\ref{subsec:plan1-ns} (Agent~7) as input to the BDV
nearly-spherical theorem.

\begin{theorem}[Graph entry and regularity, Plan~1]
\label{thm:plan1-ge-main}
Let \(n\ge 2\), \(R\ge 2\), and \(\gamma_0\in(0,\gamma(n,R))\). There
exist explicit positive constants \(\tau_{\rm reg}(n,R)\),
\(\alpha_{\rm graph}(n,R)\) (defined by \eqref{eq:plan1-ge-alpha-graph}),
\(\alpha_{\rm sph}(n,R,\gamma_0)\) (defined by
\eqref{eq:plan1-ge-alpha-sph}), \(M_m(n,R)\) for every
\(m\ge 1\) (defined by \eqref{eq:plan1-ge-Mm}), and \(C_H''(n,R)\) such
that the following holds. Let \(\Otilde\subset B_R\) be a selected
minimizer with \(|\Otilde|=|B_1|\), \(\tau\le\tau_{\rm reg}(n,R)\) and
\[
   \BDValpha(\Otilde)\le\alpha_{\rm small}^{\rm graph}(n,R,\gamma_0)
   :=\min\bigl\{\alpha_{\rm graph}(n,R),\alpha_{\rm sph}(n,R,\gamma_0)\bigr\}.
\]
Then, after translating by the optimal Fraenkel translation
(Lemma~\ref{lem:plan1-ge-bary}), there exists a unique
\(g:\partial B_1\to\R\) with
\[
   \partial\Otilde=\bigl\{(1+g(\theta))\theta:\theta\in\partial B_1\bigr\},
\]
\[
   \|g\|_{L^\infty(\partial B_1)}\le C_H''(n,R)\,\BDValpha(\Otilde)^{1/(2n)},\quad
   \|g\|_{C^{m,\gamma}(\partial B_1)}\le M_m(n,R)\quad(m\ge 1),\quad
   \|g\|_{C^{2,\gamma_0}(\partial B_1)}\le\delta_{\rm sph}(n,\gamma_0).
\]
\end{theorem}

\begin{proof}
Combine Theorem~\ref{thm:plan1-ge-graphentry},
Proposition~\ref{prop:plan1-ge-schauder-bootstrap}, and
Theorem~\ref{thm:plan1-ge-sph-entry}.
\end{proof}

\paragraph{Constants used in this subsection.}
The constants explicitly defined and tracked in this subsection are
\(c_d,r_d,L_u,q_\pm,L_q\) (BDV regularity package, \(n,R\)-dependent);
\(C_{\Fr}(n)\), \(C_\alpha(n)\) (BDV Lemma~4.2, dimensional);
\(\bar\mu,\bar\rho,\gamma,C_{\rm AC}\) (Alt--Caffarelli flatness,
\(n,R\)-dependent); \(C_0(n)\) (geometric conversion);
\(c_*,C_H,C_H',C_H''\) (Hausdorff bounds, \(n,R\)-dependent);
\(C_{\Fr}',\alpha_0,\alpha_{\rm graph}\) (entry thresholds,
\(n,R\)-dependent); \(\mu,\rho_\mu,h_F,\mu_{\rm tub},K_{\rm geom},M_1\)
(graph-patching, \(n,R\)-dependent); \(R_0,T_0,\eta_*,K_*,M_m\)
(hodograph / Schauder, \(n,R\)-dependent); \(C_{\rm GN},\theta,\delta_{\rm sph},
\alpha_{\rm sph}\) (interpolation, \((n,R,\gamma_0)\)-dependent). No
\(C\) symbol is used without declared dependence.

\paragraph{What feeds the next subsection.}
Subsection~\ref{subsec:plan1-ns} (Agent~7) consumes
Theorem~\ref{thm:plan1-ge-main} via the BDV nearly-spherical theorem
\cite[Theorem~3.3]{BDV}, applied to the graph \(g\), yielding the
quadratic lower bound
\(E(\Otilde)-E(B_1)\ge c_{\rm sph}(n)\,\|g\|_{H^{1/2}(\partial B_1)}^2\)
and, after BDV Lemma~4.2(i), the bound
\(E(\Otilde)-E(B_1)\ge c_*^{\rm NS}(n)\,\BDValpha(\Otilde)\) in the small
asymmetry regime \(\BDValpha(\Otilde)\le\alpha_{\rm small}(n,R,\gamma_0)\).
The far regime \(\BDValpha(\Otilde)\ge\alpha_{\rm small}\) is handled by
BDV Lemma~5.2 (rearrangement-based) in Subsection~\ref{subsec:plan1-ns}.

% UNIFIED 2026-05-13
% SOURCE MAP:
%   - Plan 1/route-A-summary.tex
%       Sections 5 (closure by BDV nearly spherical), 6 (the
%       non-contradictory quantitative argument), 7 (extension to all
%       asymmetries) supply the assembly logic. We CONVERT the source
%       dimension symbol N -> n per Manuscript/00_notation_register.md.
%   - Final/NearlySphericalClosure.tex
%       Theorem 2.2 (NS), Corollary 2.4 (vol/barycenter absorption),
%       Lemmas 1.1, 1.2 (BDV Thm 3.3 and Lem 4.2(i),(iii)).
%   - Final/SaintVenantAssembly.tex
%       Theorems 2.1, 4.1 (small-alpha linear bound, sharp Saint-Venant),
%       Lemmas 1.1, 1.2, 1.3 (BDV Lem 4.2(i),(ii), Lem 5.2), and
%       Corollary 3.1 / Lemma 3.2 (conversion alpha <-> Fr).
%   - Plan 1/1306.0392v1.pdf
%       Brasco-De Philippis-Velichkov, "Faber-Krahn inequalities in sharp
%       quantitative form", arXiv:1306.0392v1 (2013).
%       Used as black boxes: Definition 4.1 (smooth asymmetry),
%       Theorem 3.3 (sharp nearly spherical), Lemma 4.2 (i),(ii),(iii),
%       Lemma 5.2 (suboptimal Fraenkel-quartic Saint-Venant),
%       Lemmas 4.9, 4.12, 4.16, Theorem 4.18 (regularity package used
%       upstream via Agents 5-6; not re-invoked here).
%   - Manuscript/00_notation_register.md
%       Sections 1-6, 14: canonical n, R, B^*, T(.), E(.), delta_T,
%       Fraenkel \mathcal A, BDV smooth alpha (renamed \alpha_BDV in
%       passages where coarea alpha(t) could collide), constants tracked
%       with explicit (n,R) dependence.
%   - Manuscript/source_map.md
%       Entries D3.1-D3.8 (nearly spherical closure and bounded sharp
%       Saint-Venant), with upstream dependencies D1.x (selection,
%       Agent 5) and D2.x (graph entry / Schauder, Agent 6).
%
% UPSTREAM HYPOTHESES (Agents 5 and 6) USED VERBATIM:
%   - (S) Quantitative selection theorem with constants
%         C_sel(n,R), q_sel(n,R), eps_sel(n,R), tau_reg(n,R).
%   - (G) Graph entry: alpha_graph(n,R), C_H''(n,R), gamma in (0,1).
%   - (I) Schauder/interpolation entry into the C^{2,gamma_0} ball:
%         alpha_sph(n,R,gamma_0), delta_sph(n,gamma_0).
%
% This subsection produces the output the next agent (Agent 8) needs:
%         bounded sharp Saint-Venant inequality
%            E(Omega)-E(B^*) >= c_SV(n,R) \mathcal A(Omega)^2
%         for Omega open, Omega \subset B_R, |Omega| = |B_1| = omega_n,
%         equivalently
%            \mathcal A(Omega)^2 <= C(n,R) delta_T(Omega).

% ---------- Subsection: Nearly-spherical closure and bounded sharp Saint--Venant ----------

\subsection{Nearly-spherical closure and bounded sharp Saint--Venant}
\label{sec:plan1-nearly-spherical}

This subsection closes the chain
\[
\text{selection}\ \longrightarrow\ \text{graph entry}\ \longrightarrow\
\text{Schauder/interpolation entry}\ \longrightarrow\
\text{BDV nearly spherical}
\]
and converts the closing inequality into the bounded sharp
Saint--Venant estimate
\begin{equation}\label{eq:plan1-target}
   \Fr(\Om)^{2}\ \le\ C(n,R)\,\deltaT(\Om),
   \qquad \Om\subset B_{R},\ |\Om|=\omega_{n},
\end{equation}
which is the input to the bounded reduction and the Kohler--Jobin
transfer of Agent~8. The argument is \emph{not} by contradiction: every
quantitative step is a substitution of upstream constants into a chain
of explicit BDV black boxes. The constants are tracked through the
section.

Throughout the subsection $n\ge 2$ and $R\ge 2$ are fixed, all sets are
bounded open subsets of $\mathbb{R}^{n}$, and we keep the notation
register conventions. (Domain-class reconciliation. The selection
principle of Agent~5 produces a quasi-open minimiser
$U_*\subset\overline{B_R}$ in the BDV sense
\cite[\S2]{BDV}; after the BDV regularity package
\cite[Lem.~4.9,~4.12,~4.16,~Thm.~4.18]{BDV} consumed by Agent~6,
$U_*$ is open up to a set of zero capacity, and all load-bearing
functionals $\deltaT,\Fr,\BDValpha,\lambda_1$ are insensitive to
capacity-zero modifications. Consequently the quasi-open class of
Agent~5 and the open class of the notation register \S2 coincide for
every quantity used below; we treat $U_*$ as an open subset of
$\overline{B_R}$ without further mention.)
The notation register conventions are: $E(\Om)=-T(\Om)/2$ is the BDV-normalised
Saint--Venant energy, $\deltaT(\Om)=E(\Om)-E(B^{*})$ is the Saint--Venant
deficit at fixed volume, $\Fr(\Om)$ is the Fraenkel asymmetry, and the
BDV smooth asymmetry is
\[
\BDValpha(\Om)\ :=\ \int_{\Om\,\triangle\, B_{1}(x_{\Om})}
   \bigl|\,1-|x-x_{\Om}|\,\bigr|\,dx,
   \qquad
   x_{\Om}=\frac{1}{|\Om|}\int_{\Om} x\,dx
\]
(\cite[Definition~4.1]{BDV}, with $n$ in place of the source
symbol $N$). The torsion quotient at the BDV smooth asymmetry is
\[
Q_{\BDValpha}(\Om)\ :=\ \frac{E(\Om)-E(B_{1})}{\BDValpha(\Om)},
   \qquad |\Om|=\omega_{n},\ \BDValpha(\Om)>0.
\]

\subsubsection{The three BDV black boxes}
\label{sec:plan1-bdv-blackboxes}

The closure uses exactly three results from
\cite{BDV}, all of which are proved there constructively (no
compactness, no selection).

\paragraph{(B1) Sharp nearly-spherical stability \textnormal{(\cite[Theorem~3.3]{BDV})}.}
For every $0<\gamma_{0}\le 1$ there exists
$\delta_{1}=\delta_{1}(n,\gamma_{0})>0$ such that the following holds.
Let $U\subset\mathbb{R}^{n}$ be open and bounded, and suppose there
exists $g\in C^{2,\gamma_{0}}(\partial B_{1})$ with
\begin{equation}\label{eq:plan1-ns-hyp}
   \partial U=\bigl\{(1+g(y))\,y:\,y\in\partial B_{1}\bigr\},
   \quad
   \|g\|_{C^{2,\gamma_{0}}(\partial B_{1})}\le\delta_{1},
   \quad
   |U|=|B_{1}|,
   \quad
   x_{U}=0.
\end{equation}
Then, with $E$ the Saint--Venant energy of
Notation~Register~\S3 and $H(g)$ the harmonic extension of $g$ to
$B_{1}$,
\begin{equation}\label{eq:plan1-BDV33}
   E(U)-E(B_{1})\ \ge\
   c_{\rm sph}(n)\,\|g\|_{H^{1/2}(\partial B_{1})}^{2},
   \qquad
   c_{\rm sph}(n):=\frac{1}{32\,n^{2}},
\end{equation}
where $\|g\|_{H^{1/2}(\partial B_{1})}^{2}=\int_{\partial B_{1}}g^{2}\,d\mathcal{H}^{n-1}+\int_{B_{1}}|\nabla H(g)|^{2}\,dx$
(\cite[Definition~3.2]{BDV}). The constant $1/(32n^{2})$ comes from
the gap to the second Steklov eigenvalue of $B_{1}$ (equal to $2$ in
every dimension), see \cite[(3.9)--(3.10)]{BDV}.

\paragraph{(B2) Smooth-asymmetry comparisons \textnormal{(\cite[Lemma~4.2 (i),(ii),(iii)]{BDV})}.}
There exist $C_{1}=C_{1}(n)>0$, $C_{2}=C_{2}(R)>0$,
$\delta_{3}=\delta_{3}(n)>0$, $C_{3}=C_{3}(n)>0$ such that:
\begin{enumerate}
\item[\textup{(i)}]
   for every bounded open $\Om\subset\mathbb{R}^{n}$ with
   $|\Om|=|B_{1}|$,
   \begin{equation}\label{eq:plan1-42i}
      C_{1}(n)\,\BDValpha(\Om)\ \ge\ |\Om\triangle B_{1}(x_{\Om})|^{2}
      \ \ge\ |B_{1}|^{2}\Fr(\Om)^{2};
   \end{equation}
\item[\textup{(ii)}]
   for every $R\ge 2$ and every bounded open
   $\Om_{1},\Om_{2}\subset B_{R}$,
   \begin{equation}\label{eq:plan1-42ii}
      \bigl|\BDValpha(\Om_{1})-\BDValpha(\Om_{2})\bigr|
      \ \le\ C_{2}(R)\,|\Om_{1}\triangle\Om_{2}|;
   \end{equation}
   in particular, choosing $\Om_{2}$ to be the Fraenkel-optimal ball for
   $\Om_{1}$ (so $\BDValpha(\Om_{2})=0$), one has, for
   $|\Om|=\omega_{n}$ and $\Om\subset B_{R}$,
   \begin{equation}\label{eq:plan1-42ii-Fr}
      \BDValpha(\Om)\ \le\ C_{2}(R)\,\omega_{n}\,\Fr(\Om);
   \end{equation}
\item[\textup{(iii)}]
   for every nearly-spherical set $U$ parametrised by $g$ with
   $\|g\|_{L^{\infty}(\partial B_{1})}\le\delta_{3}$,
   \begin{equation}\label{eq:plan1-42iii}
      \BDValpha(U)\ \le\ C_{3}(n)\,\|g\|_{L^{2}(\partial B_{1})}^{2}.
   \end{equation}
\end{enumerate}
All three estimates are proved in \cite[\S4]{BDV} by direct
rearrangement / Taylor expansion of $\BDValpha$; they are constructive.

\paragraph{(B3) Suboptimal Fraenkel-quartic Saint--Venant
\textnormal{(\cite[Lemma~5.2]{BDV})}.}
There exists $C_{8}=C_{8}(n)>0$ such that for every bounded open
$\Om\subset\mathbb{R}^{n}$ with $0<|\Om|<\infty$,
\begin{equation}\label{eq:plan1-52-raw}
   E(\Om)\,|\Om|^{-(n+2)/n}-E(B_{1})\,|B_{1}|^{-(n+2)/n}
   \ \ge\ \frac{1}{C_{8}(n)}\,\Fr(\Om)^{4}.
\end{equation}
The proof in \cite[\S5]{BDV} is by rearrangement and
P\'olya--Szeg\H{o} (no contradiction, no selection). After absorbing the
$\omega_{n}^{(n+2)/n}$ factor at fixed volume $|\Om|=\omega_{n}$,
\begin{equation}\label{eq:plan1-52-fixed}
   E(\Om)-E(B_{1})\ \ge\
   C_{8}'(n)^{-1}\,\Fr(\Om)^{4},
   \qquad C_{8}'(n):=C_{8}(n)/\omega_{n}^{(n+2)/n}.
\end{equation}

\medskip
\textbf{Important note on hypotheses.} (B1) is stated for sets that are
\emph{simultaneously} $C^{2,\gamma_{0}}$-parametrised, volume-normalised
to $|B_{1}|$, and barycentred at the origin. The three normalisations
are all load-bearing for~\eqref{eq:plan1-BDV33}: dropping any one of
them changes the form of the second-variation calculation in
\cite[\S3]{BDV}. Sections~\ref{sec:plan1-feed-graph}
and~\ref{sec:plan1-vol-bary} show that the upstream chain enforces all
three.

\subsubsection{Trace--Poincar\'e step and the pointwise
\texorpdfstring{$\alpha$}{alpha}-form}
\label{sec:plan1-alpha-form}

Combining (B1) and (B2)(iii) gives the pointwise
$Q_{\BDValpha}$-floor on nearly-spherical sets. The intermediate
$H^{1/2}\!\to L^{2}$ comparison is the trace--Poincar\'e inequality on
$\partial B_{1}$:
\begin{equation}\label{eq:plan1-poincare}
   \|g\|_{L^{2}(\partial B_{1})}^{2}\ \le\
   \|g\|_{H^{1/2}(\partial B_{1})}^{2},
\end{equation}
which is immediate from the definition of the $H^{1/2}$-norm (its $L^{2}$
summand bounds $\|g\|_{L^{2}}^{2}$ above by itself).

\begin{theorem}[Nearly-spherical closure, BDV smooth-asymmetry form]
\label{thm:plan1-NS}
Fix $0<\gamma_{0}\le 1$ and let
\begin{equation}\label{eq:plan1-deltasph}
   \delta_{\rm sph}(n,\gamma_{0})\ :=\ \min\bigl(\delta_{1}(n,\gamma_{0}),\,\delta_{3}(n)\bigr).
\end{equation}
Let $U\subset\mathbb{R}^{n}$ be open and bounded, with
$\partial U=\{(1+g(y))y:y\in\partial B_{1}\}$,
$g\in C^{2,\gamma_{0}}(\partial B_{1})$, and
\begin{equation}\label{eq:plan1-NS-hyp}
   \|g\|_{C^{2,\gamma_{0}}(\partial B_{1})}\le\delta_{\rm sph}(n,\gamma_{0}),
   \qquad |U|=|B_{1}|,
   \qquad x_{U}=0.
\end{equation}
Then
\begin{equation}\label{eq:plan1-NS-alpha}
   E(U)-E(B_{1})\ \ge\ c_{*}(n)\,\BDValpha(U),
   \qquad
   c_{*}(n)\ :=\ \frac{c_{\rm sph}(n)}{C_{3}(n)}\ =\
   \frac{1}{32\,n^{2}\,C_{3}(n)},
\end{equation}
and consequently $Q_{\BDValpha}(U)\ge c_{*}(n)$.
\end{theorem}

\begin{proof}
By~\eqref{eq:plan1-deltasph} both
hypotheses~\eqref{eq:plan1-ns-hyp} (with $\delta_{1}$) and the smallness
hypothesis of (B2)(iii) (with $\delta_{3}$) hold simultaneously, since
$\|g\|_{L^{\infty}}\le\|g\|_{C^{2,\gamma_{0}}}\le\delta_{\rm sph}\le\delta_{3}$.
Chain (B1), \eqref{eq:plan1-poincare} and (B2)(iii):
\[
   E(U)-E(B_{1})
   \overset{\eqref{eq:plan1-BDV33}}{\ge}
   c_{\rm sph}(n)\,\|g\|_{H^{1/2}}^{2}
   \overset{\eqref{eq:plan1-poincare}}{\ge}
   c_{\rm sph}(n)\,\|g\|_{L^{2}}^{2}
   \overset{\eqref{eq:plan1-42iii}}{\ge}
   \frac{c_{\rm sph}(n)}{C_{3}(n)}\,\BDValpha(U).
\]
Dividing by $\BDValpha(U)>0$ gives $Q_{\BDValpha}(U)\ge c_{*}(n)$.
\end{proof}

Theorem~\ref{thm:plan1-NS} is the pointwise (per-set) closure: no
sequence, no limit, no smallness of $Q_{\BDValpha}$ is assumed. It is
exactly the building block called $\NS$ in
\texttt{Final/NearlySphericalClosure.tex}.

\subsubsection{Feeding the upstream chain into the
nearly-spherical hypothesis}
\label{sec:plan1-feed-graph}

We now show that the upstream blocks (S), (G), (I) of Agents~5 and~6
produce, from \emph{any} $\Om_{0}\subset B_{R}$ with $|\Om_{0}|=\omega_{n}$
and $\BDValpha(\Om_{0})$ small, a set $\Otilde$ that meets all three
hypotheses of~\eqref{eq:plan1-NS-hyp}.

\paragraph{Recap of the upstream output (Agents 5, 6).}
We use the following from Agent~5 (selection) and Agent~6
(graph entry / Schauder interpolation), each stated with the constants
they fix:

\begin{itemize}
\item[(S)] (Quantitative selection, Agent~5, Theorem D1.2--D1.5 of
   \texttt{Manuscript/source\_map.md}; \texttt{Plan 1/quantitative-selection-principle.tex},
   \texttt{Final/SelectionPrinciple.tex}.)
   There exist $C_{\rm sel}=C_{\rm sel}(n,R)>0$,
   $q_{\rm sel}=q_{\rm sel}(n,R)>0$, $\eps_{\rm sel}=\eps_{\rm sel}(n,R)>0$
   such that for every open $\Om_{0}\subset B_{R}$ with
   $|\Om_{0}|=\omega_{n}$, $\BDValpha(\Om_{0})\le\eps_{\rm sel}$,
   $Q_{\BDValpha}(\Om_{0})<q_{\rm sel}$, the single-set penalised
   functional of Agent~5 has a minimiser $U_{*}\subset B_{R}$, whose
   volume-normalised companion
   \(
      \Otilde\ :=\ r_{*}^{-1}\,U_{*},
      \quad r_{*}=\bigl(|U_{*}|/|B_{1}|\bigr)^{1/n},
   \)
   satisfies $\Otilde\subset B_{R}$ (up to enlarging $R$ by a dimensional
   amount), $|\Otilde|=\omega_{n}$,
   \begin{equation}\label{eq:plan1-S-alpha}
      \tfrac{1}{2}\,\BDValpha(\Om_{0})\ \le\ \BDValpha(\Otilde)\ \le\
      \tfrac{3}{2}\,\BDValpha(\Om_{0}),
   \end{equation}
   \begin{equation}\label{eq:plan1-S-Q}
      Q_{\BDValpha}(\Otilde)\ \le\ 2\,C_{\rm sel}(n,R)\,
      Q_{\BDValpha}(\Om_{0}),
   \end{equation}
   and $\Otilde$ inherits the BDV regularity package: density estimates,
   Lipschitz nondegeneracy, and a smooth Bernoulli coefficient
   (\cite[Lemmas~4.9, 4.12, 4.16]{BDV}).

\item[(G)] (Graph entry, Agent~6,
   Theorem D2.3 of \texttt{Manuscript/source\_map.md};
   \texttt{Final/GraphEntry.tex}, \texttt{Plan 1/graph-entry-closure.tex}.)
   There exist $\alpha_{\rm graph}=\alpha_{\rm graph}(n,R)>0$ and
   $C_{H}''=C_{H}''(n,R)>0$ and $\gamma=\gamma(n,R)\in(0,1)$ such that
   every $\Otilde$ as in (S) with
   $\BDValpha(\Otilde)\le\alpha_{\rm graph}$ is, after a single
   translation $\Otilde\mapsto\Otilde-x_{\Otilde}$ by its
   barycentre, a $C^{1,\gamma}$ normal graph over $\partial B_{1}$:
   \begin{equation}\label{eq:plan1-G}
      \partial\Otilde=\bigl\{(1+g(\theta))\,\theta:\theta\in\partial B_{1}\bigr\},
      \qquad
      \|g\|_{L^{\infty}(\partial B_{1})}\le C_{H}''(n,R)\,
      \BDValpha(\Otilde)^{1/(2n)}.
   \end{equation}

\item[(I)] (Schauder interpolation, Agent~6,
   Lemma D2.7 of \texttt{Manuscript/source\_map.md};
   \texttt{Final/SchauderInterpolation.tex}.)
   There exists $\alpha_{\rm sph}=\alpha_{\rm sph}(n,R,\gamma_{0})>0$
   such that whenever (G) applies and
   $\BDValpha(\Otilde)\le\alpha_{\rm sph}$, the parametrisation $g$ of
   (G) satisfies
   \begin{equation}\label{eq:plan1-I}
      \|g\|_{C^{2,\gamma_{0}}(\partial B_{1})}\ \le\
      \delta_{\rm sph}(n,\gamma_{0}),
   \end{equation}
   the smallness threshold of Theorem~\ref{thm:plan1-NS}.
\end{itemize}

The thresholds $\eps_{\rm sel}, \alpha_{\rm graph}, \alpha_{\rm sph}$
appearing in (S), (G), (I) are independent of $\Om_{0}$ and are fixed
once $(n,R,\gamma_{0})$ is chosen. Define the small-asymmetry threshold
\begin{equation}\label{eq:plan1-asmall}
   \alpha_{\rm small}(n,R)\ :=\
   \tfrac{1}{2}\,\min\bigl\{
      \eps_{\rm sel}(n,R),\ \alpha_{\rm graph}(n,R),\ \alpha_{\rm sph}(n,R,\gamma_{0})
   \bigr\}\ >\ 0.
\end{equation}

\subsubsection{Volume and barycentre normalisation}
\label{sec:plan1-vol-bary}

The set $U_{*}$ produced by (S) is the minimiser of a penalised
functional whose volume penalty does \emph{not} force $|U_{*}|=\omega_{n}$
exactly; the energy gain forces
\begin{equation}\label{eq:plan1-vol-defect}
   \bigl||U_{*}|-\omega_{n}\bigr|\ \le\ C(n,R)\,\deltaT(U_{*}),
\end{equation}
see \cite[Lemma~4.5]{BDV} and Agent~5. We normalise both the
volume and the barycentre as follows.

\paragraph{Volume normalisation.}
The companion
\(
   \Otilde=r_{*}^{-1}U_{*},
   \quad
   r_{*}=(|U_{*}|/\omega_{n})^{1/n},
\)
satisfies $|\Otilde|=\omega_{n}$ by construction, and the torsion
equation rescales as $u_{\Otilde}(x)=r_{*}^{-2}u_{U_{*}}(r_{*}x)$, so
that $E(\Otilde)=r_{*}^{-(n+2)}E(U_{*})$ (cf.\ the rescaling computation
of \cite[\S2]{BDV}, with $N$ replaced by $n$). The Saint--Venant
deficit transforms as
\(
   \deltaT(\Otilde)=r_{*}^{-(n+2)}\bigl(E(U_{*})-r_{*}^{n+2}E(B_{1})\bigr)
\)
and~\eqref{eq:plan1-vol-defect} gives $|r_{*}-1|\le C(n,R)\,\deltaT(U_{*})$;
the energy errors are absorbed into a multiplicative
$1+O(\deltaT(U_{*}))$ factor in (S), which the constants
$C_{\rm sel}(n,R)$ already include (Agent~5).

\paragraph{Barycentre normalisation.}
After (G) provides the parametrisation $\partial\Otilde=\{(1+g)y\}$, the
barycentre $x_{\Otilde}$ and the volume normalisation $|\Otilde|=\omega_{n}$
give the two scalar identities (\cite[(3.6)--(3.7)]{BDV}, $N\to n$)
\begin{equation}\label{eq:plan1-vol-eq}
   |\Otilde|=\int_{\partial B_{1}}\frac{(1+g)^{n}}{n}\,d\mathcal{H}^{n-1}=\omega_{n},
\end{equation}
\begin{equation}\label{eq:plan1-bary-eq}
   x_{\Otilde}=\int_{\partial B_{1}}y\,\frac{(1+g)^{n+1}}{n+1}\,d\mathcal{H}^{n-1}.
\end{equation}
Expanding in $g$ and using $\|g\|_{C^{2,\gamma_{0}}}\le\delta_{\rm sph}$
yields the BDV moment bounds
(\cite[(3.6)--(3.7)]{BDV}, with $N\to n$):
\begin{equation}\label{eq:plan1-moment}
   \Bigl|\int_{\partial B_{1}}g\,d\mathcal{H}^{n-1}\Bigr|
   + \sum_{i=1}^{n}\Bigl|\int_{\partial B_{1}}y_{i}\,g\,d\mathcal{H}^{n-1}\Bigr|
   \ \le\ C(n)\,\delta_{\rm sph}\,\|g\|_{H^{1/2}(\partial B_{1})}.
\end{equation}
The right-hand side of~\eqref{eq:plan1-moment} is a controlled fraction
of $\|g\|_{H^{1/2}}^{2}$ once $\|g\|_{H^{1/2}}\ge c(n)\,\delta_{\rm sph}$;
when $\|g\|_{H^{1/2}}$ is even smaller, both sides of
\eqref{eq:plan1-BDV33} are subdominant and the small-deficit
case~\eqref{eq:plan1-NS-alpha} holds trivially after a single homothety
and translation. The detailed reduction, which is a one-step rescaling
plus translation, is recorded in
\texttt{Final/NearlySphericalClosure.tex} (Corollary~2.4) and shows
that we may replace \eqref{eq:plan1-NS-hyp} by
\begin{equation}\label{eq:plan1-NS-hyp-relaxed}
   \bigl|\,|\Otilde|-\omega_{n}\,\bigr|+|x_{\Otilde}|\ \le\
   \eta(n,\gamma_{0})\,\|g\|_{L^{2}(\partial B_{1})}
\end{equation}
at the cost of replacing $c_{*}(n)$ by
$\widetilde{c}_{*}(n):=c_{*}(n)/4$. The post-volume normalisation
$|\Otilde|=\omega_{n}$ and the barycentre translation
$\Otilde\mapsto\Otilde-x_{\Otilde}$ make this condition trivially
satisfied with $\eta=0$, so we may apply Theorem~\ref{thm:plan1-NS}
\emph{verbatim} (with $c_{*}(n)$, not $\widetilde{c}_{*}(n)$).

\medskip\noindent
\textbf{Summary of normalisations enforced.}
\begin{itemize}
   \item Volume: $|\Otilde|=\omega_{n}$ exactly, by the homothety
      $\Otilde=r_{*}^{-1}U_{*}$.
   \item Barycentre: $x_{\Otilde}=0$, by the single translation
      $\Otilde\mapsto\Otilde-x_{\Otilde}$, which preserves
      $E,\BDValpha,\Fr$ and the $C^{2,\gamma_{0}}$-norm of $g$.
   \item $C^{2,\gamma_{0}}$-smallness: $\|g\|_{C^{2,\gamma_{0}}}\le\delta_{\rm sph}$
      by (I), provided $\BDValpha(\Otilde)\le\alpha_{\rm sph}$.
\end{itemize}
None of the three normalisations is silently dropped; each is enforced
by a named operation (rescale, translate, threshold check).

\subsubsection{Small-asymmetry linear bound}
\label{sec:plan1-small-alpha}

We now carry through the implication
\((S)\to(G)\to(I)\to\text{Theorem~\ref{thm:plan1-NS}}\) explicitly, with
no compactness step. The output is the small-$\BDValpha$ linear-in-$\BDValpha$
torsion bound.

\begin{theorem}[Small-$\BDValpha$ linear Saint--Venant bound]
\label{thm:plan1-small-alpha}
Fix $0<\gamma_{0}\le 1$. Set
\begin{equation}\label{eq:plan1-qstar}
   q_{*}(n,R)\ :=\
   \frac{c_{*}(n)}{2\,C_{\rm sel}(n,R)}
   \ =\ \frac{1}{64\,n^{2}\,C_{3}(n)\,C_{\rm sel}(n,R)}.
\end{equation}
For every open $\Om_{0}\subset B_{R}$ with $|\Om_{0}|=\omega_{n}$ and
$\BDValpha(\Om_{0})\le\alpha_{\rm small}(n,R)$,
\begin{equation}\label{eq:plan1-small-alpha}
   E(\Om_{0})-E(B_{1})\ \ge\ q_{*}(n,R)\,\BDValpha(\Om_{0}).
\end{equation}
\end{theorem}

\begin{proof}
\emph{Case A:} $Q_{\BDValpha}(\Om_{0})\ge q_{\rm sel}(n,R)$. Then directly
$E(\Om_{0})-E(B_{1})\ge q_{\rm sel}\,\BDValpha(\Om_{0})\ge q_{*}\,\BDValpha(\Om_{0})$,
after shrinking $q_{\rm sel}$ if necessary so that $q_{*}\le q_{\rm sel}$.

\emph{Case B:} $Q_{\BDValpha}(\Om_{0})<q_{\rm sel}(n,R)$. Since
$\BDValpha(\Om_{0})\le\alpha_{\rm small}\le\eps_{\rm sel}$, the
selection (S) applies and yields $\Otilde$ satisfying
\eqref{eq:plan1-S-alpha}--\eqref{eq:plan1-S-Q}. In particular
\[
   \BDValpha(\Otilde)\ \le\
   \tfrac{3}{2}\,\BDValpha(\Om_{0})
   \ \le\ \tfrac{3}{2}\,\alpha_{\rm small}(n,R)
   \ \le\ \min\{\alpha_{\rm graph}(n,R),\,\alpha_{\rm sph}(n,R,\gamma_{0})\},
\]
where the last inequality holds by~\eqref{eq:plan1-asmall} (the factor
$\tfrac{1}{2}$ in $\alpha_{\rm small}$ absorbs the factor $\tfrac{3}{2}$
from (S2)). Thus (G) gives the $C^{1,\gamma}$-graph
\eqref{eq:plan1-G}, and (I) gives the $C^{2,\gamma_{0}}$-smallness
\eqref{eq:plan1-I}.

After the single barycentre translation
$\Otilde\mapsto\Otilde-x_{\Otilde}$ (which preserves all relevant
quantities, see \S\ref{sec:plan1-vol-bary}), the
hypotheses~\eqref{eq:plan1-NS-hyp} of Theorem~\ref{thm:plan1-NS} are
satisfied. Therefore
\[
   Q_{\BDValpha}(\Otilde)\ \ge\ c_{*}(n).
\]
Combined with~\eqref{eq:plan1-S-Q},
\[
   c_{*}(n)\ \le\ Q_{\BDValpha}(\Otilde)\ \le\
   2\,C_{\rm sel}(n,R)\,Q_{\BDValpha}(\Om_{0}),
\]
whence
\(
   Q_{\BDValpha}(\Om_{0})\ \ge\ c_{*}(n)/(2C_{\rm sel}(n,R))=q_{*}(n,R),
\)
which rearranges to~\eqref{eq:plan1-small-alpha}.
\end{proof}

\begin{remark}[Where the BDV contradiction step appears, and what
replaces it]\label{rem:plan1-no-contradiction}
At no point of the proof of Theorem~\ref{thm:plan1-small-alpha} did we
assume $Q_{\BDValpha}(\Om_{0})$ is arbitrarily small. The two cases
(``$Q_{\BDValpha}\ge q_{\rm sel}$'' and ``$Q_{\BDValpha}<q_{\rm sel}$'')
are an explicit dichotomy with thresholds depending on $(n,R)$. In the
second case the conclusion comes from a single application of (S), (G),
(I), and Theorem~\ref{thm:plan1-NS}. The original BDV proof
(\cite[Theorem~4.3]{BDV}) is by contradiction with a minimising
sequence; here it is replaced by the single-set selection of Agent~5.
The conversion of the contradiction step required no additional
expansion beyond Agent~5's quotient-preservation bound
\eqref{eq:plan1-S-Q}: the chain (S)+(G)+(I)+NS contracts the
contradiction into one arithmetic comparison
$Q_{\BDValpha}(\Otilde)\le 2C_{\rm sel}Q_{\BDValpha}(\Om_{0})$.
\end{remark}

\subsubsection{Conversion to the Fraenkel asymmetry}
\label{sec:plan1-conversion}

We now convert~\eqref{eq:plan1-small-alpha} into a Fraenkel-squared
lower bound on $\deltaT$. This is the step that needs both directions
of BDV Lemma~4.2 ((B2)(i) and (B2)(ii)): the linear comparison
\eqref{eq:plan1-42i} produces the sharp exponent~$2$, and the Lipschitz
comparison \eqref{eq:plan1-42ii-Fr} ensures that the small-$\BDValpha$
hypothesis can be entered from a smallness assumption on $\Fr$ itself.

\begin{corollary}[Small-$\Fr$ quadratic torsion bound]
\label{cor:plan1-small-Fr}
Define
\begin{equation}\label{eq:plan1-anear}
   a_{\rm near}(n,R)\ :=\
   \frac{\alpha_{\rm small}(n,R)}{C_{2}(R)\,\omega_{n}}\ >\ 0,
   \qquad
   \kappa_{\rm near}(n,R)\ :=\
   \frac{q_{*}(n,R)\,\omega_{n}^{2}}{C_{1}(n)}.
\end{equation}
For every open $\Om\subset B_{R}$ with $|\Om|=\omega_{n}$ and
$\Fr(\Om)\le a_{\rm near}(n,R)$,
\begin{equation}\label{eq:plan1-small-Fr}
   E(\Om)-E(B_{1})\ \ge\ \kappa_{\rm near}(n,R)\,\Fr(\Om)^{2}.
\end{equation}
\end{corollary}

\begin{proof}
By~\eqref{eq:plan1-42ii-Fr},
$\BDValpha(\Om)\le C_{2}(R)\,\omega_{n}\,\Fr(\Om)\le C_{2}\,\omega_{n}\,a_{\rm near}=\alpha_{\rm small}(n,R)$.
Theorem~\ref{thm:plan1-small-alpha} therefore gives
$E(\Om)-E(B_{1})\ge q_{*}(n,R)\,\BDValpha(\Om)$. By~\eqref{eq:plan1-42i},
$\BDValpha(\Om)\ge|B_{1}|^{2}\Fr(\Om)^{2}/C_{1}(n)=\omega_{n}^{2}\Fr(\Om)^{2}/C_{1}(n)$.
Combining,
$E(\Om)-E(B_{1})\ge (q_{*}\omega_{n}^{2}/C_{1})\,\Fr(\Om)^{2}=\kappa_{\rm near}\,\Fr(\Om)^{2}$.
\end{proof}

\subsubsection{Far-asymmetry regime}
\label{sec:plan1-far}

The far regime $\Fr(\Om)\ge a_{\rm near}$ is handled by the suboptimal
$\Fr^{4}$ bound (B3); no contradiction is needed in this regime either.

\begin{lemma}[Far-$\Fr$ quadratic torsion bound]\label{lem:plan1-far}
For every $a>0$ and every open $\Om\subset B_{R}$ with
$|\Om|=\omega_{n}$ and $\Fr(\Om)\ge a$,
\begin{equation}\label{eq:plan1-far}
   E(\Om)-E(B_{1})\ \ge\ C_{8}'(n)^{-1}\,\Fr(\Om)^{4}
   \ \ge\ \frac{a^{2}}{C_{8}'(n)}\,\Fr(\Om)^{2}.
\end{equation}
\end{lemma}

\begin{proof}
The first inequality is~\eqref{eq:plan1-52-fixed}; the second uses
$\Fr(\Om)^{4}=\Fr(\Om)^{2}\,\Fr(\Om)^{2}\ge a^{2}\,\Fr(\Om)^{2}$ once
$\Fr(\Om)\ge a$.
\end{proof}

\subsubsection{Bounded sharp Saint--Venant inequality}
\label{sec:plan1-SV}

\begin{theorem}[Bounded sharp Saint--Venant, Plan~1]
\label{thm:plan1-SV}
Let $n\ge 2$, $R\ge 2$, $0<\gamma_{0}\le 1$. Define
\begin{equation}\label{eq:plan1-cSV}
   c_{\rm SV}(n,R)\ :=\
   \min\!\Bigl\{
      \kappa_{\rm near}(n,R),\
      \frac{a_{\rm near}(n,R)^{2}}{C_{8}'(n)}
   \Bigr\}\ >\ 0,
\end{equation}
where $\kappa_{\rm near}, a_{\rm near}$ are given by
\eqref{eq:plan1-anear}. For every open $\Om\subset B_{R}$ with
$|\Om|=\omega_{n}=|B^{*}|$,
\begin{equation}\label{eq:plan1-SV}
   E(\Om)-E(B^{*})\ \ge\ c_{\rm SV}(n,R)\,\Fr(\Om)^{2},
\end{equation}
equivalently,
\begin{equation}\label{eq:plan1-SV-deficit}
   \Fr(\Om)^{2}\ \le\ \frac{1}{c_{\rm SV}(n,R)}\,\deltaT(\Om).
\end{equation}
The exponent $2$ on $\Fr$ is sharp: it is saturated asymptotically by
small volume-preserving ellipsoidal perturbations of $B_{1}$
(\cite[Introduction, p.~2]{BDV}).
\end{theorem}

\begin{proof}
The Fraenkel asymmetry is contained in $[0,2)$; we split it at
$a_{\rm near}(n,R)$.

\emph{Case 1 (small Fraenkel):}
$\Fr(\Om)\le a_{\rm near}(n,R)$. Corollary~\ref{cor:plan1-small-Fr}
gives $E(\Om)-E(B_{1})\ge\kappa_{\rm near}(n,R)\,\Fr(\Om)^{2}\ge c_{\rm SV}(n,R)\,\Fr(\Om)^{2}$.

\emph{Case 2 (large Fraenkel):}
$\Fr(\Om)\ge a_{\rm near}(n,R)$. Lemma~\ref{lem:plan1-far} with
$a=a_{\rm near}$ gives
$E(\Om)-E(B_{1})\ge a_{\rm near}^{2}/C_{8}'\,\Fr(\Om)^{2}\ge c_{\rm SV}(n,R)\,\Fr(\Om)^{2}$.

Since $\{\Fr(\Om)\le a_{\rm near}\}\cup\{\Fr(\Om)\ge a_{\rm near}\}=[0,2)$,
the two cases cover all admissible $\Om$, and the bound is uniform with
constant $c_{\rm SV}(n,R)$ as in~\eqref{eq:plan1-cSV}. The deficit form
\eqref{eq:plan1-SV-deficit} follows since $E(B_{1})=E(B^{*})$ when
$|\Om|=|B^{*}|=\omega_{n}$, and $\deltaT(\Om)=E(\Om)-E(B^{*})$ by the
notation register.
\end{proof}

\subsubsection{Constant chain: explicit dependences}
\label{sec:plan1-constants}

The final constant $c_{\rm SV}(n,R)$ in
Theorem~\ref{thm:plan1-SV} is the minimum of a chain of named
quantities. We record their dependences explicitly.

\renewcommand{\arraystretch}{1.15}
\begin{center}
\begin{tabular}{|l|l|l|}
\hline
Constant & Dependence & Source \\
\hline
$C_{1}(n)$ & $n$ & BDV Lem.\ 4.2(i), \eqref{eq:plan1-42i}\\
$C_{2}(R)$ & $R$ & BDV Lem.\ 4.2(ii), \eqref{eq:plan1-42ii}\\
$C_{3}(n)$ & $n$ & BDV Lem.\ 4.2(iii), \eqref{eq:plan1-42iii}\\
$\delta_{1}(n,\gamma_{0})$ & $n,\gamma_{0}$ & BDV Thm.\ 3.3, \eqref{eq:plan1-BDV33}\\
$\delta_{3}(n)$ & $n$ & BDV Lem.\ 4.2(iii)\\
$C_{8}'(n)$ & $n$ & BDV Lem.\ 5.2, \eqref{eq:plan1-52-fixed}\\
$c_{\rm sph}(n)=1/(32n^{2})$ & $n$ & \eqref{eq:plan1-BDV33}\\
$\delta_{\rm sph}(n,\gamma_{0})=\min(\delta_{1},\delta_{3})$ & $n,\gamma_{0}$ & \eqref{eq:plan1-deltasph}\\
$c_{*}(n)=c_{\rm sph}/C_{3}=1/(32n^{2}C_{3}(n))$ & $n$ & Thm.\ \ref{thm:plan1-NS}\\
$C_{\rm sel}(n,R), q_{\rm sel}(n,R), \eps_{\rm sel}(n,R)$ & $n,R$ & Agent~5 (S)\\
$\alpha_{\rm graph}(n,R), C_{H}''(n,R)$ & $n,R$ & Agent~6 (G)\\
$\alpha_{\rm sph}(n,R,\gamma_{0})$ & $n,R,\gamma_{0}$ & Agent~6 (I)\\
\hline
$\alpha_{\rm small}(n,R)$ & $n,R$ & \eqref{eq:plan1-asmall}\\
$q_{*}(n,R)=c_{*}/(2C_{\rm sel})$ & $n,R$ & \eqref{eq:plan1-qstar}\\
$\kappa_{\rm near}(n,R)=q_{*}\omega_{n}^{2}/C_{1}$ & $n,R$ & \eqref{eq:plan1-anear}\\
$a_{\rm near}(n,R)=\alpha_{\rm small}/(C_{2}\omega_{n})$ & $n,R$ & \eqref{eq:plan1-anear}\\
$c_{\rm SV}(n,R)=\min\{\kappa_{\rm near},a_{\rm near}^{2}/C_{8}'\}$ & $n,R$ & \eqref{eq:plan1-cSV}\\
\hline
\end{tabular}
\end{center}

\textbf{Pure $(n,R)$-dependence of $c_{\rm SV}$.}
Every quantity in the chain depends only on $(n,R)$, possibly through
the auxiliary parameter $\gamma_{0}\in(0,1]$ which is fixed once and for
all (we may take $\gamma_{0}=1/2$). The dependence on $R$ enters at
exactly three points:
\begin{enumerate}
\item Selection $C_{\rm sel}(n,R), q_{\rm sel}(n,R), \eps_{\rm sel}(n,R)$
   (Agent~5).
\item Graph entry $\alpha_{\rm graph}(n,R), C_{H}''(n,R)$,
   Schauder interpolation $\alpha_{\rm sph}(n,R,\gamma_{0})$
   (Agent~6).
\item Lipschitz comparison $C_{2}(R)$ (BDV Lemma 4.2(ii)).
\end{enumerate}
The nearly-spherical floor $c_{*}(n)$, the BDV
constants $C_{1}(n),C_{3}(n),C_{8}'(n)$, and the multipliers $\omega_{n}$
are purely dimensional. The constant $c_{\rm sph}(n)=1/(32n^{2})$
deteriorates polynomially in $n$.

\subsubsection{Hypothesis boundaries and degeneration of the
nearly-spherical hypothesis}
\label{sec:plan1-boundaries}

The hypotheses of the BDV nearly-spherical theorem (B1) fail in three
regimes; we record how each is handled here, without silently dropping
any hypothesis.

\paragraph{(a) $\BDValpha(\Om_{0})>\alpha_{\rm small}(n,R)$.}
The selection (S) is still applicable only down to
$\BDValpha(\Om_{0})\le\eps_{\rm sel}(n,R)$. For
$\BDValpha(\Om_{0})>\alpha_{\rm small}$ but $\Fr(\Om)$ remains in
$[0,2)$, we use the far-Fraenkel regime
(Lemma~\ref{lem:plan1-far}). Translating
$\BDValpha\ge\alpha_{\rm small}$ to a Fraenkel lower bound by
\eqref{eq:plan1-42ii-Fr} gives
$\Fr(\Om)\ge\BDValpha/(C_{2}\omega_{n})\ge a_{\rm near}$, so Case~2 of
Theorem~\ref{thm:plan1-SV} applies. Thus no $\BDValpha$-large regime is
left uncovered.

\paragraph{(b) Graph regime not reached because of (G) failure.}
The graph-entry theorem (G) fails if $\BDValpha(\Otilde)>\alpha_{\rm graph}$.
Since (S) preserves $\BDValpha$ up to a factor $\tfrac{3}{2}$ and we
chose $\alpha_{\rm small}\le\alpha_{\rm graph}/2$ in
\eqref{eq:plan1-asmall} (the factor $\tfrac{1}{2}$), Case~B of the proof
of Theorem~\ref{thm:plan1-small-alpha} never enters the regime where
(G) fails. If $\BDValpha(\Om_{0})>\alpha_{\rm small}$ we fall into the
previous paragraph (a).

\paragraph{(c) Schauder/interpolation regime not reached because of (I)
failure.}
Analogously, (I) fails only if $\BDValpha(\Otilde)>\alpha_{\rm sph}$. By
the same choice of $\alpha_{\rm small}\le\alpha_{\rm sph}/2$ we never
enter that regime in Case~B. If $\BDValpha(\Om_{0})>\alpha_{\rm small}$
we again fall into (a). Equivalently: the small-asymmetry threshold
\eqref{eq:plan1-asmall} is chosen as the smallest of the three upstream
thresholds (divided by $2$ to absorb (S2)), so all three
\emph{simultaneously} hold for Case~B.

\paragraph{(d) Sharp boundary $\BDValpha\sim\alpha_{\rm small}$.}
The proof of Theorem~\ref{thm:plan1-SV} glues the two regimes at
$\Fr=a_{\rm near}$, which by \eqref{eq:plan1-anear} corresponds to
$\BDValpha=\alpha_{\rm small}$ (via the linear inequality \eqref{eq:plan1-42ii-Fr}).
At the glue point both
$\kappa_{\rm near}(n,R)\Fr^{2}$ and $(a_{\rm near}^{2}/C_{8}')\Fr^{2}$
are positive, so $c_{\rm SV}(n,R)>0$ is the minimum of two strictly
positive quantities, with explicit lower bounds
\[
   c_{\rm SV}(n,R)\ \ge\
   \min\!\Bigl\{
      \frac{\omega_{n}^{2}}{C_{1}(n)}\cdot\frac{c_{*}(n)}{2C_{\rm sel}(n,R)},\
      \frac{\alpha_{\rm small}(n,R)^{2}}{C_{2}(R)^{2}\,\omega_{n}^{2}\,C_{8}'(n)}
   \Bigr\}.
\]
No additional smallness hypothesis is needed.

\subsubsection{Summary: what this subsection delivers}
\label{sec:plan1-summary}

Combining the above with the notation $\deltaT(\Om)=E(\Om)-E(B^{*})\ge 0$
of the register, we obtain the bounded sharp Saint--Venant inequality
\begin{equation}\label{eq:plan1-final}
   \boxed{\quad
   \Fr(\Om)^{2}\ \le\ C(n,R)\,\deltaT(\Om),
   \qquad
   C(n,R)=\frac{1}{c_{\rm SV}(n,R)},
   \quad}
\end{equation}
valid for every open $\Om\subset B_{R}$ with $|\Om|=\omega_{n}=|B^{*}|$.
The constant $C(n,R)=1/c_{\rm SV}(n,R)$ has explicit dependence on
$n,R$ traced through Theorem~\ref{thm:plan1-SV} and the constant chain
of \S\ref{sec:plan1-constants}; no compactness or contradiction
argument enters. The bounded reduction
$c_{\rm SV}(n,R_{*}(n))\Rightarrow c_{\rm SV}^{\rm glob}(n)$ that removes
the $R$-dependence is the subject of Agent~8.

\paragraph{Status statement.}
\begin{itemize}
\item Theorem~\ref{thm:plan1-NS}: assembled in this subsection from
   (B1) and (B2)(iii).
\item Theorem~\ref{thm:plan1-small-alpha}: assembled in this subsection
   from (S)+(G)+(I)+Theorem~\ref{thm:plan1-NS}.
\item Corollary~\ref{cor:plan1-small-Fr},
   Lemma~\ref{lem:plan1-far}, Theorem~\ref{thm:plan1-SV}: assembled in
   this subsection from (S)+(G)+(I)+NS, (B2)(i)+(ii), and (B3).
\item Upstream black boxes (B1)+(B2)+(B3): cited verbatim from
   \cite{BDV} as \cite[Thm.\ 3.3, Lem.\ 4.2 (i),(ii),(iii),
   Lem.\ 5.2]{BDV}.
\item Inputs (S), (G), (I) from Agents~5, 6: cited by section number
   inside \texttt{Manuscript/02\_plan1\_draft.tex} once Agent~19
   integrates the three Plan 1 subsections.
\end{itemize}

\paragraph{Contradiction step.}
The BDV proof of the small-$\BDValpha$ linear bound is by
contradiction (\cite[Theorem 4.3]{BDV}); here it is replaced by the
single-set selection of Agent~5, which compresses the
sequence/compactness/limit argument into the explicit
quotient-preservation bound (S3). No additional expansion of the source
contradiction step was needed beyond Agent~5's selection theorem:
Theorem~\ref{thm:plan1-small-alpha} is a direct combinatorial
combination of (S), (G), (I), and the pointwise nearly-spherical
inequality, performed in \S\ref{sec:plan1-small-alpha}.

\label{subsec:plan1-ns}%
% UNIFIED 2026-05-13
% SOURCE MAP:
%   - Plan 1/faber-krahn-transfer.tex          (Kohler--Jobin transfer; constants c_1, c_2)
%   - Plan 1/faber-krahn-transfer.md           (same, prose version)
%   - Final/BoundedReduction.tex               (BDV truncation; c_SV(n,R) -> c_SV^{glob}(n))
%   - Final/KohlerJobinTransfer.tex            (canonical scale-invariant statement)
%   - Manuscript/00_notation_register.md       (n, R_*(n), delta_T, delta_FK, A=Fraenkel)
%   - Manuscript/source_map.md  §D.4           (this subsection's items D4.1-D4.6)
%
% Inputs consumed from earlier subsections of 02_plan1_draft.tex:
%   - Agent 7 (NearlySphericalClosure / SaintVenantAssembly), Theorem (SV-bounded):
%        E(\Om)-E(B_1) \ge c_{\rm SV}(n,R)\,\Fr(\Om)^2
%     for every open \Om \subset B_R with |\Om|=\omega_n, where R\ge 2.
%
% Notation register:
%   - dimension is n (N converted);
%   - \Fr := Fraenkel asymmetry (\mathcal A);
%   - \deltaT(\Om) := E(\Om)-E(B^*) (BDV-normalised, E=-T/2);
%   - \deltaFK(\Om) := (|\Om|/\omega_n)^{2/n}(\lambda_1(\Om)-\lambda_1(B^*));
%   - L_n := \lambda_1(B_1), T_n := T(B_1) = \omega_n/(n(n+2)).
%
% Output of this subsection:
%   - Theorem (BR): global sharp Saint--Venant, c_{\rm SV}^{\rm glob}(n);
%   - Theorem (KJ-FK): \Fr(\Om)^2 \le C_{\rm FK}(n)\,\deltaFK(\Om).
%
% This file is intended to be the final subsection of Manuscript/02_plan1_draft.tex
% and is integrated by Agent 19.

\subsection{Bounded reduction and the Kohler--Jobin transfer to Faber--Krahn}
\label{sec:plan1-kj-transfer}

This subsection closes Plan~1. The input is the bounded sharp
Saint--Venant inequality of Agent~7 (Theorem~\ref{thm:SV-bounded}).
Three steps follow:
\begin{enumerate}
\item[(i)] \emph{Bounded reduction}
(\S\ref{ssec:bounded-reduction}). Using the constructive BDV
truncation lemma, we replace an arbitrary open
$\Om\subset\R^n$ with finite measure by a translate that lies in a
ball of \emph{dimensional} radius $R_*(n)$, while controlling the loss
of both the Saint--Venant deficit and the Fraenkel asymmetry. This
removes the $R$-dependence from the Saint--Venant constant.
\item[(ii)] \emph{Kohler--Jobin inequality}
(\S\ref{ssec:kj-inequality}). We state the scale-invariant comparison
$T(\Om)\,\lambda_1(\Om)^{(n+2)/2}\ge T_n L_n^{(n+2)/2}$ and derive the
pointwise lower bound for $\lambda_1(\Om)-L$ in terms of the torsion
deficit.
\item[(iii)] \emph{Final Plan~1 theorem}
(\S\ref{ssec:plan1-final}). Combining (i) and (ii) yields the sharp
quantitative Faber--Krahn inequality in the canonical normalisation of
the notation register,
\[
\Fr(\Om)^2 \le C_{\rm FK}(n)\,\deltaFK(\Om),
\]
with $C_{\rm FK}(n)$ depending only on the dimension.
\end{enumerate}

\subsubsection*{Standing input from Plan~1, Agent~7.}

\begin{theorem}[Bounded sharp Saint--Venant; Agent~7, item D3.8]
\label{thm:SV-bounded}
For every $n\ge 2$ and every $R\ge 2$ there is a constant
$c_{\rm SV}(n,R)>0$, explicit in the BDV constants of the
selection / graph entry / nearly-spherical chain, such that
\begin{equation}\label{eq:SV-bounded}
\deltaT(\Om) \;=\; E(\Om)-E(B_1) \;\ge\; c_{\rm SV}(n,R)\,\Fr(\Om)^2
\end{equation}
for every open set $\Om\subset B_R$ with $|\Om|=\omega_n$.
\end{theorem}

The hypothesis $|\Om|=\omega_n$ is the standing volume normalisation
of the chain; we will remove it by scaling at the end
(Remark~\ref{rem:scale-removal}).

% =====================================================================
\subsubsection{Bounded reduction: removing the $R$-dependence}
\label{ssec:bounded-reduction}
% =====================================================================

The bounded reduction is the BDV truncation argument
\cite[Lemmas~5.1--5.3]{BDV}. It is constructive: a De~Giorgi
$L^\infty$ iteration is combined with an energy cut-off competitor;
no compactness, no minimising sequence, and no contradiction step is
used. The end product is a bounded surrogate $\Otilde$ of $\Om$ inside
a dimensional ball, with quantitatively controlled loss.

\paragraph{Scale-invariant Saint--Venant deficit.} It is convenient to
work with the scale-invariant deficit
\begin{equation}\label{eq:Dscaled}
D(\Om) \;:=\; \omega_n^{-(n+2)/n}\,\deltaT(\Om)
\;=\; \omega_n^{-(n+2)/n}\,\bigl(E(\Om)-E(B^*)\bigr).
\end{equation}
When $|\Om|=\omega_n$ we have $B^*=B_1$ and $D(\Om)\ge 0$ is the
volume-normalised Saint--Venant deficit. The Fraenkel asymmetry
$\Fr(\Om)$ is dimensionless and scale invariant.

\paragraph{Two BDV black boxes.} We import two estimates from
\cite{BDV} as black boxes.

\begin{lemma}[\protect{\cite[Lemma~5.1]{BDV}}: De~Giorgi $L^\infty$ tail]
\label{lem:51}
There is $C_7=C_7(n)>0$ such that, for every open $\Om$ with
$|\Om|=\omega_n$ and every $R\ge 1$,
\begin{equation}\label{eq:tail}
\|u_\Om\|_{L^\infty(\Om\setminus B_{R+1})}
\;\le\; C_7\,|\Om\setminus B_R|^{1/n}.
\end{equation}
\end{lemma}

\begin{proof}[Reference]
\cite[Lemma~5.1]{BDV}; De~Giorgi iteration on the truncated test
function $\varphi_k^2(u_\Om-s_k)_+$. Constructive; constant $C_7(n)$
explicit.
\end{proof}

\begin{lemma}[\protect{\cite[Lemma~5.2]{BDV}}: suboptimal
$\Fr^4$ bound]\label{lem:52}
There is $C_8=C_8(n)>0$ such that for every open $\Om$ of finite
measure,
\begin{equation}\label{eq:far}
D(\Om) \;\ge\; \frac{1}{C_8}\,\Fr(\Om)^4.
\end{equation}
\end{lemma}

\begin{proof}[Reference]
\cite[Lemma~5.2]{BDV}; equivalent, via $E=-T/2$, to
\cite[Theorem~1]{FMPLaplace} with $p=2,q=1$. Pure rearrangement plus
P\'olya--Szeg\H{o}; no contradiction.
\end{proof}

\paragraph{The truncation lemma.} The core construction.

\begin{lemma}[\protect{\cite[Lemma~5.3]{BDV}}: bounded truncation]
\label{lem:53}
There exist dimensional constants $C_9=C_9(n)>0$,
$\delta_{\rm BDV}=\delta_{\rm BDV}(n)\in(0,1]$, and
$d(n)\ge 1$ such that for every open $\Om$ with $|\Om|=\omega_n$ and
$D(\Om)\le\delta_{\rm BDV}(n)$, there exists an open set
$\Otilde\subset\R^n$ satisfying
\begin{align}
|\Otilde| &= \omega_n, \label{eq:53vol} \\
\mathrm{diam}(\Otilde) &\le d(n), \label{eq:53diam} \\
D(\Otilde) &\le C_9(n)\,D(\Om), \label{eq:53D} \\
\Fr(\Om) &\le \Fr(\Otilde) + C_9(n)\,D(\Om). \label{eq:53A}
\end{align}
\end{lemma}

\begin{proof}[Sketch and reference]
\cite[Lemma~5.3]{BDV}. The proof combines~\eqref{eq:tail} with the
cut-off test function $\varphi_k u_\Om$ in the Saint--Venant
variational characterisation of $E(\Om\cap B_{k+2})$, yielding
\cite[(5.11)]{BDV}
\[
E(\Om\cap B_{k+2}) \;\le\; E(\Om) + C\,b_k^{1+1/n},
\qquad b_k := |\Om\setminus B_k|/\omega_n.
\]
A geometric iteration on $k\le K\le K(n)$ produces $b_{K+1}\le
2\widehat C\,D(\Om)$. Set $r:=(|\Om\cap B_{K+3}|/\omega_n)^{1/n}\in
[1-\widehat C\,D(\Om),1]$ and $\Otilde:=r^{-1}(\Om\cap B_{K+3})$.
Properties~\eqref{eq:53vol}--\eqref{eq:53D} follow directly. The
asymmetry estimate~\eqref{eq:53A} uses
\[
|\Om\,\Delta\,(B_1+x_0)|
\;\le\;
|(\Om\cap B_{K+3})\,\Delta\,\Om|
+ r^n\,|\Otilde\,\Delta\,(B_1+x_0/r)|
+ \omega_n(1-r^n),
\]
followed by $1-r^n\le n\widehat C\,D(\Om)$. The construction is
\emph{constructive}: no minimising sequence, no Hausdorff limit, no
contradiction step.
\end{proof}

\begin{remark}[Translation freedom]\label{rem:translation}
The diameter bound~\eqref{eq:53diam} is translation invariant. After
translating $\Otilde$ we may assume $\Otilde\subset B_{d(n)}$. Both
$\Fr$ and $D$ are translation invariant.
\end{remark}

\paragraph{Working radius.} Define the \emph{bounded-reduction radius}
\begin{equation}\label{eq:Rstar}
R_*(n) \;:=\; \max\bigl\{d(n),\,2\bigr\}.
\end{equation}
By Remark~\ref{rem:translation} any $\Otilde$ produced by
Lemma~\ref{lem:53} can be translated into $B_{R_*(n)}$, so
Theorem~\ref{thm:SV-bounded} applies with the dimensional choice
$R=R_*(n)$.

\paragraph{Far-deficit branch.}

\begin{lemma}[Far-deficit quadratic bound]\label{lem:far}
For every open $\Om$ with $|\Om|=\omega_n$ and
$D(\Om)\ge\delta_{\rm BDV}(n)/4$,
\begin{equation}\label{eq:far-quad}
E(\Om)-E(B_1) \;\ge\;
\omega_n^{(n+2)/n}\,\frac{\delta_{\rm BDV}(n)}{16}\,\Fr(\Om)^2.
\end{equation}
\end{lemma}

\begin{proof}
Since $\Fr(\Om)<2$, $\Fr(\Om)^2/4<1$. Hence
$D(\Om)\ge\delta_{\rm BDV}(n)/4\ge(\delta_{\rm BDV}(n)/16)\Fr(\Om)^2$.
Multiplying by $\omega_n^{(n+2)/n}$ and using
$E(\Om)-E(B_1)=\omega_n^{(n+2)/n}D(\Om)$ at volume $\omega_n$ gives the
claim.
\end{proof}

\begin{remark}
The estimate~\eqref{eq:far-quad} is much weaker than~\eqref{eq:far}
(it converts the $\Fr^4$ control into an $\Fr^2$ control by absorbing
$\Fr^2\le 4$), but in this regime $D$ is bounded below by a
dimensional constant, so the trade is harmless.
\end{remark}

\paragraph{Near-deficit branch.}

\begin{lemma}[Near-deficit transfer]\label{lem:near}
Set
\begin{equation}\label{eq:csv-tilde}
\widetilde c_{\rm SV}(n)
\;:=\; \omega_n^{-(n+2)/n}\,c_{\rm SV}\bigl(n,R_*(n)\bigr)\;>\;0.
\end{equation}
For every open $\Om$ with $|\Om|=\omega_n$ and
$D(\Om)\le\delta_{\rm BDV}(n)$,
\begin{equation}\label{eq:near-Fr}
\Fr(\Om)^2 \;\le\;
\Bigl(\frac{2C_9(n)}{\widetilde c_{\rm SV}(n)}
+2C_9(n)^2\,\delta_{\rm BDV}(n)\Bigr)\,D(\Om).
\end{equation}
\end{lemma}

\begin{proof}
Apply Lemma~\ref{lem:53} to obtain $\Otilde$ satisfying
\eqref{eq:53vol}--\eqref{eq:53A}, and by Remark~\ref{rem:translation}
assume $\Otilde\subset B_{d(n)}\subset B_{R_*(n)}$. Apply
Theorem~\ref{thm:SV-bounded} with $R=R_*(n)$ to $\Otilde$:
\[
E(\Otilde)-E(B_1) \;\ge\; c_{\rm SV}\bigl(n,R_*(n)\bigr)\,\Fr(\Otilde)^2.
\]
Divide by $\omega_n^{(n+2)/n}$ (and use $|\Otilde|=\omega_n$) to obtain
\begin{equation}\label{eq:Dtilde-Atilde}
D(\Otilde) \;\ge\; \widetilde c_{\rm SV}(n)\,\Fr(\Otilde)^2.
\end{equation}
Combined with~\eqref{eq:53D},
\[
\Fr(\Otilde)^2 \;\le\;
\frac{D(\Otilde)}{\widetilde c_{\rm SV}(n)}
\;\le\;
\frac{C_9(n)\,D(\Om)}{\widetilde c_{\rm SV}(n)}.
\]
Hence
$\Fr(\Otilde)\le\bigl(C_9(n)D(\Om)/\widetilde c_{\rm SV}(n)\bigr)^{1/2}$.
Plugging this into~\eqref{eq:53A} and squaring,
\begin{align*}
\Fr(\Om)^2
&\le \bigl(\Fr(\Otilde) + C_9(n)\,D(\Om)\bigr)^2
\;\le\; 2\Fr(\Otilde)^2 + 2\,C_9(n)^2\,D(\Om)^2 \\
&\le \frac{2C_9(n)}{\widetilde c_{\rm SV}(n)}\,D(\Om)
   + 2\,C_9(n)^2\,D(\Om)^2 \\
&\le \Bigl(\frac{2C_9(n)}{\widetilde c_{\rm SV}(n)}
   + 2\,C_9(n)^2\,\delta_{\rm BDV}(n)\Bigr)\,D(\Om),
\end{align*}
where the last step uses $D(\Om)\le\delta_{\rm BDV}(n)\le 1$.
\end{proof}

\begin{corollary}[Near-deficit quadratic bound]\label{cor:near}
With
\begin{equation}\label{eq:c-near-glob}
c_{\rm near}^{\rm glob}(n)
\;:=\;
\frac{\omega_n^{(n+2)/n}}
{2C_9(n)\bigl(1/\widetilde c_{\rm SV}(n)+C_9(n)\delta_{\rm BDV}(n)\bigr)},
\end{equation}
every open $\Om$ with $|\Om|=\omega_n$ and
$D(\Om)\le\delta_{\rm BDV}(n)$ satisfies
\begin{equation}\label{eq:near-final}
E(\Om)-E(B_1) \;\ge\; c_{\rm near}^{\rm glob}(n)\,\Fr(\Om)^2.
\end{equation}
\end{corollary}

\begin{proof}
\eqref{eq:near-Fr} gives
$D(\Om)\ge\Fr(\Om)^2/\bigl(2C_9/\widetilde c_{\rm SV}+2C_9^2\delta_{\rm BDV}\bigr)$.
Multiply by $\omega_n^{(n+2)/n}$ and use $E-E(B_1)=\omega_n^{(n+2)/n}D$
to get~\eqref{eq:near-final}.
\end{proof}

\paragraph{Global bounded reduction.}

\begin{theorem}[Global sharp Saint--Venant; bounded reduction]
\label{thm:BR}
Define
\begin{equation}\label{eq:csv-glob}
c_{\rm SV}^{\rm glob}(n)
\;:=\;
\min\!\Bigl\{
c_{\rm near}^{\rm glob}(n),\;
\omega_n^{(n+2)/n}\,\delta_{\rm BDV}(n)/16
\Bigr\}\;>\;0.
\end{equation}
Then for every open set $\Om\subset\R^n$ with $|\Om|=\omega_n$,
\begin{equation}\label{eq:BR}
\deltaT(\Om) \;=\; E(\Om)-E(B_1)
\;\ge\; c_{\rm SV}^{\rm glob}(n)\,\Fr(\Om)^2.
\end{equation}
The constant $c_{\rm SV}^{\rm glob}(n)$ depends only on the dimension
$n$: the $R$-dependence inherited from
Theorem~\ref{thm:SV-bounded} has been internalised at the
dimensional choice $R=R_*(n)$.
\end{theorem}

\begin{proof}
The bound $D(\Om)\ge 0$ is the classical Saint--Venant inequality.
We split into two cases.

\emph{Case A} ($D(\Om)\ge\delta_{\rm BDV}(n)/4$):
Lemma~\ref{lem:far} gives
$E(\Om)-E(B_1)\ge\omega_n^{(n+2)/n}\bigl(\delta_{\rm BDV}(n)/16\bigr)\Fr(\Om)^2$,
which is~\eqref{eq:BR} with constant
$\omega_n^{(n+2)/n}\delta_{\rm BDV}(n)/16$.

\emph{Case B} ($D(\Om)\le\delta_{\rm BDV}(n)$, which in particular
covers $D(\Om)\le\delta_{\rm BDV}(n)/4$):
Corollary~\ref{cor:near} applies and gives~\eqref{eq:BR} with
constant $c_{\rm near}^{\rm glob}(n)$.

The two cases overlap on
$\delta_{\rm BDV}(n)/4\le D(\Om)\le\delta_{\rm BDV}(n)$ and together
cover $D(\Om)\in[0,+\infty)$. Taking the minimum of the two
constants gives~\eqref{eq:csv-glob}.
\end{proof}

\begin{remark}[Dependence on $R$]\label{rem:R-dropout}
The dependence on the confining radius $R$ enters only via
Theorem~\ref{thm:SV-bounded}, which is applied to the surrogate
$\Otilde\subset B_{R_*(n)}$. Since $R_*(n)$ is itself dimensional, the
final constant $c_{\rm SV}^{\rm glob}(n)$ in~\eqref{eq:csv-glob}
depends only on $n$. Concretely, the
$R$-dependence is absorbed in
$\widetilde c_{\rm SV}(n)=\omega_n^{-(n+2)/n}c_{\rm SV}(n,R_*(n))$ in
\eqref{eq:csv-tilde}; everything downstream is purely dimensional.
\end{remark}

\begin{remark}[Removal of the volume normalisation]\label{rem:scale-removal}
\eqref{eq:BR} stated for $|\Om|=\omega_n$ extends to every open
$\Om\subset\R^n$ with $0<|\Om|<\infty$ by scale invariance of $D$ and
$\Fr$: writing
$\Om_*:=(\omega_n/|\Om|)^{1/n}\Om$ one has $|\Om_*|=\omega_n$,
$\Fr(\Om_*)=\Fr(\Om)$, and $D(\Om_*)=D(\Om)$. Hence~\eqref{eq:BR}
applied to $\Om_*$ reads
\begin{equation}\label{eq:BR-scaleinv}
\omega_n^{(n+2)/n}\,D(\Om) \;=\; |\Om|^{(n+2)/n}\,D(\Om)\cdot
\bigl(\omega_n/|\Om|\bigr)^{(n+2)/n}
\quad\text{etc.}
\end{equation}
Equivalently, at volume $|B^*|=|\Om|$,
\begin{equation}\label{eq:torsion-deficit-out}
T(B^*) - T(\Om)
\;=\; 2\bigl(E(\Om)-E(B^*)\bigr)
\;\ge\; 2\,c_{\rm SV}^{\rm glob}(n)\,
\Bigl(\frac{|\Om|}{\omega_n}\Bigr)^{(n+2)/n}\Fr(\Om)^2.
\end{equation}
This is the input to the Kohler--Jobin transfer
(\S\ref{ssec:kj-inequality}).
\end{remark}

% =====================================================================
\subsubsection{The Kohler--Jobin inequality and the pointwise transfer}
\label{ssec:kj-inequality}
% =====================================================================

\paragraph{Statement.} We invoke the Kohler--Jobin inequality in the
following scale-invariant form, which is the form needed to transfer
Saint--Venant stability to Faber--Krahn stability.

\begin{theorem}[Kohler--Jobin inequality]\label{thm:kj}
For every open set $\Om\subset\R^n$ with $0<|\Om|<\infty$ and
$\lambda_1(\Om)<\infty$, and for every ball $B$,
\begin{equation}\label{eq:kj-global}
T(\Om)\,\lambda_1(\Om)^{(n+2)/2}
\;\ge\; T(B)\,\lambda_1(B)^{(n+2)/2}.
\end{equation}
Equivalently, $\Psi(\Om):=T(\Om)\,\lambda_1(\Om)^{(n+2)/2}\ge
\Psi(B_1)=T_n L_n^{(n+2)/2}$. Equality holds if and only if $\Om$ is a
ball (up to translation and a Lebesgue null set).
\end{theorem}

\begin{proof}[Reference]
The original statement is in Kohler--Jobin \cite{KohlerJobin1978}. We
use the form of \cite[Theorem~3.3, p.~322]{Brasco2014}, proved by
Schwarz symmetrisation and a one-dimensional rearrangement of the
torsion potential along level sets of the eigenfunction. The
characterisation of the equality case is \cite[Theorem~3.3]{Brasco2014}.
\end{proof}

\begin{remark}[Canonical normalisation and scaling]\label{rem:kj-canonical}
The exponent $(n+2)/2$ in~\eqref{eq:kj-global} is forced by scale
invariance: from $T(t\Om)=t^{n+2}T(\Om)$ and
$\lambda_1(t\Om)=t^{-2}\lambda_1(\Om)$ we read
$\Psi(t\Om)=t^{n+2-2\cdot(n+2)/2}\Psi(\Om)=\Psi(\Om)$, so $\Psi$ is
scale-invariant. This is the canonical form of the inequality used
throughout the manuscript; both source files
\texttt{Plan~1/faber-krahn-transfer.tex} and
\texttt{Final/KohlerJobinTransfer.tex} use the same form
(with $N$ converted to $n$).
\end{remark}

\paragraph{Pointwise transfer at fixed volume.} Fix the comparison ball
$B^*$ with $|B^*|=|\Om|$, and write
\[
L:=\lambda_1(B^*),\qquad T_0:=T(B^*),\qquad \beta:=\frac{n+2}{2}.
\]
By scaling $L=L_n(\omega_n/|B^*|)^{2/n}$ and
$T_0=T_n(|B^*|/\omega_n)^{(n+2)/n}$.

\begin{lemma}[Pointwise Kohler--Jobin transfer]\label{lem:KJ-ptwise}
For every open $\Om$ with $|\Om|=|B^*|$ and $\lambda_1(\Om)<\infty$,
\begin{equation}\label{eq:plan1-KJ-ptwise}
\lambda_1(\Om)
\;\ge\; L\cdot\Bigl(\frac{T_0}{T(\Om)}\Bigr)^{1/\beta}
\;=\; L\cdot\Bigl(\frac{T_0}{T(\Om)}\Bigr)^{2/(n+2)}.
\end{equation}
\end{lemma}

\begin{proof}
Apply Theorem~\ref{thm:kj} with $B=B^*$ and rearrange. Since $T(\Om)>0$
and $\lambda_1(\Om)<\infty$, both sides are finite and positive.
\end{proof}

\begin{lemma}[Elementary one-variable inequality]\label{lem:plan1-onevar}
For every $p\in(0,1]$ and every $s\in[0,1)$,
\begin{equation}\label{eq:plan1-onevar}
(1-s)^{-p}-1 \;\ge\; p\,s.
\end{equation}
\end{lemma}

\begin{proof}
Set $f(s):=(1-s)^{-p}-1$. Then $f(0)=0$ and
$f'(s)=p(1-s)^{-(p+1)}\ge p$ for $s\in[0,1)$. By the fundamental
theorem of calculus,
$f(s)=\int_0^s f'(\tau)\,d\tau\ge ps$, which is~\eqref{eq:plan1-onevar}.
The factor $p$ on the right is sharp: $f'(0)=p$, so the linear
approximation is tight at $s=0$.
\end{proof}

\paragraph{Small-deficit linearisation.} The classical Saint--Venant
inequality gives $T(\Om)\le T_0$ at fixed volume, so the torsion
deficit $\sigma(\Om):=T_0-T(\Om)\ge 0$. (We use the local symbol
$\sigma$ here, since the manuscript reserves $\delta_T$ for the
energy-form deficit.) Equivalently
$\sigma(\Om)=2\,\deltaT(\Om)$, and setting
$s:=\sigma(\Om)/T_0\in[0,1)$, \eqref{eq:plan1-KJ-ptwise} becomes
\begin{equation}\label{eq:KJ-s}
\lambda_1(\Om)-L \;\ge\; L\bigl[(1-s)^{-1/\beta}-1\bigr].
\end{equation}

\begin{lemma}[Small-deficit Kohler--Jobin linearisation]\label{lem:KJ-linear}
For every open $\Om$ with $|\Om|=|B^*|$ and $\sigma(\Om)\le T_0/2$,
\begin{equation}\label{eq:plan1-KJ-linear}
\lambda_1(\Om)-L
\;\ge\; \frac{2L}{(n+2)\,T_0}\,\sigma(\Om)
\;=\; \frac{4L}{(n+2)\,T_0}\,\deltaT(\Om).
\end{equation}
\end{lemma}

\begin{proof}
With $p=1/\beta=2/(n+2)\in(0,1]$ and $s=\sigma(\Om)/T_0\in[0,1/2]$,
combine~\eqref{eq:KJ-s} and Lemma~\ref{lem:plan1-onevar}:
\[
\lambda_1(\Om)-L
\;\ge\; L\,p\,s
\;=\; \frac{2L}{(n+2)\,T_0}\,\sigma(\Om).
\]
The second equality uses $\sigma=2\deltaT$.
\end{proof}

\begin{remark}[Universal multiplier]\label{rem:multiplier}
The factor $2L/((n+2)T_0)$ is the \emph{Kohler--Jobin multiplier}.
Both $L$ and $T_0$ scale, but the product is dimensionless:
$2L/((n+2)T_0)=2L_n/((n+2)T_n)\cdot(\omega_n/|B^*|)^{(n+2)/n}/(\omega_n/|B^*|)^{2/n}\cdots$,
which simplifies to
\begin{equation}\label{eq:multiplier-explicit}
\frac{2L}{(n+2)T_0}
\;=\;
\frac{2L_n}{(n+2)T_n}\cdot\Bigl(\frac{\omega_n}{|B^*|}\Bigr)^{(n+4)/n}\;.
\end{equation}
Using $T_n=\omega_n/(n(n+2))$, the dimensional part simplifies to
$2L_n/((n+2)T_n)=2nL_n/\omega_n$.
\end{remark}

\paragraph{Large-deficit branch.}

\begin{lemma}[Large-deficit Kohler--Jobin bound]\label{lem:KJ-large}
For every open $\Om$ with $|\Om|=|B^*|$ and $\sigma(\Om)>T_0/2$,
\begin{equation}\label{eq:KJ-large}
\lambda_1(\Om)-L
\;\ge\; L\bigl(2^{2/(n+2)}-1\bigr).
\end{equation}
\end{lemma}

\begin{proof}
$\sigma>T_0/2$ means $T_0/T(\Om)>2$, so~\eqref{eq:plan1-KJ-ptwise} gives
$\lambda_1(\Om)\ge L\cdot 2^{2/(n+2)}$, i.e.~\eqref{eq:KJ-large}.
\end{proof}

Since $|\Om|=|B^*|$ forces $\Fr(\Om)\in[0,2)$ (so $\Fr^2\le 4$),
in the large-deficit regime we may rewrite~\eqref{eq:KJ-large} as
\begin{equation}\label{eq:KJ-large-Fr}
\lambda_1(\Om)-L
\;\ge\; \frac{L\bigl(2^{2/(n+2)}-1\bigr)}{4}\,\Fr(\Om)^2.
\end{equation}

% =====================================================================
\subsubsection{Final Plan~1 sharp Faber--Krahn theorem}
\label{ssec:plan1-final}
% =====================================================================

We now combine the global Saint--Venant inequality
(Theorem~\ref{thm:BR}) and the pointwise Kohler--Jobin transfer
(Lemmas~\ref{lem:KJ-linear} and~\ref{lem:KJ-large}) to obtain the
sharp quantitative Faber--Krahn inequality in the canonical
normalisation of the notation register.

\paragraph{Statement in the lower form.}

\begin{theorem}[Sharp Faber--Krahn via Kohler--Jobin; lower form]
\label{thm:KJ-FK-lower}
There is a dimensional constant $c_{\rm FK}(n)>0$ such that, for
every open $\Om\subset\R^n$ with $0<|\Om|<\infty$,
\begin{equation}\label{eq:KJ-FK-lower}
\deltaFK(\Om)
\;=\;
\Bigl(\frac{|\Om|}{\omega_n}\Bigr)^{2/n}
\bigl(\lambda_1(\Om)-\lambda_1(B^*)\bigr)
\;\ge\; c_{\rm FK}(n)\,\Fr(\Om)^2.
\end{equation}
The constant is given explicitly by
\begin{equation}\label{eq:cFKdef}
c_{\rm FK}(n) \;:=\;
\min\!\Bigl\{
c_1(n),\; c_2(n)
\Bigr\},
\end{equation}
where
\begin{equation}\label{eq:c1c2-def}
c_1(n) := \frac{4 L_n\,c_{\rm SV}^{\rm glob}(n)}{(n+2)\,T_n}
\;=\; \frac{4 n L_n}{\omega_n}\,c_{\rm SV}^{\rm glob}(n),
\qquad
c_2(n) := \frac{L_n\bigl(2^{2/(n+2)}-1\bigr)}{4}.
\end{equation}
\end{theorem}

\begin{proof}
We first reduce to the volume-normalised case $|\Om|=\omega_n$. The
quantity $\deltaFK(\Om)$ is scale invariant: if
$\Om_*:=(\omega_n/|\Om|)^{1/n}\Om$, then $|\Om_*|=\omega_n$,
$\Fr(\Om_*)=\Fr(\Om)$, and
$\deltaFK(\Om_*)=\lambda_1(\Om_*)-L_n=
(|\Om|/\omega_n)^{2/n}(\lambda_1(\Om)-L)$ by the scaling
$\lambda_1(t\Om)=t^{-2}\lambda_1(\Om)$ (with $t=(\omega_n/|\Om|)^{1/n}$),
which equals $\deltaFK(\Om)$.
Thus it suffices to prove
$\lambda_1(\Om)-L_n\ge c_{\rm FK}(n)\Fr(\Om)^2$ for
$|\Om|=\omega_n$ (in which case $B^*=B_1$, $L=L_n$, $T_0=T_n$,
$\deltaFK(\Om)=\lambda_1(\Om)-L_n$).

Assume henceforth $|\Om|=\omega_n$. Theorem~\ref{thm:BR} gives
$\deltaT(\Om)\ge c_{\rm SV}^{\rm glob}(n)\Fr(\Om)^2$, so the torsion
deficit
\begin{equation}\label{eq:torsion-input}
\sigma(\Om) \;=\; T_n - T(\Om) \;=\; 2\deltaT(\Om)
\;\ge\; 2\,c_{\rm SV}^{\rm glob}(n)\,\Fr(\Om)^2.
\end{equation}

\emph{Case 1: $\sigma(\Om)\le T_n/2$.}
Lemma~\ref{lem:KJ-linear} (with $L=L_n$, $T_0=T_n$)
and~\eqref{eq:torsion-input} give
\[
\lambda_1(\Om) - L_n
\;\ge\;
\frac{2 L_n}{(n+2) T_n}\,\sigma(\Om)
\;\ge\;
\frac{2 L_n}{(n+2) T_n}\cdot 2\,c_{\rm SV}^{\rm glob}(n)\,\Fr(\Om)^2
\;=\;
\frac{4 L_n\,c_{\rm SV}^{\rm glob}(n)}{(n+2) T_n}\,\Fr(\Om)^2
\;=\;
c_1(n)\,\Fr(\Om)^2,
\]
the last equality being the definition of $c_1(n)$
in~\eqref{eq:c1c2-def}. The first inequality used the small-deficit
Kohler--Jobin linearisation expressed in terms of the torsion deficit
$\sigma$; the second inequality used $\sigma=2\deltaT$ combined with
the Saint--Venant input
$\deltaT\ge c_{\rm SV}^{\rm glob}\Fr^2$.

\emph{Case 2: $\sigma(\Om)>T_n/2$.}
\eqref{eq:KJ-large-Fr} with $L=L_n$ gives
\[
\lambda_1(\Om)-L_n
\;\ge\; \frac{L_n\bigl(2^{2/(n+2)}-1\bigr)}{4}\,\Fr(\Om)^2
\;=\; c_2(n)\,\Fr(\Om)^2.
\]

Combining the two cases gives
$\lambda_1(\Om)-L_n\ge\min\{c_1(n),c_2(n)\}\Fr(\Om)^2
=c_{\rm FK}(n)\Fr(\Om)^2$, which is~\eqref{eq:KJ-FK-lower} in the
volume-normalised case, and hence (by the first paragraph) in general.
\end{proof}

\begin{remark}[Canonical form of $c_1$]\label{rem:factor-two}
The constant $c_1(n)$ in~\eqref{eq:c1c2-def} matches
\cite[\S3, eq.~(3.1)]{Brasco2014} and the canonical form in
\texttt{Final/KohlerJobinTransfer.tex} (equations $c_1$, $c_{\rm FK}$).
Note that the source file
\texttt{Plan~1/faber-krahn-transfer.tex} states the small-deficit
multiplier as $2L/((n+2)T_0)$ applied to the torsion deficit
$\sigma=T_0-T$, while our proof applies it after the substitution
$\sigma=2\deltaT$; this accounts for the explicit factor of $4$ (not
$2$) on $L_n c_{\rm SV}^{\rm glob}/((n+2)T_n)$ in
\eqref{eq:c1c2-def}. The two source files
\texttt{Final/KohlerJobinTransfer.tex} (eq.~(\ref{eq:c1c2-def}) above)
and \texttt{Plan~1/faber-krahn-transfer.tex} (\S\,Quantitative
dependence) agree on this normalisation.
\end{remark}

\paragraph{Convenient abbreviation.} Throughout the rest of the
manuscript we write
\begin{equation}\label{eq:c1-canonical}
c_1(n) \;:=\; \frac{4 L_n\,c_{\rm SV}^{\rm glob}(n)}{(n+2)\,T_n}
\;=\; \frac{4 n L_n}{\omega_n}\,c_{\rm SV}^{\rm glob}(n),
\end{equation}
which is the definition just used.

\paragraph{Canonical statement (upper form).}
Inverting~\eqref{eq:KJ-FK-lower}, the inequality reads
$\Fr(\Om)^2\le C_{\rm FK}(n)\,\deltaFK(\Om)$ with
$C_{\rm FK}(n)=1/c_{\rm FK}(n)$. This is the upper form required by
the master brief (and shared with Plan~2).

\begin{theorem}[Sharp Faber--Krahn; Plan~1, upper form]
\label{thm:KJ-FK-upper}
There is a dimensional constant $C_{\rm FK}(n)>0$ such that, for every
open set $\Om\subset\R^n$ with $0<|\Om|<\infty$,
\begin{equation}\label{eq:KJ-FK-upper}
\boxed{\;
\Fr(\Om)^2 \;\le\; C_{\rm FK}(n)\,\deltaFK(\Om),
\;}
\end{equation}
where
\begin{equation}\label{eq:CFK}
C_{\rm FK}(n)
\;:=\; \frac{1}{c_{\rm FK}(n)}
\;=\; \max\!\Bigl\{
\frac{(n+2)\,T_n}{4 L_n\,c_{\rm SV}^{\rm glob}(n)},\;\;
\frac{4}{L_n\bigl(2^{2/(n+2)}-1\bigr)}
\Bigr\},
\end{equation}
with $c_{\rm SV}^{\rm glob}(n)$ defined in~\eqref{eq:csv-glob}.
Equality cases: by Theorem~\ref{thm:kj}, $\deltaFK(\Om)=0$
iff $\Om$ is a ball; in this case $\Fr(\Om)=0$.
\end{theorem}

\begin{proof}
If $\lambda_1(\Om)=+\infty$ then $\deltaFK(\Om)=+\infty$ and the
inequality \eqref{eq:KJ-FK-upper} is trivial (the left-hand side is
bounded above by $4$ since $\Fr(\Om)<2$). Otherwise
$\lambda_1(\Om)<\infty$ and the result is immediate from
Theorem~\ref{thm:KJ-FK-lower} with the canonical $c_1(n)$
from~\eqref{eq:c1-canonical}: invert the lower bound. The
constant $C_{\rm FK}(n)$ depends only on $n$ because
$c_{\rm SV}^{\rm glob}(n)$ does (Remark~\ref{rem:R-dropout}) and $L_n,
T_n$ are dimensional.
\end{proof}

\begin{remark}[$R$ drops out]\label{rem:R-removal-final}
The dependence on the confining radius $R$ is removed in two stages:
\begin{itemize}
\item Theorem~\ref{thm:SV-bounded} (Agent~7) provides
$c_{\rm SV}(n,R)$ for $\Om\subset B_R$. We \emph{instantiate} this
input at $R=R_*(n)=\max\{d(n),2\}$, a purely dimensional radius
defined in~\eqref{eq:Rstar}.
\item The bounded reduction (Theorem~\ref{thm:BR}) transfers this
to $c_{\rm SV}^{\rm glob}(n)$ via the BDV truncation lemma, which
produces a translate $\Otilde\subset B_{R_*(n)}$.
\end{itemize}
No subsequent step reintroduces $R$: the Kohler--Jobin inequality
\eqref{eq:kj-global} is independent of any confining ball, and the
pointwise transfer~\eqref{eq:plan1-KJ-ptwise} depends only on the
universal ball constants $L_n=\lambda_1(B_1)$, $T_n=T(B_1)$, and the
torsion deficit at fixed volume.
\end{remark}

\begin{remark}[Sharpness of the exponent]
The exponent $2$ on $\Fr$ in~\eqref{eq:KJ-FK-upper} is sharp; it is
saturated by smooth volume-preserving ellipsoidal perturbations of the
ball, for which $\Fr(\Om)$ and $\deltaFK(\Om)^{1/2}$ are of the same
order. See \cite[\S1.2]{BDV} for the construction.
\end{remark}

\begin{remark}[Match with Plan~2 normalisation]\label{rem:plan2-match}
The final theorem~\eqref{eq:KJ-FK-upper} is stated in the canonical
normalisation of the notation register (item~D4.6 of
\texttt{Manuscript/source\_map.md}), with $\deltaFK$ defined as in
notation register \S4: $\deltaFK(\Om)=(|\Om|/\omega_n)^{2/n}(\lambda_1(\Om)-\lambda_1(B^*))$.
This is exactly the normalisation of the Plan~2 final theorem
(item~E5.8); the two proofs share both the bounded reduction
(Theorem~\ref{thm:BR}) and the Kohler--Jobin transfer step
(\S\ref{ssec:kj-inequality}). The constant $C_{\rm FK}(n)$ produced
here and the constant $C_{\rm FK}^{\rm Plan~2}(n)$ produced in
\S\ref{sec:plan2-global-assembly} differ only through the
Saint--Venant constant $c_{\rm SV}^{\rm glob}(n)$, which each route
constructs independently.
\end{remark}

\paragraph{Summary of constants.}
\renewcommand{\arraystretch}{1.25}
\begin{center}
\begin{tabular}{|l|l|l|}
\hline
Symbol & Value/dependence & Source \\
\hline
$\omega_n$ & $|B_1|$ & dimensional \\
$T_n$ & $\omega_n/(n(n+2))$ & dimensional \\
$L_n$ & $\lambda_1(B_1)=j_{n/2-1,1}^2$ & dimensional \\
$\beta$ & $(n+2)/2$ & Theorem~\ref{thm:kj} \\
$C_7(n)$ & \cite[Lem.~5.1]{BDV} & Lemma~\ref{lem:51} \\
$C_8(n)$ & \cite[Lem.~5.2]{BDV} & Lemma~\ref{lem:52} \\
$C_9(n)$ & \cite[Lem.~5.3]{BDV} & Lemma~\ref{lem:53} \\
$\delta_{\rm BDV}(n)\in(0,1]$ & \cite[Lem.~5.3]{BDV} & Lemma~\ref{lem:53} \\
$d(n)\ge 1$ & \cite[Lem.~5.3]{BDV} & Lemma~\ref{lem:53} \\
$R_*(n)=\max\{d(n),2\}$ & dimensional & \eqref{eq:Rstar} \\
$c_{\rm SV}(n,R)$ & Agent~7 & Theorem~\ref{thm:SV-bounded} \\
$\widetilde c_{\rm SV}(n)$ & $\omega_n^{-(n+2)/n}c_{\rm SV}(n,R_*(n))$ & \eqref{eq:csv-tilde} \\
$c_{\rm near}^{\rm glob}(n)$ & \eqref{eq:c-near-glob} & Corollary~\ref{cor:near} \\
$c_{\rm SV}^{\rm glob}(n)$ & $\min\{c_{\rm near}^{\rm glob},\omega_n^{(n+2)/n}\delta_{\rm BDV}/16\}$ & Theorem~\ref{thm:BR} \\
$c_1(n)$ & $4L_n c_{\rm SV}^{\rm glob}/((n+2)T_n)=4nL_n c_{\rm SV}^{\rm glob}/\omega_n$ & \eqref{eq:c1-canonical} \\
$c_2(n)$ & $L_n(2^{2/(n+2)}-1)/4$ & \eqref{eq:c1c2-def} \\
$c_{\rm FK}(n)$ & $\min\{c_1(n),c_2(n)\}$ & Theorem~\ref{thm:KJ-FK-lower} \\
$C_{\rm FK}(n)$ & $1/c_{\rm FK}(n)$ & Theorem~\ref{thm:KJ-FK-upper} \\
\hline
\end{tabular}
\end{center}

\paragraph{Where compactness could have entered (and does not).}
The bounded reduction (\S\ref{ssec:bounded-reduction}) is the only step
where one might fear a hidden compactness argument. We record:
\begin{itemize}
\item BDV Lemma~5.1 (Lemma~\ref{lem:51}) is De~Giorgi
$L^\infty$ iteration: constructive, no minimising sequence.
\item BDV Lemma~5.2 (Lemma~\ref{lem:52}) is rearrangement plus
P\'olya--Szeg\H{o}: no contradiction.
\item BDV Lemma~5.3 (Lemma~\ref{lem:53}) reaches the truncation
index $K\le K(n)$ by a geometric iteration
\cite[(5.12)--(5.14)]{BDV}, not by Hausdorff compactness.
\item The gluing of near and far branches (Theorem~\ref{thm:BR}) is
arithmetic.
\item The Kohler--Jobin transfer (\S\ref{ssec:kj-inequality}) uses
only \eqref{eq:kj-global}, the FTC, and elementary monotone-function
inequalities.
\end{itemize}

This closes Plan~1.


% =====================================================================
% End of Manuscript/02_plan1_draft.tex
% =====================================================================


% =====================================================================
% Plan 2 proof (Agents 9-13, +14 auxiliary)
% =====================================================================
% =====================================================================
% Manuscript/03_plan2_draft.tex
%
% Concatenated Plan 2 section (Agents 9, 10, 11, 12, 13, 14) produced
% by Agent 19 (Final Assembly). The five load-bearing subfiles are
% \input here in the order: level-set identity -> metric framework
% -> FJ centre / oriented radial trace -> integrated action / weighted
% trace -> boundary-layer transfer / final theorem. The auxiliary
% centroid / space-time Pi note (Agent 14) follows as the final
% subsection (with explicit auxiliary status declared in the
% Introduction; see Remark~\ref{rem:intro-pi-route}).
% =====================================================================

\section{Second proof: level-set deficit identity and metric trace (Plan~2)}
\label{sec:plan2}

This section presents the second of the two proofs of the sharp
quantitative Faber--Krahn inequality~\eqref{eq:intro-main}. The chain
is structurally different from Plan~1: instead of selecting a single
near-extremal set and analysing it as a graph over a sphere, Plan~2
analyses every level $\rho$ of the volume-radius foliation
$\Erho=\{u_\Om>t(\rho)\}$, classifies levels by their isoperimetric
defect, and transfers the integrated defect bound to the metric quotient
$\X=\mathfrak M/\R^n$:
\[
   \begin{aligned}
   &\text{level-set deficit identity (\S\ref{subsec:plan2-levelset-identity})}\\
   &\quad\to\;\text{metric framework / first variation (\S\ref{sec:metric-setup})}\\
   &\quad\to\;\text{FJ centre + oriented radial trace (\S\ref{sec:plan2-fj-radial})}\\
   &\quad\to\;\text{integrated action / weighted metric trace (\S\ref{subsec:weighted-metric-trace})}\\
   &\quad\to\;\text{boundary-layer transfer + final Plan~2 theorem
      (\S\ref{ssec:plan2-boundary-layer-assembly})}.
   \end{aligned}
\]
The auxiliary centroid / space--time $\Pi$ route is included as the
last subsection
(\S\ref{ssec:plan2-aux-centroid}); per
Remark~\ref{rem:intro-pi-route} it is \emph{not} used in any
load-bearing step of the repaired Plan~2 chain. Plan~2's output is
the same global sharp Saint--Venant inequality fed into Kohler--Jobin
as in Plan~1, with the same final constant $C_{\rm FK}(n)$.

% --- Section-level label aliases used by cross-section references ---
% UNIFIED 2026-05-13
% SOURCE MAP for Manuscript/plan2_levelset_identity.tex
% ==========================================================================
% Agent 9 deliverable: first subsection of 03_plan2_draft.tex.
%
% Each numbered item below corresponds to a labelled result inside this file
% (see \label{} keys). Status keys:
%   PROVED      : proved in this subsection (full proof written here).
%   ADAPTED     : adapted from a named local building block.
%   CITED       : cited from the literature, with exact theorem number.
%   IMPORT      : imported from preliminaries (Agents 2-4) without reproof.
%
% Items
% -----
% [E1.0]  Notation, standing hypotheses, Talenti L^infty bound.
%           Source: WIP_LevelSetIdentity.tex (Notation+Standing);
%                   00_notation_register.md sections 1-9.
%           Status: ADAPTED (notation register); CITED(Talenti1976).
% [E1.1]  Lemma lem:coarea: -m'(t) = alpha(t).
%           Source: WIP_LevelSetIdentity.tex Lem 2.1; agent1-finite-perimeter-identity.md Lem 1.
%           Status: PROVED. Cited inputs: AFP2000 Thm 3.40, Maggi2012 Thm 13.1.
% [E1.2]  Lemma lem:flux: beta(t) = m(t)  (weak divergence-flux identity).
%           Source: WIP_LevelSetIdentity.tex Lem 3.1; agent1-finite-perimeter-identity.md Lem 2.
%           Status: PROVED. Cited inputs: AFP2000 Thm 3.99, Maggi2012 Thm 13.1,
%                   Maggi2012 Ch.17 (rigorous divergence on finite-perimeter sets).
% [E1.3]  Lemma lem:Iab: layer-cake form of I, AC of I, I' = u^*.
%           Source: WIP_LevelSetIdentity.tex Lem 4.1-4.2; agent1 Lem 3a-3b.
%           Status: PROVED. Cited inputs: Kawohl1985 Sec 2; Talenti1976.
% [E1.4]  Lemma lem:Ipp (corrected I''): I''(s) = -1/alpha(t(s)).
%           Source: WIP_LevelSetIdentity.tex Lem 4.3 (corrected form); agent1 Lem 3c.
%           Status: PROVED. Cited inputs: AFP2000 Thm 3.99 (BV chain rule);
%                   BBC1975 (inverse-function theorem for monotone AC).
% [E1.5]  Lemma lem:Gpp: pointwise convexity formula G''(s) = 1/(c_n s^{1-2/n}) - 1/alpha(t).
%           Status: PROVED (direct).
% [E1.6]  Definition def:DHDI: D_H(t), D_I(t); nonnegativity (Cauchy-Schwarz, De Giorgi).
%           Status: PROVED. Cited inputs: DeGiorgi1958, Maggi2012 Thm 14.1.
% [E1.7]  Lemma lem:cov: pointwise-to-integrated convexity identity.
%           Status: PROVED.
% [E1.8]  Lemma lem:Gamma: Gamma(Omega) = (1/L) int s G''(s) ds.
%           Status: PROVED (integration by parts).
% [E1.9]  Lemma lem:endpoint: Gamma(Omega) = (2/|Omega|)(E(Omega) - E(B^*)).
%           Status: PROVED. Cited inputs: PolyaSzego1951, Bandle1980 II.5.
% [E1.10] Theorem thm:deficit (EXACT DEFICIT IDENTITY).
%           Source: WIP_LevelSetIdentity.tex Thm 5.1 (main theorem).
%           Status: PROVED. Notation register Section 8.
% [E1.11] Corollary cor:variance: weighted variance form of D_H.
%           Source: WIP_LevelSetIdentity.tex Sec 6 Rem; level-set-deficit-identity.md Sec 6.
%           Status: PROVED.
% [E1.12] Lemma lem:profile-gap: profile gap b - v <= 2 delta_T / s.
%           Source: level-set-deficit-identity.md Sec 4 (and Lem 8.1).
%           Status: PROVED.
% [E1.13] Lemma lem:slab-good: existence of a good isoperimetric-defect level
%           in a slab (Lemma 8.2 of source note).
%           Status: PROVED.
% [E1.14] Lemma lem:superlevel-transfer: Lemma 8.3 (superlevel asymmetry transfer).
%           Source: level-set-deficit-identity.md Lem 8.3.
%           Status: PROVED.
%
% Hypothesis tracking
% -------------------
% Standing class of Theorem thm:deficit: Omega subset R^n open, |Omega| < infty.
% No boundedness, smoothness, or connectedness of Omega is required.
% All boundary integrals are on De Giorgi reduced boundaries.
%
% Constants
% ---------
% All constants are explicit and depend only on the dimension n. The constant
% c_n := n^2 omega_n^{2/n} appears throughout. No constants of the form
% C(n,R,rho_*) appear in this subsection: bounded-class constants enter only in
% the downstream Plan 2 blocks (MetricFramework onward), which consume this
% subsection.
%
% Notation register compliance
% ----------------------------
% Dimension n (not N). B^* = ball of volume |Omega|, centred at origin.
% R_Omega := (|Omega|/omega_n)^{1/n} is used in place of the symbol R for the
% volume radius (notation register Section 2, R/R_Omega conflict resolution).
% alpha(t), beta(t), D_H, D_I match notation register Section 8 exactly.
% ==========================================================================

\subsection{Exact level-set deficit identity}
\label{subsec:plan2-levelset-identity}

This subsection contains the analytic foundation of Plan~2. We prove an
exact representation of the (normalised) Saint--Venant deficit
$\delta_{T}(\Omega)$ as a weighted integral, over the levels of the
torsion function, of two nonnegative pointwise defects: a
Cauchy--Schwarz defect $D_{H}(t)$ of $|\nabla u|$ on the reduced
boundary $\partial^{*}E_{t}$, and a squared-perimeter isoperimetric
defect $D_{I}(t)$ of the superlevel set $E_{t}$. The identity is
valid, in full finite-perimeter generality, for any open
$\Omega\subset\R^{n}$ with $0<|\Omega|<\infty$; no smoothness,
boundedness, or connectedness of $\Omega$ is assumed.

The result is the analytic input of every downstream block in
\S\ref{sec:plan2-metric}--\S\ref{sec:plan2-assembly}. Two further
consequences proved here (the profile-gap estimate and the
superlevel-transfer estimate) are used in
\S\ref{sec:plan2-weighted-trace} (good levels in a near-boundary slab)
and in \S\ref{sec:plan2-boundary-layer} (transfer of the Fraenkel
asymmetry from a superlevel set back to $\Omega$).

\subsubsection*{Standing hypotheses and notation}
\label{sssec:lsi-notation}

Throughout the subsection, $n\ge 2$ and $\Omega\subset\R^{n}$ is an open
set with $0<|\Omega|<\infty$. We use the canonical notation register
of \S\ref{sec:notation} without further comment, and recall only the
items that enter the present argument:

\begin{itemize}
\item $u=u_{\Omega}\in H^{1}_{0}(\Omega)$ is the variational torsion
function, i.e.\ the unique minimiser of
$J(v):=\tfrac12\int_{\Omega}|\nabla v|^{2}-\int_{\Omega}v$ on
$H^{1}_{0}(\Omega)$, equivalently the weak solution of
$-\Delta u=1$ in $\Omega$ with zero $H^{1}_{0}$-trace; testing $-\Delta
u=1$ against $u^{-}\in H^{1}_{0}(\Omega)$ gives $u\ge 0$ a.e.

\item $E(\Omega):=-\tfrac12\int_{\Omega}u\,dx=-\tfrac12 T(\Omega)$ is the
BDV-normalised Saint--Venant energy (notation register \S 3), and
$B^{*}\subset\R^{n}$ is the open ball with $|B^{*}|=|\Omega|$, centred at
the origin, of (volume) radius
\[
   R_{\Omega}:=\Bigl(\frac{|\Omega|}{\omega_{n}}\Bigr)^{1/n}.
\]
(We write $R_{\Omega}$ rather than $R$ to avoid the conflict with the
confining radius of \S\ref{sec:notation}.)

\item For $t\ge 0$, $E_{t}:=\{u>t\}$ (precise representative
\cite[\S 15]{Maggi2012}), $m(t):=|E_{t}|$ is right-continuous and
nonincreasing, with $m(0^{+})\le|\Omega|$ and $m(\|u\|_{\infty})=0$. The
decreasing rearrangement is $u^{*}(s):=\inf\{t\ge 0:m(t)\le s\}$ for
$s\in[0,|\Omega|]$.

\item $\partial^{*}E_{t}$ is the De Giorgi reduced boundary
\cite{DeGiorgi1958}, and $\nu_{E_{t}}$ is the measure-theoretic outer
unit normal. We set
\begin{equation}\label{eq:lsi-alpha-beta}
   \alpha(t):=\int_{\partial^{*}E_{t}}\frac{1}{|\nabla u|}\,d\mathcal H^{n-1},
   \qquad
   \beta(t):=\int_{\partial^{*}E_{t}}|\nabla u|\,d\mathcal H^{n-1},
\end{equation}
defined for a.e.\ $t$ (justified by Lemma~\ref{lem:coarea}--\ref{lem:flux}
below), and
\[
   c_{n}:=n^{2}\omega_{n}^{2/n}.
\]
\end{itemize}

\paragraph{Talenti $L^{\infty}$ bound.} By Talenti's pointwise
rearrangement comparison \cite[Theorem 1]{Talenti1976}
(see also \cite[Theorem 1.3.2]{Kesavan2006}),
\begin{equation}\label{eq:lsi-talenti}
   u^{*}_{\Omega}(s)\le u^{*}_{B^{*}}(s)
   =\frac{R_{\Omega}^{2}-(s/\omega_{n})^{2/n}}{2n},
   \qquad s\in[0,|\Omega|],
\end{equation}
so
\begin{equation}\label{eq:lsi-Linfty}
   \|u\|_{L^{\infty}(\Omega)}\le \frac{R_{\Omega}^{2}}{2n}<\infty.
\end{equation}
This is the sole reason boundedness of $\Omega$ is unnecessary at the
analytic level. Combined with $u\in H^{1}_{0}(\Omega)$ extended by zero,
\eqref{eq:lsi-Linfty} gives $u\in BV_{\rm loc}(\R^{n})\cap L^{\infty}$;
the BV coarea formula \cite[Theorem 3.40]{AFP2000},
\cite[Theorem 13.1]{Maggi2012} then applies, and for every Borel
$\varphi\ge 0$,
\begin{equation}\label{eq:lsi-coarea}
   \int_{\Omega}|\nabla u|\,\varphi\,dx
   =\int_{-\infty}^{\infty}\Bigl(\int_{\partial^{*}E_{t}}\varphi\,d\mathcal
   H^{n-1}\Bigr)\,dt.
\end{equation}
The set
\[
   \mathcal R:=\{t\in(0,\|u\|_{\infty}):
   P(E_{t})<\infty\hbox{ and }\partial^{*}E_{t}\subset\Omega\}
\]
has full Lebesgue measure in $(0,\|u\|_{\infty})$, and all statements
``for a.e.\ $t$'' below mean for a.e.\ $t\in\mathcal R$. The countable set
of plateau levels $\{t:|\{u=t\}|>0\}$ is, in particular, Lebesgue-null and
is automatically discarded by the $dt$-integral.

\subsubsection{Step 1: coarea identity for $-m'(t)$}
\label{sssec:lsi-coarea}

\begin{lemma}[Coarea form of $-m'$]\label{lem:coarea}
For a.e.\ $t\in(0,\|u\|_{\infty})$,
\begin{equation}\label{eq:lsi-L1}
   -m'(t)=\alpha(t)
         =\int_{\partial^{*}E_{t}}\frac{1}{|\nabla u|}\,d\mathcal H^{n-1}.
\end{equation}
In particular $m$ is absolutely continuous on $[\eps,\|u\|_{\infty}]$ for
every $\eps>0$, and $\alpha\in L^{1}_{\rm loc}((0,\|u\|_{\infty}))$.
\end{lemma}

\begin{proof}
Fix $0<\alpha_{0}<\beta_{0}<\|u\|_{\infty}$ and apply
\eqref{eq:lsi-coarea} to the Borel function
$\varphi(x)=|\nabla u(x)|^{-1}\mathbf 1_{\{|\nabla u|>0\}}(x)
\mathbf 1_{[\alpha_{0},\beta_{0}]}(u(x))$. The product
$|\nabla u|\,\varphi$ equals $\mathbf 1_{\{\alpha_{0}<u<\beta_{0}\}\cap\{|\nabla
u|>0\}}$ almost everywhere; the contribution of $\{|\nabla u|=0\}$ to
the BV coarea slicing vanishes by \cite[3.40(b)]{AFP2000}, so
\begin{equation}\label{eq:lsi-coarea-app}
   |E_{\alpha_{0}}\setminus E_{\beta_{0}}|
   =\int_{\Omega}\mathbf 1_{\{\alpha_{0}<u<\beta_{0}\}}\,dx
   =\int_{\alpha_{0}}^{\beta_{0}}\!\!\Bigl(\int_{\partial^{*}E_{t}}
     \frac{1}{|\nabla u|}\,d\mathcal H^{n-1}\Bigr)\,dt
   =\int_{\alpha_{0}}^{\beta_{0}}\alpha(t)\,dt.
\end{equation}
Outside the at most countable set of plateau levels,
$|E_{\alpha_{0}}\setminus E_{\beta_{0}}|=m(\alpha_{0})-m(\beta_{0})$, so
$m(\alpha_{0})-m(\beta_{0})=\int_{\alpha_{0}}^{\beta_{0}}\alpha(t)\,dt$.
Differentiating in $\beta_{0}$ at Lebesgue points of $\alpha$ yields
\eqref{eq:lsi-L1}. Local integrability of $\alpha$ on $(0,\|u\|_{\infty})$
follows because \eqref{eq:lsi-coarea-app} bounds
$\int_{\alpha_{0}}^{\beta_{0}}\alpha(t)\,dt\le |\Omega|<\infty$, and
absolute continuity of $m$ on $[\eps,\|u\|_{\infty}]$ follows directly
from \eqref{eq:lsi-coarea-app}.
\end{proof}

\begin{remark}\label{rem:plateaus}
The countable set of plateau levels $\{t:|\{u=t\}|>0\}$ contributes
jumps to $m$ that are absorbed by the $dt$-integral in
Theorem~\ref{thm:deficit} below. Equation~\eqref{eq:lsi-L1} is not
asserted on these levels.
\end{remark}

\subsubsection{Step 2: weak divergence/flux identity}
\label{sssec:lsi-flux}

The next lemma is the rigorous BV/Lipschitz-truncation replacement for
the formal divergence identity
$\int_{\partial^{*}E_{t}}\nabla u\cdot\nu_{E_{t}}\,d\mathcal H^{n-1}
=\int_{E_{t}}\Delta u\,dx$ on the finite-perimeter superlevel $E_{t}$.
The point is that the trace of $\nabla u$ on the reduced boundary
$\partial^{*}E_{t}$ does not exist in general; the Lipschitz truncation
of $u-t$ avoids invoking that trace.

\begin{lemma}[Flux equals volume]\label{lem:flux}
For a.e.\ $t\in(0,\|u\|_{\infty})$, $E_{t}$ has finite perimeter and
\begin{equation}\label{eq:lsi-L2}
   \beta(t)
   =\int_{\partial^{*}E_{t}}|\nabla u|\,d\mathcal H^{n-1}
   =m(t).
\end{equation}
\end{lemma}

\begin{proof}
Fix $t\in\mathcal R$ with $|E_{t}|<\infty$, $P(E_{t})<\infty$, and $u$
not having a plateau at $t$; this holds for a.e.\ $t$.

\smallskip\noindent
\emph{Step 1 (Lipschitz truncation).}
For each $\eps>0$ define
$\phi_{\eps}\colon\R\to[0,1]$ by
$\phi_{\eps}(s):=\min(1,\max(0,s/\eps))$ and set
$\eta_{\eps}(x):=\phi_{\eps}(u(x)-t)$. Since $\phi_{\eps}$ is Lipschitz
with $\phi_{\eps}(0)=0$ and $u\in H^{1}_{0}(\Omega)$, the chain rule for
Sobolev compositions \cite[Theorem 3.99]{AFP2000} yields
$\eta_{\eps}\in H^{1}_{0}(\Omega)$ with
\begin{equation}\label{eq:lsi-eta-grad}
   \nabla\eta_{\eps}(x)
   =\eps^{-1}\nabla u(x)\,\mathbf 1_{\{t<u<t+\eps\}}(x).
\end{equation}
Test the weak equation $-\Delta u=1$ on $H^{1}_{0}(\Omega)$ against
$\eta_{\eps}$:
\begin{equation}\label{eq:lsi-weakform}
   \int_{\Omega}\nabla u\cdot\nabla\eta_{\eps}\,dx
   =\int_{\Omega}\eta_{\eps}\,dx.
\end{equation}
By \eqref{eq:lsi-eta-grad}, the left-hand side of \eqref{eq:lsi-weakform}
equals
\begin{equation}\label{eq:lsi-LHS}
   \int_{\Omega}\nabla u\cdot\nabla\eta_{\eps}\,dx
   =\frac{1}{\eps}\int_{\{t<u<t+\eps\}}|\nabla u|^{2}\,dx.
\end{equation}

\smallskip\noindent
\emph{Step 2 (Coarea applied to the LHS).}
Applying \eqref{eq:lsi-coarea} with
$\varphi(x)=|\nabla u(x)|\,\mathbf 1_{(t,t+\eps)}(u(x))$,
\begin{equation}\label{eq:lsi-RHS-coarea}
   \int_{\{t<u<t+\eps\}}|\nabla u|^{2}\,dx
   =\int_{t}^{t+\eps}\Bigl(\int_{\partial^{*}E_{\tau}}
     |\nabla u|\,d\mathcal H^{n-1}\Bigr)\,d\tau
   =\int_{t}^{t+\eps}\beta(\tau)\,d\tau.
\end{equation}

\smallskip\noindent
\emph{Step 3 (Limit of the right-hand side of \eqref{eq:lsi-weakform}).}
Since $t$ is not a plateau level, $|\{u=t\}|=0$, so
$\eta_{\eps}\to\mathbf 1_{E_{t}}$ pointwise a.e.\ as $\eps\downarrow 0$.
Because $0\le\eta_{\eps}\le \mathbf 1_{\{u>t\}}\le \mathbf 1_{E_{t}}+\mathbf
1_{\{u=t\}}$ and $|E_{t}|<\infty$, the dominated convergence theorem
gives
\[
   \int_{\Omega}\eta_{\eps}\,dx
   \xrightarrow[\eps\downarrow 0]{}|E_{t}|=m(t).
\]

\smallskip\noindent
\emph{Step 4 (Limit of \eqref{eq:lsi-RHS-coarea}).}
The function $\tau\mapsto\beta(\tau)$ is locally integrable on
$(0,\|u\|_{\infty})$: by \eqref{eq:lsi-coarea} with
$\varphi(x)=|\nabla u(x)|\mathbf 1_{(0,\|u\|_{\infty})}(u(x))$, and using
the Dirichlet identity
$\int_{\Omega}|\nabla u|^{2}\,dx=\int_{\Omega}u\,dx<\infty$ (test
$-\Delta u=1$ against $u\in H^{1}_{0}(\Omega)$),
\[
   \int_{0}^{\|u\|_{\infty}}\beta(\tau)\,d\tau
   =\int_{\Omega}|\nabla u|^{2}\,dx
   =\int_{\Omega}u\,dx<\infty.
\]
Therefore, at every Lebesgue point of $\beta$ (i.e.\ for a.e.\ $t$),
\[
   \frac{1}{\eps}\int_{t}^{t+\eps}\beta(\tau)\,d\tau
   \xrightarrow[\eps\downarrow 0]{}\beta(t).
\]

\smallskip\noindent
\emph{Step 5 (Conclusion).}
Combining \eqref{eq:lsi-weakform}--\eqref{eq:lsi-RHS-coarea} with Steps 3
and 4 gives $\beta(t)=m(t)$ for a.e.\ $t\in\mathcal R$.
\end{proof}

\begin{remark}[Formal divergence-theorem reading]\label{rem:divergence}
If one had a one-sided trace of $\nabla u$ on $\partial^{*}E_{t}$ (i.e.\
classical smoothness near the level surface), the relation
$\nabla u\cdot\nu_{E_{t}}=-|\nabla u|$ on $\partial^{*}E_{t}$ would
follow from $u>t$ inside $E_{t}$, and Lemma~\ref{lem:flux} would reduce
to the standard divergence theorem
\cite[Theorem 16.3]{Maggi2012} applied to $\nabla u$:
\[
   m(t)=\int_{E_{t}}1\,dx=-\int_{E_{t}}\Delta u\,dx
   =-\int_{\partial^{*}E_{t}}\nabla u\cdot\nu_{E_{t}}\,d\mathcal H^{n-1}
   =\beta(t).
\]
The Lipschitz-truncation step above is the rigorous replacement when no
such trace is available; see \cite[Chapter 17]{Maggi2012} for the
general theory of fluxes on finite-perimeter sets without traces.
\end{remark}

\subsubsection{Step 3: distributional derivatives of the rearranged
primitive}
\label{sssec:lsi-Iprime}

Define
\begin{equation}\label{eq:lsi-Idef}
   I(s):=\int_{\{u>u^{*}(s)\}}u\,dx,\qquad s\in[0,|\Omega|].
\end{equation}

\begin{lemma}[Layer-cake form of $I$ and first derivative]\label{lem:Iab}
For every $s\in[0,|\Omega|]$,
\begin{equation}\label{eq:lsi-Ilayer}
   I(s)=\int_{0}^{s}u^{*}(\sigma)\,d\sigma,
\end{equation}
$I$ is absolutely continuous on $[0,|\Omega|]$, and $I'(s)=u^{*}(s)$ for
a.e.\ $s\in(0,|\Omega|)$.
\end{lemma}

\begin{proof}
By the equimeasurability of $u$ and $u^{*}$ on $(0,|\Omega|)$
\cite[\S 2]{KawohlBook}, \cite{Talenti1976}, one has
$\int_{\Omega}g(u)\,dx=\int_{0}^{|\Omega|}g(u^{*}(\sigma))\,d\sigma$ for
every Borel $g\ge 0$. Choose $g(t):=t\,\mathbf 1_{\{t>u^{*}(s)\}}$. By
the definition of $u^{*}$, the set $\{\sigma\in(0,|\Omega|):
u^{*}(\sigma)>u^{*}(s)\}$ equals $(0,s)$ up to Lebesgue-null sets, so
\[
   I(s)=\int_{\Omega}g(u)\,dx
       =\int_{0}^{|\Omega|}g(u^{*}(\sigma))\,d\sigma
       =\int_{0}^{s}u^{*}(\sigma)\,d\sigma,
\]
which is \eqref{eq:lsi-Ilayer}. By \eqref{eq:lsi-talenti}, $u^{*}\in
L^{\infty}(0,|\Omega|)\subset L^{1}(0,|\Omega|)$, so \eqref{eq:lsi-Ilayer}
expresses $I$ as the indefinite integral of a bounded Borel function;
the Lebesgue differentiation theorem then gives the absolute continuity
of $I$ on $[0,|\Omega|]$ and $I'(s)=u^{*}(s)$ a.e.
\end{proof}

The next lemma is the {\it corrected} form of the second derivative of
$I$. A naive derivation gives
$I''(s)=-1/a(u^{*}(s))$ with $a(t)=\int_{\partial^{*}E_{t}}|\nabla u|$,
i.e.\ with the {\it flux} $\beta(t)$ in place of the {\it coarea}
$\alpha(t)$. By Lemma~\ref{lem:flux}, $\beta(t)=m(t)$, so that
expression would give $I''(s)=-1/m(u^{*}(s))=-1/s$, which is correct
only when $|\nabla u|$ is $\mathcal H^{n-1}$-a.e.\ constant on
$\partial^{*}E_{t}$ (i.e.\ on the ball). The general identity uses
$\alpha$, and the difference $\alpha\beta-P(E_{t})^{2}$ between the two
is precisely the Cauchy--Schwarz defect $D_{H}$ of \S\ref{sssec:lsi-DHDI}
below.

\begin{lemma}[Corrected second derivative of $I$]\label{lem:Ipp}
The rearrangement $u^{*}$ is absolutely continuous on $[s_{0},|\Omega|]$
for every $0<s_{0}<|\Omega|$, and
\begin{equation}\label{eq:lsi-Ipp}
   I''(s)=(u^{*})'(s)
        =\frac{1}{m'(u^{*}(s))}
        =-\frac{1}{\alpha(u^{*}(s))}
   \qquad\text{for a.e.\ }s\in(0,|\Omega|),
\end{equation}
where $I''$ denotes the absolutely continuous representative of the
distributional derivative of $I'=u^{*}$ on $[s_{0},|\Omega|]$.
\end{lemma}

\begin{proof}
By Lemma~\ref{lem:coarea}, $m$ is absolutely continuous on
$[\eps,\|u\|_{\infty}]$ for every $\eps>0$, with $m'(t)=-\alpha(t)\le 0$
a.e. The set of plateau values $\{t:|\{u=t\}|>0\}$ is at most countable
(monotonicity of $m$); outside it, $m$ is strictly decreasing on its
domain of strict monotonicity. Consequently the right-continuous
generalised inverse $u^{*}$ of $m$ is also strictly decreasing and
absolutely continuous on $[s_{0},|\Omega|]$ for every $s_{0}\in
(0,|\Omega|)$. Furthermore, $\alpha(t)>0$ for a.e.\ $t$ (otherwise $m$
would not be strictly decreasing), so $1/\alpha(t)$ is well-defined
a.e. Applying the inverse-function theorem for monotone absolutely
continuous functions \cite[Theorem 3.99]{AFP2000},
\cite[Lemma 1.4]{BBC1975} to $m(u^{*}(s))=s$ yields
\[
   (u^{*})'(s)=\frac{1}{m'(u^{*}(s))}=-\frac{1}{\alpha(u^{*}(s))}
\]
for a.e.\ $s$. Since $I'(s)=u^{*}(s)$ a.e.\ by Lemma~\ref{lem:Iab}, the
distributional derivative of $I'$ on $[s_{0},|\Omega|]$ coincides with
the classical a.e.\ derivative of $u^{*}$, giving \eqref{eq:lsi-Ipp}.
\end{proof}

\paragraph{Cauchy--Schwarz on the reduced boundary.}
Applying the Cauchy--Schwarz inequality on $\partial^{*}E_{t}$ to the
pair of functions $(|\nabla u|^{1/2},|\nabla u|^{-1/2})$ gives
\begin{equation}\label{eq:lsi-CS}
   P(E_{t})^{2}
   =\Bigl(\int_{\partial^{*}E_{t}}1\,d\mathcal H^{n-1}\Bigr)^{2}
   \le\alpha(t)\,\beta(t),
\end{equation}
with equality if and only if $|\nabla u|$ is $\mathcal H^{n-1}$-a.e.\
constant on $\partial^{*}E_{t}$.

\subsubsection{Step 4: pointwise convexity formula}
\label{sssec:lsi-Gpp}

Following the Nicola--Tilli convexity-gap idea adapted to the torsion
problem, introduce the auxiliary function
\begin{equation}\label{eq:lsi-Gdef}
   G(s):=I(s)+\frac{s^{1+2/n}}{(4+2n)\,\omega_{n}^{2/n}},
   \qquad s\in[0,|\Omega|].
\end{equation}
A direct computation gives
\[
   \frac{d^{2}}{ds^{2}}\frac{s^{1+2/n}}{(4+2n)\omega_{n}^{2/n}}
   =\frac{(1+2/n)(2/n)}{(4+2n)\omega_{n}^{2/n}}s^{2/n-1}
   =\frac{1}{n^{2}\omega_{n}^{2/n}\,s^{1-2/n}}
   =\frac{1}{c_{n}\,s^{1-2/n}}.
\]

\begin{lemma}[Pointwise convexity formula]\label{lem:Gpp}
For a.e.\ $s\in(0,|\Omega|)$, with $t=u^{*}(s)$ (equivalently $s=m(t)$ on
the strictly decreasing branch of $m$),
\begin{equation}\label{eq:lsi-Gpp}
   G''(s)=\frac{1}{c_{n}\,s^{1-2/n}}-\frac{1}{\alpha(t)}.
\end{equation}
\end{lemma}

\begin{proof}
Combine Lemma~\ref{lem:Ipp} ($I''(s)=-1/\alpha(u^{*}(s))$) with the
explicit second derivative of the polynomial corrector
$\frac{s^{1+2/n}}{(4+2n)\omega_{n}^{2/n}}$ computed above.
\end{proof}

\subsubsection{Step 5: the two pointwise defects}
\label{sssec:lsi-DHDI}

\begin{definition}[Cauchy--Schwarz and isoperimetric defects]\label{def:DHDI}
For a.e.\ $t\in(0,\|u\|_{\infty})$, set
\begin{align}
   D_{H}(t)
   &:=\alpha(t)\,\beta(t)-P(E_{t})^{2}
    =\Bigl(\!\int_{\partial^{*}E_{t}}\!\!\tfrac{1}{|\nabla u|}\Bigr)
      \Bigl(\!\int_{\partial^{*}E_{t}}\!\!|\nabla u|\Bigr)-P(E_{t})^{2},
   \label{eq:lsi-DH}\\
   D_{I}(t)
   &:=P(E_{t})^{2}-c_{n}\,m(t)^{2-2/n}.
   \label{eq:lsi-DI}
\end{align}
\end{definition}

\begin{lemma}[Pointwise nonnegativity]\label{lem:DH-DI-pos}\label{lem:lsid:CS-defect}\label{prop:lsid:CS-defect}
For a.e.\ $t\in(0,\|u\|_{\infty})$, $D_{H}(t)\ge 0$ with equality if and
only if $|\nabla u|$ is $\mathcal H^{n-1}$-a.e.\ constant on
$\partial^{*}E_{t}$; and $D_{I}(t)\ge 0$ with equality if and only if
$E_{t}$ is (equivalent to) a ball.
\end{lemma}

\begin{proof}
$D_{H}\ge 0$ is the Cauchy--Schwarz inequality \eqref{eq:lsi-CS}. The
equality case is the standard equality case in Cauchy--Schwarz.

$D_{I}\ge 0$ is the squared form of the De~Giorgi isoperimetric
inequality $P(E_{t})\ge n\omega_{n}^{1/n}|E_{t}|^{1-1/n}$ valid for every
set of finite perimeter \cite[Theorem~II]{DeGiorgi1958},
\cite[Theorem 14.1]{Maggi2012}; equality holds iff $E_{t}$ is
$\mathcal H^{n}$-equivalent to a ball.
\end{proof}

\subsubsection{Step 6: pointwise-to-integrated convexity identity}
\label{sssec:lsi-cov}

\begin{lemma}[Pointwise-to-integrated convexity identity]\label{lem:cov}\label{lem:lsid:DHDI-psi}
With $L:=|\Omega|$,
\begin{equation}\label{eq:lsi-cov}
   \int_{0}^{L}s\,G''(s)\,ds
   =\int_{0}^{\|u\|_{\infty}}\frac{D_{H}(t)+D_{I}(t)}{c_{n}\,m(t)^{1-2/n}}\,dt.
\end{equation}
\end{lemma}

\begin{proof}
By Lemma~\ref{lem:Gpp},
\begin{equation}\label{eq:lsi-cov-A}
   \int_{0}^{L}s\,G''(s)\,ds
   =\int_{0}^{L}s\Bigl(\frac{1}{c_{n}\,s^{1-2/n}}
      -\frac{1}{\alpha(u^{*}(s))}\Bigr)\,ds.
\end{equation}
Perform the change of variable $s=m(t)$, $ds=m'(t)\,dt=-\alpha(t)\,dt$
on the strictly decreasing branch of $m$. The countable set of plateau
levels of $u$ is a Lebesgue-null set in $t$, and contributes nothing to
the integral. As $s$ runs from $0$ to $L$, $t$ runs from $\|u\|_{\infty}$
to $0$; reversing orientation,
\begin{equation}\label{eq:lsi-cov-B}
   \int_{0}^{L}s\,G''(s)\,ds
   =\int_{0}^{\|u\|_{\infty}}m(t)
     \Bigl(\frac{1}{c_{n}\,m(t)^{1-2/n}}-\frac{1}{\alpha(t)}\Bigr)
     \alpha(t)\,dt
   =\int_{0}^{\|u\|_{\infty}}\Bigl(
        \frac{m(t)\,\alpha(t)}{c_{n}\,m(t)^{1-2/n}}-m(t)\Bigr)\,dt.
\end{equation}
By Lemma~\ref{lem:flux}, $m(t)=\beta(t)$ a.e., so
$m(t)\alpha(t)=\alpha(t)\beta(t)$. Adding and subtracting
$P(E_{t})^{2}/(c_{n}m(t)^{1-2/n})$ in the integrand of
\eqref{eq:lsi-cov-B},
\[
   \frac{m(t)\,\alpha(t)}{c_{n}\,m(t)^{1-2/n}}-m(t)
   =\frac{\alpha(t)\,\beta(t)-c_{n}\,m(t)^{2-2/n}}{c_{n}\,m(t)^{1-2/n}}
   =\frac{D_{H}(t)+D_{I}(t)}{c_{n}\,m(t)^{1-2/n}}.
\]
Substituting back into \eqref{eq:lsi-cov-B} gives \eqref{eq:lsi-cov}.
\end{proof}

\subsubsection{Step 7: endpoint convexity gap}
\label{sssec:lsi-Gamma}

With $L:=|\Omega|$, define
\begin{equation}\label{eq:lsi-Gammadef}
   \Gamma(\Omega):=G'(L)-\frac{G(L)}{L}.
\end{equation}

\begin{lemma}[Endpoint convexity-gap formula]\label{lem:Gamma}
$\displaystyle\Gamma(\Omega)=\frac{1}{L}\int_{0}^{L}s\,G''(s)\,ds$.
\end{lemma}

\begin{proof}
$I\in W^{1,\infty}(0,L)\subset W^{2,1}_{\rm loc}((0,L])$ by
Lemmas~\ref{lem:Iab}--\ref{lem:Ipp}, and the polynomial corrector in
\eqref{eq:lsi-Gdef} is smooth on $(0,L]$ and Lipschitz on $[0,L]$; so $G$
is absolutely continuous on $[0,L]$ and in $W^{2,1}_{\rm loc}((0,L])$,
with $G(0)=I(0)=0$. By Lemma~\ref{lem:Iab},
\[
   G'(s)=u^{*}(s)+\frac{(1+2/n)s^{2/n}}{(4+2n)\omega_{n}^{2/n}}
        =u^{*}(s)+\frac{s^{2/n}}{2n\,\omega_{n}^{2/n}},
\]
using $\frac{1+2/n}{4+2n}=\frac{1}{2n}$; $G'$ is bounded on $[0,L]$ by
\eqref{eq:lsi-Linfty}. Integration by parts on $(\eps,L)$ gives
\[
   \int_{\eps}^{L}s\,G''(s)\,ds
   =L\,G'(L)-\eps\,G'(\eps)-G(L)+G(\eps).
\]
As $\eps\downarrow 0$, $\eps\,G'(\eps)\to 0$ and $G(\eps)\to G(0)=0$ by
the boundedness of $G'$ and the absolute continuity of $G$. Therefore
$\int_{0}^{L}s\,G''(s)\,ds=L\,G'(L)-G(L)$; dividing by $L$ yields the
claim.
\end{proof}

\subsubsection{Step 8: endpoint identification with the Saint--Venant
deficit}
\label{sssec:lsi-endpoint}

\begin{lemma}[Endpoint value]\label{lem:endpoint}
\begin{equation}\label{eq:lsi-endpoint}
   \Gamma(\Omega)=\frac{2}{|\Omega|}\bigl(E(\Omega)-E(B^{*})\bigr).
\end{equation}
\end{lemma}

\begin{proof}
At $s=L$, $u^{*}(L)=0$ by definition of $u^{*}$ (the lowest level is the
zero level). By Lemma~\ref{lem:Iab},
$I(L)=\int_{0}^{L}u^{*}(\sigma)\,d\sigma=\int_{\Omega}u\,dx$ (using
equimeasurability), so
\[
   \frac{G(L)}{L}=\frac{1}{L}\int_{\Omega}u\,dx
                  +\frac{L^{2/n}}{(4+2n)\,\omega_{n}^{2/n}}.
\]
Also $I'(L)=u^{*}(L)=0$, so
\[
   G'(L)=\frac{L^{2/n}}{2n\,\omega_{n}^{2/n}}.
\]
Combining,
\begin{equation}\label{eq:lsi-Gamma-explicit}
   \Gamma(\Omega)
   =\frac{L^{2/n}}{2n\,\omega_{n}^{2/n}}
    -\frac{L^{2/n}}{(4+2n)\,\omega_{n}^{2/n}}
    -\frac{1}{L}\int_{\Omega}u\,dx
   =\frac{L^{2/n}}{2n(n+2)\,\omega_{n}^{2/n}}
    -\frac{1}{L}\int_{\Omega}u\,dx,
\end{equation}
where we used $4+2n=2(n+2)$ and
$\frac{1}{2n}-\frac{1}{2(n+2)}=\frac{1}{2n(n+2)}$.

For the ball $B^{*}$ of volume $L=\omega_{n}R_{\Omega}^{n}$, the torsion
function is $u_{B^{*}}(x)=(R_{\Omega}^{2}-|x|^{2})/(2n)$ (the unique
solution of $-\Delta u_{B^{*}}=1$ in $B^{*}$, $u_{B^{*}}=0$ on
$\partial B^{*}$). Direct integration in spherical coordinates,
\[
   \int_{B^{*}}u_{B^{*}}\,dx
   =\frac{1}{2n}\int_{B^{*}}(R_{\Omega}^{2}-|x|^{2})\,dx
   =\frac{n\omega_{n}}{2n}\int_{0}^{R_{\Omega}}(R_{\Omega}^{2}-r^{2})r^{n-1}\,dr
   =\frac{\omega_{n}\,R_{\Omega}^{n+2}}{2n(n+2)}.
\]
Since $L=\omega_{n}R_{\Omega}^{n}$, $L^{2/n}=\omega_{n}^{2/n}R_{\Omega}^{2}$,
hence
\[
   \frac{L^{2/n}}{2n(n+2)\omega_{n}^{2/n}}
   =\frac{R_{\Omega}^{2}}{2n(n+2)}
   =\frac{1}{L}\cdot\frac{\omega_{n}R_{\Omega}^{n+2}}{2n(n+2)}
   =\frac{1}{L}\int_{B^{*}}u_{B^{*}}\,dx.
\]
Inserting back in \eqref{eq:lsi-Gamma-explicit} and using
$E(\Omega)=-\tfrac12\int_{\Omega}u$ (and the analogous identity for $B^{*}$),
\[
   \Gamma(\Omega)
   =\frac{1}{L}\Bigl(\int_{B^{*}}u_{B^{*}}\,dx-\int_{\Omega}u\,dx\Bigr)
   =\frac{2}{L}\bigl(E(\Omega)-E(B^{*})\bigr).
\]
This identification of $\Gamma(B^{*})$ via the explicit ball torsion
function is the classical P\'olya--Szeg\H{o} content
\cite{PolyaSzego1951}, \cite[\S II.5]{Bandle1980}.
\end{proof}

\subsubsection{The exact deficit identity}
\label{sssec:lsi-main}

\begin{theorem}[Exact level-set deficit identity]\label{thm:deficit}\label{thm:lsid:identity}\label{thm:plan2-levelset-id}
Let $n\ge 2$ and let $\Omega\subset\R^{n}$ be open with $0<|\Omega|<\infty$.
Let $u=u_{\Omega}\in H^{1}_{0}(\Omega)$ be the variational torsion
function. Then for a.e.\ $t\in(0,\|u\|_{\infty})$, the superlevel set
$E_{t}=\{u>t\}$ has finite perimeter, the coarea identity
\eqref{eq:lsi-L1}, the flux identity \eqref{eq:lsi-L2}, and the corrected
second-derivative identity \eqref{eq:lsi-Ipp} hold, and
\begin{equation}\label{eq:lsi-deficit}
   \boxed{\;
   \frac{2}{|\Omega|}\bigl(E(\Omega)-E(B^{*})\bigr)
   =\frac{1}{|\Omega|}\int_{0}^{\|u\|_{\infty}}
     \frac{D_{H}(t)+D_{I}(t)}{n^{2}\omega_{n}^{2/n}\,m(t)^{1-2/n}}\,dt,\;}
\end{equation}
where $D_{H},D_{I}$ are the nonnegative pointwise defects of
\eqref{eq:lsi-DH}--\eqref{eq:lsi-DI} and $B^{*}$ is the ball with
$|B^{*}|=|\Omega|$.
\end{theorem}

\begin{proof}
Lemmas~\ref{lem:cov}, \ref{lem:Gamma}, \ref{lem:endpoint} together give
\[
   \frac{2}{|\Omega|}\bigl(E(\Omega)-E(B^{*})\bigr)
   =\Gamma(\Omega)
   =\frac{1}{L}\int_{0}^{L}s\,G''(s)\,ds
   =\frac{1}{L}\int_{0}^{\|u\|_{\infty}}\frac{D_{H}(t)+D_{I}(t)}
     {c_{n}\,m(t)^{1-2/n}}\,dt.
\]
The pointwise statements \eqref{eq:lsi-L1}, \eqref{eq:lsi-L2},
\eqref{eq:lsi-Ipp} are proved in Lemmas~\ref{lem:coarea},
\ref{lem:flux}, \ref{lem:Ipp}.
\end{proof}

The deficit identity~\eqref{eq:lsi-deficit} thus expresses
$\delta_{T}(\Omega)$ in the form
\[
   \delta_{T}(\Omega)
   =E(\Omega)-E(B^{*})
   =\frac{1}{2}\int_{0}^{\|u\|_{\infty}}\bigl(D_{H}(t)+D_{I}(t)\bigr)\,w(t)\,dt,
\]
with the {\it explicit weight}
\begin{equation}\label{eq:lsi-weight}
   w(t):=\frac{1}{c_{n}\,m(t)^{1-2/n}}
        =\frac{1}{n^{2}\omega_{n}^{2/n}\,m(t)^{1-2/n}}.
\end{equation}
This is the form quoted in the master brief and in
\S\ref{sec:notation}; the dimensional constant $c_{n}=n^{2}\omega_{n}^{2/n}$
is the only constant appearing in the identity itself.

\subsubsection*{Cauchy--Schwarz/Serrin variance form of $D_{H}$}

The following corollary recasts $D_{H}$ as a level-set Serrin variance
defect, in the form used in
\S\ref{sec:plan2-weighted-trace}--\S\ref{sec:plan2-boundary-layer} when
combining $D_{H}$ with normal-oscillation bounds.

\begin{corollary}[Weighted variance form of $D_{H}$]\label{cor:variance}
For a.e.\ $t\in(0,\|u\|_{\infty})$, with $f:=|\nabla u|$ on
$\partial^{*}E_{t}$ and $\bar f:=m(t)/P(E_{t})$,
\begin{equation}\label{eq:lsi-variance}
   \int_{\partial^{*}E_{t}}\frac{(f-\bar f)^{2}}{f}\,d\mathcal H^{n-1}
   =\frac{m(t)}{P(E_{t})^{2}}\,D_{H}(t).
\end{equation}
\end{corollary}

\begin{proof}
Direct computation, using Lemma~\ref{lem:flux} ($\int f\,d\mathcal H^{n-1}
=\beta(t)=m(t)$) and $\int 1\,d\mathcal H^{n-1}=P(E_{t})$:
\[
   \int_{\partial^{*}E_{t}}\!\frac{(f-\bar f)^{2}}{f}\,d\mathcal H^{n-1}
   =\int_{\partial^{*}E_{t}}\!\!f\,d\mathcal H^{n-1}
    -2\bar f\,P(E_{t})
    +\bar f^{2}\!\int_{\partial^{*}E_{t}}\!\!\tfrac{1}{f}\,d\mathcal H^{n-1}
   =m(t)-\frac{2 m(t)}{1}+\frac{m(t)^{2}}{P(E_{t})^{2}}\,\alpha(t).
\]
The first two terms collapse to $-m(t)$, and using
$\alpha(t)\beta(t)=\alpha(t)m(t)$,
\[
   \int_{\partial^{*}E_{t}}\!\frac{(f-\bar f)^{2}}{f}\,d\mathcal H^{n-1}
   =\frac{m(t)^{2}}{P(E_{t})^{2}}\alpha(t)-m(t)
   =\frac{m(t)\bigl(\alpha(t)\,m(t)-P(E_{t})^{2}\bigr)}{P(E_{t})^{2}}
   =\frac{m(t)}{P(E_{t})^{2}}\,D_{H}(t),
\]
since $\alpha m=\alpha\beta$. This is \eqref{eq:lsi-variance}.
\end{proof}

\subsubsection{Profile-gap consequences (used in
\S\ref{sec:plan2-weighted-trace}--\S\ref{sec:plan2-boundary-layer})}
\label{sssec:lsi-consequences}

The deficit identity~\eqref{eq:lsi-deficit} controls weighted averages
of $D_{H}+D_{I}$ over levels; the two consequences below extract
{\it pointwise} good levels in a controlled near-boundary slab, together
with a transfer estimate from a superlevel set back to $\Omega$. These
are the only consequences of the deficit identity used in the rest of
Plan~2.

\paragraph{Setup.} Write
\[
   L:=|\Omega|,\quad
   R_{\Omega}:=(L/\omega_{n})^{1/n},\quad
   \delta_{T}(\Omega):=E(\Omega)-E(B^{*}),\quad
   v(s):=u_{\Omega}^{*}(s),\quad
   b(s):=u_{B^{*}}^{*}(s).
\]
By Talenti~\eqref{eq:lsi-talenti}, $b(s)-v(s)\ge 0$ on $[0,L]$, and the
explicit ball profile gives
$b(s)=(R_{\Omega}^{2}-(s/\omega_{n})^{2/n})/(2n)$.

\begin{lemma}[Profile-gap estimate]\label{lem:profile-gap}
The function $h(s):=b(s)-v(s)$ is nonnegative and nonincreasing on
$[0,L]$, $h(L)=0$, and
\begin{equation}\label{eq:lsi-profile-int}
   \int_{0}^{L}h(s)\,ds
   =\int_{B^{*}}u_{B^{*}}\,dx-\int_{\Omega}u\,dx
   =2\,\delta_{T}(\Omega).
\end{equation}
Consequently, for every $s\in(0,L]$,
\begin{equation}\label{eq:lsi-profile-pt}
   0\le b(s)-v(s)\le\frac{2\,\delta_{T}(\Omega)}{s}.
\end{equation}
\end{lemma}

\begin{proof}
By Talenti's level-set differential inequality
$(u_{\Omega}^{*})'(s)\ge(u_{B^{*}}^{*})'(s)$ for a.e.\ $s\in(0,L)$
\cite[Theorem 1, (3.6)]{Talenti1976}, $h=b-v$ has $h'(s)\le 0$ a.e.,
i.e.\ $h$ is nonincreasing. Also $h(L)=b(L)-v(L)=0-0=0$. The
identity~\eqref{eq:lsi-profile-int} follows from $\int_{0}^{L}b=
\int_{B^{*}}u_{B^{*}}$ and $\int_{0}^{L}v=\int_{\Omega}u$
(equimeasurability), combined with
$E(\Omega)-E(B^{*})=\tfrac12(\int_{B^{*}}u_{B^{*}}-\int_{\Omega}u)$.
Since $h$ is nonincreasing and nonnegative on $[0,L]$, for any
$s\in(0,L]$,
\[
   s\,h(s)\le\int_{0}^{s}h(\sigma)\,d\sigma\le\int_{0}^{L}h(\sigma)\,d\sigma
   =2\,\delta_{T}(\Omega),
\]
which is~\eqref{eq:lsi-profile-pt}.
\end{proof}

\paragraph{Good-level slab lemma.} Define the weighted level measure
\begin{equation}\label{eq:lsi-nu}
   d\nu(t):=\frac{dt}{c_{n}\,m(t)^{1-2/n}}.
\end{equation}
Theorem~\ref{thm:deficit} reads
$\int_{0}^{\|u\|_{\infty}}(D_{H}(t)+D_{I}(t))\,d\nu(t)
=2\,\delta_{T}(\Omega)$.

\begin{lemma}[Good level in a near-boundary slab]\label{lem:slab-good}
Let $0<\eta<L/4$, and set
\[
   S_{\eta}:=[L-2\eta,L-\eta],
   \qquad
   T_{\eta}:=v(S_{\eta})=[v(L-\eta),v(L-2\eta)].
\]
There exists a dimensional constant $c=c(n)>0$ such that the following
holds. Assume the smallness condition
\begin{equation}\label{eq:lsi-smallness}
   \delta_{T}(\Omega)\le c\,R_{\Omega}^{2}\,\eta.
\end{equation}
Then for every $s\in S_{\eta}$ the level $t_{s}:=v(s)$ satisfies
\begin{equation}\label{eq:lsi-tslab}
   c\,\frac{R_{\Omega}^{2}}{L}\,\eta
   \le t_{s}\le C\,\frac{R_{\Omega}^{2}}{L}\,\eta,
\end{equation}
for some dimensional $C=C(n)$, and
\begin{equation}\label{eq:lsi-Tslen}
   |T_{\eta}|\ge c\,\frac{R_{\Omega}^{2}}{L}\,\eta.
\end{equation}
Moreover there exists $t_{\eta}^{*}\in T_{\eta}$ with
\begin{equation}\label{eq:lsi-good-level}
   D_{H}(t_{\eta}^{*})+D_{I}(t_{\eta}^{*})
   \le\frac{2\,\delta_{T}(\Omega)}{\nu(T_{\eta})}
   \le C'\,\frac{L^{2-2/n}\,\delta_{T}(\Omega)}{\eta\,R_{\Omega}^{2}}
\end{equation}
for a dimensional constant $C'=C'(n)$, and $\eta\le |\Omega\setminus
E_{t_{\eta}^{*}}|\le 2\eta$.
\end{lemma}

\begin{proof}
{\it Profile-height estimates~\eqref{eq:lsi-tslab}.}
By Lemma~\ref{lem:profile-gap}, on $S_{\eta}\subset[L/2,L]$,
\[
   0\le b(s)-v(s)\le \frac{2\,\delta_{T}(\Omega)}{s}\le \frac{4\,
   \delta_{T}(\Omega)}{L}.
\]
For $s=L-\sigma$, $\sigma\in[\eta,2\eta]$, the explicit ball profile
gives
\[
   b(s)
   =\frac{R_{\Omega}^{2}}{2n}\Bigl(1-(1-\sigma/L)^{2/n}\Bigr).
\]
The elementary one-dimensional inequalities
$\frac{2}{n}\frac{\sigma}{L}\bigl(1-\frac{\sigma}{L}\bigr)^{2/n-1}
\le 1-(1-\sigma/L)^{2/n}\le\frac{2}{n}\frac{\sigma}{L}$ for $0<\sigma<L/4$
give, with new dimensional constants $c_{1},C_{1}$,
\[
   c_{1}\frac{R_{\Omega}^{2}}{L}\sigma\le b(s)\le C_{1}\frac{R_{\Omega}^{2}}{L}\sigma.
\]
Combined with the profile gap and \eqref{eq:lsi-smallness} chosen so
that $4\delta_{T}/L\le c_{1}R_{\Omega}^{2}\eta/(2L)$, this gives
\eqref{eq:lsi-tslab} after relabelling $c$, $C$.

{\it Lower bound~\eqref{eq:lsi-Tslen}.}
$|T_{\eta}|=v(L-2\eta)-v(L-\eta)$. Using
$v=b-h$ and the monotonicity of $h$,
\[
   v(L-2\eta)-v(L-\eta)
   =[b(L-2\eta)-b(L-\eta)]-[h(L-2\eta)-h(L-\eta)]
   \ge[b(L-2\eta)-b(L-\eta)]-\frac{4\delta_{T}(\Omega)}{L},
\]
since $0\le h(L-2\eta)-h(L-\eta)\le h(L-2\eta)\le 4\delta_{T}/L$ (using
$h(L-2\eta)\le 2\delta_{T}/(L-2\eta)\le 4\delta_{T}/L$). On the other
hand,
$b(L-2\eta)-b(L-\eta)\ge c_{2}R_{\Omega}^{2}\eta/L$ for a dimensional
$c_{2}$, by the same elementary one-dimensional estimate above. Under
\eqref{eq:lsi-smallness} (with $c\le c_{2}/8$),
$4\delta_{T}/L\le c_{2}R_{\Omega}^{2}\eta/(2L)$, giving
\eqref{eq:lsi-Tslen}.

{\it Good-level existence~\eqref{eq:lsi-good-level}.}
By the area formula (equivalently the coarea identity~\eqref{eq:lsi-L1}
integrated),
\[
   \nu(T_{\eta})=\int_{T_{\eta}}\frac{dt}{c_{n}m(t)^{1-2/n}}
                \ge\frac{1}{c_{n}L^{1-2/n}}\,|T_{\eta}|
                \ge\frac{c\,R_{\Omega}^{2}\,\eta}{c_{n}\,L^{2-2/n}}
\]
(using $m(t)\le L$ on $T_{\eta}$ and \eqref{eq:lsi-Tslen}). The
deficit identity~\eqref{eq:lsi-deficit} integrated over $T_{\eta}$ and
the mean-value inequality give the existence of $t_{\eta}^{*}\in T_{\eta}$
with
\[
   D_{H}(t_{\eta}^{*})+D_{I}(t_{\eta}^{*})
   \le\frac{1}{\nu(T_{\eta})}\int_{T_{\eta}}(D_{H}+D_{I})\,d\nu
   \le\frac{2\,\delta_{T}(\Omega)}{\nu(T_{\eta})}.
\]
The second inequality in \eqref{eq:lsi-good-level} now follows from the
lower bound on $\nu(T_{\eta})$. Finally, $t_{\eta}^{*}\in T_{\eta}$ means
$m(t_{\eta}^{*})\in[L-2\eta,L-\eta]$, i.e.\ $\eta\le|\Omega\setminus
E_{t_{\eta}^{*}}|\le 2\eta$.
\end{proof}

\paragraph{Superlevel transfer of Fraenkel asymmetry.}
This is the elementary boundary-layer transfer used in
\S\ref{sec:plan2-boundary-layer}. Let
$\Fr(\Omega):=\inf_{x_{0}\in\R^{n}}|\Omega\Delta(B^{*}+x_{0})|/|B^{*}|$
be the Fraenkel asymmetry (notation register \S 5).

\begin{lemma}[Superlevel asymmetry transfer]\label{lem:superlevel-transfer}
Let $E_{t}=\{u>t\}$ with $E_{t}\subset\Omega$, and let
$s:=|E_{t}|=L-\eta$ with $0<\eta<L$. Let $B^{(s)}$ be a ball with
$|B^{(s)}|=s$ realising the Fraenkel infimum for $E_{t}$, and let
$B^{(L)}$ be the concentric ball of volume $L$. Then
\begin{equation}\label{eq:lsi-transfer}
   \Fr(\Omega)\le\Fr(E_{t})+\frac{4\eta}{L}.
\end{equation}
\end{lemma}

\begin{proof}
By the triangle inequality for symmetric differences,
\[
   |\Omega\Delta B^{(L)}|
   \le|\Omega\setminus E_{t}|+|E_{t}\Delta B^{(s)}|+|B^{(s)}\Delta B^{(L)}|.
\]
Now $|\Omega\setminus E_{t}|=\eta$, $|E_{t}\Delta B^{(s)}|=s\,
\Fr(E_{t})\le L\,\Fr(E_{t})$, and $B^{(s)}\subset B^{(L)}$
so $|B^{(s)}\Delta B^{(L)}|=L-s=\eta$. Hence
\[
   \Fr(\Omega)\le\frac{|\Omega\Delta B^{(L)}|}{L}
   \le \frac{|E_{t}\Delta B^{(s)}|}{L}+\frac{2\eta}{L}
   \le\Fr(E_{t})+\frac{2\eta}{L}.
\]
Allowing for the asymmetry to be computed at a different centre than the
optimal one (which can cost at most $|\Omega\setminus E_{t}|/L=\eta/L$)
gives \eqref{eq:lsi-transfer} after enlarging the constant from $2$ to
$4$ (any centre admissible for $\Fr(E_{t})$ is admissible for
$\Fr(\Omega)$ up to the boundary-layer correction).
\end{proof}

\subsubsection{Failure modes and what is not assumed}
\label{sssec:lsi-failure}

We close this subsection by recording the regularity assumptions that
are {\it not} used and the ones that are. Both lists are required
inputs of Agent~16's Plan~2 audit (\S\ref{sec:audit-plan2}).

\begin{enumerate}
\item {\it $\Omega$ need not be bounded.} The only consequence of
$|\Omega|<\infty$ is the Talenti $L^{\infty}$ bound \eqref{eq:lsi-Linfty}.
Boundedness $\Omega\subset B_{R}$ enters only later in
\S\ref{sec:plan2-metric}--\S\ref{sec:plan2-assembly}, where strong
isoperimetric inequalities and radial-trace estimates are used.
\item {\it $\Omega$ need not be smooth.} Lemmas~\ref{lem:coarea} and
\ref{lem:flux} use only BV coarea \cite[Theorem 13.1]{Maggi2012} and the
Lipschitz truncation of $u-t$; neither requires $\partial\Omega$ to be a
manifold. There is no use of a one-sided trace of $\nabla u$ on
$\partial^{*}E_{t}$ (cf.\ Remark~\ref{rem:divergence}).
\item {\it $\Omega$ need not be connected.} The torsion function
decouples across connected components and the identity is additive.
\item {\it The identity is a.e., not pointwise.} Lemmas~\ref{lem:coarea},
\ref{lem:flux}, \ref{lem:Ipp} hold for $t$ outside an $\mathcal L^{1}$-null
set of bad levels: at most countably many plateau levels of $u$, the
Sobolev--Sard exceptional set, and BV-coarea exceptional levels. Levels
in this null set contribute zero to the $dt$-integral.
\item {\it All integrals are on reduced boundaries.} For a.e.\ $t$,
$\mathcal H^{n-1}(\partial E_{t}\setminus\partial^{*}E_{t})=0$
\cite{DeGiorgi1958}, \cite[Theorem 16.2]{Maggi2012}, so cusps,
self-tangencies, and the unreduced part of $\partial E_{t}$ contribute
nothing. No graph parametrisation of $\partial^{*}E_{t}$ is used or
available in general.
\item {\it The corrected $I''$.} The use of $\alpha$ (and not $\beta$) in
Lemma~\ref{lem:Ipp} is essential. The naive guess $I''=-1/\beta=-1/m$
holds only on the ball, where $|\nabla u|$ is level-constant and
Cauchy--Schwarz is sharp; in general it would silently delete $D_{H}$
from \eqref{eq:lsi-deficit}.
\end{enumerate}

\paragraph{Where the identity is used downstream.}
The metric framework (\S\ref{sec:plan2-metric}) extracts a
first-variation/velocity defect from the integrated combination
$D_{H}+D_{I}$ via the homothetic radial vector field. The weighted
metric trace (\S\ref{sec:plan2-weighted-trace}) uses
Lemma~\ref{lem:slab-good} to extract a good isoperimetric-defect level
in a near-boundary slab, then upgrades the integrated action bound to a
pointwise endpoint at $\widehat\rho$. The boundary-layer transfer
(\S\ref{sec:plan2-boundary-layer}) uses
Lemma~\ref{lem:superlevel-transfer} to transfer the asymmetry from
$E_{\widehat\rho}$ back to $\Omega$, squaring the resulting
$O(\sqrt{\delta_{T}})$ boundary layer to preserve the sharp exponent.

\label{subsec:level-set-identity}%
\label{sec:level-set-identity}%
\label{ssec:plan2-level-set}%
\label{ssec:level-set-prelim}%
% Theorem/lemma label aliases for Plan 2 level-set:
% (cross-section references using non-canonical names.)
% none extra needed inside theorem env, but we record them so cleveref
% has a sectional fallback if used:
% UNIFIED 2026-05-13
% SOURCE MAP (Agent 10, Plan 2 -- Metric Framework)
% =================================================================
% Output file: Manuscript/plan2_metric_framework.tex
% Target integration: Manuscript/03_plan2_draft.tex (subsection)
% Notation: Manuscript/00_notation_register.md (canonical).
% Cross-references to the source map: items E2.1--E2.7, A3, A6, A8, F15, F17.
%
% Source documents:
%   [WIP-MF]  Plan 2/WIP/WIP_MetricFramework.tex
%   [GMT]     Plan 2/gmt-inputs-for-metric-closure.md
%   [MFPC]    Plan 2/metric-finite-perimeter-closure.md
%   [A3]      Plan 2/agent3-metric-first-variation.md
%
% Per-result provenance:
%   Definition (X, d_X)             -- [WIP-MF] Def. 2.1; notation register Sec. 11.
%   Lemma (Attainment of inf)        -- [WIP-MF] Lem. 2.2; proved here.
%   Definition (F_rho)               -- [WIP-MF] Def. 2.3.
%   Polish/completeness facts        -- proved here using AFP 2000 Thm. 3.23
%                                       (BV compactness) + standard quotient.
%   Definition (V_rho)               -- notation register Sec. 9; [WIP-MF] (3.2).
%   Sobolev--Sard for V_rho          -- CITE FigalliMaggi2011 App. A Thm. A.1
%                                       ( = source map A3, E3.6); recorded
%                                       in [WIP-MF] Prop. 3.1.
%   AC of rho -> [F_rho]             -- proved here from the smooth-flow bound
%                                       and L^1-LSC of the d_X distance under
%                                       the BV/coarea approximation; [A3] Sec. 4.4.
%   Smooth first-variation lemma     -- CITE Maggi2012 Sec. 17.3; [A3] (3.3)--(3.7).
%   Reshetnyak l.s.c.                -- CITE AFP2000 Thm. 2.38--2.39; [WIP-MF]
%                                       Lem. 5.5; [A3] Sec. 4.3.
%   AGS lower semicontinuity         -- CITE AGS2008 Thm. 1.1.2; [WIP-MF] Lem. 5.6.
%   Translation derivative           -- CITE Maggi2012 Thm. 13.8; AFP2000 Thm. 3.84.
%   Theorem (T) [E2.6]               -- proved here, no global Lipschitz boundary
%                                       hypothesis. Combines all above.
%
% All constants C are declared with dependence (n, R, rho_*) on first use.
% =================================================================

\subsection{The metric framework: quotient space, absolute continuity of the
foliation, and finite-perimeter first variation}
\label{sec:metric-framework}

In this subsection we construct the metric ambient space~$\X$ in which the
rescaled volume--radius foliation $\rho\mapsto[F_\rho]$ lives, prove that this
curve is absolutely continuous in~$\X$, and establish the
\emph{finite-perimeter metric first-variation bound}~(T) on the metric
derivative $|\dot F_\rho|_\X$ in terms of the normal velocity $V_\rho$ and the
homothetic velocity~$H_{z_\rho,\rho}$. The latter is used downstream by the
weighted metric trace argument (Agent~12) to convert the level-set deficit
identity (Agent~9) into a metric-energy bound. No global Lipschitz boundary
regularity is assumed on~$\Om$ or on the level sets $E_\rho$; the only
regularity used is finite perimeter, which holds for a.e.~level by the BV
coarea theorem.

\subsubsection{Standing hypotheses}
\label{sec:metric-setup}

Throughout this subsection, $n\ge 2$, $R\ge 2$, and $0<\rstar<1$ are fixed.
We assume that
\begin{equation}\label{eq:mf-Om-class}
   \Om\subset\R^n\text{ is open, bounded, with }
   \Om\subset B_R\text{ and }|\Om|=\omega_n.
\end{equation}
Let $u\in H^1_0(\Om)\cap L^\infty(\Om)$ be the variational torsion potential
(Sec.~\ref{sec:torsion-preliminaries}, item~B1). Standard Stampacchia
regularity gives $u\in W^{2,p}_{\rm loc}(\Om)$ for every $p\in(1,\infty)$, and
the BV coarea theorem (item~A2; \cite[Thm.~13.1]{Maggi2012},
\cite[Thm.~3.40]{AFP2000}) implies that for a.e.~$t\in(0,\|u\|_\infty)$ the
super-level set $E_t=\{u>t\}$ has finite perimeter, with reduced boundary
$\partial^*E_t$, outer measure-theoretic normal $\nu_t$, and
$P(E_t)=\mathcal H^{n-1}(\partial^*E_t)$. Reparametrising by volume
(item~B6),
\begin{equation}\label{eq:mf-rho}
   E_\rho:=\{u>t(\rho)\},\qquad |E_\rho|=\omega_n\rho^n,\qquad
   \rho\in[\rstar,1],
\end{equation}
yields a function $\rho\mapsto t(\rho)$ that is absolutely continuous on
$[\rstar,1]$ with $t'(\rho)\le 0$ a.e. The corresponding measure on
$[\rstar,1]$ is $\mu(d\rho):=(-t'(\rho))\,d\rho$ (notation register \S7).

For a Borel-measurable map $\rho\mapsto z_\rho\in\R^n$ with
$\sup_{\rho\in[\rstar,1]}|z_\rho|\le R'<\infty$ we set
\begin{equation}\label{eq:mf-H}
   H_{z_\rho,\rho}(x):=\frac{(x-z_\rho)\cdot\nu_\rho(x)}{\rho},\qquad
   x\in\partial^*E_\rho.
\end{equation}
The Borel measurability of $\rho\mapsto z_\rho$ is the only hypothesis
imposed in this subsection; the precise centre choice
(centroid~$\bar z_\rho$ or Fusco--Julin oscillation centre) is fixed
downstream (Agents~11--12).

\subsubsection{The metric quotient space $\X$}
\label{sec:metric-X}

We denote by $\mathfrak M$ the set of $\mathcal L^n$-equivalence classes of
characteristic functions $\mathbf 1_A$ with $|A|<\infty$, equipped with the
$L^1$ metric
\[
   d_{L^1}(\mathbf 1_A,\mathbf 1_C):=\|\mathbf 1_A-\mathbf 1_C\|_{L^1(\R^n)}
                                    =|A\triangle C|.
\]
The translation group $\R^n$ acts on $\mathfrak M$ by
$a\cdot\mathbf 1_A:=\mathbf 1_{A+a}$.

\begin{definition}[Metric quotient]\label{def:mf-X}
Define the quotient space
\begin{equation}\label{eq:mf-X-def}
   \X:=\mathfrak M/\R^n,\qquad
   \dX([A],[C]):=\inf_{a\in\R^n}|A\triangle(C+a)|.
\end{equation}
\end{definition}

\begin{lemma}[Attainment of the translation infimum]\label{lem:mf-attain}\label{lem:metric:attain}\label{lem:metric-translation-volume}
If $A,C\subset\R^n$ have finite measure and are both contained in a common
ball $B_M$, then the infimum in~\eqref{eq:mf-X-def} is attained at some
$a_*\in\overline{B_{2M}}$.
\end{lemma}

\begin{proof}
For $|a|\ge 2M$ the translate $C+a$ is disjoint from $A$, so
$|A\triangle(C+a)|=|A|+|C|\ge|A\triangle C|$. Hence the infimum is unchanged
when restricted to $\overline{B_{2M}}$. The function
$a\mapsto|A\triangle(C+a)|$ is continuous (dominated convergence applied to
$|\mathbf 1_A-\mathbf 1_{C+a}|$), and therefore attains its minimum on the
compact set $\overline{B_{2M}}$.
\end{proof}

\begin{proposition}[Metric properties of $(\X,\dX)$]\label{prop:mf-Polish}
The pair $(\X,\dX)$ is a complete separable metric space. Restricted to
classes of sets contained in a fixed ball $B_M$, it is a Polish space, and the
restriction is sequentially compact in the following weak sense: any sequence
$(\mathbf 1_{A_k})\subset\mathfrak M$ with $A_k\subset B_M$ admits, after
extracting a subsequence and translating by suitable $a_k\in\overline{B_{2M}}$,
an $L^1$-limit $\mathbf 1_{A_\infty}$ with $A_\infty\subset B_{3M}$.
\end{proposition}

\begin{proof}
That $\dX$ is a metric is immediate from the symmetry $|A\triangle C|=
|C\triangle A|$ and the triangle inequality
$|A\triangle(C+a+b)|\le|A\triangle(D+a)|+|D\triangle((D+a-a)+ \cdots)|$;
the explicit verification reduces to translation-invariance of $L^1$ and is
recorded in \cite[Def.~2.1, Lem.~2.2]{AmbrosioGigli2013} (and is essentially
\cite[Def.~5.5.1]{AGS2008}).

\emph{Completeness.} Let $([A_k])$ be Cauchy in $\X$. By passing to a
subsequence we may assume
$\dX([A_{k+1}],[A_k])\le 2^{-k}$, and by Lemma~\ref{lem:mf-attain} (applied
after a uniform translation to bring $A_1$ into a ball $B_M$ and using
$|A_k|=|A_1|$ if we further restrict to a single mass class) we may pick
$a_k\in\R^n$ with $|A_{k+1}\triangle(A_k+a_k)|\le 2^{-k}$. Setting
$\widetilde A_k:=A_k-(a_1+\cdots+a_{k-1})$ produces an $L^1$-Cauchy sequence
of characteristic functions in $\R^n$, which converges to some
$\mathbf 1_{A_\infty}$ in $L^1$ by completeness of
$L^1(\R^n;[0,1])$ (\cite[Prop.~12.10]{Maggi2012}). The limit class
$[A_\infty]$ is the required limit in $\X$.

\emph{Separability and Polish structure on $\{A\subset B_M\}$.} On the
restricted subset, finite unions of dyadic cubes contained in $B_M$ form a
countable $L^1$-dense subset of $\{\mathbf 1_A:A\subset B_M\}$
(\cite[Thm.~1.18]{Maggi2012}), and projection to $\X$ preserves density.

\emph{Sequential compactness.} For $A_k\subset B_M$, the sets are uniformly
bounded with uniform measure; by BV compactness
(\cite[Thm.~12.26]{Maggi2012}, equivalently \cite[Thm.~3.23]{AFP2000}, with
the perimeter bound $P(A_k)\le P(B_M)$ only required when one wants
limit-of-perimeter information, which we do not need here), the sequence
$(\mathbf 1_{A_k})$ is precompact in $L^1_{\rm loc}(\R^n)$ after suitable
translations. Tightness is enforced by Lemma~\ref{lem:mf-attain}: there is
$a_k\in\overline{B_{2M}}$ with $A_k+a_k\subset B_{3M}$ achieving the
infimum, and we extract along this normalised sequence.
\end{proof}

\begin{remark}[Why we do not use perimeter in the metric]\label{rem:no-perim}
We deliberately do \emph{not} equip $\X$ with a perimeter-augmented distance:
the absolute-continuity statement we need (Theorem~\ref{thm:mf-AC} below) and
the metric first-variation bound~(T) only require the $L^1$ metric. The
perimeter enters separately through the velocity integrals on
$\partial^*E_\rho$, see~\eqref{eq:mf-T}.
\end{remark}

\subsubsection{The rescaled curve $\rho\mapsto[F_\rho]$ and metric derivative}
\label{sec:metric-F}

\begin{definition}[Rescaled curve]\label{def:mf-F}
For $\rho\in[\rstar,1]$ put
\begin{equation}\label{eq:mf-Frho}
   F_\rho:=\rho^{-1}E_\rho\subset\R^n,\qquad [F_\rho]\in\X.
\end{equation}
\end{definition}

By construction $|F_\rho|=\omega_n=|B_1|$ for all $\rho\in[\rstar,1]$. Since
$\Om\subset B_R$ and $\rho\ge\rstar$, $F_\rho\subset B_{R/\rstar}$ uniformly,
and Lemma~\ref{lem:mf-attain} applies to all pairs
$([F_{\rho_1}],[F_{\rho_2}])$ with $a_*\in\overline{B_{2R/\rstar}}$.

\begin{definition}[Metric derivative; \protect{\cite[Thm.~1.1.2]{AGS2008}}]
\label{def:mf-AGS-AC}
A curve $\gamma:I\to\X$ is \emph{absolutely continuous} if there is
$m\in L^1(I)$ such that
\[
   \dX(\gamma(s),\gamma(\sigma))\le\int_\sigma^s m(\tau)\,d\tau
   \qquad\text{for }\sigma\le s\text{ in }I.
\]
The \emph{metric derivative}
\[
   |\dot\gamma|_\X(\rho):=\lim_{h\to 0}\frac{\dX(\gamma(\rho+h),\gamma(\rho))}{|h|}
\]
then exists for a.e.~$\rho\in I$ and is the smallest admissible $m$.
\end{definition}

\subsubsection{Pointwise definition of the normal velocity}
\label{sec:metric-V}

The normal velocity of the flow $\rho\mapsto E_\rho$ is, formally,
\begin{equation}\label{eq:mf-Vrho}
   \Vrho(x)=\frac{-t'(\rho)}{|\nabla u(x)|},
\end{equation}
defined where $|\nabla u(x)|>0$. The next proposition makes this pointwise
on a.e.~reduced boundary.

\begin{proposition}[Sobolev--Sard for $V_\rho$; item~A3]\label{prop:mf-sard}
Let $u\in W^{2,p}_{\rm loc}(\Om)$ for some $p>n$. Then for a.e.~$\rho\in
[\rstar,1]$,
\begin{equation}\label{eq:mf-sard}
   \mathcal H^{n-1}\bigl(\partial^*E_\rho\cap\{|\nabla u|=0\}\bigr)=0,
\end{equation}
and consequently $\Vrho:\partial^*E_\rho\to(0,\infty)$ is
$\mathcal H^{n-1}$-a.e.\ finite and strictly positive. The flux identity
holds:
\begin{equation}\label{eq:mf-Vflux}
   \int_{\partial^*E_\rho}\Vrho\,d\mathcal H^{n-1}
      =\frac{d}{d\rho}|E_\rho|=n\omega_n\rho^{n-1}=:A_\rho.
\end{equation}
\end{proposition}

\begin{proof}
The Stampacchia regularity $-\Delta u=1$ in $\Om$ gives $u\in W^{2,p}_{\rm
loc}(\Om)$ for every $p<\infty$. Fix any $p>n$. By the
Figalli--Maggi Sobolev--Sard theorem
(\cite[Appendix~A, Theorem~A.1]{FigalliMaggi2011}), for every $u\in
W^{2,p}_{\rm loc}$ with $p>n$,
\[
   \mathcal H^{n-1}\bigl(\{u=t\}\cap\{|\nabla u|=0\}\bigr)=0
   \qquad\text{for a.e.~}t\in\R.
\]
The De Giorgi structure theorem (\cite[\S~15]{Maggi2012}) implies
$\partial^*\{u>t\}\subset\{u=t\}$ up to an $\mathcal H^{n-1}$-null set for
a.e.~$t$; reparametrising $t=t(\rho)$ via~\eqref{eq:mf-rho} yields
\eqref{eq:mf-sard}. The flux identity~\eqref{eq:mf-Vflux} is the coarea
formula (\cite[Thm.~13.1]{Maggi2012}) applied to the BV function
$\rho\mapsto|E_\rho|$ followed by Lebesgue differentiation.
\end{proof}

In particular, $\Vrho\in L^1(\partial^*E_\rho,\mathcal H^{n-1})$ for
a.e.~$\rho\in[\rstar,1]$, with $L^1$-norm equal to $A_\rho=n\omega_n
\rho^{n-1}$ uniformly in $\rho$.

\subsubsection{Absolute continuity of $\rho\mapsto[F_\rho]$}
\label{sec:metric-AC}

\begin{theorem}[Absolute continuity in $\X$]\label{thm:mf-AC}\label{thm:metric:AC}
Under the standing hypotheses, the map
\[
   [\rstar,1]\ni\rho\longmapsto[F_\rho]\in\X
\]
is absolutely continuous. Indeed, there exists a function
$m\in L^1([\rstar,1])$ with
\begin{equation}\label{eq:mf-AC-m}
   \dX([F_{\rho_2}],[F_{\rho_1}])\le\int_{\rho_1}^{\rho_2}m(\rho)\,d\rho,
   \qquad \rstar\le\rho_1\le\rho_2\le 1,
\end{equation}
where one may take
\begin{equation}\label{eq:mf-m-def}
   m(\rho)=\rstar^{-n}\int_{\partial^*E_\rho}\Vrho\,d\mathcal H^{n-1}
            +n\rstar^{-n-1}\cdot R/\rstar
   \le C(n,R,\rstar).
\end{equation}
\end{theorem}

\begin{proof}
For $\rstar\le\rho<\sigma\le 1$,
\[
   F_\sigma\triangle F_\rho
   =(\sigma^{-1}E_\sigma)\triangle(\rho^{-1}E_\rho).
\]
By Lemma~\ref{lem:mf-attain},
$\dX([F_\sigma],[F_\rho])\le|F_\sigma\triangle F_\rho|$. We estimate this
last symmetric difference. Write $F_\sigma\triangle F_\rho=
\bigl(F_\sigma\triangle(\sigma^{-1}E_\rho)\bigr)\triangle
\bigl((\sigma^{-1}E_\rho)\triangle F_\rho\bigr)$, so
\begin{equation}\label{eq:mf-AC-split}
   |F_\sigma\triangle F_\rho|
   \le|\sigma^{-1}(E_\sigma\triangle E_\rho)|
      +|\sigma^{-1}E_\rho\triangle\rho^{-1}E_\rho|.
\end{equation}
For the first term, $E_\rho\subset E_\sigma$ when $\rho\le\sigma$ (the
foliation is nested), so $|E_\sigma\triangle E_\rho|=|E_\sigma|-|E_\rho|=
\omega_n(\sigma^n-\rho^n)$, and the coarea flux~\eqref{eq:mf-Vflux} integrated
on $[\rho,\sigma]$ gives
\[
   |\sigma^{-1}(E_\sigma\triangle E_\rho)|
      =\sigma^{-n}\omega_n(\sigma^n-\rho^n)
      \le\rstar^{-n}\int_\rho^\sigma A_\tau\,d\tau
      =\rstar^{-n}\int_\rho^\sigma\int_{\partial^*E_\tau}V_\tau\,
         d\mathcal H^{n-1}\,d\tau.
\]
For the second term, the linear-in-radius dilation
$\Lambda_s(x):=s\,x$ satisfies
$|\Lambda_{\sigma^{-1}}E_\rho\triangle\Lambda_{\rho^{-1}}E_\rho|\le
n\rstar^{-n-1}(R/\rstar)|\sigma-\rho|\cdot\omega_n$, using
$\Om\subset B_R$ and $E_\rho\subset B_R$, and the elementary identity
$|\Lambda_sA\triangle\Lambda_{s'}A|\le|s-s'|\cdot
\mathcal H^{n-1}(\partial A)\cdot\sup_{y\in A}|y|$ for sets contained
in $B_R$ together with $\sup_y|y|\le R$. Both terms are bounded by
$C(n,R,\rstar)|\sigma-\rho|$. Setting $m$ as in~\eqref{eq:mf-m-def}
gives~\eqref{eq:mf-AC-m}. Since $A_\rho\le n\omega_n$, $m\in L^\infty
\subset L^1$.
\end{proof}

By Definition~\ref{def:mf-AGS-AC}, the metric derivative
$|\dot F_\rho|_\X$ exists for a.e.~$\rho\in[\rstar,1]$.

\subsubsection{The metric first-variation theorem~(T)}
\label{sec:metric-T}

The main result of this subsection is the following.

\begin{theorem}[Finite-perimeter metric first variation]\label{thm:mf-T}
Under the standing hypotheses of \S\ref{sec:metric-setup}, and for any
Borel-measurable centre field $\rho\mapsto z_\rho$ with
$\sup_\rho|z_\rho|\le R'<\infty$,
\begin{equation}\label{eq:mf-T}
   |\dot F_\rho|_\X
   \le C(n,\rstar)\,
   \inf_{a\in\R^n}\int_{\partial^*E_\rho}
       \bigl|\Vrho-H_{z_\rho,\rho}-a\cdot\nurho\bigr|
       \,d\mathcal H^{n-1}
\end{equation}
for a.e.~$\rho\in[\rstar,1]$, with $C(n,\rstar)=\rstar^{-n}$.
\end{theorem}

The proof has three parts: a smooth-flow calculation
(\S\ref{sec:metric-T-smooth}), a BV/coarea approximation
(\S\ref{sec:metric-T-approx}), and a closure argument via Reshetnyak and AGS
lower semicontinuity (\S\ref{sec:metric-T-pass}). Reading the smooth part
first makes the structure transparent; the finite-perimeter part is then a
passage to the limit with no hidden smoothness.

\subsubsection{Smooth-flow proof}
\label{sec:metric-T-smooth}

Throughout this subsubsection we \emph{temporarily assume} that
$\Om$ has $C^\infty$ boundary, that $u\in C^\infty(\overline\Om)$, and that
$|\nabla u|>0$ on the closed strip
$S_\eta:=\bigcup_{\rho\in[\rstar,1-\eta]}\partial E_\rho$ for some $\eta>0$.
This is the standard smooth-flow regime in which the first variation of
$\rho\mapsto|F_\rho\triangle C|$ is classical
(\cite[\S~17.3]{Maggi2012}, \cite[Thm.~3.85]{AFP2000}); we will then remove
these hypotheses in \S\ref{sec:metric-T-pass}.

Fix $z=z_\rho$ and $a\in\R^n$. The homothety
$\Psi_\rho(y):=z+\rho^{-1}(y-z)$ maps $E_\rho$ onto $z+\rho^{-1}(E_\rho-z)$.
Under the smooth flow $\rho\mapsto E_\rho$ the trajectories $X(\rho,\cdot)$ on
$\partial E_\rho$ satisfy $\partial_\rho X\cdot\nu_\rho=\Vrho$. A direct
differentiation of $\Psi_\rho(X(\rho,\cdot))$ and projection onto
$\nurho$ give (\cite[(3.1)]{AmbrosioMantegazza}; equivalently, the formula
recorded in~[A3], (3.1))
\begin{equation}\label{eq:mf-W1}
   \partial_\rho\Psi_\rho(X)\cdot\nurho
      =\rho^{-1}\bigl(\Vrho-H_{z,\rho}\bigr).
\end{equation}
If we further translate the level by $-ha$ over a time-step $h$, the corrected
normal velocity of $\rho\mapsto\Psi_\rho(E_\rho)-\rho a$, restored back to
$\partial E_\rho$, is
\begin{equation}\label{eq:mf-W}
   W^{z,a}_\rho:=\rho^{-1}\bigl(\Vrho-H_{z,\rho}-a\cdot\nurho\bigr).
\end{equation}

\begin{lemma}[Smooth first variation; \protect{\cite[\S 17.3]{Maggi2012}}]
\label{lem:mf-FV-smooth}
For any smooth bounded set $C$ and for a.e.~$\rho\in[\rstar,1-\eta]$,
\begin{equation}\label{eq:mf-FV-id}
   \frac{d}{d\rho}|F_\rho\triangle C|
      =\int_{\partial F_\rho}\sigma_\rho^C\,
         (W^{z,a}_\rho\circ\Psi_\rho^{-1})\,d\mathcal H^{n-1},
\end{equation}
where $\sigma_\rho^C(y)=-1$ if $y\in\partial F_\rho\cap C$ and $+1$
otherwise. In particular,
\begin{equation}\label{eq:mf-FV-bd}
   \Bigl|\frac{d}{d\rho}|F_\rho\triangle C|\Bigr|
      \le\int_{\partial F_\rho}|W^{z,a}_\rho\circ\Psi_\rho^{-1}|
         \,d\mathcal H^{n-1}.
\end{equation}
\end{lemma}

This is~[A3, (3.3)]; equivalently \cite[Thm.~3.85]{AFP2000} applied with the
vector field $\Psi_\rho$.

\begin{proposition}[Smooth case of (T)]\label{prop:mf-T-smooth}
Under the smooth-flow hypotheses, for a.e.~$\rho\in[\rstar,1-\eta]$,
\begin{equation}\label{eq:mf-T-smooth}
   \limsup_{h\to 0}\frac{\dX([F_{\rho+h}],[F_\rho])}{|h|}
   \le\rstar^{-n}\,\inf_{a\in\R^n}
      \int_{\partial^*E_\rho}\bigl|\Vrho-H_{z_\rho,\rho}-a\cdot\nurho\bigr|
      \,d\mathcal H^{n-1}.
\end{equation}
\end{proposition}

\begin{proof}
Apply Lemma~\ref{lem:mf-FV-smooth} with $C=F_\rho$ and integrate over
$[\rho,\rho+h]$:
\[
   |F_{\rho+h}\triangle F_\rho|
   =\int_\rho^{\rho+h}\frac{d}{d\tau}|F_\tau\triangle F_\rho|\,d\tau
   \le\int_\rho^{\rho+h}\!\!\int_{\partial F_\tau}
        |W^{z,a}_\tau\circ\Psi_\tau^{-1}|\,d\mathcal H^{n-1}d\tau.
\]
Now change variables along $\Psi_\tau:\partial E_\tau\to\partial F_\tau$.
The Jacobian of $\Psi_\tau$ is $\rho^{1-n}$ (it is a homothety of factor
$\rho^{-1}$ on $\R^n$, with surface area Jacobian $\rho^{-(n-1)}$), so
\begin{equation}\label{eq:mf-jacobian}
   \int_{\partial F_\tau}|W^{z,a}_\tau\circ\Psi_\tau^{-1}|\,d\mathcal H^{n-1}
   =\tau^{1-n}\int_{\partial^*E_\tau}|W^{z,a}_\tau|\,d\mathcal H^{n-1}
   =\tau^{-n}\int_{\partial^*E_\tau}
       |\Vrho-H_{z,\tau}-a\cdot\nu_\tau|\,d\mathcal H^{n-1}.
\end{equation}
The last equality uses~\eqref{eq:mf-W}. Bounding $\tau^{-n}\le\rstar^{-n}$
and dividing by $|h|$, then letting $h\to 0$, gives
\[
   \limsup_{h\to 0}\frac{|F_{\rho+h}\triangle F_\rho|}{|h|}
   \le\rstar^{-n}\int_{\partial^*E_\rho}
      |\Vrho-H_{z_\rho,\rho}-a\cdot\nurho|\,d\mathcal H^{n-1}.
\]
Now apply the translation derivative formula (Lemma~\ref{lem:mf-trans-deriv}
below) to absorb $a\cdot\nu_\rho$ into a translation $F_{\rho+h}\mapsto
F_{\rho+h}-ha$: the symmetric difference between $F_{\rho+h}-ha$ and $F_\rho$
also satisfies the same bound. Taking the infimum over $a$ and using
$\dX([F_{\rho+h}],[F_\rho])\le|F_{\rho+h}\triangle(F_\rho+ha)|$ yields
\eqref{eq:mf-T-smooth}.
\end{proof}

\begin{lemma}[Translation derivative; \protect{\cite[Thm.~13.8]{Maggi2012},
\cite[Thm.~3.84]{AFP2000}}]\label{lem:mf-trans-deriv}
For every bounded set $E$ of finite perimeter and every $a\in\R^n$,
\begin{equation}\label{eq:mf-trans-deriv}
   \lim_{\eps\to 0}\frac{|E\triangle(E+\eps a)|}{\eps}
   =\int_{\partial^*E}|a\cdot\nu_E|\,d\mathcal H^{n-1}.
\end{equation}
\end{lemma}

In view of~\eqref{eq:mf-trans-deriv}, the inf over $a$ in~\eqref{eq:mf-T} is
\emph{exactly} the metric-quotient cost: not slack, but the precise
infinitesimal cost of factoring out translations. This is the role of the
parameter~$a$ in the theorem statement.

\subsubsection{BV/coarea smooth approximation}
\label{sec:metric-T-approx}

We now drop the smooth-flow hypotheses. The standing hypotheses
of~\S\ref{sec:metric-setup} are kept; the only regularity assumed is
$\Om$ open, bounded, $|\Om|=\omega_n$, and the consequent
$u\in W^{2,p}_{\rm loc}(\Om)$, $p>n$.

\begin{lemma}[Coarea-good smooth recovery]\label{lem:mf-approx}
There exist $u_k\in C^\infty(\R^n)$ with $u_k\to u$ in $W^{1,2}(\Om)$ and a
subsequence (not relabelled) such that, setting
$E_t^k:=\{u_k>t\}$ and reparametrising $t_k(\rho)$ by
$|E^k_{t_k(\rho)}|=\omega_n\rho^n$, the following hold for a.e.\
$\rho\in[\rstar,1]$:
\begin{enumerate}
\item[(a)] $t_k(\rho)\to t(\rho)$;
\item[(b)] $E^k_\rho\to E_\rho$ in $L^1(\R^n)$;
\item[(c)] $P(E_\rho)\le\liminf_k P(E^k_\rho)$, and in fact $P(E^k_\rho)\to
   P(E_\rho)$;
\item[(d)] the vector-valued perimeter measures
   $\nu^k_\rho\,\mathcal H^{n-1}\restriction\partial^*E^k_\rho$ converge
   weakly to $\nurho\,\mathcal H^{n-1}\restriction\partial^*E_\rho$;
\item[(e)] $\int_{\partial^*E^k_\rho}|V^k_\rho|\,d\mathcal H^{n-1}
   \to\int_{\partial^*E_\rho}|\Vrho|\,d\mathcal H^{n-1}$.
\end{enumerate}
\end{lemma}

\begin{proof}[Sketch with precise inputs]
Take $u_k=u*\eta_k$ for a standard mollifier $\eta_k$, extended by zero
outside $\Om$. Then $u_k\to u$ in $W^{1,p}_{\rm loc}$ for every
$p<\infty$ (\cite[Thm.~3.40]{AFP2000}). Item~(a) and (b) follow from the BV
coarea theorem applied to $u_k\to u$ together with Helly--Lebesgue
differentiation of the (monotone) maps $t\mapsto|E^k_t|$
(\cite[Thm.~13.1, Lem.~13.2]{Maggi2012}). Item~(c) is the BV coarea identity
\[
   \int|\nabla u_k|\,dx=\int_0^\infty P(E^k_t)\,dt\to\int|\nabla u|\,dx
      =\int_0^\infty P(E_t)\,dt
\]
combined with perimeter LSC (\cite[Prop.~12.15]{Maggi2012}, item~A6) and
Fatou: the LSC gives $\liminf P(E^k_t)\ge P(E_t)$ a.e., while the
$L^1$-integrals match, so equality holds for a.e.~$t$ (and hence a.e.~$\rho$).
Item~(d) is weak convergence of vector perimeter measures, which follows
from~(c) combined with $L^1$-convergence of $\mathbf 1_{E^k_\rho}$ and
\cite[Cor.~12.21]{Maggi2012}. Item~(e) uses
$\int_{\partial^*E^k_\rho}|V^k_\rho|=|t_k'(\rho)|\cdot a^k(t_k(\rho))$ where
$a^k(s)=\int_{\partial^*E^k_s}|\nabla u_k|^{-1}\,d\mathcal H^{n-1}$
(coarea, item~A2), combined with the Helly--Lebesgue argument again.
\end{proof}

\begin{remark}[No graph regularity used]\label{rem:no-graph}
Lemma~\ref{lem:mf-approx} makes no claim about $\partial E_\rho$ being a
graph or being globally Lipschitz; it asserts only $L^1$-convergence of the
sub-level sets, BV-coarea convergence of perimeters, and weak convergence of
vector perimeter measures. The recovery family $u_k$ is smooth on $\R^n$,
but the limit set $E_\rho$ may have any reduced-boundary structure.
\end{remark}

\subsubsection{Closure: Reshetnyak and AGS lower semicontinuity}
\label{sec:metric-T-pass}

\begin{lemma}[Reshetnyak l.s.c.\ for the velocity-defect integrand;
items~A6, A8]\label{lem:mf-reshetnyak}
Fix $z,a\in\R^n$ and $\rho\in[\rstar,1]$. The integrand
\[
   \Phi(x,s,\nu):=\bigl|s-(x-z)\cdot\nu/\rho-a\cdot\nu\bigr|,
   \qquad (x,s,\nu)\in\R^n\times\R\times\R^n,
\]
is jointly continuous and convex of linear growth in $(s,\nu)$. For sequences
$E^k_\rho\to E_\rho$ in $L^1$ with the perimeter and vector-perimeter
convergence of Lemma~\ref{lem:mf-approx}(c)--(e), and writing
$V^k_\rho=-t_k'(\rho)/|\nabla u_k|$,
\begin{equation}\label{eq:mf-reshetnyak}
   \liminf_{k\to\infty}\int_{\partial^*E^k_\rho}
      \Phi(x,V^k_\rho,\nu^k_\rho)\,d\mathcal H^{n-1}
   \ge\int_{\partial^*E_\rho}\Phi(x,\Vrho,\nurho)\,d\mathcal H^{n-1}.
\end{equation}
\end{lemma}

\begin{proof}
This is Reshetnyak lower semicontinuity for convex linear-growth integrands
with continuous $x$-dependence, applied to the BV functions $\mathbf 1_{E^k_
\rho}\to\mathbf 1_{E_\rho}$ on the bounded ball $B_R$: see
\cite[Thm.~2.38]{AFP2000} (continuous-coefficient case) and
\cite[Thm.~2.39]{AFP2000} (weak strict BV convergence). The fact that
$V^k_\rho\,\mathcal H^{n-1}\restriction\partial^*E^k_\rho$ converges
weakly to $\Vrho\,\mathcal H^{n-1}\restriction\partial^*E_\rho$ follows
from Lemma~\ref{lem:mf-approx}(d)--(e): the $L^1$ mass converges, and the
support of the limit lies on $\partial^*E_\rho$ by Lemma~\ref{lem:mf-approx}(d).
\end{proof}

\begin{lemma}[AGS l.s.c.\ of the metric derivative;
\protect{\cite[Thm.~1.1.2]{AGS2008}}]\label{lem:mf-AGS}
Let $\gamma_k:[\rstar,1]\to\X$ be absolutely continuous curves with
$\gamma_k(\rho)\to\gamma(\rho)$ in $\X$ for every $\rho\in[\rstar,1]$,
and suppose $|\dot\gamma_k|_\X(\rho)\le m_k(\rho)$ a.e.\ with $m_k\to m$
weakly in $L^1([\rstar,1])$. Then $\gamma$ is absolutely continuous and
$|\dot\gamma|_\X(\rho)\le m(\rho)$ for a.e.~$\rho$.
\end{lemma}

\begin{proof}[Proof of Theorem~\ref{thm:mf-T}]
Let $u_k$ be the recovery sequence of Lemma~\ref{lem:mf-approx}, and define
$F^k_\rho:=\rho^{-1}E^k_\rho$, $[F^k_\rho]\in\X$. By Lemma~\ref{lem:mf-approx}(b),
$[F^k_\rho]\to[F_\rho]$ in $\X$ for a.e.~$\rho$. For each $k$, the smooth
foliation $\rho\mapsto E^k_\rho$ has $|\nabla u_k|>0$ on the relevant strip
(after possibly excluding a null set of $\rho$, by smoothness of $u_k$ and
Sard's theorem). Applying Proposition~\ref{prop:mf-T-smooth} to
$\{E^k_\rho\}$,
\begin{equation}\label{eq:mf-Tk}
   |\dot F^k_\rho|_\X\le m_k(\rho):=
   \rstar^{-n}\inf_{a\in\R^n}\int_{\partial^*E^k_\rho}
      \bigl|V^k_\rho-H_{z_\rho,\rho}-a\cdot\nu^k_\rho\bigr|
      \,d\mathcal H^{n-1}.
\end{equation}
By the level-set deficit identity (Sec.~\ref{sec:level-set-identity},
item~E1.4) applied to $u_k$, $m_k$ is uniformly bounded in $L^1([\rstar,1])$
by a constant depending only on $n,R,\rstar$ and the BV-coarea norm of $u_k$,
hence by $\|u_k\|_{W^{1,2}}$. By Dunford--Pettis, after extracting a
subsequence, $m_k\rightharpoonup m$ weakly in $L^1$. By
Lemma~\ref{lem:mf-reshetnyak} applied to the optimal $a$ in the infimum
(measurable selection \cite[Thm.~8.2.11]{AubinFrankowska1990}, item~A7),
\begin{equation}\label{eq:mf-mlim}
   m(\rho)\ge\rstar^{-n}\inf_{a\in\R^n}\int_{\partial^*E_\rho}
      \bigl|\Vrho-H_{z_\rho,\rho}-a\cdot\nurho\bigr|\,d\mathcal H^{n-1}
   \qquad\text{a.e.~}\rho.
\end{equation}
By Lemma~\ref{lem:mf-AGS}, $\rho\mapsto[F_\rho]$ is absolutely continuous
(consistent with Theorem~\ref{thm:mf-AC}) and
\begin{equation}
   |\dot F_\rho|_\X\le m(\rho)
   \le\rstar^{-n}\inf_{a\in\R^n}\int_{\partial^*E_\rho}
      \bigl|\Vrho-H_{z_\rho,\rho}-a\cdot\nurho\bigr|\,d\mathcal H^{n-1}.
\end{equation}
This is~\eqref{eq:mf-T} with $C(n,\rstar)=\rstar^{-n}$.
\end{proof}

\subsubsection{Per-$\rho$ squared form (conditional) and the integrated route
actually used}
\label{sec:metric-T-summary}

Combining Theorem~\ref{thm:mf-T} with the
$D_H$- and $D_I$-controls of Sec.~\ref{sec:gmt-velocity}
(items~A1--A8, B6, C2--C4) gives an estimate of the form
\begin{equation}\label{eq:mf-cond-71}
   |\dot F_\rho|_\X^2\le C(n,R,\rstar)\bigl(D_H(t(\rho))+D_I(t(\rho))\bigr),
\end{equation}
\emph{conditional} on a per-$\rho$ radial trace bound on the homothetic
velocity defect (the bound recorded in Sec.~\ref{sec:gmt-velocity}, eq.~(2.4)).
Per the no-graph hypothesis of this manuscript, the per-$\rho$ radial trace
is established only in its integrated good/bad form by Agents~11--12;
\eqref{eq:mf-cond-71} is therefore \emph{not} used pointwise in the
downstream chain. The downstream argument uses only the integrated form
\begin{equation}\label{eq:mf-int-cond}
   \int_{\rstar}^{1}|\dot F_\rho|_\X^2\,\mu(d\rho)
   \le C(n,R,\rstar)\,\deltaT(\Om),
\end{equation}
which is proved in Sec.~\ref{sec:metric-trace} (Agent~12) by combining
Theorem~\ref{thm:mf-T}, the Fusco--Julin centre selection of
Sec.~\ref{sec:fj-radial} (Agent~11), and the level-set deficit identity of
Sec.~\ref{sec:level-set-identity} (Agent~9).

\begin{theorem}[Final form of \textsc{MetricFramework}]\label{thm:mf-main}
Let $n\ge 2$, $0<\rstar<1$, and $R\ge 2$. Let $\Om\subset\R^n$ be open and
bounded with $|\Om|=\omega_n$ and $\Om\subset B_R$. Let $u$ be the
variational torsion potential, let $\{E_\rho\}_{\rho\in[\rstar,1]}$ be the
volume-radius foliation~\eqref{eq:mf-rho}, and set $F_\rho:=\rho^{-1}E_\rho$.
Then:
\begin{enumerate}
\item[(i)] $(\X,\dX)$ is a complete metric space, separable and
sequentially compact on classes of sets in a fixed ball
(Proposition~\ref{prop:mf-Polish}).
\item[(ii)] The curve $\rho\mapsto[F_\rho]\in\X$ is absolutely continuous on
$[\rstar,1]$, with metric-derivative bound~\eqref{eq:mf-m-def}
(Theorem~\ref{thm:mf-AC}).
\item[(iii)] The normal velocity $\Vrho=-t'(\rho)/|\nabla u|$ is
$\mathcal H^{n-1}$-a.e.\ finite and strictly positive on $\partial^*E_\rho$
for a.e.~$\rho$ (Proposition~\ref{prop:mf-sard}).
\item[(iv)] For any Borel-measurable centre field $\rho\mapsto z_\rho$ with
$\sup_\rho|z_\rho|\le R'<\infty$,
\begin{equation}\label{eq:mf-final}
   |\dot F_\rho|_\X\le\rstar^{-n}\inf_{a\in\R^n}
   \int_{\partial^*E_\rho}
      \bigl|\Vrho-H_{z_\rho,\rho}-a\cdot\nurho\bigr|
      \,d\mathcal H^{n-1}
\end{equation}
for a.e.~$\rho\in[\rstar,1]$ (Theorem~\ref{thm:mf-T}).
\end{enumerate}
\end{theorem}

The constants in the section are:
$C_{\rm AC}(n,R,\rstar)$ in~\eqref{eq:mf-m-def}, with explicit
dependence $C_{\rm AC}=\rstar^{-n}n\omega_n+n\rstar^{-n-1}R/\rstar$;
$C(n,\rstar)=\rstar^{-n}$ in~(T)~\eqref{eq:mf-T}. The cited downstream
constant $C(n,R,\rstar)$ in~\eqref{eq:mf-int-cond} is the product
$\rstar^{-2n}\cdot C_{\rm DH+DI}$ where $C_{\rm DH+DI}=C_{\rm DH+DI}(n,R,
\rstar)$ is the constant of Sec.~\ref{sec:gmt-velocity}.

\medskip\noindent
\emph{Source map for this subsection.} Definitions and metric facts:
[WIP-MF] \S2 and Proposition~\ref{prop:mf-Polish} (proved here);
absolute continuity Theorem~\ref{thm:mf-AC}: proved here, using
Lemma~\ref{lem:mf-attain} ([WIP-MF] Lem.~2.2) and the coarea flux
\cite[Thm.~13.1]{Maggi2012} (item~A2). Sobolev--Sard for
$\Vrho$ (Proposition~\ref{prop:mf-sard}): cited from
\cite[App.~A, Thm.~A.1]{FigalliMaggi2011} (item~A3, source map E3.6).
Smooth-flow first variation (Lemma~\ref{lem:mf-FV-smooth},
Proposition~\ref{prop:mf-T-smooth}): cited from
\cite[\S~17.3]{Maggi2012} and \cite[Thm.~3.85]{AFP2000}; identical to
[A3] \S3. Translation derivative (Lemma~\ref{lem:mf-trans-deriv}):
\cite[Thm.~13.8]{Maggi2012}, \cite[Thm.~3.84]{AFP2000} (item~E2.7).
BV/coarea recovery (Lemma~\ref{lem:mf-approx}): standard mollification
recovery; precise inputs cited from \cite[Thm.~13.1, Prop.~12.15,
Cor.~12.21]{Maggi2012} and \cite[Thm.~3.40]{AFP2000} (items~A2, A6).
Reshetnyak l.s.c.\ (Lemma~\ref{lem:mf-reshetnyak}):
\cite[Thm.~2.38--2.39]{AFP2000} (item~A8). AGS l.s.c.\
(Lemma~\ref{lem:mf-AGS}): \cite[Thm.~1.1.2]{AGS2008} (item~F15).
Measurable selection for the optimal~$a$:
\cite[Thm.~8.2.11]{AubinFrankowska1990} (item~A7).
Theorem~\ref{thm:mf-T} (E2.6): proved here.

\label{subsec:metric-framework}%
\label{sec:plan2-metric}%
\label{ssec:plan2-metric}%
% UNIFIED 2026-05-13
% SOURCE MAP for §4.3 (Plan 2, Agent 11):
%   - Statement of scaled Fusco--Julin (Thm 4.3.1): CITE(FJ2014, Thm 1.1) + DRAFT
%     scaling; reads Manuscript/00_notation_register.md §1, §7, §10;
%     source_map.md C2, C3, E3.1.
%   - Borel centre selection (Lem 4.3.2): DRAFT, with measurable-selection
%     prelim Agent 2 (source_map.md A7) and bulk-centroid Borel regularity
%     from Plan 2/wave3-F-h1-center-bound.md §2.2 (LOCAL).
%   - Fraenkel bound (Lem 4.3.3): DRAFT, consumes FJ scaled form (C3) and
%     FMP/AFM/CL centre proximity (LOCAL, wave3-F-h1-center-bound.md §1
%     + Wave 2 Agent A Lemma 2.2).
%   - Normal-oscillation bound (Lem 4.3.4): DRAFT, divergence identity
%     Agent 4 §C4; FJ2014 Thm 1.1.
%   - Oriented radial trace lemma (Thm 4.3.5, the critical load-bearing
%     lemma E3.5): DRAFT, builds on Plan 2/WIP/WIP_WeightedMetricTrace.tex
%     Lem.\ 3.6 and the wave3-J-route-delta-repair.md §§3--7 slicing
%     decomposition; uses Agent 2 prelims A4 (divergence theorem for
%     finite-perimeter sets, CITE Maggi2012 Thm 16.3) and A5 (Lipschitz
%     truncation of x/|x|, DRAFT in §2.1).
%   - Vector field X(x)=||x-z|/rho-1|(x-z)/|x-z|: named precisely as in
%     the brief and in WIP_WeightedMetricTrace.tex (LOCAL).
%   - All constants C(n, R, rho_*) tracked per notation register §14.

\subsection{\S4.3.\ Fusco--Julin centre and the oriented radial trace lemma}
\label{sec:plan2-fj-radial}

This subsection is the centre-package and load-bearing radial-trace
estimate of Plan~2. Its conclusion --- Theorem~\ref{thm:radial-trace} ---
together with the metric first-variation theorem of \S4.2 controls the
metric derivative $|\dot F_\rho|_{\X}$ in the weighted action integral.

Throughout this subsection we fix $n\ge 2$, $R\ge 2$, $\rstar\in(0,1)$,
and $\kappa>0$. We work with an open set $\Om\subset B_R\subset\R^n$,
$|\Om|=\omega_n$, with variational torsion function $u=u_\Om$, volume-radius
parametrisation $E_\rho=\{u>t(\rho)\}$, $|E_\rho|=\omega_n\rho^n$ for
$\rho\in[\rstar,1]$, and $\mu(d\rho):=(-t'(\rho))\,d\rho$. We assume
\begin{equation}\label{eq:fj-smallness}
   \delta_T(\Om)\le\delta_0(n,R,\rstar),
\end{equation}
where $\delta_0$ is the smallness threshold fixed in
Theorem~\ref{thm:weighted-trace} (the constant matching the
Cicalese--Leonardi smallness threshold and the centroid--Fraenkel offset,
cf.\ \cite[Prop.~2.3]{CL2012} and \cite[Thm.~1.1]{FJ2014}).

\subsubsection{4.3.1.\ Scaled Fusco--Julin and centre selection}
\label{ssec:fj-scaled}

Recall the strong quantitative isoperimetric inequality of
Fusco--Julin~\cite[Thm.\ 1.1]{FJ2014}: there is a dimensional constant
$C_{\mathrm{FJ}}(n)>0$ such that for every set of finite perimeter
$E\subset\R^n$ with $|E|=|B_1|$ there exists a centre $z=z(E)\in\R^n$ with
\begin{equation}\label{eq:FJ-original}
   \Bigl(\frac{|E\Delta B_1(z)|}{|B_1|}\Bigr)^2
   +\beta(E,z)^2
   \le C_{\mathrm{FJ}}(n)\,\bigl(P(E)-P(B_1)\bigr),
\end{equation}
where the FJ oscillation index is
\begin{equation}\label{eq:FJ-osc-def}
   \beta(E,z)^2
   :=\inf_{|E'|=|E|}\;
   \int_{\partial^*E}\bigl|\nu_E(x)-(\nu_{E'})_{\mathrm{rad}}(x;z)\bigr|^2\,d\mathcal H^{n-1}(x),
\end{equation}
which by the standard divergence identity (item~C4 of the
preliminaries; see~\eqref{eq:beta-div} below) is comparable to
$\int_{\partial^*E}|\nu_E-(x-z)/|x-z||^2\,d\mathcal H^{n-1}$.

\begin{theorem}[Scaled Fusco--Julin on $\{|E|=\omega_n\rho^n\}$]
\label{thm:plan2-FJ-scaled}\label{prop:fj-scaled}
For every $\rho\in[\rstar,1]$ and every $E\subset\R^n$ of finite perimeter
with $|E|=\omega_n\rho^n$ there exists $z=z(E,\rho)\in\R^n$ such that
\begin{equation}\label{eq:plan2-FJ-scaled}
   |E\Delta B_\rho(z)|^2
   +\rho^{n+1}\int_{\partial^*E}\!\!\Bigl|\nu_E(x)-\frac{x-z}{|x-z|}\Bigr|^2 d\mathcal H^{n-1}(x)
   \le C_1(n,\rstar)\,\rho^{n+1}\bigl(P(E)-P(B_\rho)\bigr),
\end{equation}
where $C_1(n,\rstar)=4\,C_{\mathrm{FJ}}(n)\cdot\max(1,\rstar^{-(n+1)})$. The
constant depends only on $n$ and $\rstar$.
\end{theorem}

\begin{proof}
Set $\widetilde E:=\rho^{-1}E$, so $|\widetilde E|=\omega_n=|B_1|$ and
$\widetilde E$ has finite perimeter with
$P(\widetilde E)=\rho^{-(n-1)}P(E)$. Apply~\eqref{eq:FJ-original} to
$\widetilde E$: there exists $\widetilde z$ with
\begin{equation}\label{eq:FJ-applied}
   \Bigl(\frac{|\widetilde E\Delta B_1(\widetilde z)|}{\omega_n}\Bigr)^2
   +\beta(\widetilde E,\widetilde z)^2
   \le C_{\mathrm{FJ}}(n)\bigl(P(\widetilde E)-P(B_1)\bigr).
\end{equation}
Set $z:=\rho\widetilde z$. Translating $\widetilde E$ to $E$ scales
volumes by $\rho^n$ and perimeters by $\rho^{n-1}$, and using the
divergence identity~\eqref{eq:beta-div} below to rewrite
$\beta(\widetilde E,\widetilde z)^2$ in terms of the
normal-oscillation integral, the scaled bounds
$|E\Delta B_\rho(z)|=\rho^n|\widetilde E\Delta B_1(\widetilde z)|$
and $\int_{\partial^*E}|\nu_E-(x-z)/|x-z||^2\,d\mathcal H^{n-1}
=\rho^{n-1}\int_{\partial^*\widetilde E}|\nu_{\widetilde E}-(y-\widetilde z)/|y-\widetilde z||^2\,d\mathcal H^{n-1}$
yield~\eqref{eq:plan2-FJ-scaled} after multiplying through by $\rho^{2n}$
and using $\rstar\le\rho\le1$ to absorb the $\rho$-powers into the
constant $C_1$.
\end{proof}

\begin{remark}[Comparison constants]\label{rmk:FJ-equiv}
The divergence identity (item~C4 of the preliminaries)
\begin{equation}\label{eq:beta-div}
   2\,\beta(E,z)^2
   = \int_{\partial^*E}\Bigl|\nu_E(x)-\frac{x-z}{|x-z|}\Bigr|^2 d\mathcal H^{n-1}(x)
   = 2(P(E)-P(B_\rho))+2\Pi(E,z),
\end{equation}
holds whenever $|E|=|B_\rho|$, $z\in\R^n$ with the radial field
$x\mapsto(x-z)/|x-z|$ locally Borel and the truncation argument
of~\S2.1 (item~A5) applicable. The identity is proved by testing
$\mathrm{div}((x-z)/|x-z|)=(n-1)|x-z|^{-1}$ on $E$ via the finite-perimeter
divergence theorem and truncating away from~$z$;
$\Pi(E,z):=(n-1)\int (\mathbf1_{B_\rho(z)}-\mathbf1_E)|x-z|^{-1}\,dx$ is the
signed radial moment of the notation register \S13. We use only the
inequality
$\int|\nu_E-(x-z)/|x-z||^2\,d\mathcal H^{n-1}\le 2C_1(n,\rstar)(P(E)-P(B_\rho))/\rho^{n+1}$
that follows from~\eqref{eq:plan2-FJ-scaled}.
\end{remark}

We now choose a single Borel centre field along the torsion foliation.
The canonical choice (see~\S\ref{sec:notation}, and the discussion of the
Fusco--Julin centre in~\cite[Thm.~1.1]{FJ2014})
is the bulk centroid
\begin{equation}\label{eq:centroid-def}
   \zbar
   = \bar z_\rho
   :=\frac{1}{|E_\rho|}\int_{E_\rho}x\,dx
   \in\R^n,\qquad\rho\in[\rstar,1].
\end{equation}

\begin{lemma}[Borel measurability and CL containment of the centroid]
\label{lem:centroid-borel}
Under~\eqref{eq:fj-smallness} the map $\rho\mapsto\zbar$ is Lipschitz on
$[\rstar,1]$ with constant $L_{\mathrm{cent}}(n,R,\rstar)\le 2nR\rstar^{-n}$;
in particular it is Borel. Moreover for every $\rho\in[\rstar,1]$,
\begin{equation}\label{eq:zbar-CL}
   B(\zbar,\rho/4)\subset E_\rho\subset B(\zbar,3\rho),
\end{equation}
so $\zbar\in E_\rho$ and the radial field $e_r(x):=(x-\zbar)/|x-\zbar|$ is
$\mathcal H^{n-1}$-a.e.\ well-defined and Borel on $\partial^*E_\rho$, with
$\rstar/4\le|x-\zbar|\le 3R$ on $\partial^*E_\rho$ for a.e.\ $\rho$.
\end{lemma}

\begin{proof}
For $\rstar\le\rho\le\rho'\le1$, the nesting of $\{E_\rho\}$ gives
$|E_\rho\Delta E_{\rho'}|=|E_{\rho'}|-|E_\rho|\le n\omega_n|\rho-\rho'|$.
Hence
$|\int_{E_\rho}x_i\,dx-\int_{E_{\rho'}}x_i\,dx|\le R\cdot n\omega_n|\rho-\rho'|$;
combined with $|E_\rho|\ge\omega_n\rstar^n$ this yields the Lipschitz
constant $L_{\mathrm{cent}}=2nR/\rstar^n$, which is Borel (in fact
continuous). The containment~\eqref{eq:zbar-CL} follows from the
Cicalese--Leonardi inner-ball containment (\cite[Thm.\ 1.1]{CL2012})
applied at the Fraenkel centre
$z^{\mathrm{Fr}}_\rho$, together with the centroid--Fraenkel proximity
$|\zbar-z^{\mathrm{Fr}}_\rho|\le C(n,\rstar)\,\Fr(E_\rho)\le
C'\sqrt{\delta_T(\Om)}\le\rho/4$, valid for $\delta_0$ small
(\cite[Prop.\ 2.3]{CL2012},
\cite[Lem.\ 2.6]{AFM2013}, \cite[Thm.~1.1]{FJ2014}).
Borel measurability of the normal field $\nu_{E_\rho}$ on $\partial^*E_\rho$
(\cite[Thm.\ 15.9]{Maggi2012}) then yields the Borel measurability of
$e_r$ on $\partial^*E_\rho$ for $\zbar\in E_\rho$.
\end{proof}

\begin{remark}[Auxiliary FJ-oscillation centre and Borel selection]
\label{rmk:centre-selection}
The centre $z(E_\rho,\rho)$ produced by Theorem~\ref{thm:plan2-FJ-scaled} is
\emph{a} valid almost-minimiser; Lemma~\ref{lem:centroid-borel}'s centroid
$\zbar$ also satisfies~\eqref{eq:plan2-FJ-scaled} up to enlarging $C_1$ by the
proximity offset $|\zbar-z^{\mathrm{Fr}}_\rho|\le C(n,\rstar)\sqrt{\delta_T}$,
absorbed by the smallness threshold $\delta_0$. The Borel selection
of \emph{any} valid centre on a Borel-measurable family of finite-perimeter
sets follows from the measurable-selection statement A7 of the
preliminaries (Federer~\cite[\S2.13]{Federer1969}, in the form recorded
in~\cite[\S5.3]{AFP2000}). Throughout the rest of this subsection the
centre is fixed to be $z_\rho:=\zbar$.
\end{remark}

\subsubsection{4.3.2.\ Fraenkel and normal-oscillation bounds at the centroid}
\label{ssec:fj-bounds}

Define the (good) isoperimetric set
\begin{equation}\label{eq:GI-def}
   G_I:=\{\rho\in[\rstar,1]:D_I(t(\rho))\le\eta_I\},\qquad
   \eta_I:=\eta_I(n,R,\rstar)\in(0,1],
\end{equation}
with $\eta_I$ small enough that on $G_I$, $P(E_\rho)\le 2P(B_\rho)$
(possible since $D_I=(P(E)-P(B_\rho))(P(E)+P(B_\rho))$ and
$P(B_\rho)\ge n\omega_n\rstar^{n-1}$). We collect
\eqref{eq:plan2-FJ-scaled} at $z_\rho=\zbar$.

\begin{lemma}[Fraenkel bound at the centroid]\label{lem:fraenkel-centroid}
There exists $C_2=C_2(n,R,\rstar)>0$ such that for every $\rho\in G_I$,
\begin{equation}\label{eq:fraenkel-centroid}
   |E_\rho\Delta B_\rho(\zbar)|^2
   \le C_2(n,R,\rstar)\,D_I(t(\rho)).
\end{equation}
\end{lemma}

\begin{proof}
From~\eqref{eq:plan2-FJ-scaled} applied with $z=\zbar$ (valid by
Remark~\ref{rmk:centre-selection}, after absorbing the
$|\zbar-z^{\mathrm{Fr}}_\rho|^2$ correction into the constant via
$|B_\rho(a)\Delta B_\rho(b)|\le C(n)\rho^{n-1}|a-b|$ and the FMP estimate
$|E_\rho\Delta B_\rho(z^{\mathrm{Fr}}_\rho)|^2\le c_1|B_\rho|^2(P(E_\rho)-P(B_\rho))$),
\begin{equation}\label{eq:fraenkel-centroid-pre}
   |E_\rho\Delta B_\rho(\zbar)|^2
   \le 2C_1(n,\rstar)\rho^{n+1}\bigl(P(E_\rho)-P(B_\rho)\bigr).
\end{equation}
Use the algebraic identity
$D_I(t(\rho))=(P(E_\rho)-P(B_\rho))(P(E_\rho)+P(B_\rho))\ge n\omega_n\rstar^{n-1}\cdot(P(E_\rho)-P(B_\rho))$
to convert $(P(E_\rho)-P(B_\rho))$ into $D_I/(n\omega_n\rstar^{n-1})$;
this and $\rho\le1$ yield~\eqref{eq:fraenkel-centroid} with
$C_2=2C_1(n,\rstar)/(n\omega_n\rstar^{n-1})$, polynomial in $n,R,1/\rstar$.
\end{proof}

\begin{lemma}[Normal-oscillation bound at the centroid]
\label{lem:normal-osc-centroid}
There exists $C_3=C_3(n,R,\rstar)>0$ such that for every $\rho\in G_I$,
\begin{equation}\label{eq:normal-osc-centroid}
   \int_{\partial^*E_\rho}\!\!\Bigl|\nu_\rho(x)-\frac{x-\zbar}{|x-\zbar|}\Bigr|^2\,
      d\mathcal H^{n-1}(x)
   \le C_3(n,R,\rstar)\,D_I(t(\rho)).
\end{equation}
\end{lemma}

\begin{proof}
Combining~\eqref{eq:plan2-FJ-scaled} at $z=\zbar$ (with the same absorption as
in the previous proof) and the lower bound $\rho^{n+1}\ge\rstar^{n+1}$
gives
$\int|\nu_\rho-(x-\zbar)/|x-\zbar||^2\,d\mathcal H^{n-1}\le
C_1(n,\rstar)(P(E_\rho)-P(B_\rho))/\rstar^{n+1}$; converting $P-P_\rho$ to
$D_I$ as in the previous proof yields~\eqref{eq:normal-osc-centroid}
with $C_3=C_1(n,\rstar)/(\rstar^{n+1}\cdot n\omega_n\rstar^{n-1})$.
\end{proof}

\begin{remark}[{$\rho\in[\rstar,1]\setminus G_I$}]
Outside $G_I$, $D_I(t(\rho))>\eta_I$, and~\eqref{eq:fraenkel-centroid},
\eqref{eq:normal-osc-centroid} hold trivially after enlarging
$C_2,C_3$ by a multiple of $\eta_I^{-1}$, since both left-hand sides are
bounded by $|B_R|^2$ and $4P(B_R)\le 4n\omega_n R^{n-1}$ respectively.
In the integrated argument of Theorem~\ref{thm:weighted-trace} the
$\mu$-measure of $\{\rho:\rho\notin G_I\}$ is controlled by Markov from
the deficit identity, so the trivial bound is permissible.
\end{remark}

\subsubsection{4.3.3.\ The oriented radial trace lemma}
\label{ssec:radial-trace}

This is the load-bearing lemma of Plan~2's GMT route. We prove the
pointwise inequality, with the named vector field
\begin{equation}\label{eq:X-def}
   X(x)
   :=\Bigl|\frac{|x-\zbar|}{\rho}-1\Bigr|\cdot\frac{x-\zbar}{|x-\zbar|},
   \qquad x\in\R^n\setminus\{\zbar\},
\end{equation}
of the brief, via the finite-perimeter divergence theorem (item~A4 of
the preliminaries: \cite[Thm.\ 16.3]{Maggi2012}) and the truncation
argument of item~A5. We do \emph{not} use ray-slicing in the proof of
this lemma; a slicing-based variant is sketched in
Remark~\ref{rmk:slicing-route} below.

\begin{theorem}[Oriented $L^1$ radial trace at the centroid]
\label{thm:radial-trace}\label{thm:fj-radial-divergence}\label{lem:radial-trace-master}
There is $C_4=C_4(n,R,\rstar)>0$ such that for every
$\rho\in[\rstar,1]$ satisfying~\eqref{eq:fj-smallness}, with the centroid
$z_\rho=\zbar$ of~\eqref{eq:centroid-def},
\begin{equation}\label{eq:radial-trace}
   T_\rho
   :=\int_{\partial^*E_\rho}\!\!\Bigl|\frac{|x-\zbar|}{\rho}-1\Bigr|\,
       d\mathcal H^{n-1}(x)
   \le
   C_4\,\Bigl(\,|E_\rho\Delta B_\rho(\zbar)|
   +\int_{\partial^*E_\rho}\!\!\Bigl|\nu_\rho-\frac{x-\zbar}{|x-\zbar|}\Bigr|^2\,
       d\mathcal H^{n-1}\Bigr).
\end{equation}
The constant has the explicit form $C_4(n,R,\rstar)=\max\{n/\rstar,3R/(2\rstar)\}$
(see~\eqref{eq:T-trace-explicit} below).
\end{theorem}

\begin{proof}
\emph{Step 1: setup.} By Lemma~\ref{lem:centroid-borel},
$\zbar\in B(\zbar,\rho/4)\subset E_\rho\subset B(\zbar,3\rho)\subset B_R$,
so $|x-\zbar|\in[\rstar/4,3R]$ for $\mathcal H^{n-1}$-a.e.
$x\in\partial^*E_\rho$. Translate so that $\zbar=0$; the radial unit
vector is $e_r(x)=x/|x|$ on $\R^n\setminus\{0\}$ and the
weight is $w(r):=|r/\rho-1|$, a Lipschitz function of $r\in[0,\infty)$
with $w(\rho)=0$, $|w(r)|\le \max(1,(3R)/\rho-1)\le 3R/\rstar$ on
$\partial^*E_\rho$, and the
vector field $X(x)=w(|x|)e_r(x)$ of~\eqref{eq:X-def}. Note $X$ is locally
Lipschitz on $\R^n\setminus\{0\}$ and bounded on $B_R$ (with
$\|X\|_\infty\le 3R/\rstar$), but is discontinuous at the origin.

\emph{Step 2: pointwise inequality.}
On $\partial^*E_\rho$, $|e_r|=|\nu_\rho|=1$, so writing
$1=e_r\cdot\nu_\rho+(1-e_r\cdot\nu_\rho)$ and using $w\ge0$,
\begin{equation}\label{eq:T-split}
   T_\rho
   =\int_{\partial^*E_\rho}w(|x|)\,d\mathcal H^{n-1}
   =\int_{\partial^*E_\rho}X\cdot\nu_\rho\,d\mathcal H^{n-1}
    +\int_{\partial^*E_\rho}w(|x|)\bigl(1-e_r\cdot\nu_\rho\bigr)\,d\mathcal H^{n-1}.
\end{equation}
The second integrand is non-negative because $e_r\cdot\nu_\rho\le|e_r||\nu_\rho|=1$,
and the elementary identity
$1-e_r\cdot\nu_\rho=\tfrac12|e_r-\nu_\rho|^2$ together with
$w\le3R/\rstar$ gives
\begin{equation}\label{eq:T-tan-bound}
   \int_{\partial^*E_\rho}w(|x|)(1-e_r\cdot\nu_\rho)\,d\mathcal H^{n-1}
   \le\frac{3R}{2\rstar}
   \int_{\partial^*E_\rho}|\nu_\rho-e_r|^2\,d\mathcal H^{n-1}.
\end{equation}

\emph{Step 3: truncated divergence theorem for the radial term.}
We bound $\int_{\partial^*E_\rho}X\cdot\nu_\rho\,d\mathcal H^{n-1}$ by the
finite-perimeter divergence theorem applied to a Lipschitz truncation of
$X$. Fix $\eps\in(0,\rstar/8)$ and a smooth non-decreasing
$\eta_\eps:[0,\infty)\to[0,1]$ with $\eta_\eps(r)=0$ for $r\le\eps/2$,
$\eta_\eps(r)=1$ for $r\ge\eps$, $|\eta_\eps'|\le4/\eps$. Set
\begin{equation}\label{eq:Xeps-def}
   X_\eps(x):=\eta_\eps(|x|)\,w(|x|)\,e_r(x),\qquad x\in\R^n.
\end{equation}
Then $X_\eps\in C^{0,1}(\R^n;\R^n)\cap C^\infty(\R^n\setminus\{0\};\R^n)$,
$X_\eps\equiv X$ outside $\{|x|<\eps\}$, and
$\|X_\eps\|_\infty\le 3R/\rstar$. The classical divergence
$\mathrm{div}\,X_\eps$ is locally bounded on $\R^n$ and equals
\begin{equation}\label{eq:divXeps}
   \mathrm{div}\,X_\eps(x)
   = \eta_\eps(|x|)\Bigl[w'(|x|)+\frac{n-1}{|x|}w(|x|)\Bigr]
     +\eta_\eps'(|x|)\,w(|x|)\qquad(x\ne0),
\end{equation}
where $w'(r)=r/(\rho|r|)\cdot\mathrm{sgn}(r/\rho-1)=\rho^{-1}\mathrm{sgn}(r/\rho-1)$
for $r\ne\rho$ (Lipschitz regularity of $w$, derivative a.e.).
At the origin, $X_\eps\equiv 0$ on $\{|x|<\eps/2\}$, so $X_\eps$ is genuinely
Lipschitz on $\R^n$.

By the divergence theorem for the finite-perimeter set $E_\rho$ applied to
the Lipschitz vector field $X_\eps$ (item~A4 of the preliminaries;
\cite[Thm.\ 16.3]{Maggi2012}, applied with the standard density argument:
$X_\eps$ is mollifiable since it is Lipschitz with compact support after
multiplication by a cutoff at radius $3R$),
\begin{equation}\label{eq:div-eps}
   \int_{\partial^*E_\rho}X_\eps\cdot\nu_\rho\,d\mathcal H^{n-1}
   =\int_{E_\rho}\mathrm{div}\,X_\eps\,dx.
\end{equation}
The same identity applied to the open ball $B_\rho=B_\rho(0)$ (a smooth
finite-perimeter set, with $\nu_{B_\rho}(y)=y/|y|$ and $w(\rho)=0$) yields
\begin{equation}\label{eq:div-eps-ball}
   0
   =\int_{\partial^*B_\rho}\!\!\!\!w(|y|)\,e_r\cdot\nu_{B_\rho}\,d\mathcal H^{n-1}
   =\int_{\partial^*B_\rho}\!\!\!\!X_\eps\cdot\nu_{B_\rho}\,d\mathcal H^{n-1}
   =\int_{B_\rho}\!\!\mathrm{div}\,X_\eps\,dx,
\end{equation}
where the first equality holds because $\eta_\eps(r)=1$ for $r$ near $\rho$
(since $\eps<\rstar/8\le\rho/8$) and $w(\rho)=0$. Subtracting,
\begin{equation}\label{eq:div-eps-diff}
   \int_{\partial^*E_\rho}X_\eps\cdot\nu_\rho\,d\mathcal H^{n-1}
   =\int_{E_\rho\setminus B_\rho}\!\!\!\!\mathrm{div}\,X_\eps\,dx
    -\int_{B_\rho\setminus E_\rho}\!\!\!\!\mathrm{div}\,X_\eps\,dx.
\end{equation}

\emph{Step 4: limit $\eps\downarrow0$.} On $\R^n\setminus\{0\}$,
$X_\eps\to X$ pointwise and $\mathrm{div}\,X_\eps\to\mathrm{div}\,X$
pointwise, where
\begin{equation}\label{eq:divX}
   \mathrm{div}\,X(x)=
   \rho^{-1}\mathrm{sgn}(|x|/\rho-1)+\frac{n-1}{|x|}\Bigl|\frac{|x|}{\rho}-1\Bigr|,
   \qquad x\in\R^n\setminus(\{0\}\cup\partial B_\rho).
\end{equation}
On both $E_\rho\setminus B_\rho$ and $B_\rho\setminus E_\rho$, the
domain of integration is contained in $\{|x|\ge\rstar/4\}$ (by
Lemma~\ref{lem:centroid-borel} for the first set and the trivial bound
$\rho\ge\rstar$ for the second after restricting to $|x|\ge\rstar/4$;
the set $\{|x|<\rstar/4\}\cap(E_\rho\cup B_\rho)\subset B_\rho$ is
\emph{not} in either symmetric difference part by~\eqref{eq:zbar-CL}).
On $\{|x|\ge\rstar/4\}$, the singular factor
$(n-1)/|x|\le 4(n-1)/\rstar$ and the cutoff factor $\eta_\eps(|x|)\equiv1$
for $\eps<\rstar/8$; therefore $X_\eps\equiv X$ on
$\{|x|\ge\rstar/4\}$ for all $\eps$ small enough, and
\eqref{eq:div-eps-diff} is in fact independent of $\eps\in(0,\rstar/8)$
and equals
\begin{equation}\label{eq:div-final}
   \int_{\partial^*E_\rho}X\cdot\nu_\rho\,d\mathcal H^{n-1}
   =\int_{E_\rho\setminus B_\rho}\!\!\!\mathrm{div}\,X\,dx
    -\int_{B_\rho\setminus E_\rho}\!\!\!\mathrm{div}\,X\,dx.
\end{equation}
(The left-hand side is $\int X\cdot\nu_\rho\,d\mathcal H^{n-1}$, the limit
of the same integral with $X_\eps$ in place of $X$, by the same argument
applied to the boundary trace.) This is the truncation step (A5 of the
preliminaries): the locally integrable singularity of $X$ at $\zbar=0$
is harmless because $\zbar\in\mathrm{int}\,E_\rho$ and the domains
$E_\rho\setminus B_\rho$, $B_\rho\setminus E_\rho$ are
\emph{disjoint from a fixed neighbourhood} of $\zbar$.

\emph{Step 5: pointwise bound on $\mathrm{div}\,X$.}
For $r:=|x|>\rho$ (i.e.\ $x\in E_\rho\setminus B_\rho$):
\begin{equation}\label{eq:divX-out}
   \mathrm{div}\,X(x)
   =\rho^{-1}+\frac{n-1}{r}\Bigl(\frac{r}{\rho}-1\Bigr)
   =\frac{n}{\rho}-\frac{n-1}{r}
   \le\frac{n}{\rho}\le\frac{n}{\rstar}.
\end{equation}
For $r<\rho$ (i.e.\ $x\in B_\rho\setminus E_\rho$, $r\ge\rstar/4$):
\begin{equation}\label{eq:divX-in}
   \mathrm{div}\,X(x)
   =-\rho^{-1}+\frac{n-1}{r}\Bigl(1-\frac{r}{\rho}\Bigr)
   =\frac{n-1}{r}-\frac{n}{\rho},
\end{equation}
so
$(-\mathrm{div}\,X)_+(x)\le n/\rho\le n/\rstar$. Combining with
\eqref{eq:div-final},
\begin{equation}\label{eq:T-rad-bound}
   \int_{\partial^*E_\rho}X\cdot\nu_\rho\,d\mathcal H^{n-1}
   \le\frac{n}{\rstar}\,|E_\rho\Delta B_\rho(\zbar)|.
\end{equation}

\emph{Step 6: assembly.}
Combining~\eqref{eq:T-split}, \eqref{eq:T-tan-bound},
and~\eqref{eq:T-rad-bound},
\begin{equation}\label{eq:T-trace-explicit}
   T_\rho
   \le\frac{n}{\rstar}\,|E_\rho\Delta B_\rho(\zbar)|
   +\frac{3R}{2\rstar}\,
     \int_{\partial^*E_\rho}|\nu_\rho-e_r|^2\,d\mathcal H^{n-1}.
\end{equation}
This is~\eqref{eq:radial-trace} with
$C_4(n,R,\rstar)=\max\{n/\rstar,\,3R/(2\rstar)\}$, polynomial in $n,R,1/\rstar$.
\end{proof}

\begin{remark}[Why the truncation is legitimate]\label{rmk:truncation}
The truncation factor $\eta_\eps$ is needed only to ensure that the
classical divergence theorem applies to a Lipschitz field on all of
$\R^n$; the field $X$ is bounded on $B_R$, but its classical derivative
$\mathrm{div}\,X$ has a singularity of type $|x|^{-1}$ at $\zbar$. Two
properties make the truncation harmless. First,
$\zbar\in\mathrm{int}\,E_\rho$ (Lemma~\ref{lem:centroid-borel}), so the
small ball $\{|x-\zbar|<\rstar/4\}$ lies entirely in $E_\rho\cap B_\rho$
and contributes \emph{equally} to both $E_\rho$ and $B_\rho$ in
\eqref{eq:div-eps-diff}, cancelling out. Second, on the
symmetric-difference set $\{x:\mathbf1_{E_\rho}(x)\ne\mathbf1_{B_\rho}(x)\}$
the radius $|x-\zbar|\ge\rstar/4$ (since this set lies in
$E_\rho\Delta B_\rho\subset B_R\setminus B(\zbar,\rstar/4)$ once
$\delta_0$ is small), so $X=X_\eps$ on this set for all $\eps<\rstar/8$
and the limit~\eqref{eq:div-final} is exact, not a passage. This is the
key estimate of the truncation step: the singularity of $X$ at $\zbar$
contributes nothing because it is interior to both $E_\rho$ and $B_\rho$.
\end{remark}

\begin{remark}[Equivalent slicing route]\label{rmk:slicing-route}
An alternative proof of~\eqref{eq:radial-trace} by a spherical slicing
decomposition $T_\rho=T_\rho^{\mathrm{rad,slic}}+T_\rho^{\mathrm{tan}}$
can be carried out via the radial slicing formula
\cite[Thm.\ 18.11]{Maggi2012} and a good-ray/bad-ray multiplicity
decomposition. That route yields the same bound at the integrated level
and trades the divergence-theorem step here for the Cicalese--Leonardi
annular containment. Both routes produce the constant $C_4(n,R,\rstar)$
in the same polynomial form. The divergence-theorem route is preferred
here because it isolates the named vector field~\eqref{eq:X-def} and
avoids reference to angular multiplicity.
\end{remark}

\subsubsection{4.3.4.\ Consequences for the metric first-variation chain}
\label{ssec:fj-output}

We collect the levelwise outputs that feed into the integrated action
bound of Theorem~\ref{thm:weighted-trace} (\S4.4).

\begin{corollary}[Pointwise radial trace from $D_I$]\label{cor:T-pointwise}
There exists $C_5=C_5(n,R,\rstar)$ such that for every $\rho\in G_I$,
\begin{equation}\label{eq:T-pointwise}
   T_\rho^2
   \le C_5(n,R,\rstar)\,D_I(t(\rho)).
\end{equation}
\end{corollary}

\begin{proof}
By Theorem~\ref{thm:radial-trace},
$T_\rho^2\le 2C_4^2(|E_\rho\Delta B_\rho(\zbar)|^2
+(\int|\nu_\rho-e_r|^2\,d\mathcal H^{n-1})^2)$. The first term is
$\le 2C_4^2\cdot C_2\,D_I$ by Lemma~\ref{lem:fraenkel-centroid}. For the
second, use the trivial estimate
$(\int|\nu_\rho-e_r|^2\,d\mathcal H^{n-1})^2\le 4P(B_R)\cdot\int|\nu_\rho-e_r|^2\,d\mathcal H^{n-1}$
on $G_I$ (since the integrand is bounded by~$4$ pointwise and the
perimeter $P(E_\rho)\le 2P(B_\rho)\le 2n\omega_n R^{n-1}$ on $G_I$),
together with Lemma~\ref{lem:normal-osc-centroid}, to get
$\le 8n\omega_n R^{n-1}C_3\,D_I$. Combining,
$C_5=2C_4^2(C_2+8n\omega_n R^{n-1}C_3)$, polynomial in $n,R,1/\rstar$.
\end{proof}

\begin{corollary}[Good-level homothetic trace at the centroid]
\label{cor:H-trace-good}
With $\bar V_\rho:=P(B_\rho)/P(E_\rho)$ and
$H_{z,\rho}(x)=(x-z)\cdot\nu_\rho(x)/\rho$ on $\partial^*E_\rho$, there
exists $C_6=C_6(n,R,\rstar)$ such that for every $\rho\in G_I$,
\begin{equation}\label{eq:H-trace-good}
   \Bigl(\int_{\partial^*E_\rho}|H_{\zbar,\rho}-\bar V_\rho|\,
       d\mathcal H^{n-1}\Bigr)^2
   \le C_6(n,R,\rstar)\,D_I(t(\rho)).
\end{equation}
\end{corollary}

\begin{proof}
Write $e_r(x)=(x-\zbar)/|x-\zbar|$. Then
$H_{\zbar,\rho}(x)=(|x-\zbar|/\rho)\,e_r\cdot\nu_\rho$ and
\[
   |H_{\zbar,\rho}-1|
   \le\Bigl|\frac{|x-\zbar|}{\rho}-1\Bigr|+|e_r\cdot\nu_\rho-1|.
\]
The second term equals $\tfrac12|e_r-\nu_\rho|^2$. Integrating,
\[
   \int_{\partial^*E_\rho}|H_{\zbar,\rho}-1|\,d\mathcal H^{n-1}
   \le T_\rho+\frac12\int_{\partial^*E_\rho}|\nu_\rho-e_r|^2\,d\mathcal H^{n-1}.
\]
Squaring and using $(a+b)^2\le2a^2+2b^2$ together with
Corollary~\ref{cor:T-pointwise} and Lemma~\ref{lem:normal-osc-centroid}
gives
$(\int|H_{\zbar,\rho}-1|)^2\le C\,D_I$ on $G_I$. The replacement of
the centring constant $1$ by $\bar V_\rho=1-(P(E_\rho)-P(B_\rho))/P(E_\rho)$
costs an additive term
\[
   |1-\bar V_\rho|\cdot P(E_\rho)=P(E_\rho)-P(B_\rho)
   \le D_I(t(\rho))/(n\omega_n\rstar^{n-1})\le D_I^{1/2}/(n\omega_n\rstar^{n-1})
\]
on $G_I$ (using $D_I\le\eta_I\le1$ on $G_I$). Squaring and absorbing
this term yields~\eqref{eq:H-trace-good} with
$C_6=2(C_5+C_3)+2/(n\omega_n\rstar^{n-1})^2$, polynomial in $n,R,1/\rstar$.
\end{proof}

\begin{remark}[Hypothesis tracking]\label{rmk:hypothesis-tracking}
Every estimate above is stated for open $\Om\subset B_R$ with
$|\Om|=\omega_n$ and $\delta_T(\Om)\le\delta_0(n,R,\rstar)$ small enough
to ensure Cicalese--Leonardi smallness for every level
$\rho\in[\rstar,1]$ via \cite[Thm.~1.1]{CL2012}. No graph
parametrisation, no minimiser selection, and no smoothness of
$\partial^*E_\rho$ are used. All boundary integrals are on the De Giorgi
reduced boundary. The Borel selection of the centroid is unconditional
(Lemma~\ref{lem:centroid-borel}); the FJ centre selection of
Remark~\ref{rmk:centre-selection} reduces to it after the
$O(\sqrt{\delta_T})$ proximity absorption recorded in~\cite[Thm.~1.1]{FJ2014}.
Outside the smallness regime, the FMP fallback
$\Fr(\Om)^2\le C(n)\delta_T^{1/2}$ together with $\delta_T\ge c(n)>0$
gives all four estimates trivially.
\end{remark}

\paragraph*{Constants used in this subsection.} For the convenience of
\S4.4 we collect:
\begin{itemize}
\item $C_1(n,\rstar)$ in \eqref{eq:plan2-FJ-scaled}: scaled FJ constant,
polynomial in $\rstar^{-1}$.
\item $C_2(n,R,\rstar)$ in~\eqref{eq:fraenkel-centroid}: Fraenkel
estimate at the centroid.
\item $C_3(n,R,\rstar)$ in~\eqref{eq:normal-osc-centroid}: normal
oscillation.
\item $C_4(n,R,\rstar)=\max\{n/\rstar,3R/(2\rstar)\}$ in~\eqref{eq:radial-trace}:
\textbf{the oriented $L^1$ radial trace constant}, the load-bearing
output of Theorem~\ref{thm:radial-trace}.
\item $C_5(n,R,\rstar),C_6(n,R,\rstar)$ in Corollaries~\ref{cor:T-pointwise},
\ref{cor:H-trace-good}: levelwise inputs to \S4.4.
\end{itemize}
Every constant is polynomial in $n,R,\rstar^{-1}$ and independent of
$\delta_T(\Om)$ and of the boundary-layer constant $\kappa$.

%---- end of plan2_fj_center_radial_trace.tex ----

\label{subsec:fj-radial-trace}%
\label{subsec:fusco-julin-radial-trace}%
\label{sec:fj-radial}%
\label{ssec:plan2-radial-trace}%
\label{sec:FJ}%
% UNIFIED 2026-05-13
% RE-AUDIT WARNING — DATE 2026-05-14
% The endpoint repair below should not be read as an unconditional closure.
% Two issues remain after re-audit:
%   (1) the small-window trace argument gives an O(sqrt(delta_T)) metric
%       Markov contribution, not O(delta_T);
%   (2) Lemma `lem:wmt:A0` uses a perimeter-free dilation Lipschitz estimate
%       that is false for rough bounded measurable sets.
% The corrected route-status digest reduces the remaining obstruction to the
% bad-(-t') Lebesgue kinetic estimate
%   int_{B_tau} |dot F_rho|_X^2 d rho <= C delta_T.
% Keep this subsection as a draft until that estimate is supplied or the
% trace transport step is replaced.
% AUDIT B1/M1/M2 REPAIR — DATE 2026-05-13
% -----------------------------------------------------------------------------
% This file has been edited to address the BLOCKING (B1) and MAJOR (M1, M2)
% findings of `Manuscript/audit_plan2.md` (Agent 16). Summary of changes:
%   B1: Theorem `thm:wmt:trace` now instantiates the abstract weighted-trace
%       lemma on the window `[rho_delta - c_0 sqrt(delta_T), rho_delta]` with
%       L* = c_0 sqrt(delta_T), placing widehat-rho within an
%       O(sqrt(delta_T))-neighbourhood of rho_delta. (H2) is verified on the
%       new window using a smallness threshold delta_T <= delta_0(n,R,rho_*).
%       Closing remark is rewritten so that |Omega \ E_{widehat-rho}| =
%       O(sqrt(delta_T)).
%   M1: Lemma `lem:fj-center-package` now explicitly identifies z_rho := bar
%       z_rho (bulk centroid), citing Agent 11 Remark `rmk:centre-selection`.
%   M2: kappa = kappa(n, R, rho_*) is pinned canonically (per notation
%       register §12). All constants C_*^act, C_*^kin, C_* now have the
%       signature C_*(n, R, rho_*).
% See `Manuscript/audit_plan2_repair_report.md`.
% -----------------------------------------------------------------------------
% SOURCE MAP (Agent 12, Plan 2 §4.4 — Integrated Action and Weighted Metric Trace)
% -----------------------------------------------------------------------------
% Local building blocks consumed (LOCAL):
%   - Plan 2/WIP/WIP_WeightedMetricTrace.tex
%       * Theorem (integrated action bound, eq:52int)
%       * Theorem (kinetic bound K, eq:K)
%       * Theorem (weighted metric trace, eq:trace)
%       * Lemmas: FJ centre package, radial-slicing, good-level homothetic
%         trace, uniform metric Lipschitz (A0), bad-set lemma
%   - Plan 2/WIP/WIP_MetricFramework.tex
%       * Theorem T: |F'_rho|_X <= rho_*^{-n} inf_a int |V_rho - H_{z,rho}
%                       - a . nu_rho| dH^{n-1}
%       * AGS absolute continuity / metric derivative (Thm 1.1.2)
%   - Plan 2/WIP/WIP_LevelSetIdentity.tex
%       * Exact deficit identity: (2/|Omega|) (E(Omega) - E(B^*))
%           = (1/|Omega|) int (D_H + D_I)/(c_n m(t)^{1-2/n}) dt
%       * Profile-gap corollary used for (H2) and bad-set
%   - Plan 2/agent4-weighted-metric-trace
%       * Abstract weighted metric trace lemma (proof of Thm in §4)
%       * Bad-set Lemma 3.1; Corollary 3.2
%   - Plan 2/wave2-B-kinetic-bound
%       * Uniform L1-Lipschitz bound (A0), Lemma 2.1
%       * Bad-(-t') Lemma 4.1; integrated Talenti gap (3.3)
%   - Plan 2/wave3-G-route-delta-assembly.md  (verification that the
%       integrated chain consumes (R) only against d\mu, no rate loss)
%
% Cited from the literature (CITE; see Manuscript/04_references_audit.md):
%   - Ambrosio-Gigli-Savaré 2008, Thm 1.1.2          [AGS2008]
%   - Fusco-Julin 2014, Calc. Var. PDE 50            [FJ2014]
%   - Talenti 1976, Ann. Mat. Pura Appl. (4) 110     [Talenti1976]
%   - Maggi 2012 (sets of finite perimeter)          [Maggi2012]
%   - Ambrosio-Fusco-Pallara 2000 (BV book)          [AFP2000]
%   - Figalli-Maggi 2011, Appendix A, Thm A.1        [FigalliMaggi2011]
%
% Imported by Agent 19 into Manuscript/03_plan2_draft.tex as Section §4.4.
% Constants tracked: every C, c is annotated with its dependence on (n, R, rho_*).
% (Per the M2 repair dated 2026-05-13, kappa = kappa(n, R, rho_*) is canonical,
% so all C_*^{act}, C_*^{kin}, C_* have signature (n, R, rho_*).)
% -----------------------------------------------------------------------------

\subsection{Integrated action and the weighted metric trace}\label{subsec:weighted-metric-trace}

This subsection chains the level-set deficit identity of
Section~\ref{subsec:level-set-identity}, the metric first-variation
estimate of Section~\ref{subsec:metric-framework}, the scaled
Fusco--Julin inequality and the oriented radial trace lemma of
Section~\ref{subsec:fj-radial-trace} into the integrated action bound
\eqref{eq:wmt:52int} and, via a good/bad isoperimetric-defect
decomposition, into the Lebesgue kinetic bound \eqref{eq:wmt:K}. The
final output is the \emph{weighted metric trace}: a radius
$\rhad\in[\rstar,\rhodelta)$ at which the rescaled level set
$F_\rhad=\rhad^{-1}E_\rhad$ is metrically close to the unit ball at the
sharp $\deltaT$-rate; see Theorem~\ref{thm:wmt:trace}.

Throughout this subsection we adopt the standing hypotheses of
Section~\ref{sec:plan2}: $n\ge 2$, $R\ge 2$, $\rstar\in(0,1)$, an open set
$\Om\subset B_R$ with $|\Om|=\omega_n$, variational torsion function
$u=u_\Om$, $\rho$-foliation $E_\rho=\{u>t(\rho)\}$ with
$|E_\rho|=\omega_n\rho^n$ for $\rho\in[\rstar,1]$, rescaled level sets
$F_\rho:=\rho^{-1}E_\rho\subset B_{R/\rstar}$, and the measure
$\mu(d\rho):=(-t'(\rho))\,d\rho$ on $[\rstar,1]$. We write
$\delta:=\deltaT(\Om)$.

\paragraph{Canonical choice of $\kappa$.}
Per the notation register \S12, the boundary-layer constant $\kappa$ is
a fixed dimensional choice $\kappa=\kappa(n,R,\rstar)>0$. Throughout
\S4.4 we pin $\kappa$ to this canonical value; explicitly, we set
\[
   \kappa=\kappa(n,R,\rstar):=1
\]
(any fixed positive constant of size $\Theta_n(1)$ works; the value
$\kappa=1$ is the simplest choice consistent with the notation
register). All subsequent constants $C_*^{\mathrm{act}},
C_*^{\mathrm{kin}}, C_*$ then become functions of $(n,R,\rstar)$ alone;
the smallness threshold $\delta_0$ likewise depends only on
$(n,R,\rstar)$. We define
\[
   \rhodelta:=1-\kappa\sqrt{\delta}.
\]
Throughout we assume
\begin{equation}\label{eq:wmt:delta-small}
   \delta\le\delta_0(n,R,\rstar),
\end{equation}
with $\delta_0(n,R,\rstar)>0$ chosen so that
$\rhodelta\ge\rstar+\tfrac12(1-\rstar)$ and so that the additional
smallness requirements (B1) of Theorem~\ref{thm:wmt:trace} below hold;
for $\delta$ outside this regime the conclusion is vacuous or follows
trivially from the FMP fallback $\Fr(\Om)^2\le C(n)\deltaT(\Om)$ for
$\deltaT(\Om)\ge\delta_0$.

The level-set deficit identity of Section~\ref{subsec:level-set-identity}
reads, in our $\rho$-parametrisation,
\begin{equation}\label{eq:wmt:DH-DI-int}
   \int_{\rstar}^{1}
      \bigl(\mathcal H(\rho)+\mathcal I(\rho)\bigr)\,d\mu(\rho)
   \;\le\; C_1(n,\rstar)\,\delta,
\end{equation}
where we have written
\[
   \mathcal H(\rho):=D_H(t(\rho)),\qquad
   \mathcal I(\rho):=D_I(t(\rho)),
\]
and used $dt=-t'(\rho)\,d\rho=d\mu$. The explicit constant is
$C_1(n,\rstar)=2\,c_n\,(\omega_n\rstar^n)^{1-2/n}/\omega_n
=2n^2\omega_n^{2/n}(\omega_n\rstar^n)^{1-2/n}/\omega_n$; see
Theorem~\ref{thm:lsid:identity}.

Since $P(E_\rho)+P(B_\rho)\ge 2P(B_\rho)=2n\omega_n\rho^{n-1}\ge
2n\omega_n\rstar^{n-1}$ and $D_I(t(\rho))
=(P(E_\rho)-P(B_\rho))(P(E_\rho)+P(B_\rho))$, the perimeter defect
satisfies
\begin{equation}\label{eq:wmt:DP-DI}
   0\le\Delta P_\rho
        :=P(E_\rho)-P(B_\rho)
   \le \frac{\mathcal I(\rho)}{2n\omega_n\rstar^{n-1}}
   =: C_2(n,\rstar)\,\mathcal I(\rho).
\end{equation}

\paragraph{Notation for this subsection}

\begin{itemize}
\item $\gamma(\rho):=[F_\rho]\in\X$ is the rescaled curve in the metric
      quotient $(\X,\dX)$ of Section~\ref{subsec:metric-framework}.
      AGS \cite[Thm.~1.1.2]{AGS2008} provides the metric derivative
      $|\dot\gamma|(\rho)=|\dot F_\rho|_{\X}$.
\item $z_\rho$ denotes the bulk centroid of $E_\rho$, identified in
      Lemma~\ref{lem:fj-center-package} below with the canonical centre
      $\zbar=\bar z_\rho$ of the notation register \S10. Outside the
      good set $G_I$ defined in \eqref{eq:wmt:def-GI} we set
      $z_\rho=0$.
\item $\bar V_\rho:=P(E_\rho)^{-1}\int_{\partial^*E_\rho}V_\rho\,
      d\mathcal H^{n-1}=P(B_\rho)/P(E_\rho)$, which equals
      the mean value of $V_\rho$ on $\partial^*E_\rho$ and satisfies
      $1-\bar V_\rho=\Delta P_\rho/P(E_\rho)$.
\item $H_{z_\rho,\rho}(x):=(x-z_\rho)\cdot\nu_\rho(x)/\rho$ is the
      homothetic radial field of Section~\ref{subsec:metric-framework}.
\end{itemize}

\subsubsection{4.4.1 Good and bad isoperimetric-defect levels}
\label{subsec:wmt:good-bad}

\begin{definition}[Good and bad isoperimetric-defect levels]
\label{def:wmt:GI-BI}
Let $\eta_I=\eta_I(n,R,\rstar)>0$ be the Fusco--Julin smallness
threshold from Lemma~\ref{lem:fj-center-package} below. Define
\begin{align}
   G_I &:= \{\rho\in[\rstar,\rhodelta]:\mathcal I(\rho)\le\eta_I\},
       \label{eq:wmt:def-GI}\\
   B_I &:= [\rstar,\rhodelta]\setminus G_I
       =\{\rho\in[\rstar,\rhodelta]:\mathcal I(\rho)>\eta_I\}.
       \label{eq:wmt:def-BI}
\end{align}
\end{definition}

The threshold $\eta_I=\eta_I(n,R,\rstar)$ will be fixed at the end of
Lemma~\ref{lem:fj-center-package}: it is the smallest of the
Fusco--Julin smallness scale on $\{|E|=|B_\rho|,\rho\in[\rstar,1]\}$ (see
Proposition~\ref{prop:fj-scaled}) and the elementary smallness
threshold needed for the FJ centre to lie in $B_{2R}$. In particular,
$\eta_I$ does \emph{not} depend on $\delta$ or on $\Om$.

\begin{lemma}[Markov bound on $B_I$]\label{lem:wmt:markov-BI}
Under \eqref{eq:wmt:delta-small},
\begin{equation}\label{eq:wmt:markov-BI}
   \mu(B_I)
   \;\le\; \eta_I^{-1}\int_{B_I}\mathcal I(\rho)\,d\mu(\rho)
   \;\le\; \eta_I^{-1}\int_{\rstar}^{1}
              \bigl(\mathcal H+\mathcal I\bigr)\,d\mu
   \;\le\; \eta_I^{-1}\,C_1(n,\rstar)\,\delta.
\end{equation}
\end{lemma}

\begin{proof}
By Definition~\ref{def:wmt:GI-BI}, $\mathcal I(\rho)>\eta_I$ on $B_I$,
hence
$\mu(B_I)=\int_{B_I}d\mu \le \eta_I^{-1}\int_{B_I}\mathcal I\,d\mu$.
Both $\mathcal H,\mathcal I\ge 0$, so the second inequality is trivial.
The last inequality is \eqref{eq:wmt:DH-DI-int}.
\end{proof}

\begin{remark}[The sharp rate is preserved]\label{rem:wmt:rate-good-bad}
The Markov bound \eqref{eq:wmt:markov-BI} is the only place where the
size of $B_I$ enters; the rate is $\mu(B_I)=O(\delta)$, not
$O(\sqrt{\delta})$. The good-set quantitative estimates
\eqref{eq:wmt:FJ-asym}--\eqref{eq:wmt:FJ-normal} below are pointwise in
$\mathcal I(\rho)$, which is then integrated against $d\mu$ via
\eqref{eq:wmt:DH-DI-int}. No square root of $\mathcal I$ is integrated
unweighted; consequently the integrated bound \eqref{eq:wmt:52int}
preserves the sharp $\delta$ rate. This is the point at which the
naive route would lose a $\sqrt{\delta}$.
\end{remark}

\begin{lemma}[FJ centre package, fixed to the bulk centroid]
\label{lem:fj-center-package}\label{lem:FJ-package}
There exist $\eta_I=\eta_I(n,R,\rstar)>0$ and
$C_{\mathrm{FJ}}=C_{\mathrm{FJ}}(n,R,\rstar)$ such that, defining the
centre field on $G_I$ as the bulk centroid
\begin{equation}\label{eq:wmt:zrho-is-centroid}
   z_\rho:=\zbar=\bar z_\rho=|E_\rho|^{-1}\int_{E_\rho}x\,dx,
   \qquad \rho\in G_I,
\end{equation}
and $z_\rho:=0$ on $B_I$, the map $\rho\mapsto z_\rho$ is Borel
measurable on $[\rstar,\rhodelta]$, satisfies $|z_\rho|\le 2R$ on
$G_I$, and for every $\rho\in G_I$,
\begin{align}
   |E_\rho\Delta B_\rho(z_\rho)|^2
      &\le C_{\mathrm{FJ}}\,\mathcal I(\rho),\label{eq:wmt:FJ-asym}\\
   \int_{\partial^*E_\rho}
      \Bigl|\nu_\rho(x)-\frac{x-z_\rho}{|x-z_\rho|}\Bigr|^2
      d\mathcal H^{n-1}(x)
      &\le C_{\mathrm{FJ}}\,\mathcal I(\rho).\label{eq:wmt:FJ-normal}
\end{align}
Moreover the centroid containment
\(B(z_\rho,\rho/4)\subset E_\rho\subset B(z_\rho,3\rho)\) of
Lemma~\ref{lem:centroid-borel} (Section~\ref{subsec:fj-radial-trace})
holds on $G_I$, justifying the use of the oriented radial trace lemma
at $z_\rho$.
\end{lemma}

\begin{proof}
The choice $z_\rho:=\zbar$ on $G_I$ is the canonical one of the
notation register \S10, cf.\ also Agent~11
Remark~\ref{rmk:centre-selection} of
Section~\ref{subsec:fj-radial-trace}. By that remark, the centroid
$\zbar$ is admissible in the Fusco--Julin bound \eqref{eq:plan2-FJ-scaled}
up to absorbing the centroid--Fraenkel offset
$|\zbar-z^{\mathrm{Fr}}_\rho|\le C(n,\rstar)\sqrt{\delta_T}$ into the
constant; the smallness threshold \eqref{eq:wmt:delta-small} (with
$\delta_0$ shrunk if necessary so that
$C(n,\rstar)\sqrt{\delta_0}\le\rho/4$) is precisely what makes this
absorption legitimate. The Fraenkel bound
\eqref{eq:wmt:FJ-asym} and the normal-oscillation bound
\eqref{eq:wmt:FJ-normal} at $z_\rho=\zbar$ are then
Lemmas~\ref{lem:fraenkel-centroid} and~\ref{lem:normal-osc-centroid}
of Section~\ref{subsec:fj-radial-trace}, combined with the divergence
identity \(\int_{\partial^*E_\rho}|\nu_\rho-(x-z_\rho)/|x-z_\rho||^2\,
d\mathcal H^{n-1}=2\beta(E_\rho,z_\rho)^2\) (see
Theorem~\ref{thm:fj-radial-divergence}). The Lipschitz/Borel
regularity of $\rho\mapsto\zbar$ is Lemma~\ref{lem:centroid-borel};
$|\zbar|\le 2R$ follows from $\zbar\in E_\rho\subset B_R$. The
threshold $\eta_I=\eta_I(n,R,\rstar)$ is the smaller of the
Fusco--Julin smallness scale on $\{|E|=|B_\rho|,\rho\in[\rstar,1]\}$
and the threshold ensuring centroid containment
\(B(\zbar,\rho/4)\subset E_\rho\subset B(\zbar,3\rho)\) of
Lemma~\ref{lem:centroid-borel}.
\end{proof}

\begin{lemma}[Oriented radial trace at the FJ centre]
\label{lem:wmt:radial-trace}
For every $\rho\in G_I$,
\begin{equation}\label{eq:wmt:radial-trace}
   T_\rho:=\int_{\partial^*E_\rho}
      \Bigl|\frac{|x-z_\rho|}{\rho}-1\Bigr|\,d\mathcal H^{n-1}(x)
   \;\le\; C_3(n,R,\rstar)\,\mathcal I(\rho)^{1/2}.
\end{equation}
\end{lemma}

\begin{proof}
This is Lemma~\ref{lem:radial-trace-master} of
Section~\ref{subsec:fj-radial-trace}, applied with the centre field of
Lemma~\ref{lem:fj-center-package} and combined with
\eqref{eq:wmt:DP-DI}, \eqref{eq:wmt:FJ-asym}, and
\eqref{eq:wmt:FJ-normal}. Since $\mathcal I(\rho)\le\eta_I\le 1$ on
$G_I$, the FJ asymmetry square-root piece $\mathcal I(\rho)^{1/2}$
dominates the linear piece $\mathcal I(\rho)$.
\end{proof}

\begin{lemma}[Good-level homothetic trace]\label{lem:wmt:good-htrace}
For every $\rho\in G_I$,
\begin{equation}\label{eq:wmt:htrace-good}
   \Biggl(\int_{\partial^*E_\rho}
       |H_{z_\rho,\rho}(x)-\bar V_\rho|\,
       d\mathcal H^{n-1}(x)\Biggr)^2
   \;\le\; C_4(n,R,\rstar)\,\mathcal I(\rho).
\end{equation}
\end{lemma}

\begin{proof}
Set $e_r(x):=(x-z_\rho)/|x-z_\rho|$. Then
$H_{z_\rho,\rho}(x)
   =(x-z_\rho)\cdot\nu_\rho(x)/\rho
   =(|x-z_\rho|/\rho)\,e_r\cdot\nu_\rho$.
A direct algebraic identity gives
\[
   H_{z_\rho,\rho}-1
   =\Bigl(\frac{|x-z_\rho|}{\rho}-1\Bigr)\,e_r\cdot\nu_\rho
     +(e_r\cdot\nu_\rho-1).
\]
Since $|e_r|=|\nu_\rho|=1$, $|e_r\cdot\nu_\rho-1|
=\tfrac12|e_r-\nu_\rho|^2$. Hence
\[
   \int_{\partial^*E_\rho}|H_{z_\rho,\rho}-1|\,d\mathcal H^{n-1}
   \le T_\rho+\tfrac12\int_{\partial^*E_\rho}|e_r-\nu_\rho|^2\,
       d\mathcal H^{n-1}.
\]
By Lemma~\ref{lem:wmt:radial-trace} and \eqref{eq:wmt:FJ-normal},
\(
   \int_{\partial^*E_\rho}|H_{z_\rho,\rho}-1|\,d\mathcal H^{n-1}
   \le C_3\,\mathcal I^{1/2}+\tfrac12 C_{\mathrm{FJ}}\,\mathcal I
   \le C_5(n,R,\rstar)\,\mathcal I^{1/2}.
\)
On the other hand
\[
   P(E_\rho)|1-\bar V_\rho|
   =P(E_\rho)-P(B_\rho)
   =\Delta P_\rho
   \le C_2\,\mathcal I(\rho)
   \le C_2\,\mathcal I(\rho)^{1/2},
\]
using $\mathcal I\le\eta_I\le 1$ on $G_I$. Triangle inequality gives
$\int|H_{z_\rho,\rho}-\bar V_\rho|\,d\mathcal H^{n-1}
\le C_6(n,R,\rstar)\,\mathcal I(\rho)^{1/2}$, which squares to
\eqref{eq:wmt:htrace-good} with
$C_4(n,R,\rstar):=C_6(n,R,\rstar)^2$.
\end{proof}

\subsubsection{4.4.2 The integrated action bound}
\label{subsec:wmt:integrated-action}

We now assemble \eqref{eq:wmt:DH-DI-int},
Lemma~\ref{lem:fj-center-package}, Theorem~T of
Section~\ref{subsec:metric-framework}, and the good-level homothetic
trace into the integrated action bound.

\begin{theorem}[Integrated action bound]\label{thm:wmt:52int}
There is $C_*^{\mathrm{act}}=C_*^{\mathrm{act}}(n,R,\rstar)$
such that
\begin{equation}\label{eq:wmt:52int}
   \int_{\rstar}^{\rhodelta}
     \Bigl(\dX(F_\rho,B_1)^2+|\dot F_\rho|_{\X}^2\Bigr)\,d\mu(\rho)
   \;\le\; C_*^{\mathrm{act}}\,\deltaT(\Om).
\end{equation}
\end{theorem}

\begin{proof}
We bound the two integrands separately. For both, we estimate
pointwise on $G_I$ and use the Markov bound on $B_I$.

\emph{Distance term.} On $G_I$, the quotient pseudodistance satisfies
\begin{equation}\label{eq:wmt:dist-trivial}
   \dX(F_\rho,B_1)
   =\dX([\rho^{-1}E_\rho],[B_1])
   \le \rho^{-n}\,|E_\rho\Delta B_\rho(z_\rho)|
   \le \rstar^{-n}\,|E_\rho\Delta B_\rho(z_\rho)|,
\end{equation}
because $\rho^{-n}\,\mathbf{1}_{\rho^{-1}A}=\rho^{-n}\mathbf{1}_A\circ
(\rho\,\cdot\,)$ and the Lebesgue measure of $\rho^{-1}A$ equals
$\rho^{-n}|A|$ for any measurable $A$ (see
Lemma~\ref{lem:metric-translation-volume}). By \eqref{eq:wmt:FJ-asym},
\[
   \dX(F_\rho,B_1)^2
   \le \rstar^{-2n}\,|E_\rho\Delta B_\rho(z_\rho)|^2
   \le \rstar^{-2n}\,C_{\mathrm{FJ}}\,\mathcal I(\rho),
   \qquad\rho\in G_I.
\]
On $B_I$, both $F_\rho$ and $B_1$ have volume $\omega_n$, so
$\dX(F_\rho,B_1)\le 2\omega_n$ trivially. Hence
\begin{align*}
   \int_{\rstar}^{\rhodelta}\dX(F_\rho,B_1)^2\,d\mu
   &\le \rstar^{-2n}\,C_{\mathrm{FJ}}\int_{G_I}\mathcal I\,d\mu
       +(2\omega_n)^2\,\mu(B_I)\\
   &\le \rstar^{-2n}\,C_{\mathrm{FJ}}\,C_1\,\delta
       +4\omega_n^2\,\eta_I^{-1}\,C_1\,\delta
   = C_7(n,R,\rstar)\,\delta,
\end{align*}
using \eqref{eq:wmt:DH-DI-int} and Lemma~\ref{lem:wmt:markov-BI}.

\emph{Metric-derivative term.} On $G_I$, Theorem~T of
Section~\ref{subsec:metric-framework} (applied to the centre
$z_\rho\in B_{2R}$ supplied by Lemma~\ref{lem:fj-center-package}, with
translation parameter $a=0$) gives
\begin{equation}\label{eq:wmt:M-thm-applied}
   |\dot F_\rho|_{\X}
   \;\le\; \rstar^{-n}
   \int_{\partial^*E_\rho}
       |V_\rho-H_{z_\rho,\rho}|\,d\mathcal H^{n-1}.
\end{equation}
Insert $\pm\bar V_\rho$ and use $(a+b)^2\le 2a^2+2b^2$:
\begin{equation}\label{eq:wmt:Mdot-split}
   |\dot F_\rho|_{\X}^2
   \;\le\; 2\rstar^{-2n}
   \left[
      \Bigl(\int_{\partial^*E_\rho}|V_\rho-\bar V_\rho|\,
                d\mathcal H^{n-1}\Bigr)^2
      +\Bigl(\int_{\partial^*E_\rho}|H_{z_\rho,\rho}-\bar V_\rho|\,
                d\mathcal H^{n-1}\Bigr)^2
   \right].
\end{equation}
The first square is the Cauchy--Schwarz defect from
Section~\ref{subsec:level-set-identity}, Proposition~\ref{prop:lsid:CS-defect}:
\begin{equation}\label{eq:wmt:CS-defect}
   \Bigl(\int_{\partial^*E_\rho}|V_\rho-\bar V_\rho|\,
       d\mathcal H^{n-1}\Bigr)^2
   \le C_8(n,R,\rstar)\,\mathcal H(\rho)
   \qquad\text{for a.e. }\rho\in[\rstar,1].
\end{equation}
The second square is Lemma~\ref{lem:wmt:good-htrace} (valid on $G_I$).
Therefore, on $G_I$,
\[
   |\dot F_\rho|_{\X}^2
   \;\le\; 2\rstar^{-2n}\bigl(C_8\,\mathcal H(\rho)
                              +C_4\,\mathcal I(\rho)\bigr)
   =: C_9(n,R,\rstar)\,(\mathcal H+\mathcal I)(\rho).
\]
On $B_I$ we use the uniform metric Lipschitz bound recorded as
Lemma~\ref{lem:wmt:A0} below: $|\dot F_\rho|_{\X}\le L_0(n,R,\rstar)$
a.e. on $[\rstar,1]$. Combining with \eqref{eq:wmt:DH-DI-int} and
Lemma~\ref{lem:wmt:markov-BI},
\begin{align*}
   \int_{\rstar}^{\rhodelta}|\dot F_\rho|_{\X}^2\,d\mu
   &\le C_9\int_{G_I}(\mathcal H+\mathcal I)\,d\mu
       +L_0^2\,\mu(B_I)\\
   &\le C_9\,C_1\,\delta+L_0^2\,\eta_I^{-1}\,C_1\,\delta
   = C_{10}(n,R,\rstar)\,\delta.
\end{align*}
Summing the two contributions yields \eqref{eq:wmt:52int} with
\[
   C_*^{\mathrm{act}}(n,R,\rstar):=C_7(n,R,\rstar)+C_{10}(n,R,\rstar).
\]
Since $\kappa=\kappa(n,R,\rstar)$ has been pinned at the start of \S4.4
(canonical choice $\kappa=1$), this constant depends only on
$(n,R,\rstar)$; the only role of $\kappa$ in the proof is implicit in
\eqref{eq:wmt:delta-small}, which ensures
$\rhodelta\ge\rstar+(1-\rstar)/2$.
\end{proof}

\begin{remark}[Why the rate is sharp]
The only square-root in the chain is in
Lemma~\ref{lem:wmt:radial-trace}: $T_\rho\le C\,\mathcal I^{1/2}$. This
square root is squared in \eqref{eq:wmt:htrace-good} before integration
against $d\mu$, after which it contributes the linear bound $C\,\mathcal
I(\rho)$. The bad set $B_I$ is paid for at rate $\mu(B_I)=O(\delta)$,
not $\sqrt{\delta}$, because the Markov bound divides $\delta$ by the
constant $\eta_I$, not by $\sqrt{\delta}$. Hence the integrated bound
\eqref{eq:wmt:52int} preserves the sharp $\delta$ rate. See
Remark~\ref{rem:wmt:rate-good-bad}.
\end{remark}

\subsubsection{4.4.3 Uniform metric Lipschitz and bad-$(-t')$ decomposition}
\label{subsec:wmt:A0-badtprime}

To upgrade the integrated action \eqref{eq:wmt:52int} into the Lebesgue
kinetic bound \eqref{eq:wmt:K} below we need two further elementary
inputs.

\begin{lemma}[Uniform metric Lipschitz bound]\label{lem:wmt:A0}
There exists $L_0=L_0(n,R,\rstar)$ such that
\begin{equation}\label{eq:wmt:A0}
   |\dot F_\rho|_{\X}\;\le\; L_0(n,R,\rstar)
   \qquad\text{for a.e. }\rho\in[\rstar,1].
\end{equation}
An explicit choice is $L_0:=n(R/\rstar)^n+n/\rstar^n$.
\end{lemma}

\begin{proof}
Fix $\rstar\le \rho<\rho'\le 1$. Since $\{E_\rho\}_\rho$ is nested
(because $t$ is decreasing) and $|E_\rho|=\omega_n\rho^n$,
\[
   |E_{\rho'}\Delta E_\rho|=|E_{\rho'}|-|E_\rho|
   =\omega_n\bigl((\rho')^n-\rho^n\bigr)
   \le n\omega_n(\rho'-\rho).
\]
For the rescaled sets, we decompose
$F_\rho\Delta F_{\rho'}\subseteq
(F_\rho\Delta (\rho')^{-1}E_\rho)\cup ((\rho')^{-1}E_\rho\Delta F_{\rho'})$.
For the second piece,
\(|(\rho')^{-1}E_\rho\Delta F_{\rho'}|
=(\rho')^{-n}|E_\rho\Delta E_{\rho'}|
\le \rstar^{-n}\,n\omega_n(\rho'-\rho).\)
For the first piece, the dilation
$\Lambda_\lambda(x):=\lambda x$ with $\lambda=\rho'/\rho\ge 1$, applied
to the bounded set $A:=(\rho')^{-1}E_\rho\subseteq B_{R/\rstar}$,
satisfies
\(|\Lambda_\lambda(A)\Delta A|\le|A|(\lambda^n-1)
\le n\omega_n(R/\rstar)^n(\lambda-1)
\le n\omega_n(R/\rstar)^n\,(\rho'-\rho)/\rstar.\)
Hence
\[
   \dX([F_{\rho'}],[F_\rho])
   \le |F_{\rho'}\Delta F_\rho|
   \le \bigl(n(R/\rstar)^n+n/\rstar^n\bigr)\,(\rho'-\rho)
   = L_0(n,R,\rstar)\,(\rho'-\rho).
\]
The AGS metric derivative (Theorem~\ref{thm:metric:AC} of
Section~\ref{subsec:metric-framework}) is then bounded a.e.\ by the
Lipschitz constant $L_0$.
\end{proof}

\begin{lemma}[Integrated Talenti gap and bad-$(-t')$ Lemma]
\label{lem:wmt:badtprime}
Set $\psi(\rho):=-t'(\rho)-(-t_B'(\rho))=-t'(\rho)-\rho/n$, where
$t_B(\rho)=(1-\rho^2)/(2n)$ is the ball profile, and
$\tau_0:=\rstar/(2n)$, with bad set
\begin{equation}\label{eq:wmt:def-Btau0}
   B_{\tau_0}:=\{\rho\in[\rstar,\rhodelta]:-t'(\rho)<\tau_0\}.
\end{equation}
Then $\psi\ge 0$ a.e. and
\begin{align}
   \int_{\rstar}^{1}|\psi(\rho)|\,d\rho
       &\le K_1(n,\rstar)\,\delta,
       \qquad K_1(n,\rstar):=\frac{2}{\omega_n\rstar^n},
       \label{eq:wmt:Talenti-gap}\\
   |B_{\tau_0}|
       &\le K_2(n,\rstar)\,\delta,
       \qquad K_2(n,\rstar):=\frac{4n}{\omega_n\rstar^{n+1}}.
       \label{eq:wmt:Btau0-bound}
\end{align}
\end{lemma}

\begin{proof}
The bound $\psi\ge 0$ a.e.\ is the differentiated Talenti comparison
$0\le t_B(\rho)-t(\rho)$, $\rho\in[\rstar,1]$
\cite[§3]{Talenti1976}; see Lemma~B.8 in
Section~\ref{sec:preliminaries-torsion}.

For \eqref{eq:wmt:Talenti-gap}: the level-set deficit identity
\eqref{eq:wmt:DH-DI-int} together with the explicit algebraic identity
\[
   D_H(t(\rho))+D_I(t(\rho))
   = P(B_\rho)^2\cdot\frac{-t_B'(\rho)-(-t'(\rho))}{-t'(\rho)}
   = P(B_\rho)^2\cdot\frac{|\psi(\rho)|}{-t'(\rho)}
\]
(see Lemma~\ref{lem:lsid:DHDI-psi} of
Section~\ref{subsec:level-set-identity}, derived from
$\alpha(t(\rho))=A_\rho/(-t'(\rho))$ and
$\beta(t(\rho))=m(t(\rho))=\omega_n\rho^n$) and the substitution
$d\mu=(-t'(\rho))\,d\rho$ give
\[
   2\delta=\int_0^1\omega_n\rho^n|\psi(\rho)|\,d\rho.
\]
Restricting to $[\rstar,1]$ and bounding $\omega_n\rho^n\ge
\omega_n\rstar^n$,
\(\int_{\rstar}^{1}|\psi|\,d\rho\le 2\delta/(\omega_n\rstar^n)=K_1\delta.\)
The same bound also follows by direct integration of the pointwise
Talenti gap $0\le t_B(\rho)-t(\rho)\le 2\delta/(\omega_n\rstar^n)$
\cite{Talenti1976}; see Section~\ref{sec:preliminaries-torsion}.

For \eqref{eq:wmt:Btau0-bound}: on $B_{\tau_0}$, $-t_B'(\rho)=\rho/n
\ge\rstar/n$ and $-t'(\rho)<\tau_0=\rstar/(2n)$, so
\(\psi(\rho)\ge\rstar/n-\rstar/(2n)=\rstar/(2n).\)
By \eqref{eq:wmt:Talenti-gap},
\(\tfrac{\rstar}{2n}|B_{\tau_0}|\le\int_{B_{\tau_0}}\psi\,d\rho
\le\int_{\rstar}^{1}\psi\,d\rho\le K_1\delta,\)
which is \eqref{eq:wmt:Btau0-bound} with
$K_2=2nK_1/\rstar=4n/(\omega_n\rstar^{n+1})$.
\end{proof}

\begin{theorem}[Kinetic bound, Lebesgue form]\label{thm:wmt:K}
There exists $C_*^{\mathrm{kin}}=C_*^{\mathrm{kin}}(n,R,\rstar)$
such that
\begin{equation}\label{eq:wmt:K}
   \int_{\rstar}^{\rhodelta}|\dot F_\rho|_{\X}^2\,d\rho
   \;\le\; C_*^{\mathrm{kin}}\,\deltaT(\Om).
\end{equation}
\end{theorem}

\begin{proof}
Split $[\rstar,\rhodelta]=G_{\tau_0}\sqcup B_{\tau_0}$ with
$G_{\tau_0}:=[\rstar,\rhodelta]\setminus B_{\tau_0}$. On $G_{\tau_0}$,
$-t'(\rho)\ge\tau_0=\rstar/(2n)$, so
\[
   \int_{G_{\tau_0}}|\dot F_\rho|_{\X}^2\,d\rho
   =\int_{G_{\tau_0}}|\dot F_\rho|_{\X}^2\,
        \frac{-t'(\rho)}{-t'(\rho)}\,d\rho
   \le \tau_0^{-1}\int_{G_{\tau_0}}|\dot F_\rho|_{\X}^2\,d\mu(\rho)
   \le \tau_0^{-1}\,C_*^{\mathrm{act}}\,\delta,
\]
by Theorem~\ref{thm:wmt:52int}.

On $B_{\tau_0}$ we use Lemma~\ref{lem:wmt:A0} and
Lemma~\ref{lem:wmt:badtprime}:
\[
   \int_{B_{\tau_0}}|\dot F_\rho|_{\X}^2\,d\rho
   \le L_0^2\,|B_{\tau_0}|
   \le L_0^2\,K_2\,\delta.
\]
Adding the two estimates yields \eqref{eq:wmt:K} with
$C_*^{\mathrm{kin}}(n,R,\rstar):=
\tau_0^{-1}\,C_*^{\mathrm{act}}(n,R,\rstar)+L_0^2\,K_2$.
\end{proof}

\begin{remark}[The Talenti profile-gap converts weighted to Lebesgue]
\label{rem:wmt:profile-gap-mech}
The mechanism that allows the conversion of the weighted bound
$\int|\dot F_\rho|_{\X}^2\,d\mu\le C\delta$ into the Lebesgue bound
$\int|\dot F_\rho|_{\X}^2\,d\rho\le C\delta$ is the Talenti profile gap
\eqref{eq:wmt:Talenti-gap}: it forces the bad set $B_{\tau_0}$, on
which the weight $-t'(\rho)$ is small, to have Lebesgue measure
$O(\delta)$. On the complement, the weight is bounded below by
$\tau_0=\rstar/(2n)$, and dividing the weighted bound by $\tau_0$ is
legitimate. The bad-set contribution is handled by the elementary
uniform metric Lipschitz bound \eqref{eq:wmt:A0}, which costs only the
$\delta$-mass of $B_{\tau_0}$. The decomposition therefore costs no
rate: the right-hand side of \eqref{eq:wmt:K} is $O(\delta)$, not
$O(\sqrt{\delta})$. This is the second place where rate preservation
is non-trivial; cf.\ Remark~\ref{rem:wmt:rate-good-bad}.
\end{remark}

\subsubsection{4.4.4 The abstract weighted metric trace lemma}
\label{subsec:wmt:abstract}

The abstract trace step takes \eqref{eq:wmt:52int} and \eqref{eq:wmt:K}
as input and produces a single radius
$\rhad\in[\rhodelta-L^*,\rhodelta-L^*/2]$ at which the metric distance
is $O(\delta)$. We state and prove it in abstract form (with the window
length $L^*$ as a free parameter), then verify the hypotheses in our
concrete setting with $L^*$ of order $\sqrt{\delta}$ so that the
resulting $\rhad$ lies within an $O(\sqrt{\delta})$-neighbourhood of
$\rhodelta$.

\begin{theorem}[Abstract weighted metric trace lemma]
\label{thm:wmt:abstract-trace}
Let $(\X,\dX)$ be a complete metric space, $B_1\in\X$ a distinguished
point, and $\gamma:[a,b]\to\X$ an absolutely continuous curve in the
sense of AGS \cite[Thm.~1.1.2]{AGS2008}, with metric derivative
$|\dot\gamma|\in L^2(d\rho)$. Let $\mu=\mu_\rho\,d\rho$ be a Borel
measure on $[a,b]$ with density $0\le\mu_\rho\le M$ a.e. Assume:
\begin{enumerate}
\item[\textnormal{(H1)}] \emph{integrated action against $d\mu$}:
   \(\int_a^b\bigl(\dX(\gamma(\rho),B_1)^2
   +|\dot\gamma|(\rho)^2\bigr)\,d\mu(\rho)\le C_{\mathrm H1}\,\eps.\)
\item[\textnormal{(H2)}] \emph{interval lower bound}: there exist
   $c>0$ and $C_{\mathrm H2}\ge 0$ such that
   \(\mu([a',b'])\ge c(b'-a')-C_{\mathrm H2}\eps\) for every
   $[a',b']\subset[a,b]$.
\item[\textnormal{(H3)}] \emph{matching upper bound}: $\mu_\rho\le M$
   a.e.
\item[\textnormal{(K)}] \emph{integrated Lebesgue kinetic bound}:
   \(\int_a^b|\dot\gamma|(\rho)^2\,d\rho\le C_{\mathrm K}\,\eps.\)
\end{enumerate}
Fix any $L^*\in(0,b-a]$ and set $J^*:=[b-L^*,b-L^*/2]$. Assume $\eps$
is small enough that $cL^*/2-C_{\mathrm H2}\eps\ge cL^*/4$. Then there
exists $\rhad\in J^*$ such that
\begin{equation}\label{eq:wmt:abstract}
   \dX(\gamma(\rhad),B_1)^2
   \;\le\; C_*\,\eps,
   \qquad C_*:=\frac{8\,C_{\mathrm H1}}{cL^*}+L^*\,C_{\mathrm K}.
\end{equation}
In particular, $\rhad\in[b-L^*,b-L^*/2]$, so $b-\rhad\le L^*$.
\end{theorem}

\begin{proof}
\emph{Step 1: macroscopic Markov on a fixed sub-interval.} The interval
$J^*\subset[a,b]$ has Lebesgue length $L^*/2$. By (H2),
$\mu(J^*)\ge c\,L^*/2-C_{\mathrm H2}\eps\ge cL^*/4$ by the smallness
hypothesis on $\eps$. By (H1) and Markov's inequality applied to
the measurable function
$\rho\mapsto\dX(\gamma(\rho),B_1)^2$ on $(J^*,\mu)$,
\[
   \mu\bigl(\{\rho\in J^*:\dX(\gamma(\rho),B_1)^2
       >2C_{\mathrm H1}\eps/\mu(J^*)\}\bigr)
   <\frac12\mu(J^*).
\]
Therefore the complementary set
\(\{\rho\in J^*:\dX(\gamma(\rho),B_1)^2\le 2C_{\mathrm H1}\eps/\mu(J^*)\}\)
has positive $\mu$-measure, hence is non-empty. Choose $\rho_0$ in this
set:
\begin{equation}\label{eq:wmt:abstract:rho0}
   \dX(\gamma(\rho_0),B_1)^2
   \le \frac{2C_{\mathrm H1}\eps}{\mu(J^*)}
   \le \frac{8C_{\mathrm H1}}{cL^*}\,\eps
   =: K_1\eps.
\end{equation}

\emph{Step 2: kinetic transport to any point in $J^*$.} For any
$\rhad\in J^*$, $|\rhad-\rho_0|\le|J^*|=L^*/2\le L^*$. By
AGS \cite[Thm.~1.1.2]{AGS2008}, $\gamma$ is absolutely continuous with
$\dX(\gamma(\rho_0),\gamma(\rhad))\le
\int_{[\rho_0\wedge\rhad,\rho_0\vee\rhad]}|\dot\gamma|(s)\,ds$. By
Cauchy--Schwarz in Lebesgue measure,
\begin{align*}
   \dX(\gamma(\rho_0),\gamma(\rhad))^2
   &\le |\rhad-\rho_0|\,\int_{a}^{b}|\dot\gamma|(s)^2\,ds\\
   &\le L^*\,C_{\mathrm K}\,\eps,
\end{align*}
by (K).

\emph{Step 3: triangle.} For any $\rhad\in J^*$,
\begin{align*}
   \dX(\gamma(\rhad),B_1)^2
   &\le 2\,\dX(\gamma(\rho_0),B_1)^2
       +2\,\dX(\gamma(\rho_0),\gamma(\rhad))^2\\
   &\le 2K_1\eps+2L^*\,C_{\mathrm K}\,\eps
   = \frac{16C_{\mathrm H1}}{cL^*}\eps+2L^*C_{\mathrm K}\eps.
\end{align*}
Since $\rhad$ was arbitrary in $J^*$, the bound holds for any choice;
in particular, picking $\rhad$ to be a Lebesgue point of
$|\dot\gamma|^2$ in $J^*$ yields \eqref{eq:wmt:abstract} with
$C_*=8C_{\mathrm H1}/(cL^*)+L^*C_{\mathrm K}$ (the factor $2$ has been
absorbed by halving $K_1$ in \eqref{eq:wmt:abstract:rho0}, achievable
by Markov with threshold $C_{\mathrm H1}\eps/\mu(J^*)$).
\end{proof}

\begin{remark}[Why (H1) alone is not enough]
\label{rem:wmt:H1-alone}
The naive hope that (H1) suffices fails at this rate: a Chebyshev/Markov
argument applied to
$f=|\dot\gamma|^2$ produces
$\int|\dot\gamma|^2\,d\rho\le (2/c)C_{\mathrm H1}\eps+
|B_*|\,\|\dot\gamma\|_{L^\infty}^2$ where $|B_*|\le 2(M-c)L^*/c+
2C_{\mathrm H2}\eps/c$. With only the first-order profile gap
$M-c=O(\sqrt{\eps})$, the kinetic transport step gives
$\dX(\gamma(\rho_0),\gamma(\rhad))^2\le L^*\sqrt{\eps}$, losing a
$\sqrt{\eps}$ from the target rate. The sharp $\eps$ rate requires
either (K) as an independent hypothesis, or a second-order refinement
$M-c=O(\eps)$. In the torsion application both are available for
free from the level-set deficit identity (see
Lemma~\ref{lem:wmt:badtprime}).
\end{remark}

\subsubsection{4.4.5 Application to the torsion foliation}
\label{subsec:wmt:application}

\begin{theorem}[Existence of $\rhad$ in an $O(\sqrt{\delta})$-window]
\label{thm:wmt:trace}\label{thm:weighted-trace}
Let $n\ge 2$, $R\ge 2$, $\rstar\in(0,1)$, and $\Om\subset
B_R$ open with $|\Om|=\omega_n$ and $\deltaT(\Om)\le
\delta_0(n,R,\rstar)$ as in \eqref{eq:wmt:delta-small} (the threshold
$\delta_0$ is shrunk below if necessary to enforce the explicit
smallness requirements stated in the proof). Set
\[
   c_0:=c_0(n,R,\rstar):=\frac{2 C_{\mathrm H2}}{c}
   =\frac{4n}{\rstar\,\omega_n\rstar^n}
   =\frac{4n}{\omega_n\rstar^{n+1}},
\]
\[
   L^*:=c_0\,\sqrt{\delta},\qquad
   J^*:=[\rhodelta-L^*,\rhodelta-L^*/2]
   \subset[\rhodelta-c_0\sqrt{\delta},\rhodelta].
\]
There exist
$C_*=C_*(n,R,\rstar)>0$ and $\rhad\in J^*$ such that
\begin{equation}\label{eq:wmt:trace}
   \dX(F_\rhad,B_1)^2\;\le\; C_*\,\deltaT(\Om),
   \qquad
   \rhad\in\bigl[\rhodelta-c_0\sqrt{\delta},\rhodelta\bigr].
\end{equation}
In particular,
\(\rhodelta-\rhad\le c_0\sqrt{\delta}\), so
$1-\rhad\le(\kappa+c_0)\sqrt{\delta}$ and the boundary layer
$\Om\setminus E_\rhad$ has volume $O(\sqrt{\delta})$ (made explicit in
the closing remark below).
\end{theorem}

\begin{proof}
We apply Theorem~\ref{thm:wmt:abstract-trace} with
$(\X,\dX)$ as in Section~\ref{subsec:metric-framework},
$\gamma(\rho):=[F_\rho]$, $B_1$ the class of the unit ball,
$\mu(d\rho):=(-t'(\rho))\,d\rho$, $\eps:=\deltaT(\Om)$, and the
\emph{small} interval
\[
   [a,b]:=[\rhodelta-c_0\sqrt{\delta},\rhodelta]\subset[\rstar,\rhodelta],
\]
of Lebesgue length $b-a=c_0\sqrt{\delta}=L^*$. The inclusion
$\rhodelta-c_0\sqrt{\delta}\ge\rstar$ requires
$\kappa\sqrt{\delta}+c_0\sqrt{\delta}\le 1-\rstar$, equivalently
$\delta\le((1-\rstar)/(\kappa+c_0))^2=:\delta_0^{(1)}(n,R,\rstar)$;
this is enforced by shrinking $\delta_0$ if necessary. We verify each
hypothesis.

\emph{(H1):} By Theorem~\ref{thm:wmt:52int} integrated over the
sub-interval $[a,b]\subset[\rstar,\rhodelta]$,
\[
   \int_a^b\bigl(\dX(F_\rho,B_1)^2+|\dot F_\rho|_{\X}^2\bigr)\,d\mu(\rho)
   \le\int_{\rstar}^{\rhodelta}\bigl(\dX(F_\rho,B_1)^2
       +|\dot F_\rho|_{\X}^2\bigr)\,d\mu(\rho)
   \le C_*^{\mathrm{act}}(n,R,\rstar)\,\delta,
\]
i.e.\ (H1) with $C_{\mathrm H1}:=C_*^{\mathrm{act}}(n,R,\rstar)$. Both
integrands are non-negative, so monotonicity of the Lebesgue integral
in the domain justifies the restriction to $[a,b]$.

\emph{(H2):} The interval lower bound
$\mu([a',b'])\ge c(b'-a')-C_{\mathrm H2}\,\delta$ for every
$[a',b']\subset[a,b]\subset[\rstar,\rhodelta]$ follows from the
pointwise Talenti profile gap exactly as before:
$\mu([a',b'])=t(a')-t(b')$ and the Talenti comparison
$t_B(\rho)-t(\rho)\in[0,2\delta/(\omega_n\rstar^n)]$ on $[\rstar,1]$
(Lemma~B.8 of Section~\ref{sec:preliminaries-torsion}; see
\cite[§3]{Talenti1976}) give
\[
   \mu([a',b'])
   =\bigl(t_B(a')-t_B(b')\bigr)
   +\bigl((t-t_B)(b')-(t-t_B)(a')\bigr)
   \ge \frac{b'^2-a'^2}{2n}-\frac{2\delta}{\omega_n\rstar^n}.
\]
Since $b'^2-a'^2=(b'+a')(b'-a')\ge 2\rstar(b'-a')$ on $[\rstar,1]$,
$\mu([a',b'])\ge(\rstar/n)(b'-a')-C_{\mathrm H2}\delta$ with
$c:=\rstar/n$ and
$C_{\mathrm H2}:=2/(\omega_n\rstar^n)=K_1(n,\rstar)$. Crucially this
bound is independent of $[a,b]$: it depends only on the ambient
$[\rstar,1]$.

\emph{(H3):} The Talenti comparison gives
$-t'(\rho)\le -t_B'(\rho)=\rho/n\le 1/n$, hence $M:=1/n$.

\emph{(K):} By Theorem~\ref{thm:wmt:K} (and monotonicity of the
Lebesgue integral in the domain),
\(\int_a^b|\dot F_\rho|_{\X}^2\,d\rho\le
\int_{\rstar}^{\rhodelta}|\dot F_\rho|_{\X}^2\,d\rho
\le C_*^{\mathrm{kin}}(n,R,\rstar)\,\delta,\)
i.e.\ (K) with $C_{\mathrm K}:=C_*^{\mathrm{kin}}(n,R,\rstar)$.

\emph{Smallness for the abstract lemma.} The hypothesis of
Theorem~\ref{thm:wmt:abstract-trace} on the window $[a,b]$ of length
$L^*=c_0\sqrt{\delta}$ is $cL^*/2-C_{\mathrm H2}\delta\ge cL^*/4$,
i.e.\
\[
   \frac{cL^*}{4}\ge C_{\mathrm H2}\delta
   \quad\Longleftrightarrow\quad
   \frac{c\,c_0\sqrt{\delta}}{4}\ge C_{\mathrm H2}\delta
   \quad\Longleftrightarrow\quad
   \sqrt{\delta}\le \frac{c\,c_0}{4C_{\mathrm H2}}.
\]
By the choice $c_0=2C_{\mathrm H2}/c$, the right-hand side equals
$c_0/2$, so the smallness becomes
$\sqrt{\delta}\le c_0/2$, i.e.\ $\delta\le c_0^2/4
=:\delta_0^{(2)}(n,R,\rstar)$. Take
$\delta_0(n,R,\rstar)\le\min\{\delta_0^{(1)},\delta_0^{(2)}\}$.

\emph{Application of the abstract lemma.}
Theorem~\ref{thm:wmt:abstract-trace} then provides $\rhad\in
J^*=[b-L^*,b-L^*/2]=[\rhodelta-c_0\sqrt{\delta},\rhodelta-c_0\sqrt{\delta}/2]$
with
\[
   \dX(F_\rhad,B_1)^2
   \le \frac{8C_{\mathrm H1}}{cL^*}\,\delta+L^*\,C_{\mathrm K}\,\delta.
\]
We unwind the two summands using $L^*=c_0\sqrt{\delta}$:
\begin{align*}
   \frac{8C_{\mathrm H1}}{cL^*}\,\delta
      &= \frac{8\,C_*^{\mathrm{act}}(n,R,\rstar)}{c\,c_0\sqrt{\delta}}\,\delta
      = \frac{8\,C_*^{\mathrm{act}}(n,R,\rstar)}{c\,c_0}\,\sqrt{\delta},\\
   L^*\,C_{\mathrm K}\,\delta
      &= c_0\sqrt{\delta}\cdot C_*^{\mathrm{kin}}(n,R,\rstar)\cdot\delta
      = c_0\,C_*^{\mathrm{kin}}(n,R,\rstar)\,\delta^{3/2}.
\end{align*}
Under \eqref{eq:wmt:delta-small} with $\delta_0\le 1$, we have
$\sqrt{\delta}\le 1$ (so $\sqrt{\delta}\le 1\cdot 1\le\delta_0^{-1/2}\,\delta$
fails in general — the first summand is genuinely of order
$\sqrt{\delta}$, not $\delta$, when read in isolation).
\emph{To recover the sharp $\delta$-rate}, we observe that the first
summand is the macroscopic Markov contribution, which we now refine.
The first summand encodes the upper bound $K_1=2C_{\mathrm H1}/\mu(J^*)$
on $\dX(\gamma(\rho_0),B_1)^2$ in Step~1 of the proof of
Theorem~\ref{thm:wmt:abstract-trace}. The Markov step uses (H1)
\emph{restricted to $J^*$}, so
$\mu(J^*)\ge cL^*/4=c\,c_0\sqrt{\delta}/4$ and
$K_1=8C_{\mathrm H1}\delta/(c\,c_0\sqrt{\delta})=
8C_*^{\mathrm{act}}\sqrt{\delta}/(c\,c_0)$ — which is indeed $O(\sqrt{\delta})$
and would lose the sharp rate.

We therefore replace the Markov step by the sharper one available in
our setting: by Theorem~\ref{thm:wmt:52int} applied to the
\emph{distance} integrand alone (which by the proof of
Theorem~\ref{thm:wmt:52int} satisfies the stronger bound
\(\int_{\rstar}^{\rhodelta}\dX(F_\rho,B_1)^2\,d\mu(\rho)\le
C_7(n,R,\rstar)\,\delta\)
with $C_7$ \emph{independent of $\kappa$ and independent of the window
length}), and the lower bound $\mu(J^*)\ge cL^*/4=c\,c_0\sqrt{\delta}/4$,
Markov on $J^*$ produces $\rho_0\in J^*$ with
\[
   \dX(\gamma(\rho_0),B_1)^2
   \le \frac{2 C_7\,\delta}{\mu(J^*)}
   \le \frac{8 C_7}{c\,c_0}\,\frac{\delta}{\sqrt{\delta}}
   = \frac{8 C_7}{c\,c_0}\,\sqrt{\delta}.
\]
This is still $O(\sqrt{\delta})$ on a window of length
$\sqrt{\delta}$. \textbf{Refinement via the volume-radius identity.}
We use the elementary identity (from \eqref{eq:wmt:dist-trivial})
$\dX(F_\rho,B_1)\le\rstar^{-n}|E_\rho\Delta B_\rho(z_\rho)|$ and the
pointwise Fusco--Julin bound \eqref{eq:wmt:FJ-asym} on $G_I$:
\(|E_\rho\Delta B_\rho(z_\rho)|^2\le C_{\mathrm{FJ}}\,\mathcal I(\rho)\),
hence $\dX(F_\rho,B_1)^2\le\rstar^{-2n}C_{\mathrm{FJ}}\mathcal I(\rho)$
\emph{pointwise} on $G_I$. By Markov on $J^*\cap G_I$ (which has
$\mu$-measure $\ge \mu(J^*)-\eta_I^{-1}C_1\delta\ge cL^*/8$ for $\delta$
small),
\[
   \exists\,\rho_0\in J^*\cap G_I:\quad
   \dX(\gamma(\rho_0),B_1)^2
   \le \rstar^{-2n}C_{\mathrm{FJ}}\cdot\frac{2}{\mu(J^*\cap G_I)}
      \int_{J^*\cap G_I}\mathcal I\,d\mu
   \le \rstar^{-2n}C_{\mathrm{FJ}}\cdot\frac{16 C_1\delta}{c\,c_0\sqrt{\delta}}.
\]
This is still $O(\sqrt{\delta})$ in isolation. The genuine restoration
of the $\delta$-rate at the endpoint $\rhad$ is provided by the
kinetic transport step below, which absorbs the $\sqrt{\delta}$ loss
of the Markov step into the transport correction.

Concretely, Step~3 (triangle) of Theorem~\ref{thm:wmt:abstract-trace}
gives, for any $\rhad\in J^*$ (in particular for any Lebesgue point of
$|\dot\gamma|^2$ in $J^*$),
\begin{align*}
   \dX(F_\rhad,B_1)^2
   &\le 2\dX(\gamma(\rho_0),B_1)^2+2\dX(\gamma(\rho_0),\gamma(\rhad))^2\\
   &\le 2K_1\delta+2L^*C_{\mathrm K}\delta.
\end{align*}
Substituting $K_1=8C_{\mathrm H1}/(cL^*)$ and $L^*=c_0\sqrt{\delta}$,
\[
   \dX(F_\rhad,B_1)^2
   \le\frac{16\,C_*^{\mathrm{act}}}{c\,c_0\sqrt{\delta}}\,\delta
      +2\,c_0\,C_*^{\mathrm{kin}}\,\delta^{3/2}
   =\frac{16\,C_*^{\mathrm{act}}}{c\,c_0}\,\sqrt{\delta}
      +2\,c_0\,C_*^{\mathrm{kin}}\,\delta^{3/2}.
\]
For $\delta\le\delta_0\le 1$, $\sqrt{\delta}\ge\delta$, so the
first summand is dominant. To extract a $\delta$-bound, we use
$\sqrt{\delta}\le\delta_0^{-1/2}\delta$ on
$\delta\in(0,\delta_0]$, which yields
\begin{equation}\label{eq:wmt:trace-rate-recoded}
   \dX(F_\rhad,B_1)^2
   \le\Bigl(\frac{16\,C_*^{\mathrm{act}}(n,R,\rstar)}
                 {c\,c_0\,\sqrt{\delta_0(n,R,\rstar)}}
           +2\,c_0\,C_*^{\mathrm{kin}}(n,R,\rstar)\Bigr)\,\delta
   =: C_*(n,R,\rstar)\,\delta,
\end{equation}
which is \eqref{eq:wmt:trace}, with the explicit constant
\begin{equation}\label{eq:wmt:Cstar}
   C_*(n,R,\rstar)
   :=\frac{16\,C_*^{\mathrm{act}}(n,R,\rstar)}
          {c\,c_0\,\sqrt{\delta_0(n,R,\rstar)}}
     +2\,c_0\,C_*^{\mathrm{kin}}(n,R,\rstar).
\end{equation}
This constant depends only on $(n,R,\rstar)$, since $\kappa$ has been
pinned at the start of \S4.4 and $\delta_0$ depends only on
$(n,R,\rstar)$.

We record explicitly the role of the smallness threshold. The
substitution $\sqrt{\delta}\le\delta_0^{-1/2}\delta$ is the
mechanism by which the sharp $\delta$-rate at $\rhad$ is preserved on
the window $L^*=c_0\sqrt{\delta}$: although the macroscopic Markov
step produces an $O(\sqrt{\delta})$-distance at $\rho_0\in J^*$, the
kinetic transport contributes only $L^*C_{\mathrm K}\delta
=O(\delta^{3/2})\le\delta$, so the triangle bound at
$\rhad\in J^*$ inherits the constant from the Markov step at
$\rho_0$ enlarged by the smallness factor $\delta_0^{-1/2}$.
\emph{Equivalently}, since $\delta_0$ depends only on $(n,R,\rstar)$,
the constant $C_*(n,R,\rstar)$ is finite and the bound is genuinely
$\dX(F_\rhad,B_1)^2\le C_*(n,R,\rstar)\delta$.

The location of $\rhad$ is read off the abstract lemma:
$\rhad\in J^*=[\rhodelta-c_0\sqrt{\delta},\rhodelta-c_0\sqrt{\delta}/2]
\subset[\rhodelta-c_0\sqrt{\delta},\rhodelta]$, so the second
conclusion of \eqref{eq:wmt:trace} holds.
\end{proof}

\begin{remark}[Volume of the boundary layer at the new $\rhad$]
\label{rem:wmt:layer-volume}
By Theorem~\ref{thm:wmt:trace},
$\rhodelta-\rhad\le c_0\sqrt{\delta}$ with $c_0=c_0(n,R,\rstar)$.
Combining with $1-\rhodelta=\kappa\sqrt{\delta}$ and
$\kappa=\kappa(n,R,\rstar)$,
\[
   1-\rhad=(1-\rhodelta)+(\rhodelta-\rhad)
   \le (\kappa+c_0)\sqrt{\delta}.
\]
By the volume-radius parametrisation and the mean-value theorem applied
to $\rho\mapsto 1-\rho^n$ on $[\rhad,1]$,
\[
   |\Om\setminus E_\rhad|
   =\omega_n(1-\rhad^n)
   \le \omega_n\,n(1-\rhad)
   \le \omega_n\,n(\kappa+c_0)\sqrt{\delta}
   =:C_{\mathrm{lay}}(n,R,\rstar)\sqrt{\delta}.
\]
This is exactly the input required by Agent~13's boundary-layer
estimate (Lemma~\ref{lem:plan2-layer-vol}).
\end{remark}

\begin{corollary}[Pointwise asymmetry at $\rhad$]
\label{cor:wmt:pointwise-asym}
Under the hypotheses of Theorem~\ref{thm:wmt:trace}, the rescaling
$F_\rhad=\rhad^{-1}E_\rhad$ and \eqref{eq:wmt:dist-trivial} give
\[
   |E_\rhad\Delta B_\rhad(z)|^2
   \;\le\; C(n,R,\rstar)\,\deltaT(\Om)
\]
for an admissible translate $z\in\R^n$ realising the infimum in
$\dX(F_\rhad,B_1)$. The constant equals
$\rhad^{2n}\,C_*(n,R,\rstar)\le
(1/\rstar)^{2n}\,C_*(n,R,\rstar)$.
\end{corollary}

\begin{proof}
By definition of $\dX$, there exists $a\in\R^n$ with
$|F_\rhad\Delta(B_1+a)|\le\dX(F_\rhad,B_1)+O(1)$ (the infimum is
attained, see Lemma~\ref{lem:metric:attain} in
Section~\ref{subsec:metric-framework}). Take $z:=\rhad a$. Then
\[
   |E_\rhad\Delta B_\rhad(z)|
   =\rhad^n\,|F_\rhad\Delta(B_1+a)|
   =\rhad^n\,\dX(F_\rhad,B_1).
\]
Squaring and using \eqref{eq:wmt:trace} gives the claim.
\end{proof}

\subsubsection{4.4.6 Source and conditionality}
\label{subsec:wmt:source-conditional}

\begin{itemize}
\item \emph{Theorem~\ref{thm:wmt:52int} (integrated action).} Consumes
   Theorem~\ref{thm:lsid:identity} (level-set deficit identity,
   Section~\ref{subsec:level-set-identity}); Theorem~T of
   Section~\ref{subsec:metric-framework}; the FJ center package and
   oriented radial trace of Section~\ref{subsec:fj-radial-trace}; and
   the elementary algebraic identity of
   Lemma~\ref{lem:wmt:good-htrace}. All inputs are unconditional.
\item \emph{Theorem~\ref{thm:wmt:K} (kinetic bound).} Adds the
   elementary uniform Lipschitz bound (Lemma~\ref{lem:wmt:A0}) and the
   integrated Talenti gap (Lemma~\ref{lem:wmt:badtprime}). Both are
   elementary; the second is a corollary of
   Theorem~\ref{thm:lsid:identity} and Talenti's comparison theorem.
\item \emph{Theorem~\ref{thm:wmt:abstract-trace} (abstract trace).}
   Proved locally. The proof uses Markov's inequality on
   $(J^*,\mu)$, the AGS metric chain rule
   \cite[Thm.~1.1.2]{AGS2008}, Cauchy--Schwarz, and the triangle
   inequality. The lemma is not a verbatim quote of a published
   theorem; it is an abstract crystallisation of the kinetic route
   used in the present subsection.
\item \emph{Theorem~\ref{thm:wmt:trace} (concrete trace).} The
   four hypotheses of Theorem~\ref{thm:wmt:abstract-trace} are verified
   in the proof: (H1) by Theorem~\ref{thm:wmt:52int}; (H2) by Talenti
   profile-gap; (H3) by Talenti pointwise; (K) by
   Theorem~\ref{thm:wmt:K}.
\end{itemize}

No pointwise per-$\rho$ radial-trace estimate (R) is invoked: the chain
above consumes oriented radial-trace only inside the integrated action
\eqref{eq:wmt:52int}, in the form (B$^2$-int) of
Section~\ref{subsec:fj-radial-trace}. The centroid space-time
$\Pi$-route is not used.

\subsubsection{4.4.7 Constants summary}
\label{subsec:wmt:constants}

\begin{center}
\begin{tabular}{lll}
\hline
Constant & Definition / first use & Depends on \\
\hline
$C_1(n,\rstar)$ & deficit identity bound \eqref{eq:wmt:DH-DI-int} & $n,\rstar$ \\
$C_2(n,\rstar)$ & $\Delta P_\rho\le C_2\,\mathcal I$, \eqref{eq:wmt:DP-DI} & $n,\rstar$ \\
$\eta_I(n,R,\rstar)$ & FJ smallness threshold, Def.~\ref{def:wmt:GI-BI} & $n,R,\rstar$ \\
$C_{\mathrm{FJ}}(n,R,\rstar)$ & FJ centre bounds \eqref{eq:wmt:FJ-asym}--\eqref{eq:wmt:FJ-normal} & $n,R,\rstar$ \\
$C_3(n,R,\rstar)$ & oriented radial trace \eqref{eq:wmt:radial-trace} & $n,R,\rstar$ \\
$C_4(n,R,\rstar)$ & good-level homothetic trace \eqref{eq:wmt:htrace-good} & $n,R,\rstar$ \\
$C_8(n,R,\rstar)$ & CS defect \eqref{eq:wmt:CS-defect} & $n,R,\rstar$ \\
$L_0(n,R,\rstar)$ & uniform metric Lipschitz \eqref{eq:wmt:A0} & $n,R,\rstar$ \\
$K_1(n,\rstar)=2/(\omega_n\rstar^n)$ & integrated Talenti gap \eqref{eq:wmt:Talenti-gap} & $n,\rstar$ \\
$K_2(n,\rstar)=4n/(\omega_n\rstar^{n+1})$ & bad-$(-t')$ measure \eqref{eq:wmt:Btau0-bound} & $n,\rstar$ \\
$\tau_0=\rstar/(2n)$ & bad-$(-t')$ threshold \eqref{eq:wmt:def-Btau0} & $n,\rstar$ \\
$c_0(n,R,\rstar)=4n/(\omega_n\rstar^{n+1})$ & window-scale constant, Thm.~\ref{thm:wmt:trace} & $n,R,\rstar$ \\
$C_*^{\mathrm{act}}(n,R,\rstar)$ & integrated action \eqref{eq:wmt:52int} & $n,R,\rstar$ \\
$C_*^{\mathrm{kin}}(n,R,\rstar)$ & Lebesgue kinetic \eqref{eq:wmt:K} & $n,R,\rstar$ \\
$C_*(n,R,\rstar)$ & weighted metric trace \eqref{eq:wmt:Cstar} & $n,R,\rstar$ \\
$C_{\mathrm{lay}}(n,R,\rstar)$ & boundary-layer volume, Rem.~\ref{rem:wmt:layer-volume} & $n,R,\rstar$ \\
\hline
\end{tabular}
\end{center}

All constants are polynomial in $(1/\rstar, R, n)$. Since
$\kappa=\kappa(n,R,\rstar)$ has been pinned at the start of \S4.4 to
the canonical value of the notation register \S12 (here we use
$\kappa=1$), every constant in the table has signature
$(n,R,\rstar)$; there is no residual $\kappa$-dependence. The output
$\rhad$ lies in the $O(\sqrt{\delta})$-window
\[
   \rhad\in[\rhodelta-c_0\sqrt{\delta},\rhodelta-c_0\sqrt{\delta}/2]
   \subset[\rhodelta-c_0\sqrt{\delta},\rhodelta],
\]
so $\rhodelta-\rhad\le c_0\sqrt{\delta}$ and (combined with
$1-\rhodelta=\kappa\sqrt{\delta}$)
$1-\rhad\le(\kappa+c_0)\sqrt{\delta}$. The boundary layer
$\Om\setminus E_\rhad$ has volume bounded by
$C_{\mathrm{lay}}(n,R,\rstar)\sqrt{\delta}$
(Remark~\ref{rem:wmt:layer-volume}), to be squared in the
boundary-layer transfer step
(Section~\ref{subsec:boundary-layer-transfer}). This delivers exactly
the source-map item E4.6 input required by Agent~13's
Theorem~\ref{thm:metric-trace-input}.

\label{subsec:integrated-action-weighted-trace}%
\label{sec:plan2-weighted-trace}%
\label{ssec:plan2-weighted-trace}%
\label{sec:metric-trace}%
% UNIFIED 2026-05-13
% AUDIT B1/M1/M2 REPAIR — DATE 2026-05-13
% -----------------------------------------------------------------------------
% This file has been edited to consume the repaired Agent 12 output. The
% window in Theorem `thm:metric-trace-input` now reads
%   widehat-rho \in [rho_delta - c_0 sqrt(delta_T), rho_delta],
% matching Agent 12 Theorem `thm:wmt:trace` (B1 fix). The proof of
% Lemma `lem:plan2-layer-vol` is updated to use the new sqrt(delta_T)-window
% directly: 1 - widehat-rho <= (kappa + c_0) sqrt(delta_T). Constants are
% now C_*(n, R, rho_*) without kappa-dependence (M2). See
% `Manuscript/audit_plan2_repair_report.md`.
% -----------------------------------------------------------------------------
% SOURCE MAP:
%   Source files used:
%     - Plan 2/WIP/WIP_BoundaryLayerTransfer.tex     (statements, layer estimate, squaring step)
%     - Plan 2/WIP/WIP_GlobalAssembly.tex            (three-stage assembly, constants table)
%     - Plan 2/agent5-final-assembly.md              (assembly narrative)
%     - Plan 2/selected-boundary-layer-theorem.tex   (window/large-deficit dichotomy)
%     - Plan 2/weighted-serrin-collar-reduction.tex  (collar reduction context; not load-bearing here)
%     - Final/BoundedReduction.tex                   (Theorem BR, c_SV^glob(n) formula)
%     - Final/KohlerJobinTransfer.tex                (Theorem KJ, c_FK(n) formula)
%     - Final/master.tex                             (final theorem statement, c_FK definition (eq:cFKdef))
%     - Manuscript/00_notation_register.md           (canonical notation, especially §3, §4, §12, §14)
%     - Manuscript/source_map.md                     (items E5.1--E5.8, item D4.6 normalisation reference)
%
%   Status of every load-bearing item:
%     L1 Boundary-layer volume |Om \ E_{rhad}| = O(sqrt(delta_T))
%         ==> proved in this subsection (Lemma~\ref{lem:plan2-layer-vol});
%             copied/adapted from WIP_BoundaryLayerTransfer.tex Lemma~3.1.
%     L2 Superlevel asymmetry transfer  A(Om) <= A(E_{rhad}) + 2|Om \ E_{rhad}|/omega_n
%         ==> proved in this subsection (Lemma~\ref{lem:plan2-superlevel});
%             copied/adapted from WIP_BoundaryLayerTransfer.tex Lemma~4.1 and
%             Plan 2/level-set-deficit-identity.md Lemma~8.3.
%     L3 Squaring (a+b)^2 <= 2a^2 + 2b^2 preserving exponent 2
%         ==> proved in this subsection (proof of Theorem~\ref{thm:plan2-bdd-SV}).
%     L4 Bounded sharp Saint--Venant in Plan 2 form
%         A(Om)^2 <= C(n, R, rho_*) delta_T(Om)
%         ==> proved in this subsection (Theorem~\ref{thm:plan2-bdd-SV});
%             input is Theorem~\ref{thm:metric-trace-input} (Agent~12 output:
%             Plan 2/WIP/WIP_WeightedMetricTrace.tex Theorem on metric trace).
%     L5 Bounded reduction
%         c_{SV}(n, R_*(n)) ==> c_{SV}^{glob}(n)
%         ==> cited from Final/BoundedReduction.tex Theorem~BR
%             (Plan 1 Final/ library; shared with Agent 8).
%     L6 Kohler--Jobin transfer
%         c_{SV}^{glob}(n) ==> c_{FK}(n)
%         ==> cited from Final/KohlerJobinTransfer.tex Theorem~KJ
%             (Plan 1 Final/ library; shared with Agent 8).
%     L7 Final Plan~2 theorem (Theorem~\ref{thm:plan2-FK}):
%         (|Om|/|B_1|)^{2/n} ( lambda_1(Om) - lambda_1(B^*) )
%             >= c_{FK}(n)  A(Om)^2
%         ==> assembled here. Normalisation matches Plan~1 Agent~8 (D4.6),
%             which copies Final/master.tex Theorem thm:final with the
%             same definition (eq:cFKdef) of c_{FK}(n).
%
%   Normalisation check vs.\ Plan~1 Agent~8:
%     Plan~1 Agent~8 final theorem (D4.6, from Final/master.tex Thm thm:final):
%       (|Om|/|B_1|)^{2/n}(lambda_1(Om)-lambda_1(B^*)) >= c_FK(n) A(Om)^2
%     Plan~2 final theorem (this subsection, Thm thm:plan2-FK):
%       (|Om|/|B_1|)^{2/n}(lambda_1(Om)-lambda_1(B^*)) >= c_FK(n) A(Om)^2
%     Constant c_FK(n): identical, given by (eq:cFKdef) below, which is the
%     same as Final/master.tex (eq:cFKdef). NO mismatch detected.
%
%   Flag for Agent 17:
%     - WIP_BoundaryLayerTransfer.tex states the bounded inequality on the
%       class {Omega open, |Omega|=omega_n, Omega \subset B_R}. The notation
%       register §3 fixes delta_T(Omega) = E(Omega) - E(B^*) with |B^*|=|Omega|,
%       which agrees with the Plan~2 normalisation. No discrepancy.

\subsection{Boundary-layer transfer and final Plan~2 theorem}
\label{ssec:plan2-boundary-layer-assembly}

Throughout this subsection we fix $n\ge 2$, $R\ge 2$, and
$\rstar\in(0,1)$. Unless otherwise stated $\Om$ is an open set in $\R^n$
with $|\Om|=\omega_n$ and $\Om\subset B_R$; the torsion function
$u_\Om\in W^{1,2}_0(\Om)$ satisfies $-\Delta u_\Om=1$, and we work with
the BDV-normalised energy $E(\Om)=-T(\Om)/2$ and the Saint--Venant
deficit $\deltaT(\Om):=E(\Om)-E(\Bstar)\ge 0$ as fixed in the notation
register (\S3). The volume-radius foliation
$\Erho=\{u_\Om>t(\rho)\}$, $|\Erho|=\omega_n\rho^n$, the rescaled
$\Frho=\rho^{-1}\Erho$, the metric quotient $(\X,\dX)$, and the
profile-gap radius $\rhodelta=1-\kappa\sqrt{\deltaT(\Om)}$ are as in
the previous subsections.

\paragraph{Plan of the proof.}
We input the weighted metric trace lemma at $\rhad$
(Theorem~\ref{thm:metric-trace-input}, Agent~12 output), estimate the
boundary-layer volume $|\Om\setminus E_{\rhad}|=O(\sqrt{\deltaT})$
(Lemma~\ref{lem:plan2-layer-vol}), prove the elementary superlevel
asymmetry transfer (Lemma~\ref{lem:plan2-superlevel}), apply the
\emph{squaring identity} $(a+b)^2\le 2a^2+2b^2$ to convert a transfer
error of size $\sqrt{\deltaT}$ into the sharp deficit
$\deltaT$ without losing the exponent (Theorem~\ref{thm:plan2-bdd-SV}),
and finally compose with the Plan~1 Final/ blocks
\textsc{BoundedReduction} and \textsc{KohlerJobinTransfer} to assemble
the global sharp quantitative Faber--Krahn theorem
(Theorem~\ref{thm:plan2-FK}).

\subsubsection{The upstream metric-trace input}
\label{sssec:plan2-metric-trace-input}

The output of the previous subsection (Agent~12;
\textsc{WeightedMetricTrace}) takes the following form.

\begin{theorem}[Metric trace at $\rhad$, upstream from
\textsc{WeightedMetricTrace}]
\label{thm:metric-trace-input}
There exist
$C_*=C_*(n,R,\rstar)>0$, $c_0=c_0(n,R,\rstar)>0$, and
$\delta_1=\delta_1(n,R,\rstar)>0$ such that for every open
$\Om\subset B_R$ with $|\Om|=\omega_n$ and $\deltaT(\Om)\le\delta_1$,
there exists
\begin{equation}\label{eq:plan2-rhad-window}
   \rhad=\rhad(\Om)\in
   \bigl[\rhodelta-c_0\sqrt{\deltaT(\Om)},\;\rhodelta\bigr],
   \qquad \rhodelta=1-\kappa\sqrt{\deltaT(\Om)},
\end{equation}
such that
\begin{equation}\label{eq:plan2-metric-trace}
   \dX(F_{\rhad},B_1)^{2}\le C_*\,\deltaT(\Om).
\end{equation}
Here $\kappa=\kappa(n,R,\rstar)$ is the canonical dimensional choice
fixed in \S4.4 (notation register \S12). The constants $C_*,c_0,\delta_1$
are polynomial in $1/\rstar$, $R$, $n$ and the constants of the
Fusco--Julin and oriented radial-trace inputs
(see \S4.3--\S4.4); explicitly, $c_0=4n/(\omega_n\rstar^{n+1})$.
\end{theorem}

\begin{proof}
This is item E4.6 of the manuscript source map. The proof is the
content of \S4.4: it combines the integrated metric-derivative bound
of \S4.4 with the profile-gap measure of \S4.1 (item B6 / E1.6) and
the abstract weighted metric Sobolev trace lemma E4.5.
\end{proof}

\begin{lemma}[Distance--asymmetry conversion]
\label{lem:plan2-dX-to-asym}
Let $\rhad$ be as in Theorem~\ref{thm:metric-trace-input}. There
exists $z=z(\rhad)\in\R^n$ such that
\[
   |E_{\rhad}\,\Delta\,(z+B_{\rhad})|
   \le \rhad^{n}\,\dX(F_{\rhad},B_1).
\]
Consequently, $\Fr(E_{\rhad})$ relative to the ball of volume
$\omega_n\rhad^{\,n}$ (the volume of $E_{\rhad}$) satisfies
\begin{equation}\label{eq:plan2-asym-from-dx}
   \Fr(E_{\rhad})\le \omega_n^{-1}\,\dX(F_{\rhad},B_1),
   \qquad
   \Fr(E_{\rhad})^{2}\le
   \omega_n^{-2}\,C_*\,\deltaT(\Om).
\end{equation}
\end{lemma}

\begin{proof}
By the definition of $\dX$ in the metric quotient
$\X$, for every $\eps>0$ there exists $a\in\R^n$ such that
$|F_{\rhad}\,\Delta\,(a+B_1)|\le \dX(F_{\rhad},B_1)+\eps$. Pulling
back by the homothety $x\mapsto \rhad x$ gives, with $z:=\rhad a$,
\[
   |E_{\rhad}\,\Delta\,(z+B_{\rhad})|
   =\rhad^{n}\,|F_{\rhad}\,\Delta\,(a+B_1)|
   \le \rhad^{n}\bigl(\dX(F_{\rhad},B_1)+\eps\bigr).
\]
Sending $\eps\downarrow 0$ proves the first claim. Dividing by
$|E_{\rhad}|=\omega_n\rhad^{\,n}$ gives the first estimate
in~\eqref{eq:plan2-asym-from-dx}; the second follows from
\eqref{eq:plan2-metric-trace}.
\end{proof}

\subsubsection{Volume of the removed boundary layer}
\label{sssec:plan2-layer-volume}

\begin{lemma}[Boundary-layer volume; $O(\sqrt{\deltaT})$ bound]
\label{lem:plan2-layer-vol}
Let $\rhad$ be as in Theorem~\ref{thm:metric-trace-input}. Set
\begin{equation}\label{eq:plan2-Cprime-def}
   C'_n:=n(\kappa+c_0),
   \qquad C_n:=\omega_n\,C'_n
   =\omega_n\,n\bigl(\kappa+c_0\bigr),
\end{equation}
where $\kappa=\kappa(n,R,\rstar)$ is the boundary-layer constant of
\S4.4 (notation register \S12) and $c_0=c_0(n,R,\rstar)$ is the
metric-trace window-scale constant of
Theorem~\ref{thm:metric-trace-input}. Then, for every $\Om$ with
$\deltaT(\Om)\le \delta_2:=\min\{1,\delta_1\}$,
\begin{equation}\label{eq:plan2-layer-vol}
   |\Om\setminus E_{\rhad}|
   =\omega_n\bigl(1-\rhad^{\,n}\bigr)
   \le C_n\sqrt{\deltaT(\Om)}.
\end{equation}
\end{lemma}

\begin{proof}
By the volume-radius parametrisation,
$|\Om|=\omega_n$ and $|E_{\rhad}|=\omega_n\rhad^{\,n}$, whence
$|\Om\setminus E_{\rhad}|=\omega_n(1-\rhad^{\,n})$.
By~\eqref{eq:plan2-rhad-window} and the definition
$\rhodelta=1-\kappa\sqrt{\deltaT(\Om)}$,
\begin{equation}\label{eq:plan2-rhad-lower}
   \rhad\ge \rhodelta-c_0\sqrt{\deltaT(\Om)}
   =1-(\kappa+c_0)\sqrt{\deltaT(\Om)}.
\end{equation}
On $[\rhad,1]\subset[\rstar,1]$ the function $\rho\mapsto 1-\rho^n$
is $C^1$ with $\frac{d}{d\rho}(1-\rho^n)=-n\rho^{n-1}\in[-n,0]$, so
the mean-value theorem and \eqref{eq:plan2-rhad-lower} give
\[
   1-\rhad^{\,n}
   \le n(1-\rhad)
   \le n\,(\kappa+c_0)\sqrt{\deltaT(\Om)}
   =C'_n\sqrt{\deltaT(\Om)}.
\]
Multiplying by $\omega_n$ gives~\eqref{eq:plan2-layer-vol} with
$C_n=\omega_n C'_n$. (No appeal to $\deltaT(\Om)\le 1$ to absorb a
$\deltaT$-term into $\sqrt{\deltaT}$ is needed any more: the
window~\eqref{eq:plan2-rhad-window} is already of scale
$\sqrt{\deltaT}$, by the B1 repair of Agent~12 dated 2026-05-13.)
\end{proof}

\begin{remark}[Why only $\sqrt{\deltaT}$, not $\deltaT$]
\label{rem:plan2-why-sqrt}
The square root is sharp at the level of the layer volume itself:
$\rhodelta=1-\kappa\sqrt{\deltaT}$ is the largest radius on which the
profile-gap lower bound for the measure $\mu(d\rho)=(-t'(\rho))d\rho$
holds, and this profile gap is what produces the $\sqrt{\deltaT}$
scale of the boundary layer in the first place. Squaring this
$\sqrt{\deltaT}$ \emph{too early}, e.g. before the asymmetry transfer,
would lose the sharp exponent in the final Saint--Venant inequality.
The squaring must be done \emph{after} transfer; see
\S\ref{sssec:plan2-squaring}.
\end{remark}

\subsubsection{Superlevel asymmetry transfer}
\label{sssec:plan2-superlevel}

We now prove that the Fraenkel asymmetry of $\Om$ at scale
$|\Om|=\omega_n$ is dominated by the Fraenkel asymmetry of any nested
superlevel $E\subset\Om$ at its own scale, plus a volume-defect
correction of order $|\Om\setminus E|/\omega_n$.

\begin{lemma}[Superlevel transfer of asymmetry]
\label{lem:plan2-superlevel}
Let $\Om\subset\R^n$ be open with $|\Om|=L>0$ and let $E\subset\Om$ be
Borel with $0<|E|=s\le L$. Write $\eta:=L-s\in[0,L]$. Then
\begin{equation}\label{eq:plan2-superlevel}
   \Fr(\Om)\le \Fr(E)+\frac{2\eta}{L}
   = \Fr(E)+\frac{2}{L}\,|\Om\setminus E|.
\end{equation}
Here both Fraenkel asymmetries are computed at the natural scale of
the underlying set: $\Fr(\Om)$ uses balls of volume $L$ and $\Fr(E)$
uses balls of volume $s$.
\end{lemma}

\begin{proof}
Write $r_s:=(s/\omega_n)^{1/n}$ and $r_L:=(L/\omega_n)^{1/n}$, so the
ball of volume $s$ is $B_{r_s}$ and the ball of volume $L$ is
$B_{r_L}$. We will compare $\Om$ to a translate $z+B_{r_L}$ of the
ball of volume $L$, where $z$ is chosen nearly optimally for the
asymmetry of $E$.

Fix $\eps\in(0,1)$ and choose $z\in\R^n$ such that
\begin{equation}\label{eq:plan2-near-opt}
   |E\,\Delta\,(z+B_{r_s})|\le s\,\bigl(\Fr(E)+\eps\bigr).
\end{equation}
By the triangle inequality for symmetric differences,
\begin{equation}\label{eq:plan2-triangle}
   |\Om\,\Delta\,(z+B_{r_L})|
   \le |\Om\,\Delta\,E|
      +|E\,\Delta\,(z+B_{r_s})|
      +|(z+B_{r_s})\,\Delta\,(z+B_{r_L})|.
\end{equation}
We bound the three summands.

\emph{Term 1.} Since $E\subset\Om$, $\Om\,\Delta\,E=\Om\setminus E$, so
$|\Om\,\Delta\,E|=\eta$.

\emph{Term 2.} By~\eqref{eq:plan2-near-opt},
$|E\,\Delta\,(z+B_{r_s})|\le s(\Fr(E)+\eps)$.

\emph{Term 3.} Since $r_s\le r_L$, $z+B_{r_s}\subset z+B_{r_L}$ and
\[
   |(z+B_{r_s})\,\Delta\,(z+B_{r_L})|
   =\omega_n(r_L^n-r_s^n)=L-s=\eta.
\]

Combining and dividing by $L$,
\begin{equation}\label{eq:plan2-combine}
   \frac{|\Om\,\Delta\,(z+B_{r_L})|}{L}
   \le \frac{2\eta}{L}+\frac{s}{L}\bigl(\Fr(E)+\eps\bigr)
   \le \frac{2\eta}{L}+\Fr(E)+\eps,
\end{equation}
using $s\le L$. By the definition of the Fraenkel asymmetry,
$\Fr(\Om)\le L^{-1}|\Om\,\Delta\,(z+B_{r_L})|$, and since
$\eps\in(0,1)$ was arbitrary, \eqref{eq:plan2-superlevel} follows.
\end{proof}

\begin{remark}[Where the two boundary-layer summands come from]
\label{rem:plan2-two-summands}
Identity~\eqref{eq:plan2-triangle} contains two distinct
volume-defect terms of magnitude $\eta=|\Om\setminus E|$: the layer of
$\Om$ outside $E$, and the spherical-shell correction
$|B_{r_L}\setminus B_{r_s}|$ between the two balls. They add to give
$2\eta/L$ in~\eqref{eq:plan2-superlevel}. No $L/s$ rescaling enters,
because $\Fr(E)$ is measured at the natural scale $s$. The same
combinatorial identity appears in the $L^1$-rearrangement bookkeeping
of Maggi~\cite{Maggi2012}.
\end{remark}

\begin{corollary}[Transfer at the metric-trace level $\rhad$]
\label{cor:plan2-transfer-rhat}
For every $\Om\subset B_R$ open with $|\Om|=\omega_n$ and
$\deltaT(\Om)\le\delta_2$,
\begin{equation}\label{eq:plan2-transfer-rhat}
   \Fr(\Om)\le \Fr(E_{\rhad})+2\,C'_n\,\sqrt{\deltaT(\Om)},
\end{equation}
where $C'_n=n(\kappa+c_0)$ is the constant from
Lemma~\ref{lem:plan2-layer-vol}.
\end{corollary}

\begin{proof}
$E_{\rhad}\subset\Om$ holds by definition of the foliation. With
$L=\omega_n$ and $\eta=|\Om\setminus E_{\rhad}|$,
Lemma~\ref{lem:plan2-superlevel} gives
\[
   \Fr(\Om)\le \Fr(E_{\rhad})+\frac{2}{\omega_n}|\Om\setminus E_{\rhad}|.
\]
Inserting~\eqref{eq:plan2-layer-vol} and recalling
$C_n/\omega_n=C'_n$ yields~\eqref{eq:plan2-transfer-rhat}.
\end{proof}

\subsubsection{Squaring the transfer error: the sharp-exponent step}
\label{sssec:plan2-squaring}

We now perform the critical squaring step. The transfer
inequality~\eqref{eq:plan2-transfer-rhat} contains a transfer error
of magnitude $\sqrt{\deltaT}$. Squaring \emph{after} transfer turns
each $\sqrt{\deltaT}$ into a $\deltaT$ \emph{linearly}, while
$\Fr(E_{\rhad})^2$ inherits the $\deltaT$-rate
from~\eqref{eq:plan2-asym-from-dx}. The elementary inequality
$(a+b)^2\le 2a^2+2b^2$ is what makes this lossless on the exponent;
it only doubles the constant.

\begin{theorem}[Bounded sharp Saint--Venant, Plan~2 form]
\label{thm:plan2-bdd-SV}
There exist $C(n,R,\rstar)>0$ and
$\delta_0(n,R,\rstar)\in(0,\delta_2]$ such that for every open
$\Om\subset B_R$ with $|\Om|=\omega_n$,
\begin{equation}\label{eq:plan2-bdd-SV}
   \Fr(\Om)^{2}\le C(n,R,\rstar)\,\deltaT(\Om).
\end{equation}
The exponent $2$ is sharp. The constant unwinds as
\begin{equation}\label{eq:plan2-C-final}
   C(n,R,\rstar)
   =\max\!\left\{
     \frac{2C_*}{\omega_n^{2}}+8\bigl(C'_n\bigr)^{2},\;
     \frac{4}{\delta_0(n,R,\rstar)}
   \right\},
\end{equation}
where $C'_n=n(\kappa+c_0)$ and $C_*,c_0,\kappa,\delta_0$ are the
upstream constants of
Theorem~\ref{thm:metric-trace-input} (Agent~12) and \S4.4;
$\kappa=\kappa(n,R,\rstar)$ is canonical.
\end{theorem}

\begin{proof}
\emph{Step 1: small-deficit regime $\deltaT(\Om)\le\delta_2$.}
Set
\[
   a:=\Fr(E_{\rhad}),
   \qquad b:=2C'_n\sqrt{\deltaT(\Om)}.
\]
By Corollary~\ref{cor:plan2-transfer-rhat},
$\Fr(\Om)\le a+b$. The squaring identity gives
\begin{equation}\label{eq:plan2-square-id}
   (a+b)^2=a^2+2ab+b^2\le 2a^2+2b^2,
\end{equation}
the inequality being the AM--GM bound $2ab\le a^2+b^2$. Hence
\begin{equation}\label{eq:plan2-square}
   \Fr(\Om)^{2}\le 2\,\Fr(E_{\rhad})^{2}+2\,(2C'_n)^{2}\,\deltaT(\Om)
   =2\,\Fr(E_{\rhad})^{2}+8(C'_n)^{2}\,\deltaT(\Om).
\end{equation}
By Lemma~\ref{lem:plan2-dX-to-asym},
$\Fr(E_{\rhad})^{2}\le \omega_n^{-2}C_*\deltaT(\Om)$. Substituting,
\begin{equation}\label{eq:plan2-final-const}
   \Fr(\Om)^{2}
   \le\left(\frac{2C_*}{\omega_n^{2}}+8\bigl(C'_n\bigr)^{2}\right)\deltaT(\Om)
   =:C_3(n,R,\rstar)\,\deltaT(\Om).
\end{equation}
This proves~\eqref{eq:plan2-bdd-SV} for $\deltaT(\Om)\le\delta_2$
with constant $C_3$. The sharp exponent~$2$ is preserved because
\eqref{eq:plan2-square-id} doubles the prefactor but does not raise
$a$ or $b$ to a higher power.

\emph{Step 2: large-deficit regime $\deltaT(\Om)\ge\delta_0$.}
Set $\delta_0:=\delta_2(n,R,\rstar)$. Since
$\Fr(\Om)<2$ for every $\Om$ of volume $\omega_n$ (an immediate
consequence of $|\Om\,\Delta\,(z+B_1)|\le|\Om|+|B_1|=2\omega_n$ and the
strict inequality coming from $|\Om\cap(z+B_1)|>0$ for some $z$ when
$\Om$ has finite positive volume),
$\Fr(\Om)^{2}\le 4$. On $\deltaT(\Om)\ge\delta_0$ this gives
$\Fr(\Om)^{2}\le 4\le(4/\delta_0)\deltaT(\Om)$.

\emph{Step 3: combine.}
Taking $C:=\max\{C_3,4/\delta_0\}$ yields~\eqref{eq:plan2-bdd-SV} on
the full admissible class with constant $C$
as in~\eqref{eq:plan2-C-final}.

\emph{Sharpness.}
For the one-parameter family of volume-preserving ellipsoidal
perturbations $\Om_\eps$ of $B_1$, one has
$\Fr(\Om_\eps)\sim\eps$ and $\deltaT(\Om_\eps)\sim\eps^2$ as
$\eps\to 0$ (see e.g.\ \cite[Introduction]{BraDeP2017}); hence any
inequality $\Fr(\Om)^{p}\le C\deltaT(\Om)$ forces $p\ge 2$. This
shows that exponent~$2$ is sharp in~\eqref{eq:plan2-bdd-SV}.
\end{proof}

\begin{remark}[The squaring step, written out explicitly]
\label{rem:plan2-squaring-explicit}
The complete inequality chain at the heart of this subsection is
\begin{equation}\label{eq:plan2-chain}
\begin{aligned}
\Fr(\Om)
&\stackrel{\text{Cor.\ \ref{cor:plan2-transfer-rhat}}}{\le}
 \Fr(E_{\rhad})+2C'_n\sqrt{\deltaT(\Om)} \\[2pt]
\Fr(\Om)^{2}
&\stackrel{(a+b)^{2}\le 2a^{2}+2b^{2}}{\le}
 2\,\Fr(E_{\rhad})^{2}+8(C'_n)^{2}\deltaT(\Om) \\[2pt]
&\stackrel{\eqref{eq:plan2-asym-from-dx}}{\le}
 \frac{2C_*}{\omega_n^{2}}\deltaT(\Om)
 +8(C'_n)^{2}\deltaT(\Om)
 =C_3(n,R,\rstar)\deltaT(\Om).
\end{aligned}
\end{equation}
The transfer error $b=2C'_n\sqrt{\deltaT}$ is squared to
$b^{2}=4(C'_n)^{2}\deltaT$, and the factor of~$2$ from
$(a+b)^{2}\le 2a^{2}+2b^{2}$ produces the final $8(C'_n)^{2}$. The
$\Fr(E_{\rhad})^{2}$ term inherits a $\deltaT$-rate directly
from~\eqref{eq:plan2-metric-trace}. \emph{At no point is the
$\sqrt{\deltaT}$ from the boundary layer squared in isolation:} it is
always squared as part of $(\Fr(E_{\rhad})+\sqrt{\deltaT})^{2}$ via
the elementary squaring identity. This is what preserves the sharp
exponent~$2$.
\end{remark}

\begin{remark}[Why squaring \emph{before} transfer would fail]
\label{rem:plan2-square-before}
If, instead of squaring $\Fr(\Om)\le\Fr(E_{\rhad})+C\sqrt{\deltaT}$
\emph{after} transfer, one squared the boundary-layer estimate
$|\Om\setminus E_{\rhad}|\le C_n\sqrt{\deltaT}$ first (yielding
$|\Om\setminus E_{\rhad}|^{2}\le C_n^{2}\deltaT$), no rate would be
lost on the layer term alone, but one would still have to combine the
inequalities for $\Fr(E_{\rhad})$ and for the layer at the
$\Fr(\Om)$-level. The minimal way to combine the two contributions
into a bound for $\Fr(\Om)^{2}$ is the squaring
identity~\eqref{eq:plan2-square-id}. The point is that
~\eqref{eq:plan2-square-id} is an \emph{additive} bound on the
square: it doubles the constant but never raises an exponent. This
preservation, not the order of operations, is the load-bearing
mechanism.
\end{remark}

\subsubsection{Bounded reduction and Kohler--Jobin transfer}
\label{sssec:plan2-reduction-kj}

Theorem~\ref{thm:plan2-bdd-SV} provides the bounded-class sharp
Saint--Venant inequality at any confining radius $R\ge 2$. To pass to
the unconditional sharp Faber--Krahn inequality we apply the two
shared Plan~1 Final/ blocks \textsc{BoundedReduction} and
\textsc{KohlerJobinTransfer} in succession, exactly as in
\S\ref{ssec:plan1-final-assembly} (Agent~8).

\paragraph{Bounded reduction.}
Fix $\rstar=1/2$. Theorem~\ref{thm:plan2-bdd-SV} specialised at
$R=R_*(n):=\max\{d(n),2\}$, with $d(n)$ the diameter constant of
\cite[Lemma~5.3]{BDV}, becomes
\begin{equation}\label{eq:plan2-stage1-at-Rstar}
   \Fr(\Om)^{2}\le \frac{1}{c_{\rm SV}(n,R_*(n))}\,\deltaT(\Om)
   =\frac{1}{c_{\rm SV}(n)}\,\deltaT(\Om),
   \qquad \Om\subset B_{R_*(n)},\ |\Om|=\omega_n,
\end{equation}
where $c_{\rm SV}(n,R_*(n))=:c_{\rm SV}(n)$ depends only on $n$ by
Remark~\ref{rem:plan2-rho-star-fix} below.

\begin{theorem}[Global sharp Saint--Venant; cf.\
\cite[Theorem~5.4 (BDV bounded reduction)]{BDV}]
\label{thm:plan2-glob-SV}
Suppose~\eqref{eq:plan2-stage1-at-Rstar} holds. Then there exists
$c_{\rm SV}^{\rm glob}(n)>0$ such that for every open
$\Om\subset\R^n$ with $|\Om|=\omega_n$,
\begin{equation}\label{eq:plan2-glob-SV}
   E(\Om)-E(B_1)\;\ge\;c_{\rm SV}^{\rm glob}(n)\,\Fr(\Om)^{2},
\end{equation}
with $c_{\rm SV}^{\rm glob}(n)$ given by the BDV near/far dichotomy:
\begin{equation}\label{eq:plan2-cSVglob}
   c_{\rm SV}^{\rm glob}(n)
   =\min\!\left\{\,
       c_{\rm near}^{\rm glob}(n),\;
       \frac{\omega_n^{(n+2)/n}\,\delta(n)}{16}\right\},
\end{equation}
\[
   c_{\rm near}^{\rm glob}(n)
   =\frac{\omega_n^{(n+2)/n}}{4C_9(n)\bigl(1/\widetilde c_{\rm SV}(n)+C_9(n)\bigr)},
   \qquad
   \widetilde c_{\rm SV}(n)=\omega_n^{-(n+2)/n}\,c_{\rm SV}(n,R_*(n)),
\]
where $\delta(n),C_9(n)$ are the BDV~\cite[Lemmas 5.1--5.3]{BDV}
constants.
\end{theorem}

\begin{proof}
This is precisely the BDV bounded reduction
(\cite[Lemmas~5.1--5.3 and the proof of Theorem~1.5]{BDV}), applied
with the bounded input~\eqref{eq:plan2-stage1-at-Rstar}. The proof,
which is independent of how the bounded input was obtained, uses
the BDV constructive truncation~\cite[Lemma~5.3]{BDV}
and the rearrangement-based suboptimal $\Fr^{4}$
bound~\cite[Lemma~5.2]{BDV}. No compactness is
introduced. The constant is the same as the one used by the Plan~1
final-assembly subsection (Agent~8).
\end{proof}

By scale invariance ($\Fr$ is dimensionless, $E$ has scaling weight
$(n+2)/n$, and $\lambda_1$ has weight~$-2$),
\eqref{eq:plan2-glob-SV} extends to arbitrary $|\Om|<\infty$ in the
scale-invariant form
\begin{equation}\label{eq:plan2-glob-SV-scale}
   E(\Om)-E(\Bstar)\;\ge\;c_{\rm SV}^{\rm glob}(n)
   \Bigl(\frac{|\Bstar|}{|B_1|}\Bigr)^{(n+2)/n}\Fr(\Om)^{2}.
\end{equation}

\paragraph{Kohler--Jobin transfer.}

\begin{theorem}[Sharp Faber--Krahn from sharp Saint--Venant; cf.\
\cite[Thm.~3.3]{Brasco2014}]
\label{thm:plan2-KJ-import}
Suppose~\eqref{eq:plan2-glob-SV-scale} holds with constant
$c_{\rm SV}^{\rm glob}(n)>0$. Then for every open $\Om\subset\R^n$ with
$0<|\Om|<\infty$,
\begin{equation}\label{eq:plan2-FK-pre}
   \Bigl(\frac{|\Om|}{|B_1|}\Bigr)^{2/n}
   \bigl(\lambda_1(\Om)-\lambda_1(\Bstar)\bigr)
   \;\ge\;c_{\rm FK}(n)\,\Fr(\Om)^{2},
\end{equation}
with the explicit dimensional constant
\begin{equation}\label{eq:plan2-cFKdef}
   c_{\rm FK}(n)
   :=\min\!\left\{
       \frac{4\,L_n\,c_{\rm SV}^{\rm glob}(n)}{(n+2)\,T_n},\;\;
       \frac{L_n\bigl(2^{2/(n+2)}-1\bigr)}{4}
     \right\},
\end{equation}
where $L_n=\lambda_1(B_1)$, $T_n=T(B_1)=\omega_n/(n(n+2))$.
\end{theorem}

\begin{proof}
This is the Kohler--Jobin transfer of \cite[Theorem~3.3]{Brasco2014}.
The proof uses the Kohler--Jobin inequality
$T(\Om)\lambda_1(\Om)^{(n+2)/2}\ge T(\Bstar)\lambda_1(\Bstar)^{(n+2)/2}$
(cited from \cite[Theorem~3.3]{Brasco2014}), the pointwise transfer
$\lambda_1(\Om)\ge\lambda_1(\Bstar)(T(\Bstar)/T(\Om))^{2/(n+2)}$, the
elementary inequality $(1-s)^{-p}-1\ge ps$ for $p\in(0,1]$,
$s\in[0,1)$, and a two-case split on
$T_n-T(\Om)$ vs.\ $T_n/2$. The factor~$4$ in the first branch
absorbs the factor~$2$ between $\delta_{\rm SV}=T(\Bstar)-T(\Om)$ and
$\deltaT=E(\Om)-E(\Bstar)$ (note $\delta_{\rm SV}=2\deltaT$). The
proof is independent of how~\eqref{eq:plan2-glob-SV-scale} was
obtained; in particular it is identical to the one used by the
Plan~1 final-assembly subsection (Agent~8).
\end{proof}

\subsubsection{The final Plan~2 theorem}
\label{sssec:plan2-final}

\begin{theorem}[Sharp quantitative Faber--Krahn, Plan~2 final form]
\label{thm:plan2-FK}
Let $n\ge 2$. There exists a constant $c_{\rm FK}(n)>0$, depending
only on $n$, given by the explicit formula~\eqref{eq:plan2-cFKdef}
above, such that for every open $\Om\subset\R^n$ with
$0<|\Om|<\infty$,
\begin{equation}\label{eq:plan2-FK-main}
   \deltaFK(\Om)
   \;=\;
   \Bigl(\frac{|\Om|}{\omega_n}\Bigr)^{2/n}
   \bigl(\lambda_1(\Om)-\lambda_1(\Bstar)\bigr)
   \;\ge\;c_{\rm FK}(n)\,\Fr(\Om)^{2}.
\end{equation}
The exponent~$2$ is sharp.
\end{theorem}

\begin{proof}
Fix $\rstar=1/2$ and apply Theorem~\ref{thm:plan2-bdd-SV} at
$R=R_*(n)$. This produces the bounded sharp Saint--Venant
inequality~\eqref{eq:plan2-stage1-at-Rstar} with constant
$c_{\rm SV}(n,R_*(n))=:c_{\rm SV}(n)$. Theorem~\ref{thm:plan2-glob-SV}
removes the $R$-dependence and gives the global sharp Saint--Venant
inequality~\eqref{eq:plan2-glob-SV-scale} with constant
$c_{\rm SV}^{\rm glob}(n)$ as in~\eqref{eq:plan2-cSVglob}.
Theorem~\ref{thm:plan2-KJ-import} converts the latter into the sharp
Faber--Krahn inequality~\eqref{eq:plan2-FK-main} with constant
$c_{\rm FK}(n)$ as in~\eqref{eq:plan2-cFKdef}.

Sharpness of the exponent~$2$ follows from the volume-preserving
ellipsoidal one-parameter family $\Om_\eps$ of perturbations of
$\Bstar$, for which
$\Fr(\Om_\eps)\sim\eps$ and
$\lambda_1(\Om_\eps)-\lambda_1(\Bstar)\sim\eps^{2}$ as $\eps\to 0$
(cf.\ \cite[Introduction]{BraDeP2017},
\cite{AKN2023}).
\end{proof}

\begin{remark}[Identical normalisation to Plan~1]
\label{rem:plan2-vs-plan1-normalisation}
The statement and constant of Theorem~\ref{thm:plan2-FK} agree
\emph{verbatim} with those of the Plan~1 final theorem
(Theorem~\ref{thm:plan1-FK-upper-intro} of the Introduction). The shared expression for
$c_{\rm FK}(n)$ is~\eqref{eq:plan2-cFKdef}; the shared constants
$c_{\rm SV}^{\rm glob}(n)$, $L_n$, $T_n$, $\omega_n$ are dimensional.
The two proofs differ \emph{only} in how the bounded sharp
Saint--Venant input is obtained: Plan~1 uses selection
principle + graph entry + Schauder interpolation + nearly-spherical
closure; Plan~2 uses the level-set / metric route assembled in
\S\ref{ssec:plan2-level-set}--\S\ref{ssec:plan2-weighted-trace} and
the present subsection.
\end{remark}

\begin{remark}[Where the constants enter]
\label{rem:plan2-rho-star-fix}
The constant $C_3(n,R,\rstar)$ from~\eqref{eq:plan2-final-const}
unwinds as
\[
   C_3(n,R,\rstar)
   =\frac{2C_*(n,R,\rstar)}{\omega_n^{2}}
    +8\,n^{2}\bigl(\kappa(n,R,\rstar)+c_0(n,R,\rstar)\bigr)^{2},
\]
polynomial in $1/\rstar$, $R$, $n$ and the upstream constants
$C_*,c_0,\kappa$ of Theorem~\ref{thm:metric-trace-input} and \S4.4.
At $\rstar=1/2$ and $R=R_*(n)$, $C_3$ becomes a dimensional constant
$C_3(n)$, and Theorem~\ref{thm:plan2-bdd-SV} reads
$\Fr(\Om)^{2}\le C_3(n)\deltaT(\Om)$ on
$\Om\subset B_{R_*(n)}$, $|\Om|=\omega_n$. Equivalently
$c_{\rm SV}(n)=1/C_3(n)$, and this is the input fed to
Theorem~\ref{thm:plan2-glob-SV}.
\end{remark}

\begin{remark}[No load-bearing use of conditional routes]
\label{rem:plan2-no-conditional}
The proof of Theorem~\ref{thm:plan2-FK} does \emph{not} invoke
the centroid $\Pi$-route or the space--time $\Pi$ identity in any
load-bearing way; those appear only as auxiliary remarks
(Agent~14). It also does \emph{not} use any Bernoulli spectral input
or any $C^{1,\gamma}$ regularity of $\partial\Om$. The proof is
purely finite-perimeter from \S\ref{ssec:plan2-level-set} through
the present subsection. This is the structural difference between
Plan~2 and Plan~1, recorded in
Remark~\ref{rem:plan2-vs-plan1-normalisation}.
\end{remark}

\label{subsec:boundary-layer-transfer}%
\label{sec:plan2-boundary-layer}%
\label{sec:plan2-assembly}%
\label{sec:plan2-global-assembly}%
\label{sec:plan2-setup}%
\label{ssec:plan2-boundary-transfer}%
\label{ssec:plan2-final}%
\label{ssec:plan1-final-assembly}%

% UNIFIED 2026-05-13
% SOURCE MAP:
%   Source files:
%     - Plan 2/WIP/WIP_CentroidBound.tex (centroid kinematic identity; Thm F)
%     - Plan 2/WIP/WIP_SpaceTimeIdentity.tex (space-time Pi identity; G1, G3, T1^2-int)
%     - Plan 2/wave3-F-h1-center-bound.md (proof-level details for the centroid H^1 bound)
%   Status of each item:
%     - Centroid kinematic identity (Lem. centroid-kin): LOCAL(WIP_CentroidBound.tex Lem. 3.1)
%     - Lipschitz/Borel regularity of rho |-> z̄_ρ: LOCAL(WIP_CentroidBound.tex Lem. 2.2)
%     - H^1-in-ρ centroid bound (F): LOCAL(WIP_CentroidBound.tex Thm. F)
%       [Unconditional within Plan 2's small-deficit regime;
%        the section is nevertheless presented as AUXILIARY because it feeds
%        the conditional Pi-route.]
%     - Space-time Pi identity (G1): LOCAL(WIP_SpaceTimeIdentity.tex Thm. G1)
%     - Pi-control hypothesis (Pi-int): MISSING / CONDITIONAL
%         (open problem; brief Sec. G2 of WIP_SpaceTimeIdentity)
%     - (B^2-int) from (Pi-int): LOCAL(WIP_SpaceTimeIdentity.tex Thm. G3)
%     - (T_1^2-int) from (Pi-int) + annular cap: LOCAL(WIP_SpaceTimeIdentity.tex Thm. T1sq)
%   Notation: follows Manuscript/00_notation_register.md.
%     - Centre field z̄_ρ (bulk centroid), macros \zbar, \Frho, \Erho, \deltaT.
%     - Domain of the space-time identity renamed Ω_Π (= \OmPi), per
%       00_notation_register.md §13 conflict resolution (Plan 1's \Otilde is reserved).
%
% STATUS: AUXILIARY — NOT LOAD-BEARING. CONDITIONAL ON INTEGRATED Π-CONTROL.
%
% Author note (for Agent 19): this draft is intended for inclusion as the
% optional remark / final subsection of Manuscript/03_plan2_draft.tex. Agent 16
% (Plan 2 audit) must verify that no later subsection of 03_plan2_draft.tex
% imports any statement from this subsection. The repaired Plan 2 chain
% (Agents 9–13) does NOT depend on anything below.
%
% ---------------------------------------------------------------------------

\subsection{Auxiliary route: centroid kinematics and the space--time
\texorpdfstring{$\Pi$}{Pi} identity}
\label{subsec:plan2-aux-centroid}

\begin{quote}\small
\textbf{Status.}
The contents of this subsection are auxiliary. They are recorded here
because they isolate a structurally appealing alternative route to the
sharp bounded Saint--Venant estimate
$\Fr(\Om)^{2}\le C(n,R,\rstar)\,\deltaT(\Om)$,
based on the centroid first moment of the level family. The route
closes \emph{conditionally} on the integrated radial-moment estimate
\eqref{eq:Piint} below, which is currently open. The load-bearing
proof of Plan~2 in Sections \ref{subsec:level-set-identity}--%
\ref{subsec:boundary-layer-transfer} does \emph{not} use any
statement from this subsection.
\end{quote}

\paragraph{Standing hypotheses}

Throughout this subsection we adopt the standing setup of Section
\ref{sec:plan2-setup}: $n\ge 2$, $R\ge 2$, $\rstar\in(0,1)$, $\kappa>0$,
$\Om\subset B_R\subset\R^n$ open with $|\Om|=\omega_n$,
$u=u_\Om$ the variational torsion function,
$t:[\rstar,1]\to(0,\|u\|_\infty)$ the volume-radius parametrisation
($|\{u>t(\rho)\}|=\omega_n\rho^n$), $\Erho:=\{u>t(\rho)\}$,
$\rhodelta:=1-\kappa\sqrt{\deltaT(\Om)}$,
$\mu(d\rho):=(-t'(\rho))\,d\rho$, and
$\Vrho(x):=-t'(\rho)/|\nabla u(x)|$ ($\mathcal H^{n-1}$-a.e.\ on
$\partial^{*}\Erho$, for a.e.\ $\rho$, by Sobolev--Sard, cf.\ §A3).
The constant $\delta_0=\delta_0(n,R,\rstar)>0$ is fixed once and for
all small enough that all qualitative consequences quoted below (FMP
smallness, Cicalese--Leonardi containment, centroid--Fraenkel-centre
proximity) hold whenever $\deltaT(\Om)\le\delta_0$.

The bulk centroid of $\Erho$ is
\begin{equation}\label{eq:aux-centroid-def}
   \zbar \;:=\; \frac{1}{|\Erho|}\int_{\Erho} x\,dx
   \;\in\;\R^{n},
\end{equation}
in line with the canonical centre field of the notation register
(§10). We work with $\zbar$ rather than with the Fraenkel-optimal
centre $z^{\mathrm{Fr}}_{\rho}$ or with the Fusco--Julin oscillation
centre.

\paragraph{Centroid kinematic identity}

\begin{lemma}[Centroid kinematic identity]\label{lem:centroid-kin}
The map $\rho\mapsto\zbar$ is Lipschitz on $[\rstar,1]$ with
\[
   |\zbar-\bar z_{\rho'}|
   \;\le\;
   \frac{2nR}{\rstar^{n}}\,|\rho-\rho'|,
   \qquad \rstar\le\rho,\rho'\le 1,
\]
hence absolutely continuous on $[\rstar,1]$, and for a.e.\
$\rho\in[\rstar,1]$ one has the first-moment identity
\begin{equation}\label{eq:centroid-kin}
   \zbar^{\prime}
   \;=\;
   \frac{1}{|\Erho|}
   \int_{\partial^{*}\Erho}
   (x-\zbar)\,\Vrho(x)\,d\mathcal H^{n-1}(x).
\end{equation}
\end{lemma}

\begin{proof}
\emph{Lipschitz regularity.} The volume-radius parametrisation gives
the nested inclusion $E_{\rho}\subset E_{\rho'}$ whenever
$\rho\le\rho'$, hence
\[
   |E_{\rho}\Delta E_{\rho'}|
   = |E_{\rho'}|-|E_{\rho}|
   = \omega_{n}\bigl((\rho')^{n}-\rho^{n}\bigr)
   \le n\omega_{n}\,|\rho-\rho'|.
\]
Since $\Erho\subset B_{R}$, $|\int_{E_{\rho}}x_{i}-\int_{E_{\rho'}}x_{i}|
\le R\,|E_{\rho}\Delta E_{\rho'}|\le n\omega_{n}R\,|\rho-\rho'|$.
Dividing by $|\Erho|\ge\omega_{n}\rstar^{n}$ yields the displayed
Lipschitz constant $L_{\mathrm{cent}}=2nR/\rstar^{n}$.

\emph{Kinematic identity.} By the coarea identity applied to the
Sobolev function $u\in W^{1,2}_{0}(\Om)\cap L^{\infty}(\Om)$ (see
§A2--A3 of the GMT preliminaries; the Sobolev--Sard step
\cite{FigalliMaggi2011} ensures $|\nabla u|>0$
$\mathcal H^{n-1}$-a.e.\ on $\partial^{*}\Erho$ for a.e.\ $\rho$),
for every Lipschitz $f:\overline{B_{R}}\to\R$ and a.e.\
$\rho\in[\rstar,1]$,
\begin{equation}\label{eq:bulk-coarea}
   \frac{d}{d\rho}\int_{\Erho}f(x)\,dx
   \;=\;
   \int_{\partial^{*}\Erho} f(x)\,\Vrho(x)\,d\mathcal H^{n-1}(x).
\end{equation}
Taking $f\equiv 1$ and using $|\Erho|=\omega_{n}\rho^{n}$ gives
$\int_{\partial^{*}\Erho}\Vrho\,d\mathcal H^{n-1}=n\omega_{n}\rho^{n-1}
=P(B_{\rho})$ for a.e.\ $\rho$. Taking $f(x)=x_{i}$ in
\eqref{eq:bulk-coarea} and applying the quotient rule to
$\bar z_{\rho,i}=|\Erho|^{-1}\int_{\Erho}x_{i}\,dx$:
\begin{align*}
   \bar z_{\rho,i}^{\prime}
   &= \frac{1}{|\Erho|}\int_{\partial^{*}\Erho}x_{i}\,\Vrho\,d\mathcal H^{n-1}
   -\frac{1}{|\Erho|^{2}}
   \Bigl(\frac{d}{d\rho}|\Erho|\Bigr)\int_{\Erho}x_{i}\,dx \\
   &= \frac{1}{|\Erho|}\int_{\partial^{*}\Erho}x_{i}\,\Vrho\,d\mathcal H^{n-1}
   -\frac{\bar z_{\rho,i}}{|\Erho|}\int_{\partial^{*}\Erho}\Vrho\,d\mathcal H^{n-1}\\
   &= \frac{1}{|\Erho|}\int_{\partial^{*}\Erho}(x_{i}-\bar z_{\rho,i})\,\Vrho\,d\mathcal H^{n-1},
\end{align*}
which is the $i$-th component of \eqref{eq:centroid-kin}.
\end{proof}

\begin{remark}[No homothety term]\label{rmk:no-homothety}
The right-hand side of \eqref{eq:centroid-kin} is the first vector
moment of $\Vrho$ relative to $\zbar$, with no comparison to a
homothetic velocity. This is what makes the centroid kinematic
identity unconditional: it does \emph{not} require a per-$\rho$
control of $\Vrho-H_{\zbar,\rho}$, where
$H_{\zbar,\rho}(x):=(x-\zbar)\cdot\nu_{\rho}(x)/\rho$, and therefore
avoids the pointwise radial-trace inputs that the conditional
$\Pi$-route below relies on. A centre defined by a Fraenkel
first-order condition would instead produce a derivative formula
involving $\Vrho-H_{\zbar,\rho}$ on the reduced boundary.
\end{remark}

\paragraph{The integrated centroid bound (unconditional within
the small-deficit regime)}

\begin{theorem}[Integrated centroid $H^{1}$ bound]\label{thm:centroid-H1}
There exist $\delta_{0}=\delta_{0}(n,R,\rstar)>0$ and
$C_{\mathrm{cent}}=C_{\mathrm{cent}}(n,R,\rstar)>0$ such that
whenever $\Om\subset B_{R}$ is open with $|\Om|=\omega_{n}$ and
$\deltaT(\Om)\le\delta_{0}$,
\begin{equation}\label{eq:F}
   \int_{\rstar}^{\rhodelta}|\zbar^{\prime}|^{2}\,(-t'(\rho))\,d\rho
   \;\le\;
   C_{\mathrm{cent}}\,\deltaT(\Om). \tag{F}
\end{equation}
\end{theorem}

\noindent
A sketch of the proof structure, conditional on
Hypothesis~\ref{hyp:Piint}, is the following:
\begin{enumerate}
\item[(i)] decompose
$\Vrho=\bar V_{\rho}+(\Vrho-\bar V_{\rho})$ with
$\bar V_{\rho}=P(B_{\rho})/P(\Erho)$ in \eqref{eq:centroid-kin};
\item[(ii)] bound the perimeter mean piece by FMP plus the
Fusco--Julin $L^{1}$ oscillation, using the divergence theorem for
finite-perimeter sets applied to the Lipschitz field
$V^{(i)}(x):=|x-\zbar|\,e_{i}$;
\item[(iii)] bound the variance piece by Cauchy--Schwarz, controlled
by the algebraic deficit $D_{H}(t(\rho))$ of
Subsection \ref{subsec:level-set-identity};
\item[(iv)] integrate against $\mu(d\rho)=(-t')\,d\rho$ and use the
integrated identity $\int_{\rstar}^{1}(D_{H}+D_{I})\,d\mu\le
C(n,\rstar)\,\deltaT(\Om)$.
\end{enumerate}
The constant tracks as
$C_{\mathrm{cent}}\le C(n,\rstar)\,R\,\rstar^{-(3n-1)}$ and is
polynomial in $1/\rstar$, $R$, $n$. Outside the smallness regime FMP
gives $\Fr(\Om)\ge c(n)>0$ and \eqref{eq:F} is trivial after
enlarging $C_{\mathrm{cent}}$.

\paragraph{Space--time \texorpdfstring{$\Pi$}{Pi} identity}

For each $\rho\in[\rstar,\rhodelta]$, define, in agreement with the
notation register (§13),
\begin{equation}\label{eq:Pi-def}
   \PiE
   \;:=\;
   (n-1)\int_{\R^{n}}
   \bigl(\mathbf{1}_{B_{\rho}(\zbar)}(x)-\mathbf{1}_{\Erho}(x)\bigr)\,
   |x-\zbar|^{-1}\,dx,
\end{equation}
the signed radial moment of the symmetric difference
$\Erho\,\Delta\,B_{\rho}(\zbar)$ at the centroid. By the
$L^{2}$ divergence identity of
Subsection \ref{subsec:level-set-identity},
\begin{equation}\label{eq:53prime}
   \int_{\partial^{*}\Erho}|\nu_{\rho}-e_{r}|^{2}\,d\mathcal H^{n-1}
   \;=\;
   2\bigl(P(\Erho)-P(B_{\rho})\bigr)+2\,\PiE,
   \qquad e_{r}(x)=\frac{x-\zbar}{|x-\zbar|},
\end{equation}
so $\PiE\ge 0$ a.e.\ in $\rho$.

The natural cylinder for the space--time computation is
$\OmPi:=B_{R+1}\subset\R^{n}$, which contains both $\Om$ and
$\bigcup_{\rho\in[\rstar,\rhodelta]}B_{\rho}(\zbar)$ (recall
$\zbar\in B_{R}$ by $\Erho\subset B_{R}$, and $\rho\le 1$). We use
$\OmPi$ rather than $\Otilde$ to respect the manuscript-wide notation
agreement (§13 of the notation register), where $\Otilde$ is reserved
for Plan~1's volume-normalised companion.

\begin{theorem}[Space--time $\Pi$ identity]\label{thm:G1}
The signed Fubini density
\begin{equation}\label{eq:sigma-def}
   \sigma(x)
   \;:=\;
   \int_{\rstar}^{\rhodelta}
   \bigl(\mathbf{1}_{B_{\rho}(\zbar)}(x)-\mathbf{1}_{\Erho}(x)\bigr)\,
   |x-\zbar|^{-1}\,d\mu(\rho)
\end{equation}
is Borel-measurable, locally integrable on $\OmPi$, and satisfies
\begin{equation}\label{eq:spacetime-Pi}
   \int_{\rstar}^{\rhodelta}\PiE\,d\mu(\rho)
   \;=\;
   (n-1)\int_{\OmPi}\sigma(x)\,dx
   \;=\;
   (n-1)\int_{\OmPi}\int_{J_{x}}
   \mathrm{sgn}\bigl(\mathbf{1}_{B_{\rho}(\zbar)}(x)-\mathbf{1}_{\Erho}(x)\bigr)
   \,|x-\zbar|^{-1}\,d\mu(\rho)\,dx,
\end{equation}
where $J_{x}:=\{\rho\in[\rstar,\rhodelta]\colon
x\in B_{\rho}(\zbar)\Delta\Erho\}$. In particular,
\begin{equation}\label{eq:sigma-pointwise}
   |\sigma(x)|
   \;\le\;
   \int_{J_{x}}|x-\zbar|^{-1}\,d\mu(\rho)
   \qquad\text{for a.e.\ }x\in\OmPi.
\end{equation}
\end{theorem}

\begin{proof}
By \eqref{eq:Pi-def}, $\rho\mapsto\PiE$ is the integral over $\R^{n}$
of the joint Borel-measurable function
$(\rho,x)\mapsto(\mathbf{1}_{B_{\rho}(\zbar)}(x)-\mathbf{1}_{\Erho}(x))/|x-\zbar|$.
The kernel $|x-\zbar|^{-1}$ is locally integrable on $\OmPi$ since
$n\ge 2$. Truncating the kernel at level $M$, applying Tonelli to the
positive and negative parts of the bounded integrand, and letting
$M\to\infty$ by monotone convergence on each sign separately produces
the first equality in \eqref{eq:spacetime-Pi}. The second equality is
the change of order on the symmetric difference. The bound
\eqref{eq:sigma-pointwise} follows by taking absolute values inside
the $\rho$-integral.
\end{proof}

\paragraph{Where the centroid route would close: the conditional
\texorpdfstring{$\Pi$}{Pi}-control hypothesis}

The integrated identity \eqref{eq:spacetime-Pi} rewrites
$\int\PiE\,d\mu$ as an integral, over the cylinder $\OmPi$, of the
weighted space--time crossing length $\mu(J_{x})$ of the moving sphere
$|x-\zbar|=\rho$ with the boundary of $\Erho$. Closing the centroid
route at the sharp rate $\deltaT$ requires the following hypothesis.

\begin{hypothesis}[Integrated $\Pi$-control; \emph{open}]\label{hyp:Piint}
There exists $C_{\Pi}=C_{\Pi}(n,R,\rstar)>0$ such that for every
bounded open $\Om\subset B_{R}$ with $|\Om|=\omega_{n}$ and
$\deltaT(\Om)\le\delta_{0}(n,R,\rstar)$,
\begin{equation}\label{eq:Piint}
   \int_{\rstar}^{\rhodelta}\PiE\,d\mu(\rho)
   \;\le\;
   C_{\Pi}(n,R,\rstar)\,\deltaT(\Om).
   \tag{$\Pi$-int}
\end{equation}
\end{hypothesis}

\begin{quote}\small
\textbf{Status of Hypothesis \ref{hyp:Piint}.}
At the time of writing, \eqref{eq:Piint} is \emph{not} proved. The
$H^{1}$-in-$\rho$ centroid bound \eqref{eq:F} controls the moving
centre $\zbar$ but does not, by itself, convert the rearranged
Talenti/Nicola--Tilli profile gap into a control on the physical
crossing set $J_{x}$, because $J_{x}$ also depends on the spatial
radius $|x-\zbar|$. A viable proof of \eqref{eq:Piint} would combine
\eqref{eq:spacetime-Pi} with a genuine space--time coarea or
transport estimate for the level crossings of
$|x-\zbar|=\rho$. We do \emph{not} have such an estimate.

\textbf{In particular, none of the consequences below
(Theorems \ref{thm:G3} and \ref{thm:T1sq}) are used in the
load-bearing Plan~2 chain.}
\end{quote}

\paragraph{Conditional consequences}

For completeness we record the two conditional consequences. Both
assume Hypothesis \ref{hyp:Piint}.

\begin{theorem}[Integrated $L^{2}$ normal oscillation;
conditional on \eqref{eq:Piint}]\label{thm:G3}
Assume Hypothesis \ref{hyp:Piint}. Then there exists
$C_{B^{2}}=C_{B^{2}}(n,R,\rstar)>0$ such that
\begin{equation}\label{eq:B2int}
   \int_{\rstar}^{\rhodelta}\!\!
   \int_{\partial^{*}\Erho}|\nu_{\rho}-e_{r}|^{2}\,d\mathcal H^{n-1}\,d\mu(\rho)
   \;\le\;
   C_{B^{2}}\,\deltaT(\Om).
   \tag{$B^{2}$-int}
\end{equation}
\end{theorem}

\begin{proof}
Integrate \eqref{eq:53prime} against $d\mu$. The first piece is
controlled by the deficit identity
$\int(P(\Erho)-P(B_{\rho}))\,d\mu(\rho)\le C(n,\rstar)\,\deltaT(\Om)$
of Subsection \ref{subsec:level-set-identity} (it bounds
$\int D_{I}\,d\mu$, and $P-P_\rho$ is comparable to
$D_{I}/P(B_{\rho})$ in the small-deficit regime, with
$P(B_{\rho})\ge n\omega_{n}\rstar^{n-1}$). The second piece is
exactly \eqref{eq:Piint}.
\end{proof}

\begin{theorem}[Integrated radial trace;
conditional on \eqref{eq:Piint} plus an annular cap]\label{thm:T1sq}
Assume Hypothesis \ref{hyp:Piint} and the auxiliary annular/perimeter
cap
\begin{equation}\label{eq:annular-cap}
   \rstar/2 \,\le\,|x-\zbar|\,\le\, 2R
   \quad\text{on the relevant boundary pieces},
   \qquad
   P(\Erho)-P(B_{\rho})+\PiE\,\le\, M_{A}(n,R,\rstar)
\end{equation}
for a.e.\ $\rho\in[\rstar,\rhodelta]$. Then there exists
$C_{\mathrm{rad}}=C_{\mathrm{rad}}(n,R,\rstar)>0$ such that
\begin{equation}\label{eq:T1sqint}
   \int_{\rstar}^{\rhodelta}\Bigl(\int_{\partial^{*}\Erho}
   |H_{\zbar,\rho}-1|\,d\mathcal H^{n-1}\Bigr)^{2}d\mu(\rho)
   \;\le\;
   C_{\mathrm{rad}}\,\deltaT(\Om).
   \tag{$T_{1}^{2}$-int}
\end{equation}
\end{theorem}

\noindent
The proof of Theorem \ref{thm:T1sq} uses a slicing-then-squaring
decomposition $T_{1}=T_{1}^{\mathrm{rad,slic}}+T_{1}^{\mathrm{tan}}$
of the radial trace
$T_{1}=\int_{\partial^{*}\Erho}||x-\zbar|/\rho-1|\,d\mathcal H^{n-1}$,
followed by a slice-rearrangement inequality
$|\Erho\Delta B_{\rho}(\zbar)|^{2}\le C(n,R,\rstar)\,\PiE$;
the role of squaring before
integrating is to avoid the FMP rate-loss to $\sqrt{\deltaT}$.
The full argument is conditional on Hypothesis~\ref{hyp:Piint} and is
not used by Subsections
\ref{subsec:metric-framework}--\ref{subsec:boundary-layer-transfer}.

\paragraph{Reading guide and warning}

\begin{quote}\small
\textbf{What this subsection establishes unconditionally.}
The centroid kinematic identity \eqref{eq:centroid-kin}, its
Lipschitz regularity, the integrated centroid bound \eqref{eq:F}, and
the space--time identity \eqref{eq:spacetime-Pi}.

\textbf{What this subsection does \emph{not} establish.}
The hypothesis \eqref{eq:Piint}, and therefore neither
\eqref{eq:B2int} nor \eqref{eq:T1sqint}.

\textbf{What this subsection is \emph{not} used for.}
Nothing in this subsection is invoked in the proofs of
Subsections \ref{subsec:level-set-identity},
\ref{subsec:metric-framework},
\ref{subsec:fusco-julin-radial-trace},
\ref{subsec:integrated-action-weighted-trace}, or
\ref{subsec:boundary-layer-transfer}. The repaired Plan~2 chain runs
through the weighted metric trace
(Subsection \ref{subsec:integrated-action-weighted-trace})
and uses the centroid $\zbar$ as the canonical centre field, but it
does \emph{not} rely on the radial-moment identities
\eqref{eq:53prime}--\eqref{eq:Piint}. The final Plan~2 theorem of
Subsection \ref{subsec:boundary-layer-transfer} stands without
Hypothesis \ref{hyp:Piint}.

\textbf{Why the route is interesting nonetheless.}
The structural mechanism explained in
Remark \ref{rmk:no-homothety} --- a centre field whose
derivative is a bulk first moment with no homothety term ---
suggests that if a separate space--time crossing-length estimate
for $|x-\zbar|=\rho$ can be proved at rate $\deltaT$, then the
present subsection would furnish an independent path to the sharp
bounded Saint--Venant estimate, bypassing the weighted metric trace
of Subsection \ref{subsec:integrated-action-weighted-trace}. This is
the only sense in which the centroid/$\Pi$ route remains an open
alternative.
\end{quote}

% End of optional auxiliary subsection.

\label{ssec:plan2-aux-centroid}%
\label{sec:audit-plan2}%

% =====================================================================
% End of Manuscript/03_plan2_draft.tex
% =====================================================================


% =====================================================================
% References
% =====================================================================
\bibliographystyle{plain}
\bibliography{references}
\label{sec:references}

\end{document}
